% 本文件由 LaTeX 排版 Agent 根据有效 translation.json 自动生成。
% author=Ambrose; work_id=3404; title=论信德

\chapter{论信德}

% structure: p0001 source=h1 target=section reason=page-title

\section{基督教信仰阐释，卷一}

% structure: p0002 source=h2 target=omitted reason=duplicate-page-title

% structure: p0003 source=h3 target=subsection reason=relative-heading-rank; base=h3

\subsection{序言}

在出发前往\Place{东方}的前夕，为帮助其叔父\Person{瓦伦斯}抵御哥特人的入侵，西罗马皇帝\Person{格拉齐安}请求\Person{圣盎博罗削}为他撰写一篇论著，以证明耶稣基督的天主性。\Person{格拉齐安}提出这一请求的目的，是要获得某种防范措施，以抵御\OriginalTerm{亚略异端}{Arianism}的腐蚀性影响；在那时（公元378年），由于\OriginalTerm{亚略异端}{Arianism}在帝国宫廷中的确立，它已在帝国东方行省取得了对\OriginalTerm{正统信仰}{Orthodoxy}的优势。为满足\Person{格拉齐安}的愿望，这位米兰主教撰写了一篇论著，该论著现在构成了\WorkTitle{论信仰}的前两卷。皇帝对此论著颇为满意，以致于当他从\Place{东方}归来，在\Person{瓦伦斯}死于\Place{哈德良堡}之后，他写信给\Person{圣盎博罗削}，请求一份新的论著副本，并进一步请求增补一篇关于圣神天主性的论述，以扩充该论著。\Person{圣盎博罗削}确实在379年扩充了这部原论著，但新增的各卷并未涉及圣神的天主性，而是处理了\OriginalTerm{亚略异端}{Arianism}教师们提出的新异议，以及那些此前已被忽略或尚未充分讨论的要点。\Person{圣盎博罗削}的这部阐释便以此方式形成了它目前的样貌。\par

如前所述，这部《阐释》的目的是证明耶稣基督的\OriginalTerm{天主性}{Divinity}，以及祂作为天主\OriginalTerm{圣子}{Son}，与天主\OriginalTerm{圣父}{Father}的\OriginalTerm{同永恒}{co-eternity}、\OriginalTerm{同等}{co-equality}和\OriginalTerm{同性同体}{consubstantiality}。作者通过不断援引圣经——包括旧约和新约——来达成此目的；而\OriginalTerm{亚略异端者}{Arians}在许多情况下，都曾强迫圣经经文套入错误的解释模式，以使它们符合其教义。\par

除了\WorkTitle{论信仰}这个书名外，\WorkTitle{论圣三}也是这部论著在后世广为人知的另一个书名；不过可以确定的是，前一个书名是\Person{圣盎博罗削}本人所指定的。\par

% structure: p0007 source=h3 target=subsection reason=relative-heading-rank; base=h3

\subsection{前言说明}

在\WorkTitle{论信仰}前四卷的注释中，有些地方借鉴了\Person{胡特尔}神父的这部论著版本（\Place{因斯布鲁克}：\OriginalTerm{瓦格纳}{Wagner}出版社）中的注释，该版本在准备这些卷次的翻译时曾被采用。这些注释以末尾标注的字母\Quote{H}加以区分。\par

正文中引用的圣经经文，均按其原文措辞翻译。这将解释为何与英文圣经及公祷书中的译文存在任何差异，而其中的一致之处则可归因于对更熟悉译本的记忆。我们认为，采用这种方式处理\Person{圣盎博罗削}的引用是恰当的，因为这些差异值得注意，而且在某些情况下，论证的关键正有赖于此。书中的参引一律指向英文圣经的章节，而非拉丁通俗译本，除非另有特别说明。\par

序言及目录摘要均基于\Person{胡特尔}神父版本。\par

% structure: p0011 source=h3 target=subsection reason=relative-heading-rank; base=h3

\subsection{圣盎博罗削的基督教信仰阐释。}

% structure: p0012 source=h3 target=omitted reason=duplicate-page-title

\Emph{作者赞扬\Person{格拉齐安}皇帝渴慕信仰教导的热忱，并谦逊地谈及自己才疏学浅。皇帝既受天主亲自教导，无需人的指教；然而，他的虔诚已为胜利铺路。作者被委任的任务颇为艰巨；在完成此任时，他将更多地依靠权威而非理智与论证，特别是尼西亚大公会议的权威。}\par

1. 我们在\BibleBook{列王}{纪}书中读到，南方的女王前来聆听\Person{撒罗满}的智慧。\BibleRef{列上 10:1}同样，\Person{希兰}王也曾遣人往见\Person{撒罗满}，为要试验他。\BibleRef{列上 5:1}同样，你神圣的陛下，因循前代榜样，也已颁旨听取我的信仰宣认。但我并非\Person{撒罗满}，好叫你仰慕我的智慧；而你的陛下也并非一国之君；乃是统治整个世界的\Person{奥古斯都}，已下令将信仰写成一部书，并非为教导你，而是为获得你的认可。\par

2. 尊贵的皇帝啊，你的陛下自幼便一直虔诚而充满爱意地持守那信仰，又何须学习呢？经上说：\BibleQuote{我还没有在母腹内形成你以前，我已认识了你；在你还没有出离母胎以前，我已祝圣了你。}因此，圣化并非源自传承，而是源自灵感；因此，你要谨守天主的恩赐。因为那无人教导于你的，天主必已赐予并启示了你。\par

3. 你神圣的陛下即将出征作战，便要求我写一部书来阐释信仰，因为陛下深知，胜利更多地取决于对统帅的信仰，而非士兵们的英勇。因为\Person{亚巴郎}率领三百一十八人出战，却带回了无数敌人的战利品；他凭着我们主的十字架和圣名标志的大能，战胜了五位君王和数支得胜大军的势力，既为邻人报了仇，又获得了胜利，并赎回了他的侄儿。同样，\Person{农}的儿子\Person{若苏厄}，当他率领全军之力仍无法战胜敌人时，\BibleRef{苏 6:6}却在看见并认识天上万军统帅的地方，凭着七只圣号角的声音取得了胜利。因此，陛下作为基督的忠仆和信仰的捍卫者，正在为胜利做准备，而陛下愿意让我将信仰付诸文字。\par

4. 诚然，我宁愿承担劝勉持守信仰的责任，而不愿对此进行辩论；因为前者意味着虔诚的宣认，而后者容易导致鲁莽的臆断。不过，既然陛下无需劝勉，而我又不能祈求免除尽忠的义务，我虽将着手一项大胆的事业，却仍保持谦逊，与其说是对信仰进行推理和辩论，不如说是汇集大量的见证。\par

5. 在各次大公会议的决议中，我将以那次大会为主要指引，那是由三百一十八位司铎，仿佛依照\Person{亚巴郎}的裁决而被指派，他们（可谓）树立了一座凯旋纪念碑，以宣示他们对普世不信者的胜利，他们所凭借的正是那信仰的勇气，且在其中众人皆意见一致。确实，在我看来，人们在其中可以看到天主之手，因为同样的数目既是我们在信仰的公会议中的权威，也是古时记录中忠诚的榜样。\par

\Emph{作者将信仰与外教人、犹太人和异端者的错误区分开来，并在解释了\Quote{天主}和\Quote{主}之名的意义后，清楚地显示了在\OriginalTerm{本质}{Essence}同一中的\OriginalTerm{位格}{Persons}差异。在分割\OriginalTerm{本质}{Essence}时，\OriginalTerm{亚略派}{Arians}不仅引入了\OriginalTerm{三个神}{three Gods}的教义，甚至推翻了\OriginalTerm{天主圣三}{Trinity}的主权。}\par

6. 这就是我们信仰的宣言：我们说天主是唯一的，既不像外教人那样将祂的圣子与祂分离，也不像犹太人那样否认祂在万世之前由圣父所生、后来由童贞玛利亚所生；更不像\Person{撒伯里乌}那样混淆圣父与圣言，声称圣父与圣子是同一个位格；也不像\Person{佛提努斯}那样认为圣子首先在童贞玛利亚的胎中存在；更不像\Person{亚略}那样相信多种不同的权能，从而如同愚昧的外教人一样，造出多个神。因为经上记载着：\BibleQuote{以色列，你听着：主你的天主是唯一的天主。}\par

7. 因为天主和主是尊威的名号，权能的名号，正如天主亲自所说的：\BibleQuote{主是我的名，} \BibleRef{出 3:15} 又如在另一处先知所宣告的：\BibleQuote{全能的主是祂的名。} 因此，祂是天主，也是主，或者因为祂统御万有，或者因为祂鉴察万物，为万有所敬畏，没有区别。\par

8. 所以，如果天主是唯一的，那么天主圣三的名号是一，权能也是一。基督自己确实说：\BibleQuote{你们去，使万民成为门徒，因圣父及圣子及圣神之名给他们授洗。} \BibleRef{玛 28:19} 注意，是在\Emph{名}里，而不是在\Emph{众名}里。\par

9. 而且，基督自己说：\BibleQuote{我与圣父原是一体。} \BibleRef{若 10:30} 祂说\Quote{一}，是为了让人知道权能和性体没有分离；但又说\Quote{\Emph{我们是}}，好让你认出圣父与圣子，因为完全的圣父生了完全的圣子，\BibleRef{玛 5:48} 并且圣父与圣子是一，并非由于位格的混淆，而是由于性体的合一。\par

10. 因此，我们说只有一个天主，而非两个或三个天主，这正是亚略派不敬的异端以其亵渎所陷入的错误。因为它说有三个神，这就分割了天主圣三的天主性；而主在说\BibleQuote{你们去，使万民成为门徒，因圣父及圣子及圣神之名给他们授洗}时，已经显示了天主圣三具有同一权能。我们宣认圣父、圣子、圣神，在完全的天主圣三内理解完全的天主性以及权能的合一。\par

11. 主说：\BibleQuote{凡一国自相纷争，必成废墟。} 而天主圣三的国并不分裂。因此，如果不分裂，它就是一；因为不为一的才会分裂。然而亚略派却要使天主圣三的国成为那种因自相纷争而容易倾覆的国。但事实上，既然它不能倾覆，就显然没有分裂。因为没有哪个合一是分裂或破碎的，所以时间和腐朽都不能胜过它。\par

\Emph{皇帝被劝勉要在信仰上显出热忱。基督的完全天主性由祂与圣父在意志和行动上的合一而证明。天主性的属性被证明属于基督，祂的各种称号证明了祂的本质合一，但有位格区别。除此之外，无法维护天主的唯一性。}\par

12. 圣经说：\BibleQuote{不是凡向我说\InnerQuote{主啊！主啊！}的人，就能进天国。} \BibleRef{玛 7:21} 因此，尊贵的君王，信仰不能仅仅是口头上的，因为经上记载：\BibleQuote{我对你殿宇的热忱把我耗尽。} 所以，让我们以忠信的精神和虔诚的心意呼求我们的主耶稣，让我们相信祂是天主，好使我们无论向圣父求什么，都能因祂的名而获得。因为圣父的旨意是要人通过圣子向祂祈求，圣子的旨意是要人向圣父祈求。\par

13. 祂服从的恩宠促成了（与我们教训的）一致，而祂权能的行动也与此相符。因为凡圣父所做的，圣子也照样做。圣子不仅做同样的事，而且以同样方式去做，但圣父的旨意是要人在祂自己所要做的事上向祂祈求，好让你明白，并非祂不能以别的方式做，而是所显示的是同一权能。因此，天主子确实应受钦崇和敬拜，祂藉着自己天主性的权能奠定了世界的根基，又藉着服从塑造了我们的情感。\par

14. 所以我们应该相信天主是善的、永恒的、完全的、全能的、真实的，正如我们在法律和先知书以及圣经其他部分中所看到的；否则就没有天主。因为既是天主，就不可能不善，因为善的完满属于天主的本性；创造时间的天主也不可能在时间之内；天主也不可能不完美，因为较小的存在显然是残缺的，它缺少某种能使其等同于较大者的东西。这就是我们信仰的教导：天主不是恶的，天主没有做不到的事，天主不存在于时间之内，天主不次于任何存在。如果我错了，让我的对手来证明吧。\par

15. 那么，既然基督是天主，祂自然就是善的、全能的、永恒的、完全的、真实的；因为这些属性属于天主性本质。因此，让我们的对手否认基督内的天主性吧——否则他们就不能拒绝将属于天主性的归于天主。\par

16. 再者，为避免任何人陷入错误，人应注意圣经所赐予我们的那些标记，好使我们能认出圣子。祂被称为圣言、圣子、天主的权能、天主的智慧。称为圣言，因为祂毫无瑕疵；称为权能，因为祂是完美的；称为圣子，因为祂由圣父所生；称为智慧，因为祂与圣父是一，在永恒上是一，在天主性上是一。但这并不是说圣父与圣子是同一位格；圣父与圣子之间有由生育而产生的明确区别；所以基督是由天主而来的天主，由永恒者而来的永恒者，由圆满而来的圆满。\par

17. 它们并非仅仅是名称，而是能力藉着工程彰显的表记，因为圣父内既充满\OriginalTerm{天主性}{Godhead}，圣子内也充满天主性，不是不同的，而是同一的。天主性并非混杂，因为它是一体\OriginalTerm{一体}{unity}：并非繁多，因为在它内没有区别。\par

18. 再者，如果所有信者们，如经上所载，都是一个灵魂一个心：\BibleRef{宗 4:32}；如果凡依附于主的人，都是一个神\BibleRef{格前 6:17}，一如宗徒所说的；如果一个男人和他的妻子是一个肉身；如果所有我们这些凡人在普遍本性的层面上，都是同一性体\OriginalTerm{性体}{substance}；如果这是圣经论受造的人所说的，他们虽多，却是一体，他们无论如何无法与天主位格\OriginalTerm{天主位格}{Divine Persons}相比，那么父与子在\OriginalTerm{天主性}{Divinity}内更是一体，在祂们内，无论是在性体上或是在意志上，都没有区别！\par

19. 否则，我们又怎能说天主是唯一的？\OriginalTerm{天主性}{Divinity}造成复数，但能力的一体性排斥数目，因为一体不是数目，而是所有数目的本源。\par

\Emph{从圣经搜集的证据证实了父与子的合一，首先，取自\BibleBook{依撒意亚}{书}的一段经文与其他经文相互比较、加以解释，为要说明在子内，除了肉身的层面外，与父的本性并没有差别；由此可知两位的\OriginalTerm{天主性}{Godhead}是同一的。这一结论由\Person{巴路克}的权威加以确证。}\par

20. 众先知的神谕证实了，圣经所宣称的父与子在天主性上的密切合一。因为万军的上主这样说：\BibleQuote{\Place{埃及}劳碌，\Place{雇士}与\Place{色巴}的贸易（所获）都归你。身材高大的士人必前来归属你，成为你的奴隶；他们必带着锁链在你后面行走，向你跪拜，向你哀求说：惟独你内有天主，除你以外没有别的神。因为你是天主，我们却不认识，以色列的天主啊！}\par

21. 请听先知的声音：他说，\Quote{在你内}，\Quote{有天主，除你以外没有别的神。} 这与亚略派\OriginalTerm{亚略派}{Arians}的教导如何相合？他们必须否认圣父或圣子的天主性，除非他们彻底相信同一\OriginalTerm{天主性}{Divinity}的合一。\par

22. 他说，\Quote{在你内}，\Quote{有天主}——因为父在子内。经上记载：\BibleQuote{那住在我内的父，自己说话，}又说\BibleQuote{我所做的工程，祂自己也做的。}\BibleRef{若 14:10} 我们又读到子在父内，说：\BibleQuote{我在父内，父在我内。}\BibleRef{若 14:10} 如果可能，让亚略派除去这本性上的亲缘与工程中的一体！\par

23. 因此，存在天主在天主内，但不是两位天主；因为经上记载只有一位天主，且存在主在主内，但不是两位主，因为同样记载：\Quote{事奉两个主}\BibleRef{玛 6:24}。法律书上说：\Quote{以色列，听！上主你的天主是唯一的天主；}\BibleRef{申 6:4} 而且，在同一约书中记载：\Quote{上主由上主降下火来。}\BibleRef{创 19:24} 所说的主，是从\Quote{由上主}降下火来。同样，你在\BibleBook{创世}{纪}也可读到：\Quote{天主说——天主造了，}\BibleRef{创 1:6} 再往下：\Quote{天主按天主的肖像造了人；}\BibleRef{创 1:26} 然而并不是两位天主，而是一位天主，造了人。因此，在这段经文中，如同在那段经文中，保持了作为与名称的一体性。因为确实，当我们读到\Quote{天主出自天主}时，我们并不是在说两位天主。\par

24. 此外，在\BibleBook{圣咏}{第四十四篇}中，你读到先知不仅称父为\Quote{天主}，也宣认子为天主，说：\BibleQuote{天主，你的御座永远常在。} 往下又说：\BibleQuote{天主，你的天主，用欢愉的油敷了你，胜过你的同伴。} 这位傅油的天主，以及那位在肉身中被傅油的天主，就是天主圣子。因为基督在受傅油上，除了在肉身中的人以外，还能有什么同伴呢？于是你看到，天主由天主傅油，祂在受傅时取了人性，因而被宣认为天主圣子；然而法律的原理并未被破坏。\par

25. 同样，当你读到\Quote{上主由上主降下火来}时，要承认天主性的一体，因为在工程中的一体不允许有超过一位个别天主的存在，正如主自己曾表明，说：\Quote{相信我：我在父内，父在我内；否则，至少为那些工程而相信我。} 在此，我们也看到，天主性的一体通过工程中的一体而得以表明。\par

26. 宗徒小心求证父与子共有一个\OriginalTerm{天主性}{Godhead}，共有一个主位\OriginalTerm{主位}{Lordship}，以免我们陷入错误，无论是外教人的还是犹太人的不信，他指示我们当遵循的准则，说：\Quote{一个天主父，万物出于祂，我们归于祂；一个主耶稣基督，万物藉着祂，我们藉着祂。} 因为，正如他称\Person{耶稣基督}为\Quote{主}时，并未否认父是主，同样，他说\Quote{一个天主父}时，也并未否认子的真天主性，如此教导，不是有超过一位天主，而是能力的源泉是单一的，因为天主性在于主位，主位也在于天主性，如经上所载：\Quote{确信上主，祂就是天主；是祂造成了我们，不是我们自己。}\par

27. 因此，\Quote{在你内}，\Quote{有天主}，是由于本性的合一；而\Quote{除你以外没有别的神}，是由于位格完整地拥有\OriginalTerm{性体}{Substance}，毫无保留，毫无差异。\par

28. 此外，圣经在\BibleBook{耶肋米亚}{书}中论及唯一的天主，且同时承认父与子。因此我们读到：\BibleQuote{祂是我们的天主，没有谁可与祂相比。祂发现了智慧的全部途径，将之赐给祂的仆人\Person{雅各伯}，和祂所钟爱的\Place{以色列}。此后，祂又出现在地上，与人往来。}\par

29. 先知论及圣子，因为正是圣子亲自与人交谈，他说：\Quote{祂是我们的天主，与祂相比，无任何其他可视为同等。} 为什么我们对这样一位伟大先知所言、无可比拟的祂还要质疑呢？既然天主性是唯一的，又怎能将另一个拿来相比较呢？这是一个身处危险中的人民的宣认；他们敬畏宗教，因此不善于争论。\par

30. 来吧，圣神，请帮助祢的众先知，祢常居于他们内，我们也信从他们。如果我们不信先知，难道要信这世上的智者吗？可是智者在哪里？经师在哪里？当我们的农夫栽种无花果时，他发现了哲学家一无所知的事，因为天主拣选了世上愚拙的，为羞辱那强壮的。难道我们要相信犹太人吗？因为天主曾在\Place{犹太}为人所知。不，他们否认那正是我们信仰的根基，因为他们既拒绝了圣子，也就不认识圣父。\par

\Emph{天主的唯一性必然蕴含于自然秩序、信德和圣洗之中。贤士的礼物表明：（1）天主性的唯一；（2）基督的天主性与人性。天主圣三唯一性的真理，在那位与\Person{沙得辣客}、\Person{默沙客}和\Person{阿贝得乃哥}同行于火窑中的天使身上得以彰显。}\par

31. 整个自然界都见证天主的唯一性，因为宇宙是一体的。信德宣告天主只有一位，因为旧约与新约的信仰是同一的。圣洁无上的恩宠见证圣神只有一位，因为只有一种圣洗，以天主圣三之名施行。众先知宣报，众宗徒聆听，那同一的天主的声音。贤士们也信奉唯一的天主，他们怀着崇敬，将黄金、乳香和没药带到基督的摇篮前，藉黄金的馈赠承认祂的王权，以乳香敬拜祂为天主。因为黄金象征王国，乳香象征天主，没药象征殡葬。\par

32. 那么，在卑微的畜棚中献上的这些奥秘礼物意味着什么，岂不是要我们在基督身上分辨天主性与肉身的区别吗？祂外表如同人，\BibleRef{斐 2:7} 祂却被朝拜为主。祂躺在襁褓中，却在星辰间闪耀；摇篮显示祂的诞生，星辰显示祂的王权；是肉身被包裹在布中，是天主性接受天使的侍奉。因此，祂本有尊威的尊严并未丧失，祂真实地取了肉身也得到证明。\par

33. 这就是我们的信德。天主愿意这样让万民认识祂，三个青年就这样相信了，\BibleRef{达 4:17} 并在被投入的火窑中感觉不到火焰，那火毁灭焚烧了不信者，\BibleRef{达 4:22} 却如甘露般无害地落在忠信者身上，\BibleRef{欧 14:5} 因为对他们而言，别人点起的火焰变得冰冷，因那痛苦在与信德的冲突中公义地失掉了力量。因为在他们当中有一位以天使的形象出现，\BibleRef{达 4:28} 安慰他们，\BibleRef{路 22:43} 为的是以天主圣三的数目赞颂一个至高权能。天主受到赞颂，天主子在天主的天使中被看见，圣洁的属神恩宠在青年们身上发言。\par

\Emph{亚略派针对基督发出的各种亵渎之言被引述。在回应这些之前，正信者被劝诫要提防哲学家的诡辩式论点，因为异端者尤其倚赖这些。}\par

34. 现在，让我们来思考亚略派关于天主子的辩论。\par

35. 他们说天主子与圣父不像。对人说出这话也是一种侮辱。\par

36. 他们说天主子在时间中有一个开端，然而祂自身却是时间及其中万有的本源与规定者。我们作为人，尚且不愿受时间限制。我们曾经开始存在，却相信我们将拥有不随时间流逝的存在。我们渴望不朽——那么，我们怎能否认天主子的永恒性呢？天主宣告祂本性就是永恒的，而非出于恩宠。\par

37. 他们声称祂是受造的。然而有谁会将自己所造的作品与作者等同，让作者看起来像是他自己所造的物呢？\par

38. 他们否认祂的良善。他们的亵渎本身就是对自己的定罪，因此不能指望得到宽恕。\par

39. 他们否认祂真是天主子，否认祂的全能，因为他们虽然承认万物是藉着圣子的服役而被造，却将万物存在的初始本源归于天主的权能。然而，权能若不是本性的圆满，又是什么呢？\par

40. 此外，亚略派否认祂在天主性上与父是一体。那么，让他们废除福音吧，让基督的声音闭口吧。因为基督亲口说过：\Quote{我与父原是一体。} \BibleRef{若 10:30} 这话不是我说，是基督说的。难道祂是骗子，会说谎吗？ \BibleRef{户 23:19} 难道祂不义，竟宣称自己从未是过的？但关于这些问题，我们将在适当的地方分别更详细地论述。\par

41. 那么，鉴于异端者说基督与圣父不像，并力图以狡黠的辩论来支持此说，我们必须引用圣经的话：\BibleQuote{你们要小心，免得有人以哲学和虚伪的妄言，按照人的传授，依据世俗的原理，而不是依据基督，把你们勾引了去；因为在基督内，真实地住有整个圆满的天主性，你们也是在祂内得到成全。}\par

42. 因为他们将自己毒害的一切力量都积蓄在辩证的辩论中，按哲学家们的判断，那种辩论无非是毫无建树之力，只以破坏为目的。然而，天主却不用辩证法来拯救祂的子民；\Quote{因为天主的国在于信德的单纯，而不在于言辞的争辩。}\par

\Emph{为了引向基督与圣父无异的证明，\Person{圣安博}引述了亚略派中较著名的领袖，并说明他们的证词何等自相矛盾，同时指出圣经提供了何种防御来驳斥他们。}\par

那么，亚略派说基督与圣父不相像；我们否认这一点。不，实际上，我们对此话心怀畏惧。然而，我不愿陛下仅凭辩论与我们的争辩就信服。让我们查考圣经、宗徒、先知，查考基督。总之，让我们查考圣父，这些人声称他们维护圣父的尊荣，如果圣子被认为次于圣父的话。但对圣子的侮辱并不能给慈善的圣父带来尊荣。如果圣子被认为次于圣父，而非与圣父同等，绝不可能取悦慈善的圣父。\par

我祈求陛下容我片刻，让我特别针对这些人说话。但我该挑选谁出来引证呢？\Person{欧诺弥}？还是\Person{亚略}与他的老师\Person{埃提乌}？因为名字虽多，不信却是一个，邪恶顽固，但在言论上自分裂；在欺骗上并无差异，却在共同的勾当中滋生纷争。他们为何不能彼此相合，我不明白。\par

亚略派排斥\Person{欧诺弥}本人，却持守他的不信，行在他的邪恶之道上。他们说他过于慷慨地出版了\Person{亚略}的著作。诚然，是错谬的慷慨挥霍！他们称赞发命的人，却否认执行的人！因此，他们现已分裂为多个派别。有的追随\Person{欧诺弥}或\Person{埃提乌}，有的追随\Person{帕拉迪}或\Person{德莫斐罗}与\Person{奥克森提}，或是这一不信形态的继承者。另有一些人则跟随不同的师傅。那么，基督被分裂了吗？\BibleRef{格前 1:13} 不是的；而是那些使祂与圣父分离的人，亲手将自己割裂了。\par

所以，既然那些彼此不合的人全都同谋反对天主的教会，我当以异端者这个通名称呼那些我必须答辩的人。因为异端，如同神话中的九头蛇，因其伤创而变得壮大，屡遭削砍，却又新生，注定要在火焰中遭遇适当的毁灭。或者，犹如某个可怕而怪异的海妖斯库拉，分裂为许多不信的形态，她以假装成为一个基督教派为掩饰，展示其诡诈；但对于那些在她不洁海峡的波涛中飘摇不定、身处信德残骸中的可怜人，她身围凶兽，以她亵渎教义的残忍毒牙撕咬他们。\par

陛下，这怪物的洞穴，如航海之人所言，遍布隐藏的兽穴，其周围一切地方，不信的礁石回荡着她黑色犬只的嚎叫，我们必须以近乎掩耳的方式经过。因为经上记载着：\BibleQuote{以荆棘围护你的耳朵；}\BibleRef{德 28:28} 又说：\BibleQuote{要提防狗，提防作恶的人；}\BibleRef{斐 3:2} 又说：\BibleQuote{对于异端者，在受过一次惩戒后，就避开他，因为知道这样的人已败坏了，犯了罪，已由自己的判决定了罪。}\BibleRef{铎 3:10} 所以，如同谨慎的舵手，让我们为信德张起风帆，采取最安全的航程，再次沿着圣经的海岸前行。\par

\Emph{基督与圣父的相像，根据圣保禄宗徒、先知及福音的权威得到确认，尤其依赖于人是按天主的肖像受造这一事实。}\par

宗徒说，基督是圣父的肖像——因为他称祂为不可见天主的肖像，一切受造物的首生者。请注意，是首生者，而非首造者，好使我们可以相信祂既因本性而为受生者，又因永恒而为首先者。在另一处，宗徒也声明，天主立了圣子为\BibleQuote{万有的承继者，并藉着祂造成了诸世界。祂是天主光荣的辉耀，是天主本体的真像。}\BibleRef{希 1:2} 宗徒称基督为圣父的肖像，而\Person{亚略}却说祂与圣父不相像。那么，如果祂毫无相似，为何称为肖像？人们不会让自己的肖像与自己不相似，而\Person{亚略}却争辩说圣父与圣子不相像，并且愿意主张圣父生了一位与自己不相似的，仿佛祂不能生养自己的同类一般。\par

先知们说：\BibleQuote{在祢的光中，我们必得见光；}又说：\BibleQuote{智慧是永恒光明的辉耀，是天主威严的无瑕明镜，是祂美善的肖像。}\BibleRef{智 7:26} 请看说的是何等伟大的名称！\BibleQuote{辉耀，}因为在圣子内，圣父的光荣清晰照耀：\BibleQuote{无瑕明镜，}因为在圣子内看见圣父：\BibleRef{若 12:45} \BibleQuote{美善的肖像，}因为这不是一个身体在另一个身体中反射所见，而是圣子内的整个[天主性]大能。\BibleQuote{肖像}一词教导我们毫无差别；\BibleQuote{真像}表明祂是圣父形体的对应；而\BibleQuote{辉耀}则宣告祂的永恒。事实上，这\BibleQuote{肖像}并非肉身的容貌，不是由色彩构成，也不是用蜡塑造，而是纯粹源自天主，发自圣父，由泉源汲出。\par

主藉着这肖像，将圣父显示给\Person{斐理伯}，说：\BibleQuote{斐理伯，谁看见了我，就是看见了父。你怎么说：把父显示给我们呢？你不信我在父内，父在我内吗？}\BibleRef{若 14:9} 是的，凡看见圣子的，就在肖像中看见了圣父。请留意说的是何种肖像。这肖像是真理，是正义，是天主的大能：不是哑巴，因为是圣言；不是无知觉的，因为是智慧；不是虚幻愚蠢，因为是大能；不是没有灵魂，因为是生命；不是死的，因为是复活。所以，你可以看到，当说到肖像时，其意义是指圣父，圣子是圣父的肖像，因为没有人能是自己的肖像。\par

我还可以从圣子的证言中写下更多；然而，为避免祂或许显得过于张扬自己，让我们查考圣父。因为圣父说过：\BibleQuote{让我们照我们的肖像，按我们的模样造人。}\BibleRef{创 1:26} 圣父对圣子说\BibleQuote{照\Emph{我们}的肖像和模样，}而你却说天主圣子与圣父不相像。\par

\Person{若望}说：\BibleQuote{可爱的诸位，我们是天主的子女，但我们将来如何，还没有显明；可是我们知道：一显明了，我们必要相似他。} \BibleRef{若一 3:2} 盲目的疯狂！无耻的顽固！我们是人，并且就我们所能而言，我们将具有天主的肖像：我们怎敢否认圣子相似天主呢？\par

因此，圣父说过：\BibleQuote{让我们照我们的肖像，按我们的模样造人。}在我所读到的，宇宙之初，圣父和圣子就已存在，而我看到一个创造。我听见说话的那一位，我承认做事的那一位：但我所读到的，是一个肖像，一个模样。这模样不属于差异，而属于统一。所以，你为自己所要求的，你便从圣子那里夺去，因为事实上，若不藉天主的肖像，你便不能在祂的肖像内。\par

\Emph{已证实圣子与圣父的相似，所以证明圣子的永恒并不困难，尽管事实上，这一点可以依据\Person{依撒意亚}先知和\Person{圣史若望}的权威来确立，藉着这权威，异端的首领们显明被驳斥。}\par

因此，很明显，圣子并非不相似圣父，所以我们更容易宣认祂也是永恒的，因为相似永恒者必然是永恒的。但如果我们说圣父是永恒的，却否认圣子的永恒，那么我们就是说圣子不相似圣父，因为暂时的不同于永恒的。先知宣称祂是永恒的，宗徒也宣称祂是永恒的；旧约和新约都充满了对圣子永恒的见证。\par

那么，让我们按照顺序引用它们。在旧约中——从众多见证中举一例——经上记载：\BibleQuote{在我以前，没有一个天主；在我以后，也必没有一个。}\BibleRef{依 43:10} 我不打算评论这处经文，而是直接问你：\Quote{这话是谁说的——是圣父还是圣子？}无论你说哪一位，你都会发现自己被说服，或者，如果你是个信徒，就会受到教导。那么，这话是谁说的，是圣父还是圣子？如果是圣子，祂说：\BibleQuote{在我以前，没有一个天主；}如果是圣父，祂说：\BibleQuote{在我以后，也必没有一个。}一位在祂以前没有任何存在，另一位在祂以后也没有任何存在；正如圣父在圣子内被认识，圣子也在圣父内被认识，因为每当你提到圣父，你也就暗示了祂的圣子，因为没有一个人是自己的父亲；而当你提到圣子时，你也承认了祂的圣父，因为没有一个人能是自己的儿子。所以，圣子不能没有圣父而存在，圣父也不能没有圣子。因此，圣父是永恒的，圣子也是永恒的。\par

\BibleQuote{在起初已有圣言，圣言与天主同在，圣言就是天主。圣言在起初就与天主同在。}\Quote{在}，注意，\Quote{与天主同在}。\Quote{在}——看，我们有了四次的\Quote{在}。亵渎者从何处找到经上记载祂\Quote{不在}呢？再者，\Person{若望}在另一处——在他的书信中——论到\BibleQuote{那在起初就有的}。\BibleRef{若一 1:1} \Quote{在}的延伸是无限的。你想象多长的时间都可以，但圣子仍然是\Quote{在}。\par

如今，在这段简短的经文中，我们的渔夫堵住了一切异端的道路。因为那\Quote{在起初}有的，不受时间限制，也不被任何起始所先存。因此，让\Person{亚略}住口吧。再者，那\Quote{与天主同在}的，并非与祂混淆或混合，而是以其作为与天主同在的圣言所具有的无玷完美而被区分；所以让\Person{撒伯里乌}保持沉默吧。而\Quote{圣言就是天主}。因此，这圣言不在于发出的言语，而在于属天卓越的指称，这样\Person{佛提努斯}的教训就被驳斥了。此外，藉着祂在起初就与天主同在的事实，证明了圣父与圣子内永恒天主性不可分割的统一，使\Person{优诺米}蒙羞受辱。最后，既然万物都被说是藉着祂而造成的，祂就被清楚地显示为旧约和新约的共同创立者；这样\OriginalTerm{摩尼教徒}{Manichæan}就找不到攻击的根据了。如此，这位好渔夫用一张网捕捉了他们所有人，使他们无力欺骗，尽管是不值得捕获的鱼。\par

\Emph{\Person{圣盎博罗削}质问异端者并展示他们的回答，即圣子确实在一切时间之前就已存在，却并非与圣父同等永恒；于此，圣人指出他们将天主性描绘为可变的，并进一步指出，每一位格都必须被信奉为永恒的。}\par

告诉我，你这个异端者——因为皇帝无比宽仁，准许我暂时向你讲话，并非我想与你辩论，也不是渴望听你争论，而是我愿意展示它们——告诉我，我说，是否曾有一段时间，全能天主不是圣父，然而仍是天主？你的回答是：\Quote{我什么也没说关于时间。}巧妙而狡猾地反驳！因为如果你将时间引入争论，你就会定罪自己，因为你必须承认曾有一段时间圣子不存在，然而圣子是时间的统御者和创造者。祂不可能在祂自己的作品之后才开始存在。因此，你必须容许祂是祂作品的统御者和创造者。\par

你回答说：\Quote{我不说}圣子在时间之前不存在；但当我称祂为\Quote{圣子}时，我便宣告祂的圣父在祂以前就存在，因为，如你所说，父亲存在于儿子之前。但这意味着什么？你否认时间在圣子之前，但你却又坚持有某种东西先于圣子的存在——某种时间内的受造物——而且你展示出某些生发的中间阶段，由此你让我们理解为从圣父的生发是一个时间中的过程。因为如果祂开始成为圣父，那么最初祂只是天主，然后祂才成为圣父。那么，天主如何是不变的呢？因为如果祂先是天主，然后才是圣父，那么祂确实因着后来附加的生发行为而经历了变化。\par

60. 但愿天主保护我们免于这种疯狂；因为我们提出这个问题，只是为了驳斥异端者的不敬。虔诚的精神确认一种不在时间内的生，并因此宣告圣父与圣子同永恒，且不主张天主曾经历任何变化。\par

61. 因此，让圣父与圣子在崇拜中结合，正如他们在\OriginalTerm{天主性}{Godhead}中结合一样；不要让亵渎分离那紧密的生的联系所结合起来的。让我们尊敬圣子，好使我们也能尊敬圣父，正如福音中所记载的。\BibleRef{若 5:23} 圣子的永恒是圣父尊威的装饰。如果圣子并非从永远就存在，那么圣父就经历了变化；但圣子是从永远就存在的，因此圣父从未改变，因为他是永远不变的。这样我们就明白，那些否认圣子永恒的人是在教导圣父是可变的。\par

\Emph{基督的永恒已由宗徒的教导得到证明，\Person{圣盎博罗削}告诫我们，论及天主性的生，不应按人类生殖的方式来设想，也不可过于好奇地探究。他拒绝处理由此产生的难题，并说，凡从我们对身体的知识中撷取的、用于谈论这神圣的生的词语，都必须以属灵的意义来理解。}\par

62. 现在请听另一个论证，清楚显示圣子的永恒。宗徒说，天主的能力和\OriginalTerm{天主性}{Godhead}是永恒的，而基督是天主的能力——因为经上记载基督是\BibleQuote{天主的能力和天主的智慧}。那么，如果基督是天主的能力，既然天主的能力是永恒的，基督也必然是永恒的。\par

63. 因此，异端者，你不能从人类生殖的习惯建立虚假的道理，也不能从我们的论述中搜集这种工作所需的东西，因为我们无法用我们有限的语言去包罗无限天主性的伟大，\Quote{祂的伟大无可测度}。如果你试图解释人的出生，你必须指出一个时间。但神圣的生成超出万物之上；它广阔深远，远超一切思想和感觉。因为经上记载：\BibleQuote{除非藉着我，没有人能到父那里去。}\BibleRef{若 14:6} 因此，凡你对于圣父所想到的——是的，即使是他的永恒——除非藉着圣子的帮助，你什么也想不到，任何理智不藉着圣子也无法上升到圣父那里。\BibleQuote{这是我的爱子}，圣父说。\BibleQuote{是}——留意这个词——祂是自有永有者，永远是。因此\Person{达味}也感动地说：\BibleQuote{上主，你的圣言永远存留在天上。}——因为那存留的，在那存在和永恒上都毫不动摇。\par

64. 你问我，如果他没有一位先他而存在的父，他怎么是子？我反过来问你，你认为子在什么时候或以什么方式被生？对我来说，关于他生成之奥秘的知识是我所无法达到的——心智失败，口舌哑然——是的，不仅是我，连天使也是如此。它超越诸德能，超越天使，超越\OriginalTerm{革鲁宾}{Cherubim}、\OriginalTerm{色辣芬}{Seraphim}，以及一切有感觉和思想者，因为经上记载：\BibleQuote{基督的平安，超乎各种思虑之上。}如果基督的平安超乎各种思虑之上，那么如此奇妙的生成怎能不超乎一切理解之上？\par

65. 因此，你也（如同天使一样）用双手掩面，因为不容你窥探那超越的奥秘！我们被允许知道圣子是被生的，却不许争论他被生的方式。我不能否认前者；后者我害怕探究，因为如果\Person{保禄}说，他被提到第三层天时所听见的话是不许说出的，\BibleRef{格后 12:2} 我们如何能解释这个从圣父并出自圣父的生成之奥秘呢？这奥秘我们既不能听见，凭理智也无法达到。\par

66. 但如果你强迫我遵守人类生殖的规则，以便你可以说圣父先于圣子存在，那么请考虑，取自地上受造物的生殖实例，是否适合彰显神圣的生成。如果我们按人间习俗来说，你无法否认，在人中，父在存在方面的变化先于子。父先成长，先进入老年，先悲伤，先哭泣。那么，如果子在时间上晚于他，他在经验上就比子更老。如果孩子出生，父母逃脱不了生育的羞耻。\par

67. 为什么如此喜欢那种拷问呢？你听见\OriginalTerm{天主子}{Son of God}的名字；要么废除它，要么承认他的真实本性。你听见说起胎——承认那确凿无疑的生之真理。说起他的心——知道这是天主的圣言。说起他的右手——承认他的能力。说起他的面——承认他的智慧。这些词语，当我们论及天主时，不可像论及身体那样去理解。圣子的生是\OriginalTerm{不可理解的}{incomprehensible}，圣父\OriginalTerm{无痛地}{impassibly}生，而且是从自己、在无限遥远的世代中，\OriginalTerm{真天主}{very God}生了真天主。圣父爱圣子，\BibleRef{若 5:20} 你却急切地探究他的位格；圣父喜悦他，你却与犹太人同伙，以恶眼看他；圣父认识圣子，你却与外邦人一起辱骂他。\BibleRef{路 23:36}\par

\Emph{不能从圣经证明圣父先于圣子存在，取自人类生殖的论证也达不到这个目的，因为它们带来了无穷的荒谬。胆敢断言基督在时间的过程中开始存在，是亵渎的极点。}\par

68. 你问我，那身为圣父者是否可能在存在上不居先。我要你告诉我，圣父何时存在，而圣子尚不存在；请证明这一点，从论证或圣经的凭据中收集。如果你依靠论证，你无疑曾受教：天主的权能是永恒的。再者，你读过圣经的话：\BibleQuote{以色列，如果你听从我，在你内将不会有新神，也不可敬拜外邦神。}这诫命的第一条标志着圣子的永恒，第二条标志着祂具有同一的本性，因此我们既不能相信祂是在圣父之后才存在，也不能设想祂是另一位神明的儿子。因为，如果祂并非永远与圣父同在，祂就是一位\Quote{新}神；如果祂与圣父没有同一的天主性，祂就是一位\Quote{外邦}神。但祂并不在圣父之后，因为祂并不是\Quote{一位新神}；祂也不是\Quote{一位外邦神}，因为祂是由圣父所生，且如经上所载，祂是\BibleQuote{在万有之上，永远可赞颂的天主。}\BibleRef{罗 9:5}\par

69. 但是，如果亚略派相信祂是一位外邦神，他们为何又要敬拜祂？因为经上记载：\BibleQuote{不可敬拜外邦神。}否则，如果他们不敬拜圣子，就让他们承认这一点，事情便了结了——以免他们以宗教的表白欺骗任何人。那么，我们看得出，这就是圣经的证明。如果你们还有别的证据可提出，那是你们自己的事了。\par

70. 现在让我们再进一步，从论证中总结真理。因为尽管论证通常甚至应向人的证据让步，然而，异端者，你仍可任意争辩。你说：\Quote{经验教导我们，生者先于被生者。}我回答：从各方面遵循我们惯常的经验，如果其余全都与此一致，我不反对你的主张得到认可；但若没有这样的一致，你怎么能在这一点上要求同意，而在所有别的点上都缺乏支持呢？可见，既然你要求惯常的，那就引出了：圣子由圣父所生时，曾是一个小孩子。你们见过祂在摇篮中的婴儿，正在哭泣。随着岁月流逝，祂从力量走向力量——因为如果祂因被生之物的软弱而软弱，祂也必陷入软弱之中，不但在诞生上，也在生命上。\par

71. 但是，或许你竟疯狂到这种地步，以致毫不退缩地将这些事归于天主子，正如你所做的，以人性的软弱为尺度来衡量祂。那么，虽然你不能拒绝给祂天主的名称，你竟意图因软弱而证明祂是一个人，这又怎样呢？当你查究圣子的位格时，如果牵连到圣父，又当怎样？当你对前者草率判断时，你正将后者包含在同一谴责之中！\par

72. 如果\OriginalTerm{天主子的诞生}{Divine Generation}受限于时间——如果我们从人的生育习惯中借用这点来假设，那么结果就会进一步：圣父在一个肉身子宫中孕育了圣子，并在十个月的历程中承受着重负。然而，常见的生育如何才能在没有另一性别的帮助下发生呢？你看到，通常的生育程序并不是开端，你却认为生育的过程总是受制于身体必然性的规律。你要求惯常的程序，我却要求性别的差异：你要求时间的限定，我却要求秩序的限定：你探究终结，我探究开端。确实，终结依赖开端，而非开端依赖终结。\par

73. 你说：\Quote{一切被生的都有一个开端，因此由于圣子是圣子，祂有一个开端，并首先在时间的限度内存在。}这算是他们亲口说的话；至于我，我承认圣子是受生的，但他们声明的其余部分令我战栗。人啊，你承认天主，却以上述诽谤减损祂的光荣吗？愿天主拯救我们脱离这种疯狂。\par

\Emph{对圣子天主性的进一步异议，皆以同样的答案回应——即它们同样可以加于圣父。圣父既不受时间、空间或任何受造之物的限制，这种限制也不应加于圣子，其玄妙的诞生不但出自圣父，也出自童贞玛利亚；因此，既然祂由圣父的诞生不涉及性别的差异或类似的事，祂由童贞的诞生亦然。}\par

74. 下一个异议如此：\Quote{若圣子不具有所有儿子皆有的属性，祂就不是圣子。}愿圣父、圣子、圣神宽恕我，因为我愿以全心的虔敬提出这个问题。确实，圣父存在，并永远常在；受造之物亦如天主所安排。那么，在这些受造物中，有哪一个不受空间、时间、受造这一事实，或某个起始因或创造者的限制呢？肯定没有。那么怎么样？它们中可有任何一个为圣父所需？这样的说法是亵渎。所以，不要将只属于受造物的事应用于天主性，或者，如果你坚持强行比较，就想想你的邪恶通向何方。天主不准我们竟然看到它的结局。\par

75. 我们坚持虔敬所给的答复。天主是全能的，因此天主圣父不需要那些事物，因为在祂内没有变化，也没有我们这类援助的余地，我们软弱，需靠此类事物支持。然而那全能的，显然是那非受造的，不受任何空间限制，并超越时间。在天主之前，什么都没有——甚至，谈论有任何事物先于天主，便是严重的罪。因此，如果你承认在天主圣父的本性中，没有任何暗示需要维系的东西，因为祂是天主，那么就可以推论：在天主圣子内，也不能设想存在任何这类的东西，任何表示开端或增长的东西，因为祂是\Quote{出自真天主的真天主。}\par

76. 既然如此，我们发现并没有普遍依循的正常规律，那么，\OriginalTerm{亚略派}{Arian}啊，你就当满足于相信圣子的诞生是奥秘的。我说，你当满足；你若不信我，至少也该尊重天主的声音，祂说：\BibleQuote{你把我同谁相比，以为相似？}\BibleRef{依 46:5} 又说：\BibleQuote{天主不像人能反悔。}\BibleRef{户 23:19} 的确，如果天主以奥秘的方式工作，祂并不用双手的劳作，也不在若干时日的进程中完成任何工作、塑造任何事物、或使之完善，\BibleQuote{因为他一发言，万有造成，他一出命，万物化生。} 那么，我们既然承认造物主以奥秘方式工作，又在祂的工程中辨识出这一点，为什么我们不该相信，祂也以奥秘的方式生了祂的圣子呢？祂既拥有无比尊威的恩宠，理当也享有奥秘诞生的光荣；理当认为祂以特别而奥秘的方式生了圣子。\par

77. 不仅基督由圣父所生是奇妙的，祂由童贞女所生也同样令人惊叹。你说前者的方式与我们人类的受孕相似，我却要指出——不，我要迫使你承认——后者也与我们的出生方式毫无相似之处。告诉我，祂是怎样由玛利亚出生的？祂在童贞女胎中的成孕，与什么法则相符？没有男人的种子，如何能有出生？童贞女如何能怀孕？在经历夫妻间的交合之前，她如何成为母亲？这里没有〔看得见的〕原因，却有圣子诞生。那么，这个出生究竟是怎样依照一种新法则实现的呢？\par

78. 那么，既然在童贞玛利亚身上，并没有依照人类出生的普遍秩序，你又怎能要求天主圣父以你所受生的方式来生呢？普遍的秩序确实是由性别差异决定的，因为这已植入我们血肉的本性之中；可是，在血肉不存在之处，你又怎能期待找到血肉的软弱呢？没有人会去质疑一个比自己更优越的：你被命令去相信，而没有质询的自由。因为经上记载说：\BibleQuote{亚巴郎相信了上主，上主就以此算为他的正义。}\BibleRef{创 15:6} 语言不仅无法说明圣子的诞生，甚至也无法说明天主的化工，因为经上又说：\BibleQuote{他的一切作为，无不真实。} 难道祂的化工是在信实中完成，祂的诞生就不是吗？是的，我们竟去质疑那看不见的，我们这些受命要去相信而非考究所见之事的人。\par

\Emph{继续讨论圣子的诞生。圣安博用与\BibleBook{希伯来}{书}作者所用相同的例子来说明其方式。借拿步高王和圣伯多禄的例子，显示出相信所启示之事的责任。借着显示给圣伯多禄的神视，显明了圣子的永恒和天主性——因此，应当相信宗徒，胜于相信那些哲学教师，后者的权威已到处被人弃绝。另一方面，亚略派则显示出与外教人相似。}\par

79. 有人会问：\Quote{圣子是以怎样的方式受生的呢？} 祂是永存者，是圣言，是永恒之光的光辉，\BibleRef{希 1:3} 因为光辉在产生的那一刻即发挥效力。这是宗徒的例证，不是我的。因此，你不要以为曾有一刻天主没有智慧，就如不可设想曾有一刻光没有辉耀一样。亚略派啊，不要用人的尺度判断天主的事，而要在你找不到人的事物的地方，相信天主的事。\par

80. 异教国王在烈火中，同三个希伯来青年一起，看见有第四位的形像，仿佛天使，\BibleRef{达 3:25} 因为他想这位天使超越众天使，便断定祂是天主之子，他虽未研读过经典，却相信了祂。亚巴郎也看见了三位，却只朝拜了一位。\BibleRef{创 18:1}\par

81. 伯多禄在山上看見梅瑟和厄里亞同天主子在一起時，對於他們的本性和光榮並沒有受騙。因為他沒有問他們，而是問基督該做什麼；雖然他準備好要向三位致敬，卻仍等待其中一位的命令。但由於他愚昧地以為三個位格應搭三個帳棚，便被天主聖父至尊的聲音所糾正，說：\BibleQuote{這是我的愛子，你們要聽從祂！}\BibleRef{玛 17:5} 這等於說：\Quote{你為什麼把你的同伴與你的主同等看待呢？} \Quote{這是我的兒子，} 並不是\Quote{梅瑟是我的兒子}，也不是\Quote{厄里亞是我的兒子}，而是\Quote{這是我的兒子。} 這位宗徒並非遲鈍，他明白這項譴責；他俯伏在地，因聖父的聲音和聖子光榮的美麗而屈服，卻又由聖子扶起，因為聖子慣於扶起跌倒的人。\BibleRef{玛 17:6} 然後他只見一位，\BibleRef{玛 17:8} 唯獨天主子，因為僕役們已經退去，好使人看出祂是唯一的主，唯獨祂有資格稱為兒子。\par

82. 那麼，那個神視的目的何在？它並非顯示基督與祂的僕役同等，而是預示了一個奧秘，好叫我們明白，法律和先知與福音一致的，揭示天主子是永恆的，就是他們所預報的那一位。因此，當我們聽說聖子出自母胎，聖言發自內心時，讓我們相信聖子不是由手所塑造，而是由聖父所生，不是工匠的作品，而是父的親子。\par

83. 因此，那位曾说\BibleQuote{这是我的爱子}的，并没有说\Quote{这是受造于时间之物}，也没有说\Quote{此存在是我所创造、所造作、我的仆人}，而是说\BibleQuote{这是我的爱子，你所见的荣耀者}。这就是\Person{亚巴郎}的天主，\Person{依撒格}的天主，\Person{雅各伯}的天主，祂曾在荆棘丛中显现给\Person{梅瑟}，\Person{梅瑟}论到祂说：\BibleQuote{那自有者派遣了我}\BibleRef{出 3:14}。在荆棘丛中或在旷野里向\Person{梅瑟}说话的不是圣父，而是圣子。\Person{斯德望}论及这位\Person{梅瑟}时说：\BibleQuote{这\Person{梅瑟}就是曾在旷野中的教会里，与天使同在的那位}\BibleRef{宗 7:38}。那么，这就是那位颁赐法律者，与\Person{梅瑟}交谈，说：\BibleQuote{我是\Person{亚巴郎}的天主，\Person{依撒格}的天主，\Person{雅各伯}的天主}。那么，这就是圣祖们的天主，这就是先知们的天主。\par

84. 所以，我们所读到的有关圣子的，你的心灵既明白了所读的，你的口舌就该宣认。哪里需要信德，哪里就该摒弃争辩；如今，即使在辩证法的学派之中，也让辩证法学说缄默吧。我所问的不是哲学家们说什么，而是想知道他们做什么。他们孤寂地坐在自己的学派中。看哪，信德战胜了论证。那些诡辩之人日复一日地被同伴遗弃；那些单纯相信的人却日见增多。不是哲学家，而是渔夫；不是辩证法学的大师，而是税吏，现在得到了信任。一类人，通过逸乐与奢侈，将世界的重担加在自己身上；另一类人，通过斋戒和克己，将之卸去；如此，忧伤现在开始比逸乐赢得更多的追随者。\par

85. 现在让我们看看亚略派与异教徒到底有多大的差别。后者呼求不同性别且能力不平等的众神；前者则确认一个同样能力不平等且天主性多样化的天主圣三。异教徒声称他们的众神曾是某一时刻开始存在的；亚略派则谎称基督是在时间的过程中开始存在的。他们岂不都是把自己的不虔浸泡在哲学的染缸里吗？然而，异教徒确实夸耀自己所崇拜的，而亚略派却主张天主之子——本身就是天主——是一个受造物。\par

\Emph{以下论点证明天主之子并非受造存在：（一）祂未命令福音应传给祂自己；（二）受造存在被置于虚幻之下；（三）圣子创造了万有；（四）我们读到祂是受生的；（五）在所有同时宣示天主性与人性共存于祂内的经文中，始终有着受生与嗣养的区别。这一切见证都得到宗徒阐释的证实。}\par

86. 我相信，神圣陛下，如今已经清楚显明，主耶稣既不与圣父不同，也不是在时间过程中开始存在的。我们还需驳斥另一亵渎之说，并证明天主之子并非受造存在。我们读到、也听到主说的话，便是我们帮助所在的那赋予生命的话语：\BibleQuote{你们往普天下去，向一切受造物宣传福音。}\BibleRef{谷 16:15}那位说\Quote{一切受造物}的，没有例外。那么，那些称基督为\Quote{受造物}的人，如何站立得住呢？如果祂是受造物，祂能命令福音传给自己吗？所以，祂不是受造物，而是造物主，把训导受造之物的任务托付给了自己的门徒。\par

87. 所以，基督并非受造物；因为正如宗徒所说，\BibleQuote{受造物被置于虚妄之下}。\BibleRef{罗 8:20}基督难道被置于虚妄之下了吗？再者，按照同一位宗徒的话，\BibleQuote{受造物一同叹息，同受产痛，直到如今}。那么，基督是否参与了这叹息和产痛呢？——祂可是将我们这些悲惨哀恸者从死亡中释放出来的那位啊！宗徒说：\BibleQuote{受造物将脱离腐朽的奴役}。\BibleRef{罗 8:21}因此，我们看到，受造物与其主之间有着巨大的差别，因为受造物受役使，而\BibleQuote{主就是那神，主的神在哪里，哪里就有自由}。\BibleRef{格后 3:17}\par

88. 是谁最先陷入这个错误，宣称那创造和造成万有的祂是受造物呢？我要问，难道主创造了祂自己吗？我们读到：\BibleQuote{万物是藉着祂而造成的；凡受造的，没有一样不是由祂而造成的。}\BibleRef{若 1:3}既然如此，祂岂是创造了祂自己吗？我们读到——谁又能否认呢？——天主以智慧造成了万有。若如此，我们又怎能设想智慧是在它自己内被造成的呢？\par

89. 我们读到圣子是受生的，正如圣父说：\BibleQuote{我在晨星之前从胎中将你生出}。我们读到\Quote{首生}的圣子，\BibleRef{哥 1:15}和\Quote{独生子}\BibleRef{若 1:14}——首生，因为祂之前无人在先；独生，因为祂之后无人在后。再者，我们读到：\BibleQuote{谁能述说祂的受生呢？}\BibleRef{依 53:8}请注意，是\Quote{受生}，而非\Quote{受造}。面对如此伟大有力的见证，还能提出什么论据来反驳呢？\par

90. 此外，天主之子指出了受生与恩宠的区别，当祂说：\BibleQuote{我升到我的父和你们的父那里去，升到我的天主和你们的天主那里去。}祂没有说\Quote{我升到我们的父那里去}，而是说\BibleQuote{我升到我的父和你们的父那里去。}这一区别正是差异的标记，因为基督的父是我们的造主。\par

91. 祂又说\BibleQuote{升到我的天主和你们的天主那里去}，因为尽管祂与圣父原是一体，而且圣父因拥有相同的本性而是祂的父，而天主却是藉着圣子的职分开始成为我们的父，不是出于本性，而是出于恩宠——然而，祂在此似乎指出基督内存在两种性体：天主性与人性——天主性来自他的圣父，人性来自他的母亲，前者先于万有，后者源自童贞女。因为首先，他以圣子的身份发言，称天主为他的父，之后，以人的身份发言，称祂为天主。\par

92. 的确，圣经中各处都有证据表明，基督称天主为祂的天主，是作为人而说的。\BibleQuote{我的天主，我的天主，祢为什么舍弃了我？}又说：\BibleQuote{从母胎中祢就是我的天主。}在前一处，祂是以人的身分受苦；在后一处，则是从母胎中出生的人。因此，当祂说\BibleQuote{从母胎中祢就是我的天主}时，意思是那位从永远是祂圣父的天主，从祂由母胎出生那一刻起，就是祂的天主。\par

93. 既然我们在福音中、在宗徒书信中、在先知书中都读到基督是\OriginalTerm{受生}{begotten}的，亚略异端者怎敢称祂是\OriginalTerm{受造}{created}或造成的呢？他们实在应当想想，他们在何处读到祂是受造的，何处读到祂是被造成的。因为已经明显指出，天主子是由天主受生，由天主所生——那么，让他们仔细思考，他们在何处读到祂是被造成的，因为祂并非被造成天主，而是生而为天主，即天主子；后来，祂才按血肉由玛利亚造成人。\par

94. \BibleQuote{但时期一满，天主就派遣了祂的儿子，由女人\Emph{所生}，生于法律之下。}\BibleQuote{\Emph{祂的}儿子，}请注意，并非众多儿子中的一位，也不是与另一位共有的儿子，而是祂自己的儿子；当宗徒说\BibleQuote{祂的儿子}时，他表明这子的生发是永恒的，是出于子的本性。宗徒随后断言，祂后来是由女人\BibleQuote{造成}的，这是为了让人明白，这造成并非指天主性，而是指祂穿上了肉身——因此，\BibleQuote{由女人造成}是指取了血肉；\BibleQuote{生于法律之下}是指遵守法律。然而，前者，即属神的生发，是在法律之前；后者，即血肉之生，则是在法律之后。\par

\Emph{\BibleRef{宗 2:36}与\BibleRef{箴 8:22}的解释，指明它们专指基督的人性。}\par

95. 因此，异端者惯常引用的圣经经文——\BibleQuote{天主立祂为主，为基督}——是无的放矢。让这些无知的人读完整段经文，并理解它。因为经上这样写道：\BibleQuote{天主已立你们所钉死的这位耶稣为主，为基督。}被钉死的，并非天主性，而是肉身。这确实是可能的，因为肉身容许被钉死。因此，不能由此推断天主子是受造之物。\par

96. 那么，让我们也来处理那段被他们用来歪曲的经文——让他们明白\BibleQuote{上主造了我}这句话的含义。并非\Quote{圣父创造}，而是\BibleQuote{上主造了我}。肉身承认其主，颂扬声言圣父：我们受造的本性承认前者，爱慕并认识后者。那么，谁能不意识到这些话是在宣报\OriginalTerm{降生成人}{Incarnation}呢？因此，圣子是在论及自己作为人的那一方面时，自称是受造的，当祂说：\BibleQuote{你们为什么图谋杀害我——一个\Emph{人}，一个向你们讲了真理的人？}祂说的是祂的\OriginalTerm{人性}{manhood}，在此人性中祂被钉死、死亡并被埋葬。\par

97. 此外，毫无疑问，作者将未来之事写成了过去；这是预言的惯常用法，即将要来之事说成仿佛已经发生或过去了一样。例如，在\BibleBook{圣咏}{第二十一篇}中你读过：\BibleQuote{巴商的肥牛围绕着我，}又说：\BibleQuote{他们瓜分了我的衣服。}福音作者指出，这是预言性地论及受难的时刻，因为对天主而言，未来之事是当下临在的，对预知万物的那位来说，它们如同已经过去并结束了；正如经上所载：\BibleQuote{是祂造成了将要成就的事。}\par

98. 祂宣称自己的位置在万世之前已被立定，这并不奇怪，因为圣经告诉我们，祂是在万世之先已被预定的。接下来这段经文揭示出，那些话如何是对\OriginalTerm{降生成人}{Incarnation}的真正预言：\BibleQuote{智慧建造了自己的房屋，竖立了七根柱子来支撑它，宰杀了祭牲，将酒调在壶中，摆好筵席，派遣了使者，以宣报的大声召集众人，说：\InnerQuote{谁若是单纯的，就转向我这里来。}}难道我们没有在福音中看到，这一切在\OriginalTerm{降生成人}{Incarnation}之后都应验了吗？基督揭示了圣晚餐的奥秘，派遣了宗徒，大声呼喊说：\BibleQuote{谁若渴，到我这里来喝吧。}\BibleRef{若 7:37} 因此，前事与后事相呼应，我们在预言中看到了\OriginalTerm{降生成人}{Incarnation}的整个事迹被简要地陈述出来。\par

99. 许多其他关于降生成人的预言段落也很容易找到，但我不想在经卷上过多耽搁，免得这篇论著显得过于冗长。\par

\Emph{若亚略异端者认为\Quote{受造}与\Quote{受生}二词意义相同，便是亵渎基督。然而，若他们承认二词意指不同，便不可将那位他们读过是受生者说成是受造之物。这一准则得到了圣\Person{保禄}见证的支持，他自称基督的仆人，并禁止敬拜受造之物。天主是纯粹而非复合的实体，在祂内没有任何受造的本性；再者，圣子不可被贬低到受造物的层次，因为圣父在祂内满心喜悦。}\par

100. 现在我将特别质问亚略异端者：他们是否认为\OriginalTerm{受生}{begotten}与\OriginalTerm{受造}{created}是相同的。若他们说二者相同，那么生发与创造便没有区别。由此可推，因为我们也都是受造的，那么我们与基督以及万物之间就没有区别了。然而，纵然他们如此癫狂，也不敢说出这样的话。\par

101. 再者——姑且依他们愚妄的说法，承认这本非真理——我问他们：倘若如他们所想，言语上本无分别，为什么他们不以其较为尊贵的名号呼求他们所崇拜的那一位呢？为什么他们不利用圣父的言词呢？为什么他们摈弃这尊荣的称号，却使用一个羞辱的名号呢？\par

102. 然而，如果——如我所想——\Quote{受造的}与\Quote{受生的}之间有所区别，那么当我们读到祂是受生的时，我们必不会以\Quote{受生}和\Quote{受造}为同一回事。因此，要么让他们承认祂是圣父所生、由童贞女所生，要么让他们说明，天主子如何能既受生又受造。一个单一的本性，尤其在天主性那里，是排斥（内在）冲突的。\par

103. 但无论如何，让我们暂且放下个人的判断：让我们请教保禄，他被天主的圣神充满，预见到这些争论，已对外邦人尤其是亚略派作出了判决，说他们因事奉受造物而不事奉造物主，已属受天主审判定罪的人。事实上，你可以读到：\BibleQuote{天主任凭他们顺从心中的情欲，以致他们彼此玷辱自己的身体；他们将天主的真理变为虚妄，去崇拜事奉受造之物，而不事奉那位造物主——祂就是永远受赞颂的天主。}\BibleRef{罗 1:24}\par

104. 因此，保禄禁止我崇拜受造物，并提醒我事奉基督是我的本分。由此可知，基督并非受造之物。宗徒自称\BibleQuote{保禄，耶稣基督的仆人}，\BibleRef{罗 1:1}这位承认其主的忠仆，也同样要我们不崇拜受造之物。那么，他如何能自己作基督的仆人，如果他以为基督是受造的人呢？因此，让这些异端者停止崇拜他们称为受造之物的那一位，或者停止称他们所伪装崇拜的那一位为受造物，以免借着崇拜者的外表，陷入更大的不虔之中。因为家贼比外敌更坏，而这些人假基督之名以凌辱基督，更增加了他们的过犯。\par

105. 我们确实还能找到比这位外邦人的导师、这蒙拣选的器皿——由迫害者之中被拣选的——更好的圣经讲解者吗？那曾迫害基督的人，如今宣认祂。无论如何，他所读的撒罗满的智慧，远超亚略所读的，并且他精通法律，因此，因为他读了，所以他不说基督是受造的，而是说祂是受生的。因为他曾读到：\BibleQuote{祂一发言，万有造成；祂一出命，万物化生。}那么，我问你们：基督是因一言而造成的吗？祂是因一命而受造的吗？\par

106. 而且，天主内\Emph{如何}能有什么受造的本性呢？事实上，天主具有\OriginalTerm{单纯的本性}{uncompounded nature}；不能有所增减，且只在祂的本性内具有那属神的本质；祂充满万有，却不与任何事物混淆；祂渗透万有，却不被任何事物渗透；在同一时刻，以祂的全部充满而亲临于天上、地上和海洋的最深处，视之无形，言之难宣，触之难测；唯凭信德以追随，唯凭虔敬以钦崇；因此，任何在属灵深义上、在显示光荣与尊崇上、在抬举权能上卓绝的称号，你都可知道它们理当归属天主。\par

107. 那么，既然圣父在圣子内得享喜悦；你们便当相信，圣子堪配圣父，祂由天主而出，正如祂亲自作证说：\BibleQuote{我由天主而来，且来到这里}；\BibleRef{若 8:42}又说：\BibleQuote{我出于天主}。\BibleRef{若 16:27}从天主那里出发而来的那一位，只能具有唯独天主才有的属性。\par

\Emph{基督是真天主，可从祂是天主的亲子、由天主受生并由天主而来，以及圣父与圣子共有一个意志和行动等事实证明。众宗徒和百夫长的见证——圣盎博罗削将其与亚略派的教导相对照——连同依撒意亚和圣若望的见证，均被引证。}\par

108. 因此，基督不仅是天主，且确实是\OriginalTerm{真天主}{very God}——出自真天主的真天主，乃至祂本身就是真理。\BibleRef{若 14:6}那么，如果我们询问祂的名号，那就是\BibleQuote{真理}；如果我们寻求祂本性的地位与尊荣，祂真是天主的亲子，确实是天主的\Emph{亲}生子；正如经上所载：\BibleQuote{祂没有怜惜自己的亲子，反倒为我们的缘故将祂交出}，\BibleRef{罗 8:32}所谓交出，即指在肉身方面而言。祂是天主的亲子，这宣示了祂的天主性；祂是真天主，表明祂是天主的亲子；祂的慈悯是祂服从的保证，祂的祭献则是我们救恩的保证。\par

109. 不过，为防止人曲解\BibleQuote{天主将祂交出}这句圣经，宗徒自己另在一处说：\BibleRef{迦 1:3}\BibleQuote{恩宠与平安，由天主我们的父及主耶稣基督赐与你们；祂曾为我们的罪过舍弃了自己}；又说：\BibleRef{弗 5:2}\BibleQuote{正如基督爱了我们，为我们舍弃了自己}。那么，既然祂既被圣父交出，又自己甘愿交出，可见圣父与圣子的行动和意志是同一的。\par

110. 那么，如果我们探究祂本性的卓越之处，我们发现这就在于受生。否认天主子是由（天主）所生，就是否认祂是天主的\Emph{本有}之子，而否认基督是天主的本有之子，就是将祂与其余人等同，不过和任何其他人一样是个儿子而已。然而，如果我们探究祂受生方式的独特性质，那就是：祂是从天主而来的。因为，在我们的经验中，出来意味着已存在之物，而所谓出来之物似乎是从隐藏的内在之处发出的；但我们在圣经的简短章节中，观察到神圣受生的特有属性，即圣子似乎不是从任何地方出来的，而是犹如出自天主的天主，出自圣父的圣子，也不是在时间过程中有起始，因为祂是由圣父所生而出来的，正如那位受生者自己所说：\BibleQuote{我从至高者的口中出来。}\BibleRef{德 24:3}\par

111. 但是，如果\OriginalTerm{亚略派}{Arians}不承认圣子的本性，如果他们不信圣经，至少让他们相信大能的工作。圣父对谁说：\BibleQuote{让我们照我们的肖像，按我们的模样造人？}\BibleRef{创 1:26}，难道不是对祂所知的真圣子说的吗？在谁身上，除了真圣子，祂能认出自己的肖像呢？养子与真正的儿子不同；而圣子也不会说：\BibleQuote{我与父原是一体}\BibleRef{若 10:30}，如果他自己不是真的，却拿自己与那真实者相比。因此，圣父说：\BibleQuote{让我们造。}说话的那位是真的；那么，创造的那位难道不是真的吗？对说话者所有的尊敬，难道应该不归于创造者吗？\par

112. 然而，若非圣父知晓祂是自己的真圣子，祂怎会将祂的旨意托付给祂，以圆满合作，并将祂的工作托付给祂，使之在现实中圆满完成呢？正如经上记载，圣子做圣父所做的工作，且圣子愿意复活谁就复活谁，因此祂在能力上平等，在意志上是自由的。因此，合一得以保持，因为天主的能力在于天主性为每一位格所固有，而自由不在于任何差异，而在于意志的合一。\par

113. 宗徒们在海上遭遇风浪颠簸，他们一看见水在主的脚下翻腾，看见祂无畏的足迹踏在水面上，而祂在汹涌的海浪中行走；那条被波浪冲击的船，在基督一进入就平静了；他们又看见风浪都听从祂——那时，虽然他们心中尚未完全相信，却相信祂是天主的真圣子，说道：\BibleQuote{你真是天主子。}\BibleRef{玛 14:33}\par

114. 同样，百夫长和与他在一起的人，当主受难时大地的根基震撼时也宣认了——而这一点，异端者啊，你竟否认。百夫长说：\BibleQuote{这人真是天主子！}\BibleRef{玛 27:54} \Quote{真是}是百夫长说的——而亚略派却说\Quote{\Emph{不是}真的}。百夫长，双手沾满鲜血，但心灵虔诚，宣告了基督受生的真实性与永恒性；而你，异端者啊，却否认其真实性，并使之成为时间中的事！但愿是你亲手犯下罪过，而不是你的灵魂！但你这手也不洁、心存杀机的，尽力想致基督于死地，因为你认为祂卑微且软弱；而且，更糟糕的是，尽管天主性不会受到任何伤害，你仍不遗余力地想要在基督身上杀死祂的荣耀，而非祂的身体。\par

115. 那么，我们不能怀疑祂是真天主，因为连执行死刑的人都相信、魔鬼也都承认祂是真天主。他们的见证我们现在虽不需要，但这见证却比你们的亵渎更大。我们引他们作见证，为使你们羞愧，同时也引天主的圣言，好使你们相信。\par

116. 主藉依撒意亚的口宣告说：\BibleQuote{在我的仆人嘴上要称呼一个新名，这新名将在全地上受赞颂，他们必祝福真天主，凡在地上起誓的，都必指真天主起誓。}\BibleRef{依 65:16} 我说这些话，正是依撒意亚看见天主的光荣时所说的，因此在福音书中明说，他看见了基督的光荣并且论及了祂。\BibleRef{若 12:41}\par

117. 但再听圣史若望在他的书信中写的：\BibleQuote{我们也知道天主子来了，赐给了我们理智，叫我们认识那真实者，我们确实是在那真实者内，就是在他的儿子，我们的主耶稣基督内。他即是真实的天主和永远的生命。}\BibleRef{若一 5:20} 若望称祂为真天主子及真天主。那么，祂若是真天主，就必定是非受造的，毫无谎言或欺诈的污点，在祂内毫无混乱，也不与圣父有任何不相似。\par

\Emph{亚略派的谬误被列举在\OriginalTerm{尼西亚信德定义}{Nicene Definition of the Faith}中，以防他们欺骗任何人。这些谬误被列举出来，并附以\OriginalTerm{绝罚}{anathema}，据说不仅在尼西亚宣布过，而且在\OriginalTerm{阿里米努姆}{Ariminum}也两次重新宣布。}\par

118. 因此，基督是\Quote{出自天主的天主，出自光明的光明，出自真天主的真天主；由圣父所生，而非受造；与圣父同一性体。}\par

119. 我们的教父们正是如此，遵循圣经的指引，宣布了信理，并进一步主张将不虔敬的学说列入他们的法令记录中，好使亚略的不信自行显露，而不是像用染料或脂粉掩盖起来。因为那些不敢公开阐述其思想的人，会给他们的思想涂上虚假的色彩。因此，如同监察官的卷册一样，亚略异端不是以名字被揭露，而是以其所遭的谴责被标记出来；为的是，那些好奇且急于了解它的人，在听到阐述之前，先知道它已被谴责，从而避免堕落，并最终使之相信。\par

120. \Quote{那些}，法令如此写道，\Quote{说天主子曾有一段时间不存在，且在祂受生之前祂并不存在，又说祂是从无中造成的，或是出于另一实体或\OriginalTerm{性体}{\GreekText{οὐσία}}，或是能改变的，或是祂内有什么转变的阴影的人——这样的人，大公及宗徒的教会宣布为可诅咒的。}\par

121. 至尊的陛下业已认同，凡发表此类教义的人理当受绝罚。并非出于人的决定，也非出于人的谋略，才有三百一十八位主教，如前文我已更为详细指出的，聚集召开大公会议，而是借着他们的数目，主耶稣以祂圣名和受难的标记，证实了祂就在他们中间，在自己的子民聚集之处。\BibleRef{玛 18:20} 三百之数，是祂十字架的标记；十八之数，是耶稣圣名的标记。\par

122. 这也是\Place{阿里米努姆}会议中第一信纲的教导，以及该会议之后的第二修正的教导。寄给\Person{君士坦丁}皇帝的信函为这信纲作证，而随后的会议则宣布了那修正。\par

\Emph{亚略被指控犯有上述错误中的第一项，并以圣若望的见证予以驳斥。异端鼻祖的可悲死亡被描述，其余亵渎的错误一一受到审查并推翻。}\par

123. 那么，\Person{亚略}说：\Quote{曾有一段时间，天主子不存在，}但圣经却说：\Quote{祂存在，}而非\Quote{祂不存在。}再者，圣\Person{若望}写道：\BibleQuote{在起初已有圣言，圣言与天主同在，圣言就是天主。这圣言在起初就与天主同在。}\BibleRef{若 1:1} 请看动词\Quote{存在}出现了多少次，而\Quote{不存在}却无处可寻。那么，我们该信谁呢？——是曾靠在\Person{基督}怀中的圣\Person{若望}，还是在自己肠子喷涌中打滚的\Person{亚略}？——如此打滚，叫我们明白\Person{亚略}在其教导中如何显出了\Person{犹达斯}的样子，遭受了相似的惩罚。\par

124. 因为\Person{亚略}的肠子也流出来了——为礼所禁，不便言明何处——他就这样在中间爆裂，头朝下跌倒，玷污了他用以否认\Person{基督}的那张污秽的嘴。他碎裂了，正如宗徒\Person{伯多禄}论及\Person{犹达斯}所说的，因为他\BibleQuote{用不义的代价买了一块田地，头朝下跌倒，在中间爆裂，所有肠子都流了出来。}\Person{亚略}的死并非偶然的方式，因为同样的恶行遭受了同样的惩罚，为使那些否认并出卖同一主的人，也承受同样的折磨。\par

125. 让我们继续探讨其他要点。\Person{亚略}说：\Quote{在祂受生之前，天主子并不存在，}但圣经却说，万物都靠着子的工而得以存立。那么，一位不存在的，又如何能赋予其他存在呢？再者，当这亵渎者使用\Quote{当……的时候}和\Quote{在……之前}这些词语时，他无疑是在使用表示时间的词语。那么，亚略派如何能否认在子存在以前时间就已存在，却又想主张在时间内受造的事物存在于子之前呢？因为正是这些词语——\Quote{当……的时候}，\Quote{在……之前}，\Quote{曾一度不存在}——宣告了时间的概念。\par

126. \Person{亚略}说天主子是从虚无中存在的。那么，祂如何是天主子呢？——祂是如何从圣父的胎中受生？——我们如何读到祂是心中丰盈所发出的圣言，除非我们应当相信，正如经上所载，祂是从圣父最内里、不可接近的圣所发出的？如今，被称为子的要么是因收养，要么是因本性，就如我们被称为子乃是因收养。\Person{基督}是天主子，是因祂真实且永存的本性。那么，那位从虚无中造了万有的，祂自己又怎能是从虚无中受造的呢？\par

127. 那不知道圣子从何而来的人，就没有圣子。因此，犹太人没有圣子，因为他们不知道祂从何而来。所以主对他们说：\BibleQuote{你们不知道我从哪里来；}\BibleRef{若 8:14} 又说：\BibleQuote{你们既不认识我是谁，也不认识我的父，}因为凡否认圣子是由圣父所出的，就不认识圣父——圣子正是由祂所出；同样，他也不认识圣子，因为他不认识圣父。\par

128. \Person{亚略}说：\Quote{［圣子是］另一实体的。}但有什么别的实体竟被高举到与天主子同等的地位，以至于单凭那实体，祂就成了天主子？再者，亚略派有何权利指责我们，因为我们在希腊语中谈论\OriginalTerm{实体}{\GreekText{οὐσία}}，或在拉丁语中谈论\Emph{实体}（Substantia），而他们自己却说天主子是另一\Quote{实体}，从而承认了一个神圣的\Emph{实体}（Substantia）？\par

129. 然而，若他们想要争辩\Quote{天主实体}或\Quote{天主性体}这些词语的使用，他们将很容易被反驳，因为圣经在希腊语中多次论及\OriginalTerm{实体}{\GreekText{οὐσία}}，在拉丁语中论及\Emph{实体}（Substantia），而且我们读到，圣\Person{伯多禄}要有分于\OriginalTerm{天主性体}{divine Nature}。但如果他们坚持圣子是另一\Quote{实体}的，他们就用自己的口自我反驳了，因为他们既承认了他们所畏惧的\Quote{实体}一词，又将圣子置于受造物的水平上，而他们却假意要将祂高举在这水平之上。\par

130. \Person{亚略}称天主子为受造物，不过\Quote{不像其他受造物}。然而，有什么受造物与其他受造物没有区别呢？人不像天使，地不像天，太阳不像水，光不像暗。因此，\Person{亚略}的这种偏向是空洞的——他不过是用拙劣的色彩伪装了他那欺人的亵渎之言，为要迷惑愚人。\par

131. \Person{亚略}宣称天主子可以改变和动摇。那么，如果祂是可变易的，祂又怎能是天主呢？因为祂亲口说过：\BibleQuote{我是，我是，我永不改变}？\par

\Emph{圣安博罗削表明他的愿望，希望有位天使飞到他面前来洁净他，就像从前色辣芬为依撒意亚所做的那样——不仅如此，更希望基督本人会降临到他、皇帝以及读者们面前；最后他祈求，愿主之杯的能力与神奇将格拉齐安及其他信友高举，他以神秘的语言描述了这杯。}\par

132. 然而，如今我必须宣认先知\Person{依撒意亚}的宣认，那是他在宣告主的言语之前所作的：\Quote{祸哉，我完了！因为我是个唇舌不洁的人，住在唇舌不洁的人民中间，竟亲眼见了\OriginalTerm{万军的主}{Lord of Sabaoth}！} 既然\Person{依撒意亚}在见了\OriginalTerm{万军的主}{Lord of Sabaoth}之后说\Quote{祸哉，我完了！}那么我，一个\Quote{唇舌不洁的人}，却不得不谈论天主的圣生，我将说什么呢？当\Person{达味}祈求在他所知道的事上，在他的口边派上守卫时，我又如何敢于开口谈论我所畏惧的事呢？哦，但愿也有一位\OriginalTerm{色辣芬}{Seraphim}，从天上的祭台用两约的钳子取一块燃烧的炭，并用它的火洁净我污秽的嘴唇！\par

133. 但是，既然当时\OriginalTerm{色辣芬}{Seraph}是在神视中降下到先知那里，而祢，主啊，却藉着奥秘的启示，亲自降生成人来到我们中间，那么求祢，不要藉任何代理人，也不要藉任何使者，而是祢亲自洗涤我的良心，清除我的隐罪，使我也能像从前不洁、如今却因祢的仁慈藉信德得以洁净的人一样，用\Person{达味}的歌词歌唱：\Quote{以色列的天主啊，我要弹琴赞美祢，我的口唇必将欢腾，我的灵魂也要一同赞美祢，因为它已蒙祢救赎。}\par

134. 因此，主啊，求祢远离那些诽谤并仇恨祢的人，来到我们这里，圣化我们至高元首\Person{格拉齐安}的耳朵，以及所有其他将读到这本小书的人的耳朵——并清洁我的耳朵，使任何地方都不再存留他们所听到的不信之瑕疵。那么，请彻底洗净我们的耳朵，不是用井水、河水，也不是用潺潺或淙淙的溪水，而是用那洁如清水的话，比任何水都清，比任何雪还纯——正是祢所说的话——\Quote{你们的罪虽似朱红，将变成雪一样的洁白。}\BibleRef{依 1:18}\par

135. 此外，还有一个杯，祢用它来涤净灵魂的密室，这杯不属于旧约的秩序，也不是装满普通葡萄树的酒——而是一个新杯，自天降下到地上，盛满了从奇妙葡萄串中榨出的酒；这葡萄串以肉躯的形状悬挂在十字架的树上，正如葡萄悬挂在葡萄树上一样。那么，正是从这个葡萄串，才有了那使人心喜悦的酒，\BibleRef{民 9:13} 它振奋悲伤者，充满芳香，将信德的狂喜、真诚的敬虔与洁净倾注在我们内。\par

136. 因此，主我的天主，求祢用这酒洁净我们至高皇帝的神魂之耳，好使他，如同常人饮了普通的酒之后，喜爱安息宁静，驱散了对死亡的恐惧，不感到伤害，不贪求属于他人的东西，也忘记了自己的；同样，他被祢的酒所陶醉，能喜爱和平，满怀信德之乐的信心，永不认识不信的死亡，表现出爱心的忍耐，不分受他人的亵渎之事，且把信德看得比亲族和子女更为重要，正如经上所说：\Quote{舍弃你所有的一切，来跟随我吧。}\BibleRef{玛 19:21}\par

137. 主耶稣，也求祢用这酒炼净我们的感官，好使我们能朝拜祢，钦崇祢，那有形与无形之物的创造者。的确，祢本身绝不会不是无形而美善的，因为祢已将无形性与美善赐予了祢手所造的工程。\par

\SourceRecord{Translated by H. de Romestin, E. de Romestin and H.T.F. Duckworth.；Nicene and Post-Nicene Fathers, Second Series；Vol. 10.；Edited by Philip Schaff and Henry Wace.；Buffalo, NY: Christian Literature Publishing Co.,；1896.；<http://www.newadvent.org/fathers/34041.htm>.}

% structure: p0001 source=h1 target=section reason=page-title

\section{论基督教信仰，卷二}

% structure: p0002 source=h3 target=omitted reason=duplicate-page-title

\Emph{列举天主子十二项名号，并分为三类。这些名号既证明子的永恒，也证明父的永恒。进而与\OriginalTerm{大司祭}{High Priest}胸牌上的十二块宝石相比，并以新的分类说明其不可分离性。再次回到大司祭胸牌的比喻，作者陈明这神秘衣袍上编织工艺与宝石之美，以及划分编织与宝石的深层含义；随后阐述所设譬喻，指出信仰必须与行动交织，并对论及子的同一信仰作简短总结。}\par

1. 尊威的皇帝陛下，我以为前一卷中已充分说明天主子是永恒的存在，与圣父没有差别，是受生而非受造：我们也已依据圣经章节证明，天主的真子即是天主，且其威严之明显记号也如此宣告。\par

2. 因此，尽管已经陈述的足以保障信仰——甚至已经丰盈外溢—— 正如江河之大小多由其源头涌流之貌来判断，然而为使我们信仰更清晰地显现于眼前，我以为泉源之水仍须分为三道。那么，首先，有表明在天主性内本质内在的明显记号；其次，有表明圣父与圣子相似的说法；最后，有表明天主性威严无可置疑之合一的说法。第一类名号为\Quote{生育}、\Quote{天主}、\Quote{圣子}、\Quote{圣言}；第二类为\Quote{光明}、\Quote{表达}、\Quote{镜子}、\Quote{肖像}；第三类为\Quote{智慧}、\Quote{能力}、\Quote{真理}、\Quote{生命}。\par

3. 这些记号如此宣示圣子的本性，使你们藉以知道圣父是永恒的，而圣子与祂没有差别；因为生育之源乃是那位\OriginalTerm{自有者}{He Who is}，且由永恒者所生，祂就是天主；从圣父出发，祂是圣子；从天主而来，祂是圣言；祂是圣父光荣的光辉，祂本体的肖像，天主的摹本，祂威严的肖像；是施恩者的恩惠，智慧者的智慧，全能者的能力，真实者的真理，生活者的生命。因此，圣父与圣子的属性相互一致，无人可以设想任何差别，或怀疑它们不属于同一威严。我们原可逐一举例说明这些名号的用法，若非受限于篇幅，力求简要。\par

4. 我们的信仰柱石，便由这十二项名号，如同十二块宝石，建立起来。因为这些宝石——\OriginalTerm{红玛瑙}{sardius}、\OriginalTerm{碧玉}{jasper}、\OriginalTerm{祖母绿}{smaragd}、\OriginalTerm{橄榄石}{chrysolite}等等——被织入圣洁的\Person{亚郎}的圣衣，即那背负\OriginalTerm{基督}{Christ}肖像，也就是真\OriginalTerm{司铎}{Priest}的预像者的圣衣；这些宝石嵌在黄金中，刻着以色列儿子的名字，十二块紧密相连，互相契合，若有人将其拆散分离，整个信仰结构便会坍塌。\par

5. 这便是我们信仰的基础——知道天主子是受生的；若祂非受生，便不是圣子。然而，仅称祂为圣子还不够，除非你更指明祂是\OriginalTerm{独生子}{Only-begotten Son}。若祂是受造物，便不是天主；若祂不是天主，便不是生命；若祂不是生命，便不是真理。\par

6. 因此，前三项记号，即\Quote{生育}、\Quote{圣子}、\Quote{独生子}，表明圣子原初即出自天主，并由其本性而然。\par

7. 随后三项——即\Quote{天主}、\Quote{生命}、\Quote{真理}——揭示祂的能力，祂以此建立并支撑受造世界。\Quote{因为，}如\Person{保禄}所说，\BibleQuote{我们生活、行动、存在，都在他内}\BibleRef{宗 17:28}。所以在前三项中显明圣子的自然权利，而后三项则彰显圣父与圣子之间运作的一致。\par

8. 天主子也被称为\Quote{肖像}、\Quote{光辉}和\Quote{表达}（\Quote{天主的}），因为这些名号显明圣父不可理解、不可测度的\OriginalTerm{威严}{Majesty}住在圣子内，以及祂在圣子内的相似表达。因此，我们看出这三项名号指涉于圣子与圣父的相似。\par

9. 我们还剩下\OriginalTerm{能力}{Power}、\OriginalTerm{智慧}{Wisdom}和\OriginalTerm{正义}{Justice}的运作，用来分别证明圣子的永恒。\par

10. 这就是那饰以宝石的袍子；是真\OriginalTerm{司铎}{Priest}的肩带；是婚宴礼服；此处有受默感的织工，他深知如何织成这件作品。这不是普通的编织物，主曾借先知论及它说：\Quote{谁将编织的技巧赐给了妇女？}也不是普通的石头——我们见到它们被称为\Quote{镶满的}石头；因为一切完美都基于不缺少任何东西。它们是彼此契合、嵌在黄金中的宝石——即属灵的类型；借我们的理智将它们结合，并以确凿的论证镶嵌。最后，圣经教导我们这些宝石何等非凡：一些人献上某种、另一些人献上另一种较次的祭品时，敬虔的首领们却把它们带来，佩戴在肩上，用它们做成\Quote{决断的胸牌}，也就是一件编织品。当信仰与行动并行时，我们便有了一件编织品。\par

11. 但愿无人以为我最初作了三重的划分，每部分包含四个，后来又作了四重的划分，每部分包含三个，是出于误导。美好事物的美妙，若从不同角度展示，便更令人喜悦。因为这些是美好的事物，司祭长袍的纹理正是它们的象征，也就是说，或是指法律，或是指教会，后者为她的新郎做了两件衣服，正如经上所载——一件是行动的，一件是精神的，交织着信德与善工的线。因此，在一处，如我们所读的，她以金子作底，然后在其上织以蓝色、紫色、朱红色和白色。此外，在另一处，她先制作蓝色和其他颜色的小花，然后附上金子，制成一件司祭长袍，为使由相同鲜艳颜色组成的各种恩宠与美善的装饰，因着排列的多样而获得新的荣耀。\par

12. 此外（为完成我们对这些预像的解释），可以肯定，精金和银子所指的是主的神谕，我们的信德借此得以稳固。\BibleQuote{主的神谕是纯净的神谕，是经火炼的银子，除去渣滓，七次炼净的。}蓝色如同我们所呼吸和吸入的空气；紫色则代表水的颜色；猩红象征火；而白麻布代表地，因其起源在地。这四种元素又构成了人的身体。\par

13. 因此，无论你是将信德与灵魂中已有的信德配合以身体的行为；或是行为在先，信德作为同伴加入，将它们呈献给天主——这便是司祭的礼服，这便是司祭的祭衣。\par

14. 因此，当她的额头因善工的美好冠冕而闪耀时，信德便是有益的。\BibleRef{雅 2:14}这信德——为把这事简明地说明——包含在以下不可推翻的原则中。如果子的起源在于虚无，祂就不是子；如果祂是受造物，祂就不是创造者；如果祂是被造的，祂就没有创造万物；如果祂需要学习，祂就没有预知；如果祂是接受者，祂就不是圆满的；如果祂进步，祂就不是天主。如果祂（与父）不相似，祂就不是（父的）肖像；如果祂是因恩宠而为子，祂就不是因性体而为子；如果祂不分享天主性，那么祂就有犯罪的倾向。\Quote{除天主性以外，没有谁是善的。}\BibleRef{谷 10:18}\par

\Emph{亚略派从圣玛尔谷福音\BibleRef{谷 10:18}那句\BibleQuote{除天主一位以外，没有谁是善的}所提出的论证，因着基督这些话的解释而被驳斥。}\par

15. 我现在必须面对的反驳，神圣的陛下，使我充满困惑，我的灵魂和身体因想到竟有人——或不如说不是人，而是外形像人，内心却充满野兽般的愚昧——他们在从主的手中领受了如此众多且重大的恩惠之后，竟能说所有美善之物的创造者自己不是善的，便感到虚弱。\par

16. 他们说，经上记载，\BibleQuote{除天主一位以外，没有谁是善的。}我承认圣经——但文字并无虚假；但愿在亚略派对它的解释中没有虚假。写下的文字是无辜的，可指责的是人们所理解的含义。我承认这话是主和救主的话——但让我们想想祂是在什么时候、对谁、以怎样的理解说这话的。\par

17. 圣子确实是以人的身份说话，并且是对一个经师说话——就是那个称圣子为\BibleQuote{善师}，却不愿承认祂是天主的人。对于他所不信的，基督进一步让他明白，好使他相信圣子，不是将其看作一位善师，而是看作善的天主，因为，如果凡提到\Quote{唯一的天主}时，圣子从不与那唯一者的完满分离，那么，当说唯一天主是善的时，独生子又怎能被排除在天主美善的完满之外呢？因此，亚略派必须要么否认圣子是天主，要么承认天主是善的。\par

18. 因此，主以受天主默感的理解，并没有说\BibleQuote{除圣父一位以外，没有谁是善的}，而是说\BibleQuote{除天主一位以外，没有谁是善的}，\Quote{父}乃是生育者的专有名字。但天主的唯一性绝不排除三个位格的天主性，因此所颂扬的是祂的性体。所以，美善属于天主的性体，而在天主的性体中，又存在着圣子——因此，这谓词所表达的，不是归于单一的位格，而是归于[完全]的唯一性[天主性]的。\par

19. 所以，主并没有否认自己的美善——祂是在责备这种门徒。因为当经师说\BibleQuote{善师}时，主回答说：\BibleQuote{你为什么称我善？}——这等于说：\Quote{称你所不相信是天主的那位为善，这是不够的。}我所寻找的门徒不是这样的人——那些只顾念我的人性，把我视为善师，却不仰望我的天主性，不相信我是善的天主的人。\par

\Emph{圣子的美善可从祂的工程得到证明，即祂在旧约下向以色列子民，以及在新约下向基督徒所施的恩惠。相信作为我们之主和审判者的美善，是与我们自己有利的。圣父为圣子所作的见证。为数不少的犹太人为圣子作证；因此亚略派显然比犹太人更坏。新娘的言语，同样宣告了基督的美善。}\par

20. 然而，我却不愿圣子仅依靠其性体的特权及祂至尊的固有权利。如果祂不配得这称号，我们就不要称祂为善；而如果祂不通过工程和仁爱的行为来配得此称号，就让祂放弃因性体而享有的权利，接受我们的审判。那位将要审判我们的，并不拒绝接受审判，为要使祂\BibleQuote{在祂所言的事上得以显明，在受审时要显为正直。}\par

21. 那么，祂岂不是善的，祂已向我显示了美好的事物？祂岂不是善的，当犹太人的六十万民众在追击者前奔逃之时，祂突然分开\Place{红海}的潮水，使一整片海水凝固？——波浪环绕着忠信者，成了他们的墙壁，却又倒流淹没了不信者。\par

22. 祂岂不是善的？在祂的命令下，海成为逃跑者脚下的干地，磐石为干渴的人涌出水来？这样，真正创造者的工程就得以被人认识，当流动的水变得坚实，磐石流出水来。为使我们承认这是基督的工程，宗徒说：\BibleQuote{那磐石就是基督。}\BibleRef{格前 10:4}\par

23. 祂岂不是善的，祂在旷野中以天降的食粮喂养了无数民众，使他们不受饥荒之苦，无需劳苦，安享休息？——以致在四十年的时间里，他们的衣服没有穿旧，鞋子没有磨破，为信众预表了将来的复活，表明伟大事业的荣耀、祂披覆我们身上的大能的华美，以及人类生命之流，都非徒然造成。\par

24. 祂岂不是善的，祂将大地提升到天上，以致如同光明的群星在天上如镜子般反映祂的荣耀，同样，宗徒、殉道者与司铎的团体，如荣耀的星辰闪耀，光照整个世界。\par

25. 那么，祂不但是善的，而且更超越。祂是善牧，不但为祂自己，也是为祂的羊，\BibleQuote{因为善牧为羊舍掉自己的性命。}是的，祂舍命为提升我们的生命——但这全凭祂天主性的权能，祂舍去性命，又再取回来：\BibleQuote{我有权舍掉我的性命，也有权再取回它。没有人从我夺去它，是我自己舍掉它。}\par

26. 你看到祂的善，因祂是自愿舍命：你看到祂的权能，因祂又取回性命——你怎能否认祂的善，当祂在福音中论到自己说：\BibleQuote{如果我是善的，为什么你的眼睛是恶的呢？}忘恩负义的可怜人，你在做什么？你否认祂的善，而你若真信的话，你对美善事物的盼望正是寄托在祂身上？你否认祂的善，而祂赐给了我们\BibleQuote{眼所未见，耳所未闻}的？\par

27. 为我的利益，我相信祂是善的，因为\BibleQuote{信赖上主是件好事。}为我的利益，我承认祂是主，因为经上记载：\BibleQuote{你们要称谢上主，因为祂是善的。}\par

28. 为我的利益，我尊崇我的审判者为善，因为主是以色列家的公义审判者。那么，如果天主子是审判者，而审判者就是公义的天主，并且天主子是审判者，[那么]那审判者兼天主子就是公义的天主。\par

29. 但是，也许你不相信别人，甚至不相信圣子。那么，请听圣父所说的：\BibleQuote{我的心灵吐露出美善的圣言。}那么，圣言是善的——\BibleQuote{圣言与天主同在，圣言就是天主。}\BibleRef{若 1:1} 因此，如果圣言是善的，而圣子就是天主的圣言，那么，无论\OriginalTerm{亚略派}{Arians}如何不悦，天主子就是天主。现在让他们至少羞愧脸红吧。\par

30. 犹太人曾常说：\Quote{祂是善的。}虽然有些人说：\Quote{祂不是，}但另有些人却说：\Quote{祂是善的。}——而你们\Emph{全都}否认祂的善。\par

31. 赦免一个人罪的那位是善的；那么，除免世罪者岂不是善的吗？因为指祂所说的正是：\BibleQuote{看，天主的羔羊，看，除免世罪者。}\par

32. 但我们为什么怀疑呢？教会历代信奉祂的善，并宣认其信德说：\BibleQuote{愿祂以口唇亲吻我，因为你的胸脯比美酒更佳美；}\BibleRef{歌 1:1} 又说：\BibleQuote{你的颈脖宛如最醇美的酒。}\BibleRef{歌 7:9} 因此，祂以自己的善，用法律和恩宠的胸脯滋养我们，以讲述天上的事理抚慰人的悲伤；而我们竟否认祂的善，当祂正是善的彰显，在祂身上表现出永恒美善的肖像，正如我们上文所指出经上记载的，祂就是那美善的无瑕映照和对应？\BibleRef{歌 7:9}\par

\Emph{因为天主是唯一的，天主子就是天主，既是善的又是真的。}\par

33. 但你们那些否认天主子之善和真天主性的人怎么想呢？虽然经上记载说，除了唯一的天主外没有别的神？\BibleRef{格前 8:4} 因为虽有所谓的神，你们竟把基督与那些被称为神而实非神者并列吗？而永恒是祂的本质，在祂以外，没有别的善而真的天主，因为天主在祂内；\BibleRef{若 17:22} 而且由父的本质推论，在祂之后没有别的真天主，因为天主是唯一的，既不像\OriginalTerm{萨培里派}{Sabellians}那样混淆父与子，也不像\OriginalTerm{亚略派}{Arians}那样分离父与子。因为父与子，作为父与子，是不同的位格，但他们的天主性不容分割。\par

\Emph{天主子之全能，藉旧约与新约的权威得以彰显。}\par

34. 既然天主子是真的、善的，那么祂无疑是全能的天主。对此还能有任何怀疑吗？我们已经引用过经上记载说：\BibleQuote{上主、全能者的名号是祂。}那么，因为圣子是主，而主是全能者，天主子就是全能者。\par

35. 但请听一段不容置疑的经文：圣经说：\BibleQuote{看，祂乘云而来，众目都要看见祂，连那些刺透祂的人也要看见，地上的万族都要因祂而哀恸。的确，阿们。我是\OriginalTerm{阿耳法}{Alpha}和\OriginalTerm{敖默加}{Omega}，那今在、昔在及将来永在的全能者上主天主说。}\BibleRef{默 1:8} 请问，他们刺透的是谁？我们盼望谁再来，岂不是子吗？因此，基督是全能的主，天主。\par

36. 神圣的陛下，请听另一段话——请听基督的话。\BibleQuote{全能的主这样说：在祂的光荣之后，祂派遣我到那掳掠你们的异民那里，因为凡触动你们的，就是触动祂眼中的瞳仁。看，我要向那掳掠你们的人挥动我的手，我要拯救你们，他们必成为那曾掳掠你们者的掠物；这样，他们必知道，是全能的主派遣了我。}显然，说话的是全能的主，而被派遣的也是全能的主。因此，全能的能力既属于圣父也属于圣子；然而，只有一位全能的天主，因为尊威是一体的。\par

37. 再者，为使您至杰出的陛下知道，那在福音中说话的，同样也在先知书中说话的，就是基督；祂藉着\Person{依撒意亚}的口，仿佛预先安排了福音，说：\BibleQuote{我自己，那说过话的，已来到了。}这就是说，我在法律中说过话的，如今临在于福音中。\par

38. 在别处，祂又说：\BibleQuote{凡圣父所有的，都是我的。}\BibleRef{若 16:25} 祂说\Quote{一切}是什么意思呢？显然，不是指受造之物，因为这一切都是藉着圣子而造成的，而是指圣父所有的一切——即永恒、主权、天主性，这些是祂的产业，因为是由圣父所生的。那么，我们不能怀疑，那拥有圣父所有的一切的，是全能的（因为经上记载：\BibleQuote{凡圣父所有的，都是我的}）。\par

\Emph{针对基督的全能所引用的某些圣经章节，得到释疑；作者也特别着力于表明，基督常常依照人性的感受说话。}\par

39. 虽然经上论到天主说：\BibleQuote{那真福而唯一全能者，}\BibleRef{弟前 5:15} 但我并不因此担忧，认为圣子与祂分离，因为圣经称天主，并非仅称圣父自己，为\Quote{唯一全能者}。圣父自己也藉着先知论及基督说：\Quote{我已将援助放在一位大能者身上。}所以，并非唯独圣父是唯一全能者；圣子也是全能者，因为在颂扬圣父时，圣子也同受颂扬。\par

40. 的确，让人指出有什么是圣子不能做的。当祂创造诸天时，谁是祂的助手？——当祂奠定世界的根基时，谁又是祂的助手？那在建立众天使和宰制者时无需助手的，难道需要助手来释放世人吗？\BibleRef{哥 1:15}\par

41. 他们说：\Quote{经上记载：\InnerQuote{我的圣父啊，若是可能，求你叫这杯离开我。}如果祂是全能的，祂怎么会怀疑可能性呢？}这意思是，因为我已证明祂是全能的，我也就证明了祂不可能怀疑可能性。\par

42. 你们说，这些话是基督的话。确实如此——但请考虑一下，祂说这些话的场合，以及祂以什么身份说话。祂已承担了人的实体，并随之有了人的感受。此外，你们在之前引用的篇章中发现：\BibleQuote{祂稍往前行，俯首至地祈祷说：圣父啊！若是可能……}因此，祂不是以天主身份，而是以人的身份说话；因为天主怎能不知任何事的可能与不可能呢？或者，为天主有什么不可能呢？\BibleQuote{因为对你来说，没有不可能的事。}\BibleRef{约 22:17}\par

43. 然而，祂怀疑谁呢——怀疑自己，还是怀疑圣父？显然是那位说\BibleQuote{求你叫这杯离开我}的，——祂是作为人，因着人性的感受而生疑。先知认为在天主没有不可能的事。先知不怀疑；难道你们认为圣子会怀疑吗？你们要把天主置于人之下吗？什么——天主竟怀疑圣父，在想到死亡时就惊恐吗？那么，基督害怕——害怕，而\Person{伯多禄}却毫无畏惧。\Person{伯多禄}说：\BibleQuote{我愿意为你舍命。}\BibleRef{若 13:37} 基督却说：\BibleQuote{我的心灵忧伤。}\BibleRef{若 12:27}\par

44. 两种记载都是真实的，而且同样自然的是，那较小的一位不知畏惧，而那较大的一位却忍受了这感受，因为前者完全不知死亡的力量，而后者，身为居住在肉身内的天主，却显示出肉身的软弱，好使那些否认降生成人奥迹的人，他们的邪恶无可推诿。祂这样说了，然而摩尼教徒不信，\Person{瓦伦廷}否认，\Person{马西昂}则断定祂只是个幻影。\par

45. 但事实上，祂实在是把自己置于与人同等的地位，正如祂在祂肉身的真实形体中所显示的，以至于祂说：\BibleQuote{但不要照我的意愿，而要照你的意愿，}\BibleRef{玛 26:39} 然而，基督的特有德能确是愿意圣父所愿意的，正如祂的特有德能是行圣父所行的一样。\par

46. 那么，在此，你们惯常用来反对我们的异议，即基于主曾说：\BibleQuote{不要照我的意愿，而要照你的意愿，}以及又说：\BibleQuote{我自天降下，不是为执行我的旨意，而是为执行派遣我来者的旨意，}\BibleRef{若 6:38} 应告一段落了。\par

\Emph{以上所引的圣经章节作为离题的缘由，藉着将这样的自由归于圣神，以及归于圣子的圣经篇章，证明了我们的主行动的自由。}\par

47. 现在，我们且更充分地解释，为何我们的主说：\Quote{若是可能，}并以此作为休战，同时我们证明祂拥有意志的自由。你们——在邪路上走得太远了——否认圣子有自由意志。再者，你们惯常贬损圣神，尽管你们不能否认经上记载说：\BibleQuote{圣神随意往哪里吹。}\BibleQuote{往哪里，}圣经说，而不是\BibleQuote{往哪里被命令。}那么，如果圣神随意往哪里吹，圣子难道不能做祂愿意做的吗？原来，正是这同一位圣子在祂的福音中说，圣神有能力随意往哪里吹。那么，圣子是否承认圣神比祂大，因为圣神有能力做那不允许祂自己的事呢？\par

48. 宗徒也说过：\BibleQuote{这一切都是同一位圣神所行的，随祂的旨意分给每个人。}\BibleRef{格前 12:11} \Quote{随祂的旨意，}注意——就是说，按照自由意志的判断，并非出于强迫。再者，圣神所分施的并非微小的恩赐，而是天主惯常所做的工作——治愈的神恩和行大能的作为。圣神既随己意分施，圣子就不能释放祂所愿意的人。但且听祂说，当祂确实照己意行事时：\Quote{我的天主，我乐意承行祢的旨意；}又说：\Quote{我要向祢献上自愿的祭献。}\par

49. 圣宗徒后来知道，耶稣有权柄随己意行事，所以当他看见耶稣在海面上行走时，便说：\BibleQuote{主，如果是祢，就让我在水面上走到祢那里去罢。}\BibleRef{玛 14:28} \Person{伯多禄}相信，如果基督命令，自然状况就可改变，水能承载人的脚步，不协调的事物能归于和谐一致。伯多禄请求基督命令，而非请求：基督没有请求，而是命令了，事就成就——而\Person{亚略}却否认这一点！\par

50. 实际上，有什么是圣父愿意而圣子不愿意，或者圣子愿意而圣父不愿意的呢？经上记载：\BibleQuote{父使祂所愿意的人复活。}\BibleRef{若 5:21} 圣子也使祂所愿意的人复活。现在，你告诉我，圣子使谁复活了，而圣父却不愿使他复活？然而，既然圣子使祂所愿意的人复活，且圣父与圣子的行动是同一的，你会看出，不仅圣子承行圣父的旨意，圣父也承行圣子的旨意。因为复活是什么？难道不是藉着基督的苦难而复活吗？而基督的苦难正是圣父的旨意。因此，圣子使谁复活，就是按圣父的旨意使他复活；所以，圣父与圣子的旨意是同一的。\par

51. 再说，圣父的旨意是什么？不就是要耶稣来到世界，洗净我们的罪过吗？听那癞病人的话：\BibleQuote{主！你若愿意，就能洁净我。}\BibleRef{玛 8:2} 基督回答说：\BibleQuote{我愿意，}健康的效果立刻随之而来。你难道看不出，圣子是祂自己旨意的主宰，而基督的旨意与圣父的旨意相同吗？的确，既然祂说过：\BibleQuote{凡父所有的一切，都是我的，}\BibleRef{若 16:15} 无疑没有任何被排除在外，圣子便拥有与圣父相同的旨意。\par

\Emph{前面提出的难题又再次着手解决。基督实实在在取了人性意志和情感，这些是一切与祂天主性不符之事的来源，因此必须归因于祂既是真天主又是真人的事实。}\par

52. 因此，行动同一的地方，意志也同一；因为在天主内，祂的意志直接产生实际效果。但天主的意志是一个，人的意志是另一个。再者，为显示生命是人性意志的目标，因为我们畏惧死亡，而基督的苦难却取决于天主的意志，即祂应为我们受苦，所以当\Person{伯多禄}要阻止祂受难时，主说：\BibleQuote{你所体会的，不是天主的事，而是人的事。}\BibleRef{玛 16:23}\par

53. 所以，祂取了我的意志，我的忧苦。我满怀信心地称之为忧苦，因为我宣认祂的十字架。祂称为自己的意志的那个意志，原是我的，因为祂作为人承受了我的忧苦，作为人说了话，所以祂说：\Quote{不要照我所愿意的，而要照祢所愿意的。}忧苦是我的，祂所承受的沉重也是我的，因为没有人临终时会欢欣踊跃。祂与我一同受苦，为我受苦，为我忧愁，为我沉重。因此，祂代替了我，并在我内忧伤了，祂自己本无理由忧伤。\par

54. 主耶稣，不是祢的创伤，而是我的创伤，使祢痛楚；不是祢的死亡，而是我们的软弱，正如先知所说：\BibleQuote{祂承受了我们的忧苦，}\BibleRef{依 53:4} 主，我们却以为祢受了苦难，其实祢不是为自己忧伤，而是为我。\par

55. 若祂为所有人忧伤，又为一个人哭泣，这有什么奇怪呢？在死亡的时刻，祂为所有人沉重，当祂要复活\Person{拉匝禄}时，也为一个人哭泣，这又有什么奇怪呢？\Emph{那时}，祂确实被爱他的姊妹的眼泪所感动，因为那些眼泪触动了祂的人心——而在这里，通过隐密的忧伤，祂成就了这样的事：正如祂的死亡结束了死亡，祂的鞭伤医治了我们的伤痕，同样，祂的忧伤也除去了我们的忧伤。\par

56. 作为人，祂疑惑；作为人，祂惊愕。祂的权能和天主性并不惊愕，而是祂的灵魂惊愕；祂惊愕是因为祂取了人性的软弱。既然祂取了灵魂，祂也取了一切灵魂的情感，因为就天主性而言，天主不可能忧伤或死亡。最后，祂呼喊说：\Quote{我的天主，我的天主，祢为什么舍弃了我？}所以，祂作为人说话，带着我的恐惧，因为当我们处于危险之中时，我们以为自己被天主舍弃了。因此，祂作为人忧伤，作为人哭泣，作为人被钉十字架。\par

57. 宗徒\Person{保禄}也这样说过：\Quote{因为他们已把基督的肉身钉在十字架上了。}宗徒\Person{伯多禄}也说：\BibleQuote{基督既在肉身上受了苦难。}\BibleRef{伯前 4:1} 因此，受苦的是肉身；天主性则远离死亡，安然无恙；身体按人性规律屈服于痛苦；如果灵魂不能死，天主性又怎能死呢？\BibleQuote{你们不要害怕那杀害肉身，}我们的主说：\BibleQuote{而不能杀害灵魂的。}\BibleRef{玛 10:28} 如果灵魂不能被杀死，天主性又怎能被杀死？\par

58. 那么，当我们读到荣耀的主被钉十字架时，不要以为祂是在祂的荣耀中被钉十字架的。\BibleRef{格前 2:8} 这是因为祂是\OriginalTerm{天主}{God}，也是人，因着祂的\OriginalTerm{天主性}{Divinity}而为天主，又因取了血肉，而成为人——\Person{耶稣基督}——因此荣耀的主被说成是被钉十字架的；因为祂拥有两个\OriginalTerm{本性}{natures}，即\OriginalTerm{人性}{humanity}和天主性，祂在祂的人性中承受了\OriginalTerm{受难}{Passion}，为要使那位受苦者既被称为荣耀的主，又被称为人子，正如经上所记载的：\BibleQuote{祂自天降下。}\BibleRef{若 3:13}\par

\Emph{\Person{基督}的这句话：\BibleQuote{父比我大}，是按照刚才确立的原则来解释的。其他类似的话也以同样的方式阐述。就祂的\OriginalTerm{天主性}{Godhead}而言，我们的主不能被称为次于父。}\par

59. 因此，正是由于祂的\OriginalTerm{人性}{humanity}，我们的主才疑惑、极其忧伤，并从死者中复活，因为跌倒的也必再起来。再者，正是因着祂的人性，祂说了那些话，我们的对手却恶意地用来攻击祂：\BibleQuote{因为父比我大。}\BibleRef{若 14:28}\par

60. 然而，当我们在另一处读到：\BibleQuote{我出于父，来到了世界上；我又离开世界，往父那里去，}\BibleRef{若 16:28} 祂怎会离去，若非藉着死亡？祂又怎样来到，若非藉着复活？再者，祂又加上句话，为要表明祂是在论及祂的\OriginalTerm{升天}{Ascension}：\BibleQuote{因此，我在事发生以前告诉了你们，好叫你们在事发生以后，能够相信。}\BibleRef{若 14:20} 因为祂是在论及祂身体的苦难和复活，而藉着那复活，那些先前怀疑的人得着引领而相信——的确，\OriginalTerm{天主}{God}永远临在每一处，不会从一个地方移到另一个地方。正如是人离去，照样是祂自己来到。再者，祂在另一处说：\BibleQuote{起来，我们从这里走吧。}\BibleRef{若 14:31} 因此，祂离去和来到，是在那与祂和我们共同的事上。\par

61. 的确，祂既然是完美而真实的\OriginalTerm{天主}{God}，怎会是较小的天主呢？然而，就祂的\OriginalTerm{人性}{humanity}而言，祂是较小的——但你仍诧异，祂以人的身份说话时，称父比祂大，而祂以人的身份时，竟称自己是虫，不是人，说：\BibleQuote{但我是虫，不是人；}又说：\BibleQuote{祂像被牵去宰杀的羊。}\BibleRef{依 53:7}\par

62. 若你在这方面宣称祂次于父，我不能否认；然而，按圣经的话说，祂不是被生为次等的，而是\BibleQuote{被贬抑成为稍逊者}，\BibleRef{希 2:9} 就是说，被\Emph{造成}次等的。祂怎样\BibleQuote{被贬抑成为稍逊者}呢？岂不是因为，\BibleQuote{祂虽具有天主的本体，却没有以与天主同等作为应把持不舍的，反而空虚了自己；}\BibleRef{斐 2:6} 确实，祂并没有舍弃祂原有的，而是取了祂所没有的，因为\BibleQuote{祂取了奴仆的形体。}\BibleRef{斐 2:6}\par

63. 再者，为要使我们知道，祂因取了一个身体而\BibleQuote{被贬抑成为稍逊者}，\Person{达味}已表明他所预言的是一个人，说：\BibleQuote{人算什么，你竟顾念他？人子算什么，你竟眷顾他？你使他稍微逊于天使。}宗徒在解释同一段经文时说：\BibleQuote{我们看见了\Person{耶稣}，因他所受的死亡之苦，接受了荣耀和尊荣的冠冕；他稍微逊于天使，好叫他在远离天主的情况下，为众人尝到死味。}\BibleRef{希 2:9}\par

64. 因此，\OriginalTerm{天主子}{Son of God}不仅被贬抑成为稍逊于父，也逊于天使。你若将此视为祂的屈辱，[我问]那么，就祂的\OriginalTerm{天主性}{Godhead}而言，难道圣子比那些事奉祂、服侍祂的天使更低吗？这样，你意图贬损祂的尊荣，反而陷入了扬天使之本性于天主子之上的亵渎中。但\BibleQuote{仆人不能大过主人。}\BibleRef{玛 10:24} 再者，即使在祂\OriginalTerm{降生成人}{Incarnation}之后，天使仍在服侍祂，为要叫你承认祂并未因取了肉身的性质而损伤其尊威，因为\OriginalTerm{天主}{God}不能承受丝毫自身的贬损，而祂从童贞女所取的，既不加增也不减损祂的\OriginalTerm{天主性能力}{divine power}。\par

65. 因此，祂既拥有圆满的\OriginalTerm{天主性}{Divinity}和光荣，\BibleRef{哥 2:9} 就祂的天主性而言，并非次等的。大和小是属于形体存在之物的区别；所谓更大，是指地位、品质，至少是年龄上的更大。当我们论及天主的种种时，这些词语便失去意义。通常，那教导并启迪他人的，被称为更大；但就天主的智慧而言，情形并非如此：它并非藉由接受他人的教导而建立，因为正是它为一切教导奠定了根基。然而，宗徒写下的话多么明智：\BibleQuote{好叫他在远离天主的情况下，为众人尝到死味，}——恐怕我们会以为那受难的是天主性，而非肉躯！\par

66. 那么，如果我们的对手找不出方法来证明[父]比[子]更大，就不要扭曲言辞以作虚假的报告，而应寻求其中的真义。因此，我问他们：在什么方面他们以为父更大？若因他是父，那么[我回答]这里不涉及年龄或时间的问题——父并非以白发为特征，子也不是以年轻为特征——而父亲的更大尊严正是依赖这些条件。然而，\Quote{父}与\Quote{子}不过是名称，一个是生育者的，一个是被生育者的——这些名称似乎使人结合而非分离；因为孝道并不带来个人价值的丧失，正如亲属关系把人联系起来，而不是把他们撕裂。\par

67. 因此，如果他们不能以自然的秩序作为任何质疑的依据，就让他们现在相信圣经的见证。如今，\OriginalTerm{圣史}{Evangelist}作证说，圣子并非因为身为圣子而低于圣父；不仅如此，他甚至宣称，圣子因其为圣子而与圣父同等，说：\BibleQuote{犹太人想要杀他，因为他不单违犯了安息日，而且又称天主为自己的父，使自己与天主同等。}\BibleRef{若 5:10}\par

68. 这并不是犹太人所说的话——是圣史作证说，当他称天主为自己的父时，他使自己与天主同等，因为圣经并没有把犹太人描述为说：\Quote{因此我们想要杀他}；而是圣史本人发言说：\Quote{因此犹太人想要杀他}。此外，他发现了原因，[说] 犹太人之所以想要杀他，是因为他作为天主违犯了安息日，又声称天主是自己的父，他不仅将违犯安息日的天主权威的尊威归于自己，而且，在论及他的父时，还将属于永恒同等的权利归于自己。\par

69. 天主圣子对这些犹太人做出了最恰当的回应，证明自己是圣子和与天主同等的。他说：\BibleQuote{凡父所作的，子也照样作。}\BibleRef{若 5:19} 因此，圣子既被宣称也被证明是与圣父同等的——这是一个真正的同等，既排除了天主性的差异，又显明了圣子及与圣子同等的圣父；因为有差异就没有同等，如果只有一位格，也没有同等，因为没有人与自己是同等的。这样，圣史就指出了，为什么基督称自己为天主圣子，也就是使自己与天主同等，是合宜的。\par

70. 因此，宗徒遵循这一启示，说：\BibleQuote{他并没有将自己与天主同等视为把持不舍的。}因为一个人没有的东西，他才想夺取。因此，与圣父同等乃是作为天主和主，他凭自己的本体所拥有的，并不是非法抢夺的战利品。所以宗徒又加上说：\BibleQuote{他取了奴仆的形体。}一个奴仆当然是与同等者相对的。因此，圣子在天主的形体中是同等的，但在取肉躯及作为人受苦时是低下的。因为同一性体怎能既低下又同等呢？而且，如果圣子较低，他怎能以同样方式作父所作的事呢？能力不同，又怎能有相同的行动呢？较低者能作出较大者那样的效果吗？在实体不同之处，能有行动的统一吗？\par

71. 因此，要承认，就基督的天主性而言，他不能被称为低于圣父。基督对亚巴郎说：\BibleQuote{我指着自己起誓。}\BibleRef{创 22:16} 现在宗徒指出，指着自己起誓的那一位不能低于任何别的。因此他说：\BibleQuote{当初天主应许亚巴郎的时候，因为没有大于自己可以指着起誓的，就指着自己起誓，说：\InnerQuote{我必多多祝福你，使你的后裔繁多。}}\BibleRef{希 6:13} 因此，基督没有大于他的，为此他指着自己起誓。而且，宗徒又正确地加上说：\BibleQuote{因为人都是指着比自己大的起誓}，因为人有一位比自己大的，但天主却没有。\par

72. 另一方面，如果我们的对手将这段经文理解为指圣父，那么记录的其余部分就不相合了。因为圣父没有显现给亚巴郎，亚巴郎也没有给天主圣父洗脚，而是给那位内里有着未来之人的肖像者洗脚。而且，天主圣子说：\BibleQuote{亚巴郎看见了我的日子，就欢喜。}\BibleRef{若 8:56} 因此，就是这一位指着自己起誓，[且]亚巴郎所看见的也是这一位。\par

73. 的确，那位在天主性上与圣父为一体的，怎能还有大于自己的呢？\BibleRef{若 10:30} 哪里有合一，哪里就没有差异，而在较大与较小之间却有区别。因此，我们面前这段圣经实例关于圣父和圣子的教导是：既不是圣父较大，圣子也没有任何在自己之上的，因为在圣父和圣子内，没有天主性的差异将他们分开，而是同一的尊威。\par

\Emph{面对圣子因受圣父派遣而至少在这一点上较低的说法，回答是：圣子也受圣神派遣，而圣神却不被认为大于圣子。此外，圣神反过来也由圣父派遣给圣子，为的是显示他们行动上的合一。因此，我们有责任仔细分辨哪些言论合宜地归于作为天主的基督，哪些归于作为人的基督。}\par

74. 对于那种常被提出的异议，即基督因受派遣而较低微，我毫不担心。因为即使他较低微，也并未因此得到证实；另一方面，他同等的尊位确实得到了证实。既然所有人都尊敬子如同尊敬父一样，\BibleRef{若 5:23} 那么可以肯定，圣子并不因受派遣而较低微。\par

75. 所以，不要计较人类语言的狭隘限制，而要注重言语的浅显意义，并相信已成的事实。请想一下，我们的主耶稣基督在\BibleBook{依撒意亚}{书}中说，祂是由圣神所派遣的。那么，圣子难道因为由圣神所派遣，就低于圣神吗？因此你们有记录，圣子声明自己是由圣父和祂的圣神所派遣的。祂说：\BibleQuote{我是元始，}\BibleRef{依48:12}\BibleQuote{且我永远生活；我的手奠定了大地的基础，我的右手使诸天稳固站立；}后面又说：\BibleQuote{我曾发言，也曾召唤；我领了他来，使他的道路亨通。你们接近我，听这些事：我从起初就没有在暗中说话。当它们被造时，我就在那里：现在，主和祂的圣神派遣了我。}\BibleRef{依48:15}在这里，创造天地的那一位亲自说，祂是由主和祂的圣神所派遣的。那么，你们看到，语言的贫乏并不减损祂使命的尊荣。因此，祂由圣父所派遣；也由圣神所派遣。\par

76. 为使你们领会，在尊威上毫无分离的差异，圣子相应地也派遣圣神，正如祂自己所说：\BibleQuote{当护慰者来时，就是我从父那里给你们派遣的、发自父的真理之神。}\BibleRef{若15:26} 关于这同一位护慰者也将由父所派遣，祂先前已教导说：\BibleQuote{但那护慰者，就是父因我的名所要派遣的圣神。}\BibleRef{若14:26} 请看他们的一体性，因为天主父所派遣的，子也派遣，父所派遣的，圣神也派遣。否则，如果亚略派不承认子是受派遣的，因为我们读到子是父的右手，那么他们自己将不得不承认有关父的，正是他们否认于子的，除非他们为自己发现另一位父或另一位子。\par

77. 所以，对那些在言语上无谓争吵的，要休战了，因为天主的国，如经上所载，不在于动听的言辞，而在于彰显的能力。我们当注意，天主性与肉体有所区别。在两者之中，说话的是同一位天主圣子，因为两种本性都在祂内；然而虽是同一位在说话，祂却并不总是以同一种方式说话。请看，在祂身上，一会儿显现天主的光荣，一会儿显露人的情感。作为天主，祂讲论天主的事，因为祂是圣言；作为人，祂讲论人的事，因为祂是在我的本性中说话。\par

78. \BibleQuote{这是从天上降下来的生活的食粮。}\BibleRef{若6:51} 这食粮就是祂的肉，正如祂自己说的：\BibleQuote{我所要赐给的食粮，就是我的肉。}\BibleRef{若7:52} 这就是那从天上降下来的，这就是父所祝圣并派遣到世界上来的。就连文字本身也教导我们，并非天主性，而是肉体需要祝圣，因为主亲自说：\BibleQuote{我为他们祝圣我自己，}\BibleRef{若17:19} 为使你们承认，祂既在肉体中为我们被祝圣，也藉着祂的天主性施行祝圣。\par

79. 这就是父所派遣的同一位，但如宗徒所说：\BibleQuote{由女人所生，属于法律之下，}\BibleRef{迦4:4} 这就是那位说：\BibleQuote{上主的神临于我身上，因此祂给我傅了油，派遣我向贫穷人传报喜讯：} 这就是那位说：\BibleQuote{我的教导不是我的，而是派遣我者。谁若愿意承行祂的旨意，就会认出这教导是出于天主，还是我凭自己讲的。}\BibleRef{若7:16} 那么，出于天主的教导是一回事；出于人的教导是另一回事；因此，当犹太人视祂为人，质疑祂的教导，并说：\Quote{这人没有学过，怎么通晓经书？} 耶稣回答说：\BibleQuote{我的教导不是我的，} 因为祂在教导时没有文字的修饰，似乎祂不是作为人，而是作为天主在教导，不是学来的，而是自己构想的教导。\par

80. 因为祂找到并策划了一切的智慧之道，正如我们在前面读到的，关于天主子曾说过：\BibleQuote{这是我们的天主，没有别的可同祂相比。祂找到了一切的智慧之道。此后祂出现在地上，与世人往来。} 那么，作为天主的祂，怎么可能没有自己的教导——祂在显现于地之前已经找到了一切的智慧之道？或者说，论到祂曾说\BibleQuote{没有可同祂相比的}，祂怎么会是低等的呢？确实，祂堪称为无可比拟的，与祂相比，没有别的可算得上——然而这样，祂也不能与父相比。但如果有人认为这是指父说的，他们就免不了陷入\Person{撒伯流}的亵渎中，将人性的摄取归于父。\par

81. 让我们继续往下看。\BibleQuote{凡凭自己讲的，是寻求自己的光荣。}\BibleRef{若7:18} 请看，父与子的一体性在此明显启示出来。说话的那位不可能不存在；但祂所讲的，却不可能仅仅出自祂自己，因为在祂内，所有的一切，自然都是从父而来。\par

82. 那么，\BibleQuote{寻求自己的光荣}这些话是什么意思呢？就是说，不是父无份的光荣——因为天主的圣言本就是祂的光荣。再者，主说：\BibleQuote{使他们看见我的光荣。}\BibleRef{若17:24} 但那圣言的光荣也就是父的光荣，正如经上记载：\Quote{主耶稣基督在天主父的光荣中。} 因此，就祂的天主性而言，天主子拥有自己的光荣，而父与子的光荣是同一的：所以，祂在光辉上并非低等，因为光荣是同一的；在天主性上也不是低等，因为天主性的圆满就在基督内。\par

83. 那么，你问，经上为何这样记载：\BibleQuote{父啊，时辰来到了，求你光荣你的子吗？} \BibleRef{若 17:1} 你说，说这些话的那位需要受光荣。至此你还有眼可见；你却没有读经上其余的话，因为接着写道：\BibleQuote{好叫你的子也光荣你。} 难道父需要什么，以至他要由子受光荣吗？\par

\Emph{论据基于子的顺服而被驳斥，且阐明天主圣三内能力、天主性和作为的合一；另外提及基督对祂母亲的顺服，祂显然不能因为顺服而被称为次等。}\par

84. 同样，我们的反对者通常质疑子的顺服，因为经上写道：\BibleQuote{祂空虚自己，取了奴仆的形体，与人相似，形状也一见如人；祂贬抑自己，听命至死，且死在十字架上。} \BibleRef{斐 2:7} 作者不但告诉我们子听命至死，而且首先说明祂是人，好让我们明白听命至死并非祂天主性的事，而是祂降生成人的事，由此祂承担了属于我们本性的职责和名称。\par

85. 因此，我们得知天主圣三的能力是合一的，正如我们在受难时及受难后所学到的：因为子藉着祂的身体受苦，这身体是祂受苦的保证；圣神倾注在宗徒们身上；灵魂被交托在父手中；进而，天主以大能的声音被宣告为父。我们得知父与子有同一形式、同一肖像、同一圣化，同一作为、同一光荣，最终，同一天主性。\par

86. 因此，只有一位天主，因为经上记载：\BibleQuote{你要敬畏上主你的天主，只事奉祂。} \BibleRef{申 6:13} 一位天主，并非如不虔敬的\Person{萨培里}所主张的那样，父与子是同一\OriginalTerm{位格}{Person}——而是因为父、子及圣神只有一个天主性。而那里有一个天主性，那里就有一个意志、一个目的。\par

87. 再者，为使你知道父存在，子存在，并且父与子的工作相同，请听从宗徒的话：\BibleQuote{愿天主我们的父，和我们的主耶稣，引领我们的道路到你们那里去。} \BibleRef{得前 3:11} 父与子都被提及，但引领是合一的，因为能力是合一的。同样，在另一处我们读到：\BibleQuote{愿我们的主耶稣基督，和爱我们，并赐给我们永远的安慰和在恩宠中的美好希望的天主我们的父，鼓励你们，并坚固你们的心。} \BibleRef{得后 2:15} 宗徒向我们呈现了多么完美的合一啊，以至于安慰的泉源不是多个，而是一个。那么，让疑惑止息吧，或者，如果疑惑不能被理性克服，就让它因我们主的仁慈心肠而软化。\par

88. 让我们回想，我们的主怎样仁慈地对待了我们，祂不仅教给了我们信德，也教给了我们品行。因为，祂取了人的形态，便服从\Person{若瑟}和\Person{玛利亚}。\BibleRef{路 2:51} 那么，祂因为服从就比所有人更卑微吗？服从是一回事，主权是另一回事，但服从并不排斥主权。那么，祂在哪方面服从父的法律呢？当然是在祂的身体上，在祂服从祂的母亲的地方。\par

\Emph{降生成人的目的和治愈效果。信德的好处，由此我们知道基督为了我们承受了一切软弱——基督，祂的天主性在祂的苦难中彰显出来；由此我们明白天主子被差遣不包含任何从属关系，我们无需担心这种信仰会令父不悦，因为父自己宣称祂在子内得享喜悦。}\par

89. 让我们同样以仁慈行事，让我们劝说反对者那对他们有好处的事，\Quote{让我们在造我们的上主面前，敬拜哀号。} 因为我们不是要推翻，而是要治愈；我们不给他们设伏，而是尽本分警告他们。仁慈善于折服那些武力和辩论都无法胜过的人。再者，我们的主曾用油和酒治愈了那个从\Place{耶路撒冷}下到\Place{耶里哥}、落入强盗手中的人；祂没有用法律的严厉药剂或预言的苛刻来对待他。\par

90. 因此，愿意痊愈的人，都要到祂那里去。让他们领受祂从父那里带下来、在天上制成的良药，这药是用那些不会枯萎的天上果实的汁液调制的。这不是地上的出产，因为自然界中根本没有这种合剂。祂怀着令人惊叹的目的取了我们的肉躯，为要显示肉身的法律已服从于心灵的法律。祂降生成人，为使祂——人类的导师——作为人而战胜。\par

91. 如果祂作为天主，只显露祂能力的臂膀，仅展示祂那不可侵犯的天主性，那对我有何益处呢？祂为什么取人性，不就是为在与我本性及软弱相同的条件下受试探吗？祂理当受试探，理当与我一同受苦，好使我晓得受试探时如何得胜，被困迫时如何逃脱。祂凭着节制之力、轻看财富、信德而得胜；祂践踏野心、逃避放纵、远避淫佚。\par

92. \Person{伯多禄}看见了这药物，便舍弃了他的网，就是说，谋利的工具和保障，弃绝了肉身的私欲，如同弃绝一艘漏水的船，这船好像接收了众多偏情的舱底污水。这真是强效的良药，不仅除去了旧伤的疤痕，甚至切断了偏情的根源与源头。信德啊，比一切宝库更丰裕；绝妙的良药啊，治愈我们的创伤和罪恶！\par

93. 我们当想想正确信仰的裨益。为我有益的是认识到，基督为了我的缘故承受了我的软弱，顺服于我肉身的偏情，为了我，也就是说，为了每一个人，祂成了罪，成了诅咒，为了我，并在我内，祂被贬抑、被服从，为了我，祂是羔羊、葡萄树、磐石、仆人、婢女的儿子，不知道审判的日子，为了我的缘故，不知晓那日子和那时辰。\par

94. 因为创造了时日的那一位，怎可能不知道那一天？那位既已宣告了审判的时日与缘由，又怎可能不知道那一天？因此，祂成为诅咒，并不是按祂的\OriginalTerm{天主性}{Godhead}，而是按祂的\OriginalTerm{肉身}{flesh}；因为有经上记载说：\Quote{凡挂在木架上的，是可咒骂的。} 所以，祂是在肉身内及按肉身被悬挂，为此缘故，那担负我们诅咒的，自己成了诅咒。祂流泪，是为使你，人啊，不致长久流泪。祂忍受凌辱，是为使你不在遭受错待时悲痛。\par

95. 多么荣耀的医治--- 得享基督的安慰！因为祂为我们的缘故，以超乎寻常的忍耐承受了这些事---而我们，为了祂圣名的光荣，竟连普通的忍耐都无法忍受！谁能在遭受攻击时，不学会宽恕？看哪，基督甚至在十字架上，还为迫害祂的人祈祷。难道你没有看出，基督的这些\Quote{软弱}（如你所喜欢称呼的）正是你的\Emph{力量}吗？你为何在关于我们医治的事上质疑祂？祂的眼泪洗涤我们，祂的哭泣洁净我们---至少在这\Emph{疑惑}中仍有力量：你若开始怀疑，就会绝望。所受的凌辱越大，所当献的感恩也越大。\par

96. 即便在嘲笑和凌辱的时刻，也当承认祂的\OriginalTerm{天主性}{Godhead}。祂被悬在十字架上，一切元素都向祂致敬。\BibleRef{玛 27:51} 太阳收回了光线，白日消逝，黑暗降临，遮盖大地，大地震动；然而那悬在十字架上的祂却毫不震动。这些征兆岂不是表明对造物主的敬畏？祂被悬在十字架上---\Person{亚略}派啊，你只看到这一点；祂赐予天主的国---你却看不到。你读到祂尝了死味，\BibleRef{路 23:43} 却无视祂也邀请强盗进入乐园。\BibleRef{若 20:11} 你注视在坟墓旁哭泣的妇女，却不看在那旁守候的天使。你读祂\Emph{说}了什么，却不读祂\Emph{做}了什么。你说主对客纳罕妇人说：\BibleQuote{我被派遣，只是为了以色列家迷失的羊。}\BibleRef{玛 15:24} 你却不说祂照那妇人所求的给她做了。\par

97. 由此你应明白，祂\Quote{被派遣}并不意味着祂是奉别人的命令而受强迫，而是祂凭\OriginalTerm{自由意志}{free will}，按照自己的判断行事，否则你就是指责祂轻视父。因为按照你的解释，如果基督来到犹太地区，只是执行父的命令，解救犹太地区的居民，并不顾及其他人，但在此之前，祂却解救了客纳罕妇人的女儿脱离疾病，那么祂显然不仅是执行他人的指示，更是自由地运用自己的判断。但哪里有按自己意志行事的自由，哪里就不会违背自己使命的条件。\par

98. 不要怕子的行为使父不悦，因为子亲自说过：\Quote{凡祂所喜悦的事，我常作，} 并说：\Quote{我所做的事业，也是祂所做的。} 那么，父怎能不悦于祂借子所做的呢？因为只有一位天主，正如经上记载的：\Quote{祂使受割损的人因着\OriginalTerm{信德}{faith}而成义，也使未受割损的人因着\OriginalTerm{信德}{faith}而成义。}\BibleRef{罗 3:30}\par

99. 要阅读全部圣经，细心留意一切，你就会发现，基督如此显示自己，为使人能在人身上认出天主。当你听见父宣告祂对子的喜悦时，不要恶意误解子因父而欢跃。\par

\Emph{天主教徒与亚略派，谁更能确保自己获得作为审判者的基督的恩待？针对 \BibleRef{咏 109:1} 的异议得以解答，显示出当父邀请子坐在自己右边时，既不意味着臣服，也不意味着偏袒，因为子坐在父的右边。天主圣三以及性体合一的真理，已由天使的三圣颂所证实。}\par

100. 然而，如果我们的对手不能以善意挽回，就让我们把他们传唤到审判者面前。那么，我们要到哪一位审判者面前呢？当然是到那位拥有审判权者面前。是到父面前吗？不，因为\Quote{父不审判任何人，祂把审判全交给了子。}\BibleRef{若 5:22} 祂交给了子，意思是，这并非出于慷慨施予，而是在\OriginalTerm{受生}{generation}的过程中给付的。所以，你看，祂何等不愿意你侮辱祂的子---甚至将子赐给你作你的审判者。\par

101. 那么，在审判之前，让我们看看谁更有理，是你还是我？诉讼中谨慎的一方，必先努力赢得审判者的好感。你尊敬人---难道不尊敬天主吗？请问，何者能赢得法官的青睐---尊敬还是轻蔑？就算我有错---其实我确实没错：难道基督会不悦于人们对祂表示的敬意？我们都是罪人---那么，谁更堪当宽恕，是那献上敬拜的，还是那显出傲慢的？\par

102. 如果推理不能打动你，至少让审判的清晰景象打动你吧！抬起你的眼睛，注视审判者，看看是谁坐在那里，与谁同坐，坐在何处。基督坐在父的右边。如果你肉眼不能看见，就听听先知的话：\BibleQuote{上主对我主说：你坐在我的右边。}\BibleRef{咏 109:1} 因此，子坐在父的右边。现在告诉我，你这认为天主的事要按世间的事来判断的人---你认为坐在右边的那位是较低的吗？父坐于子的左边，这对父是羞辱吗？父光荣子，你却将此当作侮辱！父愿这邀请成为爱与尊重的标记，你却把它当作君主的命令！基督从死者中复活，坐在天主的右边。\par

\Quote{但，}你反驳说，\Quote{圣父\Emph{说了}。}很好，现在听一段圣父没有说话，而圣子预言的话：\BibleQuote{从此以后你将看见人子坐在大能者的右边。}\BibleRef{玛 26:64} 这是指祂要收回自己的身体——对祂，圣父曾说：\Quote{坐在我右边。}如果你问的是\OriginalTerm{天主性}{Godhead}的永恒居所，当比拉多问祂是否是犹太人的君王时，祂说：\Quote{我为此而生。}因此，宗徒确实指出，我们应当相信基督坐在天主的右边，不是出于命令，也不是出于任何恩赐，而是作为天主最最亲爱的圣子。因为经上为你记载：\BibleQuote{你们寻求天上的事吧，那里有基督，坐在天主的右边；你们要体味天上的事。}\BibleRef{哥 3:2} 这就是体味天上的事——相信基督在坐席时，不是如同接受命令的仆人那样服从，而是作为爱子受到尊敬。因此，圣父是针对基督的身体说：\Quote{坐在我右边，等我使你的仇敌作你的脚凳。}\par

如果你还想曲解这句话的意思，\Quote{我要使你的仇敌作你的脚凳，}我回答说，圣父也将那些要被圣子复活和赐予生命的人带到圣子面前。因为基督说：\BibleQuote{凡不是派遣我的父所吸引的人，谁也不能到我这里来，而我在末日要叫他复活。}\BibleRef{若 6:44} 而你们却说，天主子因软弱而服从——这位圣父将人带到祂面前，好让祂在末日复活他们的圣子。请问，在你看来，这就是服从吗？当国度是为圣父预备，而圣父把人带到圣子那里，并且没有任何曲解言辞的余地，因为圣子把国度交给圣父，而没人能优先于祂？因为圣父交给圣子，圣子又交还给圣父，这清楚地证明了爱与尊重：双方的如此交付，使得接受者并非获得他人之物，交付者也没有损失。\par

再者，坐在右边并不表示优越，坐在左边也并不表示耻辱，因为\OriginalTerm{天主性}{Godhead}中没有等级，不受空间或时间的限制，这些只是我们微小人心的度量衡。没有爱的差别，没有任何东西分裂\OriginalTerm{合一}{Unity}。\par

但是何必走这么远？你们已环顾四周，看见了审判者，注意到天使宣告祂。他们赞美，你却辱骂祂！宰制者和掌权者在祂面前俯伏——你却污蔑祂的名！所有圣人都朝拜祂，但天主子不朝拜，圣神也不朝拜。\OriginalTerm{色辣芬}{seraphim}说：\BibleQuote{圣！圣！圣！} \BibleRef{依 6:3}\par

这同一个名字\Quote{圣}的三次呼喊意味着什么？如果重复三次，为什么只是一个赞美行动？如果是一个赞美行动，又为什么重复三次？为什么重复三次，难道不是说明圣父、圣子和圣神在圣善上是一体的吗？\OriginalTerm{色辣芬}{seraphim}说出这名字，不是一次，免得他排除圣子；不是两次，免得他忽略圣神；不是四次，免得他把\OriginalTerm{受造物}{created beings}也纳入（对造物主的赞美）。此外，为表明\OriginalTerm{天主圣三}{Trinity}的\OriginalTerm{天主性}{Godhead}是唯一的，他在三次\Quote{圣}之后，又以单数添上\Quote{\OriginalTerm{万军的主天主}{Lord God of Sabaoth}。}因此，圣父是圣的，圣子是圣的，圣神同样是圣的，所以\OriginalTerm{天主圣三}{Trinity}受朝拜，但不朝拜；受赞美，但不赞美。至于我，我宁愿相信色辣芬，按照天上所有总领天使和掌权者的方式朝拜。\par

\Emph{亚略派、萨培理派和摩尼教徒关于其审判者所持的邪恶可耻的意见被简短驳斥。基督对其他仇敌的警告被阐明后，圣盎博罗削表示希望自己得到更仁慈的审判。}\par

那么，让我们继续你们的指控，看看你们如何赢得审判者的欢心。现在说吧，说吧，我说，告诉祂：\Quote{我认为祢，基督，与祢的圣父不相像；} 祂必回答说：\Quote{注意，如果你能，注意，我说，并告诉我你们认为我在哪方面不同。}\par

你再说：\Quote{我断定祢是\OriginalTerm{受造物}{created being}；} 基督必回答：\Quote{如果两个人的见证是真的，你们岂不该相信我和我父，就是称我为祂圣子的吗？}\par

然后你会说：\Quote{我不承认祢[完美]的美善；} 祂必回答：\Quote{照你的\OriginalTerm{信德}{faith}，给你成就吧；这样我就不对你行美善。}\par

\Quote{祢是\OriginalTerm{全能者}{Almighty}，我不承认；} 祂转过身来回答：\Quote{那么我就不能赦免你的罪了。}\par

\Quote{祢是一个受统治的存在。}对此祂必回答：\Quote{那么，你为什么向你所认为如同奴隶般受统治的那位寻求自由和宽赦呢？}\par

我看见你的控告止于此。我不逼迫你，因为我自己也深知我的罪过。我不吝惜宽赦你，因为我自己也想获得宽免，但我想知道你们祈祷的对象。看，当我在审判者面前陈述你们的愿望时，我不揭露你们的罪过，而是看你们的祈祷和愿望如何依次陈明。\par

所以，说出那些愿望吧，所有的人都想得到它们。\Quote{主，使我成为天主的\OriginalTerm{肖像}{image}。}对此祂必回答：\Quote{什么肖像？是你们所否认的肖像吗？}\par

\Quote{使我成为\OriginalTerm{不朽}{incorruptible}的。}祂的回答必然是：\Quote{我怎能叫你不朽呢？你们称我为\OriginalTerm{受造物}{created being}，从而把我当作\OriginalTerm{可朽}{corruptible}的。死人复活时将要洁净，脱离腐朽——你们称你所看见是天主的那位可朽的吗？}\par

116. \Quote{你对我行善吧。} \Quote{你为什么向我求你所否认的呢？我原愿你成为善的，我曾说 \BibleQuote{你们应是圣的，因为我是圣的}，\BibleRef{肋 19:2} 你却立意否认我是善的？\Emph{你}呢，还寻求罪的赦免吗？不，除了天主，没有人能赦罪。\BibleRef{谷 2:7} 既然对你而言我不是真实唯一的天主，我绝不能赦免你的罪。}\par

117. 亚略和福提努斯的追随者这样说。\Quote{我否认你的天主性。} 主将回答他们说：\Quote{愚妄的人心中说：没有天主？你们想这是指谁说的？——是犹太人、外邦人，还是魔鬼？无论这是说谁，福提努斯的门徒啊，那沉默不言的还更可容忍；你竟然胆敢扬声说出来，显明你比愚人更愚。你否认我的天主性，然而我曾说过：\BibleQuote{你们是神，你们都是至高者的儿子}？你否认他是天主，却看见他周围如同神的工作。}\par

118. 萨培里派的人接着发言。\Quote{我认为祢，祢自己，同时是圣父、圣子和圣神。} 主将回答他说：\Quote{你既不听从圣父，也不听从圣子。此事还能有什么疑惑呢？圣经本身教导你，是圣父把审判权交出来，而圣子施行审判。\BibleRef{若 5:22} 你没有听从我的话：\InnerQuote{我不是独自一个，而是有圣父与我同在，祂派遣了我。}}\par

119. 摩尼教徒发言。\Quote{我认为魔鬼是我们肉身的创造者。} 主将回答他说：\Quote{那么，你在这天上做什么呢？离开，到你创造者那里去吧。\BibleQuote{我愿我的人与我同在，就是圣父赐给我的人。} \BibleRef{若 17:24} 你，摩尼教徒，自认为魔鬼的受造物；那么，赶快到他居所去，就是那火与硫磺的地方，那里的火不灭，免得刑罚永无止境。}\par

120. 我暂且不提其他异端——不是人物，而是怪象。什么审判在等待他们，对他们的判决形式将怎样？对所有这些人，祂确实要回答，更多是忧伤而非愤怒：\Quote{我的人民啊，我对你们做了什么，在哪里烦难了你们？难道不是我领你们出了 \Place{埃及}，从为奴之家带你们进入自由吗？}\par

121. 但是，仅领我们出 \Place{埃及} 得自由，救我们脱离为奴之家，还不够：比这更大的恩惠，祢已将自己赐给了我们。那么祢要说：\Quote{\BibleQuote{我不是背负了你们的一切疾苦吗？}\BibleRef{依 53:4} 难道我没有为你们交付我的身体？难道我没有寻求死亡，这死亡与我天主性无关，却是为你们的救赎所必需的吗？难道\Emph{这就是}我要得的感谢？我的圣血就换得这样的结果吗？正如我昔日藉先知的口说过：\BibleQuote{我的血有什么益处，因为我竟归于腐朽？} 难道这就是益处，让你们邪恶地否认我——你们，正是我为他们忍受这些的人？}\par

122. 至于我，主耶稣，虽然我自觉身负大罪，但仍要说：\Quote{我没有否认祢；祢可以宽恕我肉身的软弱。我承认我的过犯；我不否认我的罪。\BibleQuote{祢若愿意，就能洁净我。}\BibleRef{玛 8:2} 因这句话，那癞病人得了所求。求祢不要审判祢的仆人，我恳求祢。我所求的，不是审判，而是宽恕。}\par

\Emph{审判者的判决已颁布，反对者的抗辩已受审，判决的终局性，即无可上诉，已获证明。}\par

123. 我们从基督期待什么判词？这事我知道。难道我要说，祂\Emph{将会}给出什么判词吗？不，祂已经宣布了判决。我们手中有这判决。祂说：\BibleQuote{众人都要尊敬圣子，如同尊敬圣父一样。那不尊敬圣子的，就是不尊敬派遣祂来的圣父。} \BibleRef{若 5:23}\par

124. 如果这判决不令你满意，就向圣父上诉，撤销圣父已作出的判语。你竟说祂有一个与祂不相似的儿子。祂将回答：\Quote{那么是我说谎了，我曾对圣子说：\BibleQuote{让我们照我们的肖像，按我们的模样造人。}} \BibleRef{创 1:26}\par

125. 告诉圣父，祂创造了圣子，祂将回答：\Quote{那么，你为什么朝拜了你认为是受造者的那位呢？}\par

126. 告诉圣父，祂生了比自己次要的儿子，祂将回答：\Quote{将我们比较一下，让我们看看。}\par

127. 告诉圣父你毫不信任圣子，祂要回答：\Quote{难道我没有对你们说过：\BibleQuote{这是我的爱子，我所喜悦的，你们要听从祂。}} \BibleRef{玛 17:5} 这些字句 \InnerQuote{你们要听从祂} 是什么意思，岂不是 \InnerQuote{当祂说：\BibleQuote{凡父所有的一切，都是我的}，要听从祂} 吗？宗徒们听到了，正如经上记载：\BibleQuote{他们就俯伏在地，非常害怕。} \BibleRef{玛 17:6} 如果承认祂的人都仆倒在地，那些否认祂的人又当如何？但耶稣伸手拉起宗徒们——却要任你俯伏，使你看不见你所否认的荣耀。\par

128. 所以我们要注意，因为圣子定罪的，圣父也定罪，所以我们应尊敬圣子如同尊敬圣父，好能藉着圣子抵达圣父。\par

\Emph{圣盎博罗削婉拒对自己功德的任何称赞：无论如何，信德已获圣经的权威支持所充分辩护，而亚略派却如犹太人一样顽固，充耳不闻。他祈求他们能受感动而热爱真理；同时，他们应被视为异端者和基督的仇敌，予以躲避。}\par

129. 陛下，我已将这些论证，以简要概括的方式，粗略而非以完全解释和精确顺序的形式陈述出来。如果亚略派认为它们不完善且未完成，我确实承认它们几乎尚未开始；若他们认为还有需要提出的，我承认几乎全部都还需提出；因为不信者亟需论证，而信友则绰绰有余。的确，\Person{伯多禄}的单次宣认便足以使人对基督怀有信心：\BibleQuote{你是基督，永生天主之子。}因为知晓祂的\OriginalTerm{神圣的生育}{Divine Generation}，既无分裂也无减损，既非衍出也非受造，就已足够。\par

130. 此理，确实在一部部圣经中均已宣告，无一例外，然而不信者仍存疑惧：因为，如经上所载，\BibleQuote{这百姓的心迟钝了，耳朵沉重，眼睛紧闭，免得眼睛看见，耳朵听见，心里明白，}\BibleRef{依 6:10} 因为，就像犹太人一样，亚略派惯于掩耳，或喧哗鼓噪，每当\OriginalTerm{救恩的圣言}{Word of salvation}被听闻时。\par

131. 不信者怀疑人的话，又有什么奇怪呢，既然他们拒绝相信\OriginalTerm{天主的圣言}{Word of God}？天主子，正如你将看到福音中所载的，说：\BibleQuote{父啊，光荣祢的名吧！}就有声音从天上传来，说：\BibleQuote{我已光荣了它，我还要再光荣它。}\BibleRef{若 12:28} 这些话不信者听见了，却不相信。子说话，父回答，犹太人却说：\BibleQuote{是打雷。}另有些人说：\BibleQuote{是天使向他说话。}\BibleRef{若 12:29}\par

132. \Person{保禄}，再者，如\BibleBook{宗徒}{大事录}所载（\BibleRef{宗 22:9}），当保禄因基督的圣音领受\OriginalTerm{恩宠}{grace}的召叫时，与他同行的几个同伴中，唯有他说自己听见了基督的圣音。这样，神圣的陛下，那信的人，听见了——他听见，是为能信；而那不信的人，听不见，不，他不愿听，他无法听，免得他信了！\par

133. 至于我，诚然，巴不得他们愿意听，好能相信——以真正的爱与温良来听，如同寻求真理的人，而不攻击一切真理。因为经上写道，我们不要留意\BibleQuote{那无稽的传说和无穷的族谱，这些事只会引发争论，对于天主藉信德所立的工作毫无益处。这训令的目的就是爱，即由纯洁的心、光明的良心和真诚的信德所发出的爱；有些人离弃了这些，转而讲些虚妄的言论，他们愿作法学士，却不明白自己所说的话，也不明白自己坚决所主张的道理。}另一处，同一宗徒又说：\BibleQuote{至于那些愚拙而无知的辩论，你务要躲避。}\BibleRef{弟后 2:23}\par

134. 对这些播弄争端的人——即\OriginalTerm{异端者}{heretics}——\OriginalTerm{宗徒}{Apostle}吩咐我们远离他们。关于他们，宗徒在另一处说：\BibleQuote{有些人必离弃\OriginalTerm{信德}{faith}，听从那欺骗人的神和魔鬼的教训。}\BibleRef{弟前 4:1}\par

135. 同样，\Person{若望}说异端者是\OriginalTerm{假基督}{Antichrists}，这清楚地指向亚略派。因为这个[亚略]\OriginalTerm{异端}{heresy}是在所有其他异端之后兴起的，并汇集了所有异端的毒害。正如经上论到假基督说：\BibleQuote{牠开口亵渎\OriginalTerm{天主}{God}，亵渎他的圣名，并与他的\OriginalTerm{圣人}{saints}作战。}\BibleRef{默 13:6} 他们也是这样侮辱天主子，连他的\OriginalTerm{殉道者}{martyrs}也不放过。更有甚者，也许假基督都不会做的事，他们竟然篡改了\OriginalTerm{圣经}{Scriptures}。因此，凡说耶稣不是基督的，这人就是假基督；凡否认\OriginalTerm{世界的救主}{Saviour of the world}的，就是否认耶稣；凡否认\OriginalTerm{子}{Son}的，也否认\OriginalTerm{父}{Father}，因为经上记载：\BibleQuote{凡否认子的，也否认父。}\BibleRef{若一 2:23}\par

\Emph{\Person{圣盎博罗削}向\Person{格拉齐安}保证胜利，并宣称这在\Person{厄则克耳}的预言中早已预言。此盼望更坚立于皇帝的虔诚之上，先前的灾难乃是对东方异端的惩罚。此卷书以向\OriginalTerm{天主}{God}祈祷作结，恳求祂此刻显示仁慈，拯救信友的军队、国土与君主。}\par

136. 在这备战并战胜\OriginalTerm{蛮族}{Barbarians}的时节，我不敢再久留陛下。去吧，你的庇护，确实，在\OriginalTerm{信德}{faith}的盾牌之下，并佩带\OriginalTerm{圣神的剑}{sword of the Spirit}；去吧，去到那自古所许、并在\OriginalTerm{天主}{God}所赐\OriginalTerm{神谕}{oracles}中预言的胜利之中。\par

137. 因为\Person{厄则克耳}，在那些遥远的日子里，已经预言了我们民族的减少与哥特战争，说：\BibleQuote{人子，你要说预言，讲论\Person{哥格}说：上主这样说——在那一天，当我的百姓以色列安居之时，你难道不从你的住处，从极北之地，率领许多民族，都是骑马而来，形成庞大强盛的聚集，以及众多军队的威力，上来攻击我的百姓以色列，如密云遮盖地面，在末日吗？}\par

138. 那\Person{哥格}就是\OriginalTerm{哥特}{Goth}人，我们已看见他的到来，而上主的话应许在未来的日子里将战胜他：\BibleQuote{上主说，他们要掠夺那些曾掠夺他们的人，抢掠那些曾夺走他们财物为掠物的人。在那一天，我要将一块著名的地方，就是以色列著名的坟墓，给\Person{哥格}——}即给哥特人——\BibleQuote{那是堆积许多死人的地方，就是那些前往海滨的人，它将环绕并封闭山谷的出口，在那里，\OriginalTerm{以色列家族}{the house of Israel}必倾覆\Person{哥格}和他的所有群众，那地要称为\Person{哥格}群众的谷：以色列家族必倾覆他们，使这地得以洁净。}\par

139. 再者，神圣的陛下，我们既已投身于与\OriginalTerm{异教不信}{alien unbelief}的抗争，就毋庸置疑，我们将享有那在您內坚强的\OriginalTerm{天主教信仰}{Catholic Faith}的助佑。的确，天主义怒的原因业已显明，以致对\Place{罗马帝国}的信仰首先被动摇，正是在那里，对天主的信仰退让了。\par

140. 我无意追述那些\OriginalTerm{宣信者}{confessors}的死亡、酷刑与流放，信友的职任如何被当作礼物赏给了卖国贼。我们岂未听闻，沿整个边境——从\Place{色雷斯}，经河畔的\Place{达基亚}、\Place{默西亚}，直至\Place{潘诺尼亚}人的\Place{瓦莱里亚}全境——亵渎者的宣道与蛮族的入侵交相喧嚷？如此嗜血的邻人能带给我们什么益处？或者，\Place{罗马国家}依靠这样的防卫者如何能得安全？\par

141. 全能的天主啊，我们以自己的血、自己的流放，为\OriginalTerm{宣信者}{confessors}的死亡、\OriginalTerm{司铎}{priests}的流放、以及如此狂傲的邪恶罪愆所做的补赎，已经够了，实在绰绰有余——那些背弃信仰的人必不能得享安全，这已足够显明了。主啊，求你回转，重树你信仰的旗帜。\par

142. 引领我军前行的，并非军团的鹰帜，亦非飞鸟的占卜，而是你的圣名，\Person{主耶稣}，与对你的敬拜。此地绝非不信者之地，而是素以遣出\OriginalTerm{宣信者}{confessors}为俗的土地——\Place{意大利}；常受试探，却从未被诱离的\Place{意大利}；你的陛下久加护卫，现今又从蛮族手中解救出来的\Place{意大利}。我们的皇帝毫无动摇之心，信德坚如磐石。\par

143. 求你如今彰显你尊威的明显标记，好使那相信你是真\OriginalTerm{万军的上主}{Lord of Hosts}、\OriginalTerm{天军统帅}{Captain of the armies of heaven}的人；那相信你是天主的真能力与真智慧，既非时间之内的存有，亦非受造之物，而是正如经上所载的，\BibleQuote{天主永恒的能力与天主性}（\BibleRef{格前 1:24}）的人，能蒙你至高威能的助佑，为他所持的信德赢得胜利的桂冠。\par

\SourceRecord{Translated by H. de Romestin, E. de Romestin and H.T.F. Duckworth.；Nicene and Post-Nicene Fathers, Second Series；Vol. 10.；Edited by Philip Schaff and Henry Wace.；Buffalo, NY: Christian Literature Publishing Co.,；1896.；<http://www.newadvent.org/fathers/34042.htm>.}

% structure: p0001 source=h1 target=section reason=page-title

\section{论基督徒的\OriginalTerm{信德}{Faith}，第三卷}

% structure: p0002 source=h3 target=omitted reason=duplicate-page-title

\Emph{陈述前两卷所论问题何以在后三卷中更详尽探讨的理由。为使用寓言辩护，并以圣经的榜样为支持，因圣经中可见多种诗体寓言的比喻，尤其是塞壬，即感官享乐的象征；基督徒应受教以\Person{保禄}的话和基督的作为来逃避它们。}\par

1. 既然您最仁慈的陛下曾命我为教导您而撰写一些有关\OriginalTerm{信德}{faith}的论著，且您曾亲自召见我并鼓励我的怯懦，我如同临战前夕的人，仅写了两卷，为指出我们信德进展的某些道路与途径。\par

2. 然而，我看到某些心怀恶意之人，一心播散纷争，尚未耗尽攻击之力，而您仁慈的陛下的虔敬关切又召唤我进一步努力，因为您渴望在更多事上考验您已在少数事上检验过的人，我遂决定对已在寥寥数语中处理过的问题，再稍加详尽论述，以免人们以为，我并非在宁静与信心之中提出那些主张，而是我虽断言了它们，却因怀疑而放弃了辩护。\par

3. 再则，我们曾提及\OriginalTerm{许德拉}{Hydra}与\OriginalTerm{斯库拉}{Scylla}（见第一卷第6章46节），并以比喻的方式引入它们，为要表明，我们当如何提防不信的层出不穷，或提防在其浅滩上造成的不祥舟覆之灾；若有人以为，借自诗人传奇的此类论据的润饰是不合法的，并因无机会谤毁我的信德，便攻击我的言辞，那么，让他知道，不仅片语，甚至连完整的诗句也已被织入圣经的文本中。\par

4. 例如，那一句，\BibleQuote{我们也是祂的后裔，}\BibleRef{宗 17:28}，\Person{保禄}凭先知经验教导并加以引用，是从哪里来的呢？先知的话语方式并不避讳巨人，也不避讳\OriginalTerm{提坦}{Titans}山谷；\Person{依撒意亚}也曾提及\OriginalTerm{塞壬}{sirens}和鸵鸟的女儿。\Person{耶肋米亚}也曾预言\Place{巴比伦}将有塞壬的女儿居住其中，为要表明，巴比伦的罗网，即此世的喧嚣，好比古代情欲的故事：它们似乎在此生多岩的岸边歌唱着悦耳的曲调，却足以致命，要网罗青年的\OriginalTerm{灵魂}{soul}——那位希腊诗人自己也告诉我们，那位明智者仿佛被自己审慎的锁链捆绑，才得以逃脱。在基督来临之前，即使对强者而言，要从享乐那欺人的幻象和诱惑中拯救自己，也被视为极难的事。\par

5. 但是，如果那位诗人尚且断定，世俗享乐与放纵的诱惑能毁灭人的心智，并必然导致覆舟之灾，那么，对于圣经为我们写下的这话：\BibleQuote{不要为肉体的\OriginalTerm{私欲}{concupiscence}安排}，我们又当如何思想呢？经上又说：\BibleQuote{我痛击我的身体，使它为奴，免得我给别人报捷，自己反落选。}\BibleRef{格前 9:27}\par

6. 实在，基督为我们赢得了\OriginalTerm{救恩}{salvation}，不是藉着享乐，而是藉着禁食。况且，祂禁食不是为了为自己求得\OriginalTerm{恩宠}{grace}，而是为了教导我们。祂饥饿，也并非因身体软弱而受制，而是借着饥饿证明祂确实取了\OriginalTerm{肉身}{body}；为教导我们，祂不但取了我们的肉身，也取了肉身的软弱，正如经上所说：\BibleQuote{祂承受了我们的脆弱，担负了我们的疾苦。}\par

\Emph{基督为了我们而取的肉身的属性和遭遇，不应归于祂的\OriginalTerm{天主性}{Godhead}；在天主性方面，祂是\OriginalTerm{至高者}{Most High}。否认这一点，就等于说父\OriginalTerm{降生成人}{Incarnation}了。当我们读到天主是唯一的，且除祂以外再没有别的，或唯独祂具有不死不灭性时，这也必须理解为适用于基督，这不仅是为了避免上述的罪恶\OriginalTerm{异端}{heresy}（\OriginalTerm{父受难论}{Patripassianism}），而且因为父与子的活动被宣布为一而同一的。}\par

7. 那么，祂饥饿，是肉身的软弱，即是我们的软弱；祂哭泣，并且忧伤至死，这是出于我们的本性。为何要将我们本性的属性和事件归于天主性呢？别人告诉我们祂被\Emph{造成}，这也是肉身的属性。因此，我们确实读到：\Quote{\Place{熙雍}我们的母亲要说：\InnerQuote{祂是个人，}在她内祂成了人，至高者亲自奠定了她的基础。} \Quote{祂成了人，}请注意，并非\Quote{祂成了天主。}\par

8. 然而，那位同时是至高者和人的，是谁呢？无非是\BibleQuote{在天主与人之间的中保，也就是那身为人的基督耶稣，祂奉献了自己，为众人作赎价。}\BibleRef{弟前 2:5} 这段话确实正适用于祂的降生成人：因为我们的\OriginalTerm{救赎}{redemption}是由祂的宝血完成的，我们的宽赦是借着祂的德能而来的，我们的生命是藉祂的恩宠而保全的。祂以至高者的身份赐予，以人的身份祈求。前者是创造者的职分，后者是救赎者的职分。恩赐虽各有不同，赐恩者却是一位，因为我们的创造者成为我们的救赎者乃是合宜的。\par

9. 谁能否认我们有明显的证据表明基督是至高者呢？凡是另有想法的人，就是将降生成人的\OriginalTerm{奥迹}{mystery}归功于天主父。然而，基督是至高者，这一点已经由圣经在另一处论及\OriginalTerm{受难}{Passion}奥迹时所说的话清楚证明了：至高者发出自己的声音，大地为之震动。而且在福音中你可以读到：\BibleQuote{至于你，小孩，你要称为至高者的先知，因为你要走在上主前面，为祂预备道路。}\BibleRef{路 1:76} \Quote{至高者}是谁？就是天主子。那么，那位至高者天主就是基督。\par

10. 再者，虽然圣经处处称天主是唯一的天主，但天主子并未与这\OriginalTerm{一体性}{Unity}分离。因为那位至高者是唯一的，正如经上所说：\BibleQuote{让他们知道祢的名是上主：祢是全地唯一的至高者。}\par

因此，对手所引以为据的可耻结论，即他们从圣经论及天主的话语中得出者，便带着轻蔑与羞辱而被驳斥了：\BibleQuote{祂是那独享\OriginalTerm{不死性}{immortality}、居于不可接近的光中者}\BibleRef{弟前 6:16} 因为这些论及天主的话，这名号同等地归于圣父与圣子。\par

12. 实际上，如果在读到\Quote{天主}之名时，他们否认这也指圣子（如同圣父），他们便是在亵渎，因为他们否定了圣子的神圣主权，并且他们似乎分享了\OriginalTerm{撒伯里乌派}{Sabellians}的罪错，即教导圣父\OriginalTerm{降生成人}{Incarnation}。那么，让他们解释一下，他们如何能不将宗徒的这话理解为对圣父的亵渎：\BibleQuote{你们在祂内也一同复活了，是因相信那使祂从死者中复活的天主的大能。}\BibleRef{哥 2:12} 同时让他们从接下来的话中警醒自己正面临什么——因为接下来便是：\BibleQuote{你们从前因你们的过犯和肉身的未受割损，原是死的；但祂却使你们与祂一同生活了，恕赦了我们的一切过犯，涂抹了那反对我们、与我们对立的规条上所写的债据，把它从我们中间拿去，钉在十字架上；解除了自己的肉身。}\BibleRef{哥 2:13}\par

13. 那么，我们不可设想那位使肉身复活的圣父是唯一的天主；同样，我们也不可设想那位身体复活的圣子有什么不同。那位使肉身复活的，确实也使人生活；那位使人生活的，也赦免了罪过；那位赦免罪过的，也涂抹了债据；那位涂抹债据的，也把它钉在十字架上；那位钉在十字架上的，也解除了自己的肉身。但是，解除肉身的并不是圣父；因为\Emph{不是}圣父，而是如经上所载：圣言成了血肉。\BibleRef{若 1:14} 你们由此看出，\OriginalTerm{亚略派}{Arians}在分裂圣父与圣子时，有陷于声称圣父受苦的危险。\par

14. 然而，我们能够轻松地证明这些话论及圣子的行动，因为圣子确实亲自使自己的身体复活了，正如祂亲口说的：\BibleQuote{你们拆毁这座圣殿，三天之内，我要把它重建起来。}\BibleRef{若 2:19} 祂也亲自使我们与祂的身体一同生活：\BibleQuote{就如父唤醒死者，使他们生活，同样，子也随自己的意愿使人生活。}\BibleRef{若 5:21} 祂亲自赦免罪过，说道：\BibleQuote{你的罪赦了。}\BibleRef{路 5:20} 祂也将那记在债据上的手书钉在十字架上，因为祂被钉在十字架上，在肉身上受了苦。除天主子外，再没有谁解除肉身的，而祂正是那穿上肉身的。因此，那位完成了我们复活之工的，显然被指明为真天主。\par

\Emph{圣父与圣子不可分割，这由宗徒的话所证明，因为圣子理当是受赞颂的、唯一的掌权者、本性而\OriginalTerm{不死}{immortal}——不是靠恩宠，天使本身亦不死——并居于不可接近的光中。圣父与圣子如何同样且同等地被称为\Quote{唯一者}。}\par

15. 因此，当你读到\Quote{天主}之名时，不要分开圣父与圣子，因为圣父与圣子的\OriginalTerm{天主性}{Godhead}是一而同一的。所以当你读到\BibleQuote{受赞颂的、唯一的掌权者}\BibleRef{弟前 6:15} 也不要分开他们，因为这些话论及天主，正如你也读到：\BibleQuote{我在那使万物生活的天主前恳切求你。}\BibleRef{弟前 6:13} 基督也确实使人生活，所以天主之名恰当地归于圣父与圣子，因为他们行动的效果是一致的。让我们继续看下文：\BibleQuote{我恳切求你，}他说，\BibleQuote{在那使万物生活的天主和耶稣基督前。}\par

16. 圣言在天主内，正如经上所载：\BibleQuote{我在天主内赞美祂的圣言。}在天主内的是祂的永恒大能，即耶稣；因此，在论及天主时，宗徒见证了天主性的合一，而藉着基督之名，他则见证了\OriginalTerm{降生成人的圣事}{sacrament of the Incarnation}。\par

17. 此外，为表明他所讲的是基督的降生成人，他又加上：\BibleQuote{祂曾在\Person{般雀比拉多}前宣了好誓，}〔我恳切求你〕\BibleQuote{持守这训令，不受玷污，无可指摘，直到我们的主耶稣基督显现；在适当的时候，那受赞颂的唯一掌权者，万王之王，万主之主，必要将这显现彰显出来；祂是那独享不死性、居于不可接近的光中者，没有人看见过，也不能看见。}\BibleRef{弟前 6:13} 所以，那些话论及天主，这名号的尊荣和真理性同为〔圣父与〕圣子所共有。\par

18. 那么，这些事既然全都适用于圣子，此处又怎能没有想到圣子呢？若非如此，那就否认祂的天主性吧，这样你就可以否认对天主所应有的说法了。那位施予祝福者，祂的受赞颂性不可否认，因为\BibleQuote{那些过犯得赦免的人，是有福的。}那位赐予我们健全教导者，不能不被称为\BibleQuote{受赞颂的}，正如经上所载：\BibleQuote{这教导是按着那受赞颂的天主的光荣福音。}\BibleRef{弟前 1:11} 祂的大能不可否认，圣父论及祂说：\BibleQuote{我把救助放在一位大能者身上。}当祂自己也使其他者不死时，谁敢拒绝承认祂是不死的呢？正如论及天主的智慧所载：\BibleQuote{藉着她，我将拥有不死性。}\BibleRef{智 8:13}\par

19. 但是，祂本性的不死是一回事，我们本性的不死是另一回事。可朽坏的事物不能与神圣的事物相比。天主性是唯一不为死亡所触及的\OriginalTerm{实体}{Substance}，因此，宗徒虽知道〔人的〕灵魂和天使都是不死的，却宣称唯一天主具有不死性。其实，灵魂也会死亡：\BibleQuote{谁犯了罪，谁就该死，}\BibleRef{则 18:20} 而天使也非绝对不死，他的不死取决于造物主的旨意。\par

20. 不要匆匆弃绝这一点，因为\Person{加俾额尔}没有死，\Person{辣法耳}也没有，\Person{乌黎耳}也没有。即使就他们的本性而言，也有犯罪的可能性，虽然不是经由训练而进步的可能，因为每一个理性的受造物都受到自身之外的影响，并趋向于审判。我们所受的影响，决定了审判的赏罚、败坏或趋向成全，因此\WorkTitle{训道篇}说：\BibleQuote{因为天主必审判一切行为}。因此，每一个受造物内都蕴含着败坏与死亡的可能性，即使它[目前]不死或不犯罪；即使在任何事上它没有自陷于罪，这也不是出于它不朽本性的恩赐，而是出于训练或恩宠。所以，作为赠予的\OriginalTerm{不死不灭}{immortality}是一回事：没有变化可能的不死不灭是另一回事。\par

21. 我们是否因为基督在肉躯内为众人尝了死味，就否认基督\OriginalTerm{天主性}{Godhead}的不死不灭呢？那么加俾额尔就比基督更好了，因为加俾额尔从未死过，而基督却交付了灵魂。但仆人不能大过主人，\BibleRef{玛 10:24}，我们必须分清肉躯的软弱和天主性的永恒。基督的死亡源自肉躯，而不死不灭则是基督主权之本然。但如果天主性使得肉躯不见腐朽，而肉躯确实本性上是会腐朽的，那天主性本身又怎能死亡呢？\par

22. 如果圣子在圣父怀中，如果圣父是光，圣子也是光，因为天主是光（\BibleRef{若一 1:5}），那么圣子怎能不居于不可接近的光中呢？或者，如果我们设想在天主性的光之外，另有某种光是不可接近的光，那么这光岂不比圣父更大，以至那经上记载与圣父同在、又在圣父内的那位，竟不在那光中了么？因此，当人们只读到\Quote{\Emph{天主}}时，不要排除对圣子的思想；当他们只读到\Quote{\Emph{子}}时，也不要排除对圣父的思想。（\BibleRef{若 16:32}）\par

23. 在世上，圣子并不缺少圣父，\BibleRef{若 10:30}，而你竟认为圣父在天上缺少圣子吗？圣子在肉躯内——（当我说\Quote{祂在肉躯内}或\Quote{祂在世上}时，我说的是假如我们生活在福音所载的那些日子，因为如今我们不再按肉躯认识基督了 \BibleRef{格后 5:16}）——祂在肉躯内，却不独自，如经上所说：\BibleQuote{我并不独自，因为有父与我同在}\BibleRef{若 8:16}，你竟想圣父独自居于光中吗？\par

24. 免得你们以为这论证只是推测，请接受这权威的话语。圣经说：\BibleQuote{从来没有人见过天主；只有那在父怀中的独生子，启示了祂。}\BibleRef{若 1:18} 如果圣子在圣父怀中，圣父怎能独居呢？圣子怎能启示祂所未见的呢？因此，圣父并不独自存在。\par

25. 现在请留意圣父和圣子的\Quote{独处}是什么样的。圣父是独一的，因为没有其他的父；圣子是独一的，因为没有其他的子；天主是独一的，因为天主圣三的天主性是一体的。\par

\Emph{我们被告知，基督只是在肉躯方面是\Quote{被造}的。为了救赎人类，祂不需要任何援助方式，正如祂也不需要任何援助以完成祂的复活；而其他人，为了复活死者，则需要求助于祈祷。甚至当基督祈祷时，那祈祷也是以祂的人性身份献上的；同时，由于祂命令（某事当行），祂必须被视为天主。在这一点上，魔鬼的证词比亚略派的论证更真实。讨论最后解释，为何\Quote{大能者}的称号被赋予人子。}\par

26. 现在已经充分说明，圣父并非没有圣子而独自为天主，圣子也不能被认为是没有圣父而独自为天主，因为我们读到天主之子是\Quote{\OriginalTerm{被造}{made}}的，这是就祂的肉躯而言，并非就祂由圣父天主受生而言。\par

27. 的确，祂藉圣祖之口，说明了祂在何种意义上\Quote{\OriginalTerm{被造}{made}}，说：\BibleQuote{因为我的灵魂满溢忧苦，我的生命临近阴府。我已算作与堕入深坑的人同列；我成了一个自由的人，无助，在死者中。} 因此，我们读到的是：\Quote{我被造如人}，而不是\Quote{我被造如天主}；又说：\Quote{我的灵魂满溢忧苦}。请注意，是\Quote{我的灵魂}，不是\Quote{我的天主性}。祂之被\Quote{造}，乃在于那与阴府相关，那被算作与他人同列的部分，因为天主性不容有任何相似，可以成为与它物归类的根据。然而请注意，即使在那注定要死的肉躯内，天主性的尊威也显于基督。虽然祂被\Quote{造}如人，被\Quote{造}如肉躯，但祂在死者中成为自由的，\Quote{自由，无助}。\par

28. 但是，既然已经说过：\BibleQuote{我已将援助加在大能者身上}，圣子又怎能在此说祂无助呢？请区分这里呈现的两个本性。肉躯需要帮助，天主性却不需要。因此，祂是自由的，因为死亡的锁链无法拘禁祂。祂没有被黑暗的势力俘虏，反倒是祂在他们中间施展权能。祂是\Quote{无助}的，因为祂本人、主，并非藉任何使节或特使的职分，而是凭自己的大能，拯救了祂的子民。祂既能使别人复活，又怎可能需要任何帮助来使自己的身体复活呢？\par

29. 虽然人也曾使死者复活，但他们并非凭借自己的能力，而是因基督的名这样做。祈求是一回事，命令是另一回事；获得与赐予不同。\par

30. \Person{厄里亚}确实复活了死者，但他是祈祷了——他并未命令。\Person{厄里叟}则是先伏在死者身上，按照死者的姿势，使之复活；\BibleRef{列下 4:34}而且，再次，\Person{厄里叟}的遗体仅仅一接触，便赋予了死者生命，为使这位先知预示那一位的来临——祂虽被派遣取了罪恶的肉身的形状，却甚至在祂被埋葬后，也能使死者复活。\par

31. 再者，\Person{伯多禄}在治好\Person{艾乃阿}时，说：\BibleQuote{因纳匝肋人耶稣之名，起来行走吧！}并非因他自己的名字，而是因基督之名。但\Quote{起来}是一项命令；另一方面，这是对自己权利的信赖的一个实例，\BibleRef{谷 16:17}而非对权力的傲慢宣称，并且这命令的权威在于这名字的有效效力，而不在于它自身的能力。那么，亚略派能作何回答呢？\Person{伯多禄}以基督之名命令——这是一方面；另一方面，他们却认为圣子并未命令，而是请求。\par

32. 他们反对说，我们读到祂曾发出祈祷。\BibleRef{若 11:41}但请注意其中的区别。祂以\Emph{人子}的身份祈祷，以\Emph{圣子}的身份命令。难道你们不愿将连魔鬼都归于圣子的，也归于祂吗？你们岂非要指控自己比\Person{撒殚}更为邪恶吗？魔鬼说：\BibleQuote{若你是圣子，就命令这块石头变成饼吧！}\BibleRef{路 4:3}\Person{撒殚}说\Quote{命令}，你们却说\Quote{恳求}。魔鬼相信，藉着圣子的一句话，一种单一物质的本质可以转变为复合物质的本质；你们却以为，除非圣子提出请求，甚至祂的旨意也无法成就。再者，魔鬼认为应从祂的能力来评估圣子，\BibleRef{罗 1:4}你们却认为应从祂的软弱来评估。魔鬼的试探比亚略派的争论更可容忍。\par

33. 那么，我们若在一处发现人子被称作\Quote{大能的}，而在另一处却发现光荣的主被钉在十字架上，我们就不应感到困惑。\BibleRef{格前 2:8}有什么能力比统御天上的大能者更大呢？但这能力却在祂手中，祂统御著上座者、宰制者和天使；因为，正如经上所载，祂虽置身于野兽之中，却有天使服侍祂，为叫你们明了那属于降生成人的本性与那属于至高权力的本性之间的区别。所以，就祂的肉身而言，祂遭受野兽的攻击；就祂的天主性而言，祂却受天使的朝拜。\par

34. 那么，我们已经得知，祂成了人，并且祂之被造成必须归于祂的人性。此外，在圣经的另一处，你可以读到：\BibleQuote{祂被造成属于祂，出自\Person{达味}的后裔，}\BibleRef{罗 1:3}这就是说，就肉身而言，祂由\Person{达味}的后裔被\Quote{造成}，但祂是在万世之前由天主所生的天主。\par

\Emph{从圣经引证的一些段落，为表明\Quote{造成}未必总是与\Quote{创造}同义；由此得出的结论是，不应像犹太人那样，以圣经的字面意义为吹毛求疵的争论的依据；然而犹太人还是被证明不像异端者那样恶劣，因此先前确立的原则重新得到了确认。}\par

35. 同时，\Quote{成为}并不总意味着\Emph{创造}；因为我们读到：\BibleQuote{主，你成了我们的避难所}，以及\BibleQuote{你成了我的救援。}显然，这里并没有陈述\Emph{创造}的事实或目的，而是说天主成了我的\Quote{避难所}，并转向了我的\Quote{救援}，正如宗徒所说：\BibleQuote{祂为了我们成了来自天主的智慧、正义、圣化与救赎，}\BibleRef{格前 1:30}这就是说，基督是为了我们而被\Quote{造成}的，出于圣父，而非\Emph{受造}。再者，同一作者在后续章节中解释了他是在何种意义上说基督为我们成了智慧：\BibleQuote{然而我们宣讲的是在奥秘教义中的天主智慧，这智慧是隐藏的，是天主在世界存在以前，为使我们获得光荣而预定的，这世界的统治者没有一个认识这智慧，因为如果他们认识了，决不至于将光荣的主钉在十字架上。}当宣讲苦难的奥秘时，确实没有涉及永恒的生成过程。\par

36. 所以，主的十字架是我的智慧；主的死亡是我的救赎；因为我们是藉著祂的宝血被赎回的，正如宗徒\Person{伯多禄}所说的。\BibleRef{伯前 1:19}因此，作为人，主以祂的血救赎了我们；而作为天主，祂也赦免了罪。\BibleRef{谷 2:8}\par

37. 因此，我们不要在言辞中设下陷阱，也不要急切地在其中寻找纠缠；我们不要因为不信者曲解了书面文字的意义，便只提出字面所呈现的，而非深层的意义。犹太人正是由此走向毁灭，他们藐视深奥隐藏的意义，只追随字句的赤裸形式，因为\BibleQuote{文字叫人死，圣神却叫人活。}\BibleRef{格后 3:6}\par

38. 然而，在这两种可憎的不虔中，将仅属于人性的归于天主性，或许比将仅属于字面的归于精神更为可憎。犹太人惧怕相信人性被提升到天主内，因而失去了救赎的恩宠，因为他们拒绝了那救恩所赖以存在的基础；亚略派则将天主性的尊威贬抑到人性的软弱。尽管那些钉死主之肉身的犹太人是可憎的，但我仍认为那些相信基督的天主性被钉在十字架上的人更为可憎。因此，一位常与犹太人打交道的人说：\BibleQuote{在警告过异端者一两次后，就该远离他。}\BibleRef{铎 3:10}\par

39. 再者，这些人也不谨慎避免对圣父的侮辱，他们竟将基督为了我们而\Quote{成了}智慧这一事实，不虔敬地应用于祂那超越一切时间界限与区分的不可理解的\OriginalTerm{受生}{generation}上；因为，暂且不论对圣子的侮辱就是对圣父的冒犯，他们甚至在对圣父的攻击中大肆亵渎，关于圣父经上记载说：\BibleQuote{愿天主是真实的，而众人都是撒谎者。}\BibleRef{罗 3:4}如果他们以为这是指圣子而言，他们并没有堵塞祂的受生，反而因倚靠此经文的权威而承认他们所拒绝的，即基督是天主，是真天主。\par

40. 若我要逐一细数我们读到祂\Quote{成了}之处，那未免太冗长了；这\Quote{成了}并非出于本性，而是出于恩宠的安排。例如，梅瑟说：\BibleQuote{你成了我的助佑和保护者，为拯救我；}\BibleRef{出 15:2}达味说：\BibleQuote{求你成为我救恩的天主，作我的避难所，为拯救我；}依撒意亚说：\BibleQuote{祂成了每座卑微之城的助佑。}\BibleRef{依 25:4}当然，圣人们并非对天主说：\Quote{你被创造了，}而是说：\Quote{藉着你的恩宠，你成了我们的保护者和助佑。}\par

\Emph{为处理基于\BibleBook{若望}{福音}某段经文的一项异议，圣盎博罗削首先指出，亚略派的解释助长了摩尼教徒的论调；接着，在阐明这同一段经文的不同断句方式后，他清楚地表明，若按亚略派赋予的含义理解这段经文，不可能不侮辱圣父的天主性，并就此阐发其真正含义。}\par

41. 因此，我们无需害怕亚略派以其鲁莽的解经方式所构建的论点，他们试图证明天主圣言是\Quote{受造的，}因他们说，经上记载：\BibleQuote{在祂内受造的是生命。}\par

42. 首先，让他们明白，若他们将\BibleQuote{那受造的}这句话指向天主性，他们便陷入了摩尼教徒所提出的难题，因为这些人这样论证：\Quote{若在祂内受造的是生命，那么，在祂内还有未曾受造的，就是死亡，}因而他们可不虔敬地引入\OriginalTerm{两个本原}{two principles}。但教会谴斥这种教训。\par

43. 再者，亚略派怎能证明圣史确实说了这话？大多数精通信德的人所读的段落如下：\BibleQuote{万物是藉着祂造成的；凡受造的，没有一样不是藉着祂造成的。}另一些人则这样读：\BibleQuote{万物是藉着祂造成的；凡受造的，没有一样不是藉着祂。}然后他们接着说：\Quote{那受造的}，并把它与\Quote{在祂内}连起来；也就是说，\BibleQuote{但凡所有，都是在祂内受造的。}但\Quote{在祂内}这话是什么意思？宗徒告诉我们，他说：\BibleQuote{我们生活、行动、存在，都在祂内。}\BibleRef{宗 17:28}\par

44. 然而，任凭他们随意解读这段经文，他们无法削减天主圣言的尊威，若他们将\Quote{那受造的}这话指向祂的位格，作为主语，就无法不同时侮辱天主圣父，关于圣父经上记载说：\BibleQuote{但履行真理的，却来就光明，为显示出他的行为是在天主内完成的。}\BibleRef{若 3:21}看，这里我们读到人的行为在天主内完成，但我们却不能因此把天主性当作这些行为的主体。我们必须或者承认这些行为是藉着祂完成的，正如宗徒的声明所表明的：\Quote{万物都是藉着祂，在祂内受造的，祂在万有之先，万有都\Emph{一同存在}在祂内，}或者，如所引经文的证言教导我们的，我们应把那获得永生之果的诸德，视为在天主内完成的——贞洁、虔敬、热忱、信德，以及类似的德行，天主的旨意正是藉此表达出来。\par

45. 那么，正如这些行为是天主圣父的旨意与大能的表达，它们也是基督的，正如我们读到：\BibleQuote{在基督内受造，为行善工；}\BibleRef{弗 2:10}圣咏中又说：\BibleQuote{愿和平在你大能中成就；}又说：\BibleQuote{你以智慧造成了它们。}\Quote{你以智慧造成，}请注意——不是说\Quote{你造成了智慧；}因为既然万物都是在智慧中造成的，而基督是天主的智慧，那么这智慧显然不是依附性的，而是实体，且是永恒的实体；但如果智慧是受造的，那么它便是在比万物更恶劣的状况下受造，因为它自己无法成为智慧。因此，如果说\Quote{受造}常常涉及某种依附性事物，而非事物的本质，那么创造也可以指向某种预定的目的。\par

\Emph{\Person{撒罗满}的话，\BibleQuote{上主造了我}等，意指基督的降生成人是为了救赎圣父的创造，正如圣子自己的话所表明的。祂是\BibleQuote{开端}，可从祂德行的明显证据来理解，并显示主如何开启了诸德之路，成为它们的真正开端。}\par

46. 由此我们便可明白，关于降生成人的预言，\BibleQuote{上主造了我，作祂工程的元始，为行祂的化工，}\BibleRef{箴 8:22}意思是，主耶稣由童贞女受造，为的是救赎圣父的工程。确实，我们不能怀疑，这是指降生成人的奥迹，因为主取了我们的肉身，为将祂手中的工程从败坏的奴役中拯救出来，好能藉着祂自己肉身的苦难，推翻那握有死亡权势者。因为基督的肉身是为了受造物，但祂的天主性在受造物之前就已存在，因为祂在万有之先，万有都一同存在在祂内。\BibleRef{哥 1:16}\par

47. 因此，祂的天主性并非由于创造，而是创造因天主性而存在；正如宗徒所指明的，说万物因天主子而存在，因为我们读到：\BibleQuote{原来那为万物终向和万物根源的，既以使众子进入光荣，又使救他们的首领因受苦难而达到完全，是适当的。} \BibleRef{希 2:10} 难道他没有明确宣告，那因天主性而为万物的创造者的天主子，后来为了祂子民的救恩，甘心取了肉躯并遭受死亡吗？\par

48. 主为了做哪样的事工而从童贞女\Quote{受造}，祂在治愈那个天生的盲人时亲自表明了，说：\BibleQuote{我必须作那派遣我来者的工。} 而且在同一段圣经中，祂为使我们认为祂是在讲论降生成人的奥秘，说：\BibleQuote{当我在世界上的时候，我是世界的光。}\BibleRef{若 9:5} 因为就祂是人而言，祂暂时在这个世界上，但作为天主，祂永远存在。在另一处，祂也说：\BibleQuote{看，我同你们天天在一起，直到今世的终结。}\BibleRef{玛 28:20}\par

49. 对于\Quote{元始}一词，也没有任何置疑的余地，因为在祂在世生活时，当有人问祂\Quote{祢是谁？}祂回答说：\BibleQuote{我就是元始，正如我给你们所说的。}这不仅是指永恒天主性的本质特性，也指各种美德的可见证据，因为由此祂证明了祂是永恒的天主：祂是万物的元始，是各项美德的创立者，祂是教会的头，正如经上所载：\BibleQuote{因为祂是身体的头，就是教会的头；祂是元始，是死者中的首生者。} \BibleRef{哥 1:18}\BibleRef{弗 4:15}\par

50. 因此，很明显，\Quote{祂行径的元始}这句话，看来我们必须将其指为祂穿上肉体的奥秘，这是对降生成人的预言。因为基督降生成人的目的，是为我们铺平通往天堂的道路。注意祂如何说：\BibleQuote{我升到我的父和你们的父那里去，升到我的天主和你们的天主那里去。}\BibleRef{若 20:17} 然后，为使你们知道，全能圣父在降生成人之后，将祂的行径指派给圣子，在匝加利亚书中，天使对身穿污秽衣服的耶叔亚所说的话：\BibleQuote{万军的上主这样说：如果你遵行我的道路，谨守我的法令。}\BibleRef{匝 3:7} 那污秽衣服除了指穿上肉躯之外，还有什么意思呢？\par

51. 我们可以说，主的行径就是由基督引导的善生所采取的某些途径，祂说：\BibleQuote{我是道路、真理、生命。}\BibleRef{若 14:6} 因此，这道路就是天主的超越大能，因为基督是我们的道路，祂也是一条美善的道路，一条为信徒开启了天国的道路。此外，主的行径是正直的，如经上所载：\BibleQuote{上主，求祢使我认识祢的道路。} 贞洁是一条道路，信德是一条道路，节德是一条道路。确实，有美德的道路，也有邪恶的道路；因为经上记载：\BibleQuote{求祢察看在我内是否有邪恶的道路。}\par

52. 因此，基督是我们各种德行的元始。祂是洁德的元始，教导童贞女不要寻求人的拥抱，而要献出她们身体和心灵的全部纯真，服事圣神，而非丈夫。基督是节俭的元始，因为祂本是富有的，却成了贫穷的。\BibleRef{格后 8:9} 基督是忍耐的元始，因为祂受辱骂，不还口；受打击，不报复。基督是谦卑的元始，因为祂虽具有与天主圣父同等的权能尊威，却取了奴仆的形体。各项美德都源于祂。\par

53. 因此，为了使我们学习这些不同的美德，\Quote{有一个儿子赐给了我们，祂的元始在祂肩上。}那\Quote{元始}就是主的十字架——刚勇的元始，由此为圣洁的殉道者开辟了一条踏入神圣战争苦难的道路。\par

\Emph{刚才引用的依撒意亚经文中所含关于基督天主性和人性的预言，得到阐述，并展示其反驳各种异端的力量。}\par

54. 这个元始，依撒意亚看见了，因此他说：\BibleQuote{有一个婴孩为我们诞生了，有一个儿子赐给了我们。} 贤士也如此，当他们看见马厩中的小婴孩时，便朝拜了祂，说：\Quote{有一个婴孩诞生了；}当他们看见那颗星时，便宣示：\Quote{有一个儿子赐给了我们。}一方面，是来自地上的恩赐；另一方面，是来自天上的恩赐——而二者同是一位格，就每方面而言都是完美的，在天主性方面毫无变异，在人性方面也毫无缺损。贤士们朝拜了一位，向同一位奉献了礼物，为显示那在马槽中所看见的正是天上的主。\par

55. 请留意两个动词的含义如何不同：\Quote{有一个婴孩诞生了，有一个儿子赐给了。}祂虽由圣父所生，却非为我们而生，而是赐给我们的，因为圣子并非为了我们，而是我们为了圣子。因为事实上，祂并非为我们而生，祂在我们之前已存在，且是受造万物的创造者：祂也并非此时才首次获得生命，祂常存在，且在元始就存在；\BibleRef{若 1:1} 另一方面，那先前不存在的，却为我们诞生了。我们又看到，天使向牧人们说话时，怎样记载祂的诞生：\BibleQuote{今天在达味城中，为你们诞生了一位救主，即主基督。}\BibleRef{路 2:11} 因此，为我们诞生的，是那先前不存在的——即童贞女的孩子，由玛利亚而来的身体——因为这身体是在人受造之后才形成的，而天主性则在我们之前已存在。\par

有些抄本读作：\BibleQuote{有一个婴孩为我们诞生了，有一个儿子赐给了我们；}这就是说，那位身为天主子者，为了我们，由玛利亚所生，并赐给了我们。至于祂被\Quote{赐给}的事实，请听先知的话：\BibleQuote{求祢赐给我们祢的救恩。}但那在我们之上的，是赐下的：那从天上来的，是赐下的：正如我们关于圣神所读到的：\BibleQuote{天主的爱，借着所赐与我们的圣神，已倾注在我们心中了。}\BibleRef{罗 5:5}\par

但请注意，这段经文对众多异端，如同水泼在火上。\BibleQuote{有一个婴孩为\Emph{我们}诞生了，}不是为犹太人；\BibleQuote{为\Emph{我们}，}不是为摩尼教徒；\BibleQuote{为\Emph{我们}，}不是为马西昂派。先知说\BibleQuote{为\Emph{我们}，}就是指为那些相信的人，不是为不信者。祂确实，由于祂的怜悯，为所有人诞生了；但正是异端者的不忠，使得那位为所有人诞生者的诞生，不能使所有人受益。因为太阳受命升起，光照善人也照恶人，但那些看不见的人，却感觉不到日出的显现。\par

因此，正如这婴孩并非为所有人诞生，而是为信者诞生；同样，子也是赐给信者，而非不信者。祂是赐给我们的，不是赐给\OriginalTerm{福提诺派}{Photinians}；因为他们断言，天主子并未赐给我们，而是诞生于我们中间，且在我们中间才开始存在。祂是赐给我们的，不是赐给\OriginalTerm{撒伯里乌派}{Sabellians}，他们不愿听闻有子被赐给，坚持圣父与圣子是同一个且相同的。祂是赐给我们的，不是赐给\OriginalTerm{亚略派}{Arians}，在他们看来，子并非为救恩而被赐给，而是作为臣属与低等者被派遣；再者，祂为他们并非\Quote{\OriginalTerm{谋士}{Counsellor}}，因为他们认为祂对未来的事一无所知，也不是子，因为他们不信祂的永恒性，尽管论到天主的\OriginalTerm{圣言}{Word}，经上记载：\BibleQuote{那从起初就有的；}又说：\BibleQuote{在起初已有圣言。}\BibleRef{若 1:1} 回到我们当前要讨论的经文。圣经说：\BibleQuote{在起初，在祂创造大地之前，在祂创造深渊之前，在祂使水泉涌出之前，在一切山丘之前，祂生了我。}\par

\Emph{对前述引自\Person{撒罗满}\WorkTitle{箴言}的经文，将于下文更详细地解释。}\par

\Emph{或许}你会问，我如何引用\BibleQuote{上主造了我}这处经文，作为指向基督的\OriginalTerm{降生成人}{Incarnation}，因为宇宙的创造发生在基督降生成人之前？但请思考，圣经的惯常用法是把将要发生的事说成似乎已经过去，并暗示在基督内，\OriginalTerm{天主性}{Godhead}与\OriginalTerm{人性}{Manhood}的两性结合，以免有人否认祂的天主性或祂的人性。\par

例如，在\WorkTitle{依撒意亚}中，你可以读到：\BibleQuote{有一个婴孩为我们诞生了，有一个儿子赐给了我们；}同样，在这里（在《箴言》中），先知首先提出了血肉的\OriginalTerm{受造}{creation}，并将天主性的宣告与之结合，为叫你明白基督并非两位，而是一位，祂既在世界之前由圣父\OriginalTerm{受生}{begotten}，又在末时由童贞女\OriginalTerm{受造}{created}。这样，其含义便是：我，那位在世界之前受生者，就是那由凡女受造，为特定目的而受造的那一位。\par

再者，在宣告\BibleQuote{上主造了我}之前，祂说：\BibleQuote{我要讲述自永恒就有的事，}而在说\BibleQuote{祂生了}之前，祂又首先说：\BibleQuote{在起初，在祂创造大地之前，在一切山丘之前。}就其范围而言，\Quote{在……之前}这个介词一直回溯到无限无量的过去；因此，\BibleQuote{在亚巴郎之前，我就有，}\BibleRef{若 8:58} 显然不一定意味着\Quote{在亚当之后}，正如\BibleQuote{在晨星之前}不一定意味着\Quote{在天使之后}。但当祂说\Quote{在……之前}时，祂的意思不是祂被包含在任何一者的存在中，而是万物都被包含在祂的存在中，因为这是圣经显示天主永恒性的惯常方式。最后，在另一处经文中你可以读到：\BibleQuote{群山尚未形成，大地和世界尚未造成，从永远直到永远，祢就是天主。}\par

因此，子在一切受造物之前\OriginalTerm{受生}{begotten}；祂在万物内，并为万物的益处而\OriginalTerm{受造}{made}；由圣父受生，超越法律，\BibleRef{谷 2:28} 由玛利亚而出，属于法律之下。\BibleRef{迦 4:4}\par

\Emph{关于\Person{洗者若翰}的话语（若1:30）的观察，这段话语可能涉及天主的预定，但无论如何，根据前述思考，必须理解为指\OriginalTerm{降生成人}{Incarnation}。以\WorkTitle{卢德传}的历史为参照，神秘地阐释基督的优先性。}\par

但〔他们说〕经上记载：\BibleQuote{在我以后来的一个人，成了在我以前的，因祂本来在我以前；}\BibleRef{若 1:30} 于是他们辩论道：\Quote{看，那在先前的，却是\InnerQuote{受造}了。}让我们单独来看这些话。\BibleQuote{在我以后来的一个人。}那么，来的那一位是个人，而这个人就是\Quote{受造}的那一位。但\Quote{\OriginalTerm{人}{man}}这个词含有性别之意，而性别归属于人性，绝非归于\OriginalTerm{天主性}{Godhead}。\par

我可以论证说：那一位人（基督耶稣）在预先存在中，乃至于祂的肉身是预先被\OriginalTerm{预知}{foreknown}的，尽管祂的权能来自永恒——因为教会和\OriginalTerm{圣人们}{Saints}在世界开始之前就已被\OriginalTerm{预定}{foreordained}。但在这里，我暂搁这一论证，而强调：那受造并非涉及天主性，而是涉及\OriginalTerm{降生成人}{Incarnation}的本性，正如若翰自己所说的：\BibleQuote{这就是我所说的那一位：在我以后来的一个人，成了在我以前的。}\par

65. 那么，正如我在上文所指出的，圣经揭示了\OriginalTerm{基督}{Christ}内的双重性，为使你们理解祂同时具有\OriginalTerm{天主性}{Godhead}与\OriginalTerm{人性}{Manhood}，此处乃从肉身开始；因为圣经的习惯是不按固定规则，有时从基督的天主性开始，然后下降到\OriginalTerm{降生成人}{Incarnation}的可见标记；有时则相反，从祂的谦卑开始，上升到天主性的光辉，正如在众先知、众福音作者和圣\Person{保禄}那里屡见不鲜。在此，遵循此惯例，作者先讲论我们主的降生成人，然后宣告祂的天主性，不是为了混淆，而是为了区分人性与天主性。然而，\OriginalTerm{亚略派}{Arians}却像犹太酒贩一样，把水掺进酒里，混淆了天主性的诞生与人性的诞生，并把原本只针对肉身的卑微所说的话归于天主的威严。\par

66. 对于他们可能提出的某种异议我毫不担心，即在所引的话中有\Quote{一个\Emph{人}}——因为有些版本作\Quote{那在我以后来的}。但即使在此，也让他们注意前文。经上说：\BibleQuote{圣言成了\Emph{血肉}}。\BibleRef{若 1:14} 福音作者既说圣言成了\Emph{血肉}，便没有另提\Emph{人}。我们在提到\Quote{血肉}时理解为\Quote{人}，在提到\Quote{人}时理解为\Quote{血肉}。那么，在说了\Quote{圣言成了\Emph{血肉}}之后，此处就无需特别提及\Quote{人}，因为用\Quote{血肉}之名已经意指人了。\par

67. 稍后，圣\Person{若望}以\Quote{除免世罪的}羔羊为例；为清楚地教导你们关于他先前所谈论的那一位的降生成人，他说：\BibleQuote{这就是我曾说过的那一位：在我以后来的一个人，却在我之前被造成}，意即，我曾说祂是作为人、而非作为天主而\Quote{被造成}的。然而，为表明正是那位在万世之前、而非其他哪位成了血肉，以免我们设想有两个天主子，他补充说：\BibleQuote{因为祂先我而有}。假如\Quote{被造成}一词指的是天主性的诞生，那么作者又何必加上这句画蛇添足呢？不过，他先针对降生成人说道：\BibleQuote{在我以后来的一个人，却在我之前被造成}，然后补充说：\BibleQuote{因为祂先我而有}，因为有必要教导[基督的]天主性的永恒性；这就是为何圣\Person{若望}承认基督的优先性的原因，好使那位作为自己圣父永恒大能者，正因此缘故而被恰当地尊崇。\par

68. 然而，属灵领悟力的丰富活动使发起攻击并将亚略派逼入死角成为一项愉快的操练，因为他们试图将这段经文中的\Quote{被造成}一词理解成指[基督的]天主性，而非人性。事实上，当施洗者已宣告\BibleQuote{在我以后来的那一位，是在我以前被造成的}，也就是说，那位虽然在世间生涯中后我而来，却位居我之价值与恩宠之上，并有资格被尊为天主——他们还有什么立场可站呢？因为\Quote{后我而来}这一用语属于时间中的事件，而\Quote{他先我而有}则标示基督的永恒性；\Quote{在我以前被造成}则指祂的卓越，因为事实上，降生成人的奥迹超越了人的功德。\par

69. 再者，圣\Person{若翰}施洗者也以较不沉重的言辞教导了他所结合的理念，他说：\BibleQuote{在我以后来的一个人，我就是给祂提鞋也不配}，至少阐明了[基督]更为卓越的尊威，尽管未及祂天主性诞生的永恒性。这些话是如此全然指向降生成人，以致圣经在一部更早的书卷中，为我们提供了那神秘之鞋的一个人类预像。因为，按法律，人若死了，他与妻子的婚姻关系便转归其兄弟或另一位近亲，好使兄弟或近亲的后裔能延续家族的生命；\Person{卢德}便是如此：她虽是外邦出生，却曾拥有一个犹太民族的丈夫，这丈夫留下了一位关系很近的亲族；她在拾麦穗并以此养活自己和婆婆时被\Person{波阿次}看见并喜爱，然而\Person{波阿次}并未立即娶她为妻，直等到她先解下[那男子脚上的]鞋——按法律，她本应成为那人的妻子。\par

70. 这故事表面简单，却蕴含着深层的隐义，因为所行之事乃是更深奥事实的外在标记。的确，如果我们强解字面意思以就文句，几乎会发现这些话引发某种羞耻和恐惧，竟会使我们以为它们意在表达和传递平常肉身交合的思想；但这却是一种预像，预表那从犹太中将兴起的一位——基督按肉身即出于此——祂将以属天教导的种子，复活祂已死亲属的种子，即是说，那子民，并且法律的规条，在其属灵意义上，将婚姻的鞋归于祂，为了\OriginalTerm{教会}{Church}的订婚。\par

71. \Person{梅瑟}不是新郎，因为有话对他说：\BibleQuote{将你脚上的鞋脱下}，\BibleRef{出 3:5} 好让他给主腾出位置。\Person{农}的儿子\Person{若苏厄}也不是新郎，因为同样有话对他说：\BibleQuote{将你脚上的鞋脱下}，\BibleRef{苏 5:16} 以免因名字的相似，他竟被当成教会的配偶。除基督外别无新郎，圣\Person{若望}论到祂说：\BibleQuote{拥有新娘的是新郎}。\BibleRef{若 3:29} 因此，他们脱下鞋，但祂的鞋却不能被解开，正如圣若望所说：\BibleQuote{我连解祂鞋带也不配}。\BibleRef{若 1:27}\par

72. 那么，唯独基督是那位新郎，教会——祂的新娘——从列国来到祂面前，在婚姻中委身于祂；她从前贫穷饥饿，如今却因基督的收获而富足；在她心灵隐秘的深处，收集着丰富的庄稼和\OriginalTerm{圣言}{Word}的遗穗，以便用新鲜的食粮滋养那因丧子而疲倦、饥馑的人，就是那已死之民的母亲——在寻求新儿女的同时，她并不撇下那寡妇和穷乏者。\par

73. 那么，唯独基督是那新郎，祂甚至不吝啬把祂收获的禾捆赐给会堂。但愿会堂没有自行闭门不纳！她本有禾捆，自己可以收集。然而，因她的子民已死，她就像一个因丧子而孤苦的人，开始借着教会之手收集禾捆，好赖以维生——那欢欢喜喜来的人，必将这些禾捆带回来，正如经上记着说：\BibleQuote{他们必欢欢喜喜地来，带着自己的禾捆。}\par

74. 确实，除了基督，还有谁敢宣称教会是祂的新娘呢？惟有祂，而不是其他任何人，从\Place{黎巴嫩}呼叫她说：\BibleQuote{你從黎巴嫩來吧，我的新娘！你從黎巴嫩來吧！}\BibleRef{歌 4:8} 除了祂，教会还能论到谁而说：\BibleQuote{他的喉咙甘甜，他全然可爱。}\BibleRef{歌 5:26} 何况我们是从讨论祂脚上的鞋而开始这谈论的——除了降生成人的\OriginalTerm{天主圣言}{Word of God}，这些话还能适用于谁呢？\BibleQuote{他的双腿像一对大理石柱，置于纯金座上。}\BibleRef{歌 5:15} 因为惟独基督行走在众人的灵魂里，并在祂圣徒的心中开辟道路；在这些心中，一如在精金座基和宝石根基上，天上的\OriginalTerm{圣言}{Word}留下了不能磨灭的足迹。\par

75. 那么，我们清楚地看到，无论是那个人物还是预表，都指向\OriginalTerm{降生成人}{Incarnation}的奥迹。\par

圣安博回到主要问题上，并指出每当论及基督时说到祂\Quote{被造}（或\Quote{成为}），这必须理解为有关祂的\OriginalTerm{降生成人}{Incarnation}，或某些局限。在这个意义上，他阐释了圣经——尤其是\Person{圣保禄}著作——的几段经文。基督的永恒司祭职，预表在\Person{默基瑟德}身上。基督不仅拥有与圣父的相似，更拥有\Emph{同体}。\par

76. 因此，当论及基督时，如果说祂\Quote{被造}，或说祂\Quote{成为}，这一说法并非指祂的\OriginalTerm{天主性}{Godhead}的本体，而往往是指\OriginalTerm{降生成人}{Incarnation}——有时甚至是指某一特定职分；因为，如果将它理解成指祂的天主性，那么天主就成了受侮辱和嘲笑的对象，正如经上记着说：\BibleQuote{你却弃绝了你的基督，把祂归于虚无，你使祂流浪；} 又说：\BibleQuote{祂成了邻人的笑柄。} 请留意，是祂的邻人——而不是祂的家人，也不是那些依附着祂的人，因为\BibleQuote{那与主结合的，便是与祂成为一神。}\BibleRef{格前 6:17} 邻人并不依附于祂。再者，\Quote{祂成了笑柄，} 因为主的十字架为犹太人是绊脚石，为希腊人是愚妄：\BibleRef{格前 1:23} 然而对智慧人而言，正是借着这同一个十字架，祂便被\Emph{造成}高于诸天，高于天使，并且成为更美盟约的中保，正如祂曾是前一个盟约的中保一样。\par

77. 请注意我如何重复这个说法；我绝非试图避免它。然而，请留意祂是在什么意义上被称为\Quote{被造}的。\par

78. 首先，\BibleQuote{他既完成了净化之事，就坐在高天之上尊威者的右边，被立为远超过众天使。}\BibleRef{希 1:3} 哪里有了净化，哪里就有牺牲；哪里有牺牲，哪里也就有身体；哪里有身体，哪里就有献品；哪里有献品的职司，哪里也就有受苦的祭献。\par

79. 其次，祂是更美盟约的中保。然而，哪里有了遗嘱，立遗嘱者的死必须先发生，正如稍后一点所记的。不过，这死并不是祂永恒\OriginalTerm{天主性}{Godhead}的死，而是祂软弱之人躯体的死。\par

80. 再者，我们得知祂是如何被造得\Quote{高于诸天}的。圣经说：\BibleQuote{他是圣善的、无玷的、别于罪人的、高于诸天的；他无需像那些大司祭一样，每日必须先为自己的罪，然后为人民的罪祭献牺牲；因为他奉献了自己，只一次而永远完成了这事。}\BibleRef{希 7:26} 人若非在某一方面曾居下位，便不能说被高举；因此，基督因坐在圣父的右边，便在那方面被高举，即祂曾成为比天使微小一点，为要受苦而奉献了自己。\par

81. 最后，宗徒亲自对斐理伯人说：\BibleQuote{他取了人的形像，被人见有人的样式，就自己卑微，听命至死。}\BibleRef{斐 2:7} 请注意，在宗徒所说祂\Quote{被造}的那一方面，祂是被造为人的样式，而不是指天主性方面的主权，并且祂听命至死，以显出人所应有的服从，并获享按理属于天主性的王国。\par

82. 我们还需要援引多少经文，才能证实祂的\Quote{被造}必须被理解为指向祂的\OriginalTerm{降生成人}{Incarnation}，或某项特殊经纶呢？凡被造的，也同样是受造的，因为\BibleQuote{他一发言，它们便造成；他一出命，它们便受造。} 又说\BibleQuote{上主创造了我。} 这些话是就祂的人性而说的；我们也在第一卷书中指出，\Quote{受造}一词似乎是与降生成人有关。\par

83. 再者，宗徒亲自宣告，不应向受造之物敬拜，从而表明圣子并不是被造的，而是由天主所生的。\BibleRef{罗 1:25} 同时，他在别处也指出，在基督内有什么是被造的，以便阐明他是在什么意义上读了\WorkTitle{撒罗满书}中的话：\BibleQuote{上主创造了我。}\par

84. 现在，让我们按顺序回顾一整段经文。\BibleQuote{因此，孩子既都有同样的血肉，祂也照样取了一样的血肉，为能藉著死亡，毁灭那握有死亡的权势者。} \BibleRef{希2:14} 那么，那位愿意让我们分享祂自己血肉的，是谁呢？当然是天主子。除了藉着血肉，祂如何与我们成为一体？除了藉着肉身的死亡，祂又如何摧毁死亡的锁链？因为基督承受了死亡，便使死亡本身被置于死地。\BibleRef{格前15:54} 所以这段经文是在谈论\OriginalTerm{降生成人}{Incarnation}。\par

85. 让我们看下文：\BibleQuote{祂确实不曾取了天使的本性，而是取了亚巴郎后裔的本性。因此，祂应当在各方面相似弟兄们，好能在关于天主的事上，成为一个仁慈而忠信的首领，向上主奉献的司祭，以补赎人民的罪过；因为祂自己既受过试探，也就能援助受试探的人。为此，有分于天上召叫的诸位圣弟兄啊！你们应细想我们宣认的宗徒和大司祭——\Person{耶稣}，祂对那创造祂者是忠信的，正如\Person{梅瑟}在天主的全家一样。} 以上便是宗徒的话。\par

86. 你看到作者是在哪方面称祂为受造的了：\Quote{就在祂取了亚巴郎后裔这一点上}；这显然是在宣示一个身体的生成。的确，祂若非在身体内，如何能赎人民的罪呢？祂若非在身体内，又是在什么内受了苦呢？——正如我们上面所说：\Quote{基督在肉身内受了苦}？祂若非从司祭的族裔中取了什么，又怎能成为司祭呢？\par

87. 司祭的职责就是献上某物，并按法律藉着血进入圣所；既然天主拒绝了公牛和公山羊的血，这位大司祭确实必须藉着自己的血，经过并进入天上的至圣所，好能成为我们罪过的永远补赎。司祭与祭品，原本为一；然而司祭职和祭献，却是在人性的条件下进行的，因为祂曾被牵去如同待宰的羔羊，而祂是照\Person{默基瑟德}的品位而为司祭。\par

88. 因此，任何人看到一个人间的建制的品位时，都不可坚持其中蕴含着天主性；因为就连那位\Person{默基瑟德}——\Person{亚巴郎}曾藉他的职务献了祭——教会确实不认为是天使（像某些犹太妄谈者所主张的那样），而是一位圣善的人、天主的司祭，他作为我主的\Emph{预像}，被描述为\BibleQuote{无父、无母，无族谱，无时日之始，也无生命之终}，为要预先显示天主永生的圣子来到这个世界，祂同样降生成人，而后被生出，却没有肉身的父亲；作为天主受生，却没有母亲，也没有族谱，因为经上记载：\BibleQuote{谁将宣述祂的世代？} \BibleRef{依53:8}\par

89. 因此，我们接受这位\Person{默基瑟德}为按照基督的模子而造的天主司祭，但我们将前者视作预像，将后者视作真体。预像本是真理的阴影，我们承认前者的王权只在一座城的名下，而后者的王权则显于整个世界的和好；因为经上记载：\BibleQuote{天主在基督内使世界与自己和好} \BibleRef{格后5:19}——也就是说，[在基督内]有永恒的\OriginalTerm{天主性}{Godhead}；或者，如果父在子内，如同子在父内，那么他们在本性和运作上的合一显然是不可否认的。\par

90. 可是，我们的对手再怎么想否认，又怎能公正地否认呢？因为圣经说：\BibleQuote{住在我内的父，亲自作祂自己的事业}；又说：\BibleQuote{我所行的事业，正是祂自己行的}。 \BibleRef{若14:10} 这里不是说\Quote{祂\Emph{也}作事业}，而是应视为事业的相似而非合一；宗徒在说\Quote{我所行的事，正是祂自己行的}时，已清楚表明我们应当相信：父的事业与子的事业是同一个。\par

91. 另一方面，当祂要人明了事业的相似而非合一时，祂说：\BibleQuote{信我的人，我所行的事业，他也要行}。 \BibleRef{若14:12} 祂巧妙地在此插入了\Quote{也}这个词，准许我们有相似，却未赋予本性的合一。因此，父的事业与子的事业是同一个，无论亚略派是否乐意这样想。\par

\Emph{圣父与圣子的国度是唯一且不可分的，各自的天主性也是如此。}\par

92. 现在我要问：他们怎能设想圣父与圣子的国度为分开的呢？因为主曾说过，正如我们上文所示：\BibleQuote{凡一国自相纷争，必成废墟}。 \BibleRef{玛12:25}\par

93. 的确，正是为了驳斥亚略派仇恨的不敬教导，圣\Person{伯多禄}亲自断言圣父与圣子的主权是唯一的，他说：\BibleQuote{为此，弟兄们，你们更要尽心竭力，使你们的蒙召和被选坚定不移，因为你们这样做，就不会跌倒；这样，你们便有更丰富的恩宠，得以进入天主及我们的主救主耶稣基督的永远国度。} \BibleRef{伯前2:10}\par

94. 然而，如果有人认为此处只论及基督的主权，并且因此将这段经文理解为圣父与圣子的权柄是分开的——却仍要承认那是属于子的主权，且是永恒的，这样，不但引入了两个分离、因而可能崩溃的国度，而且，既然没有任何国度堪与天主的国度相比——他们无论怎样渴望，也无法否认那是属于子的国度——那么他们要么必须收回他们的主张，承认圣父与圣子的国度是唯一的；要么必须将管理一个较小国度的职责归于圣父——这是亵渎；要么必须承认那位他们恶毒地宣称在天主性上次等的（圣子），竟拥有同等的国度，这是自相矛盾的。\par

但此说[他们的教导]与其前提不符，不相合，不成立。那么，让他们承认，国是一，正如我们所承认并证明的，不是凭我们自己的证据，而是凭上天所赐的见证。\par

首先，从圣经的更多证言中了解，天国也是圣子的国：\\BibleQuote{我实在告诉你们：站在这里的人中，就有些人在未尝到死味以前，必要看见人子来到自己的国里。}\\BibleRef{玛 16:28}因此，毫无疑问，国是属于天主子的。\par

如今，要知道，圣子的国就是圣父的国：\\BibleQuote{我实在告诉你们：站在这里的人中，就有几个在未尝到死味以前，必要看见天主的国带着威能降来。}\\BibleRef{谷 8:39}这国确实是一，乃至赏报是一，继承者是同一的，功绩也是一，应许者也是如此。\par

这国如何能不是一呢？尤其是当圣子亲口论到自己说：\\BibleQuote{那时，义人要在他们父的国里，发光如同太阳}？\\BibleRef{玛 13:43}因为凡属于父的，凭其与父的尊威相称，也属于子，因为在同一光荣中合一。\\BibleRef{若 17:5}所以，圣经宣告这国既是圣父的国，也是圣子的国。\par

如今，要知道，当天主的国被称呼时，父或子的权柄并未搁置，因为父的国和子的国都被涵盖在\\Emph{天主}这名下，说：\\BibleQuote{几时你们看见亚巴郎、依撒格、雅各伯及众先知在天主的国里}\\BibleRef{路 13:28}。我们岂能否认众先知在圣子的国里，当临终的强盗说：\\BibleQuote{当你来为王时，请你纪念我！}主却回答：\\BibleQuote{我实在告诉你：今天你就要与我一同在乐园里}吗？\\BibleRef{路 23:42}在天主的国里，若非逃脱了永死，还能有何理解？但逃脱了永死的人，看见人子来到自己的国里。\par

那么，他所赐予的，他如何不能掌管呢？因为他说：\\BibleQuote{我要将天国的钥匙交给你}？\\BibleRef{玛 16:19}且看两者间的鸿沟。仆人开，主赐与；一位靠自己，另一位靠基督；仆人领受钥匙，主授予权能：一是赐予者之权，一是管家的职分。\par

再看一个证据，父与子的国，其治理是一。在\\BibleBook{弟茂德}{前书}中写道：\\BibleQuote{保禄，奉我们救主天主和作我们希望的基督耶稣的命，作基督耶稣的宗徒。}\\BibleRef{弟前 1:1}因此，父与子的国被明确宣告为一，正如宗徒保禄也曾断言说：\\BibleQuote{因为你们应该清楚知道：不论是犯邪淫的，行不洁的，或是贪婪的——即崇拜偶像的，在基督和天主的国内，都不得承受产业。}\\BibleRef{弗 5:5}因此，国是一，\\OriginalTerm{天主性}{Godhead}是一。\par

法律证明了\\OriginalTerm{天主性}{Godhead}的合一，它说及唯一的天主，\\BibleRef{申 6:4}，宗徒也说及基督：\\BibleQuote{因为在基督内，真实地住有整个圆满的天主性。}因为，倘若如宗徒所言，整个圆满的天主性真实地住在基督内，那么父与子必须被宣认具有同一的天主性；或者，若有人想要将子的天主性与父的天主性分开，而子却拥有整个圆满的天主性，那么，既然圆满的完美之外已无剩余，还有什么可以被认为保留下来呢？所以，天主性是一。\par

\\Emph{圣子的尊威是本有的，与父同等，且天使非分享者，而是瞻望者。}\par

我们已经确定父与子具有同一形象和模样，现在还需证明他们也具有同一尊威。我们不必远求证据，因为圣子亲口论到自己说：\\BibleQuote{当人子在自己的尊威中，与众天使一同降来时，他要坐在自己尊威的宝座上。}\\BibleRef{玛 25:31}看，圣子的尊威被宣告了！他那非受造的尊威无可否认，还缺少什么呢？因此，尊威属于圣子。\par

让我们的对手现在确认无疑：父与子的尊威是一，因为主亲口说过：\\BibleQuote{谁若以我和我的话为耻，将来人子在自己的尊威，和父及众圣天使的尊威中降来时，也要以这人为耻。}\\BibleRef{路 9:26}\\BibleQuote{及众圣天使的尊威}这话的力量，岂不是表明仆人们从崇拜其主中获得荣耀吗？\par

因此，圣子将自己的尊威归于父，也归于自己，并非让天使与父和子同等分享那尊威，而是让他们仰望那超凡的天主光荣；因为诚然，即使天使也没有本有的尊威，如同圣经论及子那样：\\BibleQuote{当他坐在自己尊威的宝座上时}，而是他们侍立面前，好能看见父和子的光荣，按照他们堪当或能够承受的视觉程度。\par

106. 此外，天主所赐的言语本身就表明了其含义，好让你明白圣父与圣子的荣耀并非天使所能共享。经上如此说：\BibleQuote{当人子在自己的光荣中，与众天使一同降来时。} 在另一处，为表明祂父的尊威与光荣，以及祂自己的尊威与光荣，原是一体，我们的主亲自说：\BibleQuote{人子将要在祂父的光荣中，同诸圣天使降来。}\BibleRef{谷 8:38} 天使奉命而来，祂则在光荣中降临：天使是祂的侍从，祂坐在宝座上；天使侍立，祂安坐——借用人类日常生活的用语，祂是审判者，天使是法庭的官员。请注意，祂并未先提圣父的天主尊威，而后才提及自己及天使的，以免看起来像是在制造一种从至高到较低本性的降序。祂先提及自己的尊威，然后才说到圣父及天使的尊威（因为圣父不可能显得低于天使）。这样，祂就不会因为将自己置于圣父与天使之间，而显得似乎构成了一种通过提升自身尊位而由天使升到圣父的升序；也不会相反地被认为展示了一种从圣父到天使的降序，导致尊位的减损。我们承认圣父与圣子有同一的\OriginalTerm{天主性}{Godhead}，当然不会像亚略派那样设想这种区别的次序。\par

\Emph{圣子与圣父同性同体。}\par

108. 现在，陛下，关于\OriginalTerm{性体}{substance}的问题，我何须告诉您，圣子与圣父同性同体呢？因为我们读过，圣子是圣父性体的肖像，好让您明白，就\OriginalTerm{天主性}{Godhead}而言，圣子与圣父并无分别。\par

109. 基于这种肖像，基督说：\BibleQuote{凡父所有，全是我的。}\BibleRef{若 16:15} 因此，我们不能否认天主有性体，因为祂并非没有性体——祂赋予了其他事物存在的基础，尽管这在受造物中与在天主内有所不同。圣子，万物藉着祂而恒存，不可能没有性体。\par

110. 因此，圣咏作者说：\BibleQuote{我的骨骸并不隐藏，它们是你暗中形成的，我的性体在地府深处。} 因为，对于祂的大能和天主性来说，创世以前所成就的一切，尽管其宏伟尚未彰显，也是无法隐藏的。所以，此处我们提到了\OriginalTerm{性体}{substance}。\par

111. 但或许有人会反对说，提及祂的性体是其降生成人的结果。我已表明，\OriginalTerm{性体}{substance}一词不止一次被使用，并且不是你所解释的继承财产之意。现在，若你愿意，让我们承认，按照奥秘的预言，基督的性体曾临在于地府——因为祂确实在阴间施展自己的大能，以那赋予自己肉身的灵魂，释放了死者的灵魂，解开死亡的束缚，赦免罪过。\BibleRef{伯前 3:19}\par

112. 并且，事实上，何障碍会使你无法将\Quote{性体}理解为祂的天主性体呢？既然天主无所不在，因此经上对祂说：\BibleQuote{我若上升于高天，你已在那里；我若下降于阴府，你也在那里。}\par

113. 而且，圣咏作者在接下来的话语中已表明，我们必须将这理解为指向天主性体，因为他如此说：\BibleQuote{你的眼睛早已看见我未受造的实体。} 因为圣子并非受造，也非天主的工程，而是出自永恒大能的圣言。祂称祂为\OriginalTerm{未受造的}{\GreekText{ἀχατέργαστον}}，意思是圣言既非受造之物，也非被创造之物，而是在没有任何受造物见证在场的情况下，由圣父所生。但除此之外，我们还有大量证据。即使我们承认这里所说的性体是肉身的性体，只要你自己不说天主子是由工程形成的某物，而是承认祂非受造的天主性。\par

114. 现在我知道，有些人断言，那奥秘的降生成人之形体是未受造的，因为其中没有任何经由与人交合而成的事，因我们的主是由童贞女所出。那么，如果许多人都根据这段经文断言，由玛利亚所生的那一位也非受造工程的产物，你这个作\Person{亚略}门徒的人，还敢认为天主的圣言也是那样的产物吗？\par

115. 然而，这是我们读到\Quote{性体}的唯一所在吗？难道在另一处不也这样说：\Quote{城门的门闩被毁，山岭倒塌，祂的性体显露出来}？怎么，难道这个词在这里也指受造物吗？我知道有些人习惯将性体解释为钱财产业。那么，如果你将此意赋予这词，山岭倒塌就是为了让人看到某人的金钱财物了。\par

116. 但让我们记得，倒塌的是\Emph{哪些}山，即有关它们经上说：\BibleQuote{假如你们有信德如芥子，你们将对这座山说：从这里移到那边去！它必会移位。}\BibleRef{玛 17:19} 那么，所谓山，指的是那些自高自大的事物。\BibleRef{格后 10:5}\par

117. 此外，在希腊文版本中，译法是：\Quote{宫殿倒塌了。} 若非撒殚的宫殿，还能是什么宫殿呢？论及这宫殿，主曾说过：\Quote{祂的国怎能站得住呢？} 因此，我们正读到的那些属于魔鬼的宫殿，就如同一座座高山；所以当这些宫殿从信徒心中倒塌时，真理便显明出来：天主子基督，出自圣父的永恒\OriginalTerm{性体}{substance}。再者，那些铜山又是什么？从中出来四辆马车。\BibleRef{匝 6:1}\par

我们看见那高耸，自高自大反对天主之智识的，被主的话打倒，当天主子说：\BibleQuote{别作声！从他身上出来！你这邪魔！} \BibleRef{谷 1:25} 关于他，先知也说：\Quote{看，我来对付你，你这败坏的山！}\par

那么，\Emph{那些}山已经倾倒了，并且显明在基督内有天主的实体，正如那些见过他的人说：\BibleQuote{你真是天主子！} \BibleRef{路 4:41} 因为他是以天主性的、而非人性的权能命令魔鬼。耶肋米亚也说：\Quote{在山上举哀，在旷野的路上捶胸，因为他们败落了，因为没有人了，他们没有听到实体的话；从飞鸟到家畜，他们都战栗，都败落了。}\par

我们也没有忽略，在另一处，为陈述人的状况之脆弱，为显示他承担了肉身的软弱和我们心灵的情感，主借他的先知的口说：\Quote{上主，求你记念我的实体是什么，} 因为那是天主子在人性的脆弱中说话。\par

论到他，圣经在上述章节中说，为揭示降生成人的奥迹：\Quote{但是你拒绝，上主，视同无物——你抛弃了你的基督。你推翻了你与你仆人订立的盟约，将他的圣洁践踏于地。} 圣经称他为\Quote{仆人}，指的是他的什么，岂不是他的肉身吗？——因为\Quote{他并没有以自己与天主同等，把持不舍，反而空虚自己，取了奴仆的形体，成为与人相似的，形状也完全与人一样。} \BibleRef{斐 2:6} 所以，正因他取了我们的本性，他是仆人，但凭他自己的权能，他是主。\par

此外，你们读到\Quote{谁曾站在上主的\Emph{真理}（\Emph{实体}）中了呢？}，以及\Quote{如果他们曾站在我的真理中，听了我的话，教训了我的百姓，我必早已使他们离开他们的邪念和恶行。}，这些话是什么意思呢？\par

\Emph{阿里乌斯派，由于他们声称子是\Quote{由另一实体}，明显承认天主内有实体。他们回避使用此术语的唯一原因是，正如\Person{尼科美底亚的优西比乌}所显明，他们不愿承认基督是真天主子。}\par

阿里乌斯派怎能否认天主的实体呢？他们怎能以为\Quote{实体}一词虽然出现在圣经许多章节中却应当被禁用，而他们自己却因说子是\OriginalTerm{异体}{\GreekText{ἑτεροούσιος}}，即另一实体，从而承认天主内有实体呢？\par

因此，他们所回避的，不是术语本身，而是其效力与后果，因为他们不愿承认天主子是真实的[天主]。虽然天主生育的过程无法用人类的语言表述，但教父们断定，他们的信德可以适当地借由使用这一术语，与\OriginalTerm{异体}{\GreekText{ἑτεροούσιος}}对抗而加以分辨，这是遵循先知的权威，他说：\BibleQuote{谁曾站在上主的\Emph{真理}（\Emph{实体}）中，看见了祂的圣言？} \BibleRef{耶 23:18} 所以阿里乌斯派，当他们使用\Quote{实体}一词以配合他们的亵渎时，便予以承认；相反，当这词依信友的虔诚热忱而被采用时，他们却加以拒绝和争辩。\par

除了不愿承认他是真天主子之外，还能有什么别的原因，使他们不愿意称子是与父\OriginalTerm{同性同体}{\GreekText{ὁμοούσιος}}呢？这在\Person{尼科美底亚的优西比乌}的信中流露出来。他写道：\Quote{如果我们说子是真天主、非受造，那么我们就是正在承认他与父\OriginalTerm{同性同体}{\GreekText{ὁμοούσιος}}。} 当这封信在\Place{尼西亚}大公会议上被宣读之后，教父们便将这个词写入了他们的信德说明中，因为他们看到这词使他们的对手胆怯；为要拿起对手已拔出的剑，斩下对手亵渎异端的头。 \BibleRef{撒上 17:51}\par

然而，他们以撒伯里乌派为借口来回避使用该词，这是徒劳的；借此他们暴露了自己的无知，因为一个存有是与另一个\OriginalTerm{同性同体}{\GreekText{ὁμοούσιος}}，而非与自身同性同体。故此，我们正确地称子与父\OriginalTerm{同性同体}{\GreekText{ὁμοούσιος}}，因为这个术语既表达了位格的区分，也表达了性体的合一。\par

他们能否认圣经中遇到\OriginalTerm{实体}{\GreekText{οὐσία}}一词吗？既然主曾说到饼，即\OriginalTerm{超性之粮}{\GreekText{ἐπιούσιος}}，而梅瑟写道\OriginalTerm{你们将成为我的特选子民}{\GreekText{ὑμεῖς} \GreekText{ἔσεσθέ} \GreekText{μοι} \GreekText{λαὸς} \GreekText{περιούσιος}} \BibleRef{出 19:6}？\OriginalTerm{实体}{\GreekText{οὐσία}}是什么意思，这个名称从何而来，难道不是来自\OriginalTerm{永存者}{\GreekText{οὖσα} \GreekText{ἀεί}}，即那永远存在的吗？因为那位\Quote{是}，且\Quote{永远长存}的，就是天主；故此，永恒的实体，永远常在的，被称为\OriginalTerm{实体}{\GreekText{οὐσία}}。饼是\OriginalTerm{超性之粮}{\GreekText{ἐπιούσιος}}，因为它从圣言的实体中汲取了持续力的实体，并将它供给心灵与灵魂，因为经上记载：\Quote{饼健壮人心}。\par

128. 因此，让我们谨守祖先的规诫，不要以粗鲁和鲁莽的妄为亵渎传承给我们的象征。我们听说过的那本封印的预言书，无论是长老、掌权者、天使、总领天使，都未曾敢展开；因为展开它的特权唯独保留给基督。\BibleRef{默 5:5} 我们中间有谁敢揭开那被宣信者们封印、并因众多见证而长久被尊为圣的司祭职之书？那些被迫揭开的人，后来却因尊重所遭受的欺骗，又重新封印了；那些不敢以亵渎之手触摸它的人，却挺身而出成为殉道者和宣信者。我们怎能否认我们宣扬其胜利之人所持守的信德？\par

\Emph{为了武装正统派以对抗亚略派的诡计，圣安博揭露了后者使用的一些欺诈性宣认，并通过各种论证表明，虽然他们有时称圣子为\Quote{天主}，但这还不够，除非他们也承认祂与圣父平等。}\par

129. 不要惧怕，不要战栗；那威吓的人反而让忠信者占了上风。欺诈之人的安抚膏油是有毒的——因此，当他们假装宣讲自己否认之事时，我们必须提防他们。从前那些轻信他们的人就这样被欺骗了，以至于当他们以为一切都是诚信时，却落入了背叛的罗网。\par

130. 他们说：\Quote{愿那说基督如同其他受造物一样，是受造物的人受诅咒。} 单纯的民众听到这话就相信了，因为经上记载：\BibleQuote{幼稚的人，有话必信。} \BibleRef{箴 14:15} 他们就这样听到并相信了，被那话的初闻所蒙蔽，好比鸟儿，因贪求信德的诱饵，没有注意到为他们张开的罗网，这样，在追求信德的过程中，却吞下了不虔欺诈的钩子。因此，主说：\BibleQuote{你们要机警如同蛇，纯朴如同鸽子。} \BibleRef{玛 10:16} 智慧被置于首位，是为了使纯朴不致受伤害。\par

131. 因为福音所指的蛇，正是那些脱去旧习，穿上新行径的人：\BibleQuote{脱去旧人和他的作为，且穿上了新人，这新人既是照造他者的肖像而更新。} \BibleRef{哥 3:9-10} 那么，让我们学习福音称为蛇的那些人的作为，蜕去旧人的皮囊，这样，我们就能如同蛇一样，知道如何保全生命并提防欺诈。\par

132. 只要说\Quote{愿那说基督是受造物的人受诅咒}就足够了。亚略派的人啊，为什么你要将毒药掺入你宣认中的良善，从而玷污整个宣认呢？因为加上\Quote{如同其他受造物一样}，你并不是在否定基督是受造物，而是在否定祂像其他受造物一样——因为你确实称祂为受造物，尽管你赋予祂超越一切受造物的尊位。再者，亚略，这不虔教义的始作俑者，曾说过圣子是完美的受造物，而非如同其他受造物一样。那么，你看，你如何继承了从你父辈传下来的言辞。否认基督是受造物就足够了：为什么还要加上\Quote{但不像其他受造物一样}？切除这坏疽的部分，以免传染扩散——这是有毒的，致命的。\par

133. 再者，你们有时说基督是天主。不，你们应当称祂为真天主，意思是，你们承认祂拥有圣父圆满的天主性——因为在天上或地上，都有被称为神的。那么，\Quote{天主}这个名称，不应仅作为一种称呼和提及的方式，而应理解为，你们确认圣子拥有与圣父相同的天主性，正如经上所说：\BibleQuote{就如父是生命之源，照样他也使子成为生命之源；} \BibleRef{若 5:26} 这就是说，祂是藉着生祂而赐给祂，如同对圣子，而不是出于恩宠，如同对缺乏者。\par

134. \BibleQuote{并且赐给他行审判的权柄，因为他是人子。} \BibleRef{若 5:27} 要特别注意这附加的话，免得你们趁机藉一句话宣讲虚妄。你们读到祂是人子；难道你们因此就否认祂接受[所赐的权柄]吗？那么，如果天主的一切专有属性没有都赐给圣子，你们就否认天主吧，因为祂说过：\BibleQuote{凡父所有的一切，都是我的；} \BibleRef{若 16:15} 为什么不承认，天主性的所有特性和属性都存在于圣子内呢？因为那位说\BibleQuote{凡父所有的一切，都是我的}的，祂排除什么为没有呢？\par

135. 为什么你们\Quote{如此执着地}、用这般恳切的言语，细数基督使死人复活、在水面上行走、医治人的疾病呢？这些能力，祂确实也赐给了祂的仆人们，让他们也能像祂一样彰显。当这些能力见于人身上时，更激起我的惊奇，因为天主赐给了他们如此大的能力。我愿意听到一些关于基督的事，是祂所独有且特有的，是受造物不能与祂共享的，因为祂是受生的，天主的独生子，真天主出自真天主，坐在圣父的右边。\par

136. 无论在哪里我读到圣父与圣子并坐，我总是发现圣子坐在右边。这是因为圣子高于圣父吗？不，我们不是这样说；而是那位因天主的爱而受尊荣的，却因人的不虔而受羞辱。圣父知道，关于圣子的疑惑必然会被散播，因此祂给了我们一个尊崇的榜样，好让我们效法，免得我们羞辱圣子。\par

\Emph{基于圣斯德望见到主站立的异象的反对意见被驳斥，并藉由同一位圣人向圣子所做的祈祷，显示出圣子与圣父的平等。}\par

137. 唯一一处记载，斯德望说他看见主耶稣\Emph{站立}在天主的右边。\BibleRef{宗 7:55}现在，你们要明白这些话语的深意，以免你们借此滋生疑问。你们或许会问：为什么我们在各处读到圣子坐在天主右边，唯独一处写祂站立？祂坐着，是作为审判生者死者的判官；祂站着，是作祂子民的辩护者。祂站着，是作为司铎，向祂的圣父献上一位优秀殉道者的祭献；祂站着，就像裁判，将要赏给一位善斗者那场大争战的奖品。\par

138. 你们也要领受圣神，使你们能明了这些事，正如斯德望领受了圣神；你们也能像那位殉道者一样说：\BibleQuote{看，我看见天开了，人子站在天主的右边。}\BibleRef{宗 7:55}那有天向他敞开的人，看见耶稣在天主的右边；那灵魂的眼目紧闭的人，看不见耶稣在天主的右边。因此，让我们宣认耶稣在天主的右边，好使天也为我们敞开。那些以别样宣认的人，便是为自己关闭了天堂之门。\par

139. 但如果有人提出反对，强调\Emph{圣子}是站着的，就让他们根据这段经文指出圣父是坐着的，因为尽管斯德望说人子站着，但他在这里并没有进一步说圣父坐着。\par

140. 然而，为更充分清楚地表明站立并不意味羞辱，反是主权，斯德望向圣子祈祷，愿自己更蒙圣父悦纳，说：\BibleQuote{主耶稣，接我的灵魂去吧！}\BibleRef{宗 7:58}再者，为显示圣父与圣子的主权是同等的，他又祈祷说：\BibleQuote{主，不要向他们算这罪债。}\BibleRef{宗 7:51}这些话，正是主在受难时以人子身份向圣父说的——也正是斯德望在自己的殉道时向天主圣子祈祷的话。当人向圣父和圣子祈求同一恩宠时，便确认了双方有同等的权能。\par

141. 否则，如果我们的对手坚持斯德望是向圣父说话，就让他们看看，按他们自己的说法，所肯定的是什么。我们确实不为他们的论证所动；然而，那些注重字面和次序的人，应当留意，\Emph{第一项}祈求是向圣子发出的。我们即使按他们对这经文的理解，也能从中证明圣父与圣子尊威的合一；因为，当祈祷同等地向圣子并圣父发出时，祈祷所赋予的同等地位便指向行动中的合一。但如果他们不承认圣子被以\Quote{主}的称谓来呼求，我们就看明，他们确实在试图否认祂是主。\par

142. 然而，既然这位伟大殉道者的荣冠已经展现，就让我们减弱争论的热切，结束今天的讲论。让我们歌颂这位圣洁的殉道者，正如在一场大争战后所常行的——这殉道者虽因敌人的打击而流血，却蒙基督赐予的荣冠为赏报。\par

\SourceRecord{Translated by H. de Romestin, E. de Romestin and H.T.F. Duckworth.；Nicene and Post-Nicene Fathers, Second Series；Vol. 10.；Edited by Philip Schaff and Henry Wace.；Buffalo, NY: Christian Literature Publishing Co.,；1896.；<http://www.newadvent.org/fathers/34043.htm>.}

% structure: p0001 source=h1 target=section reason=page-title

\section{\WorkTitle{基督信仰阐释}，卷四}

% structure: p0002 source=h3 target=omitted reason=duplicate-page-title

\Emph{惊奇之处不在于人未能认识基督，而在于他们未曾听从圣经的话语。基督确实未被认识，即使天使也须由启示才得认识，祂的前驱亦然。以下描述基督凯旋升天的荣耀，其卓越胜过某些先知的被提升。最后，借由对此场合与天使对话的阐释，证明圣子的全能，以驳斥亚略派。}\par

1. 陛下，思考人们为何如此歧误，或者说——唉！——为何众人竟对天主子信奉种种异见，令人惊奇之处似乎完全不在于人的知识在处理超性事物时遭受挫败，而在于它未曾服从圣经的权威。\par

2. 的确，若人们以其世俗智慧未能领悟天主圣父与主耶稣基督的奥秘——在祂内蕴藏着一切智慧与知识的宝藏\BibleRef{哥 2:3}，这奥秘甚至连天使也须藉启示才能得知——这又有何可惊奇的呢？\par

3. 因为谁能仅凭想象、而不靠信德，去追随主耶稣呢？祂一时从最高天降至地下的阴影，一时又从阴府升起，回到天上的居所；祂在瞬间\OriginalTerm{自我空虚}{self-emptied}，好能居于我们中间，却又丝毫不损祂原有的本体，圣子永远在圣父内，圣父也在圣子内？\par

4. 连基督的前驱，尽管仅是在代表会众的层面上，也对他产生了疑惑；这位前驱本是奉派行在主面前的人，最后竟派遣使者去问：\BibleQuote{祢就是要来的那一位，还是我们该等候别人？}\BibleRef{玛 11:3}\par

5. 天使们也因这属天的奥秘而惊愕不已。因此，当主复活时，连高天也承受不住祂从死者中复活的荣耀——祂的肉身不久前还被局限在狭窄的坟墓内——以致天上的万军也疑惑惊奇。\par

6. 因为一位征服者来了，身佩奇妙的战利品，上主在祂的圣殿中，天使和总领天使行在祂面前，惊叹于从死亡夺来的掠物；他们虽知从肉身并不能给天主增添什么，因为万有都低于天主，但目睹十字架的凯旋标志——其上\BibleQuote{王权担在祂肩上}——以及永恒征服者所携带的战利品，他们仿佛觉得，那曾从他们中间出去的祂，这些门扉无法容祂通过，尽管它们实在永不能容纳祂的伟大——于是他们在祂归来时，为祂寻求更宽阔更高大的通道——祂在\OriginalTerm{自我空虚}{self-emptied}之后，竟然丝毫未减损其尊位。\par

7. 然而，理当为新征服者预备一条新路——因为征服者本人在身形上似乎总比他人更高更大；但因\OriginalTerm{正义之门}{Gates of Righteousness}——即那敞开天堂的旧约与新约之门——是永恒的，它们并未改变，而是被举扬起来，因为进入的并非单独一人，而是整个世界，亦即全人类的救赎主。\par

8. \Person{哈诺客}曾被提升，\Person{厄里亚}被接去，但仆人不能大过主人。因为\BibleQuote{没有人升过天，除了那自天降下的人子；}\BibleRef{若 3:13}至于\Person{梅瑟}，虽然他的遗体在地上从未被人见过，我们也未在经上读到过他居于天上的光荣中，除非是在主以自己复活的保证，冲破了地狱的束缚、举扬了虔敬者灵魂之后。因此，哈诺客曾被提升，厄里亚被接去；两者都是仆人，都在肉身内，但他们并非在从死者中复活之后，也未曾带着死亡的战利品和十字架的凯旋行列，被天使看见。\par

9. 因此，当（天使们）远远望见万有之主、死亡的首位且唯一的征服者来临时，便吩咐他们的首领，将门扉举起，敬拜地说：\BibleQuote{城门，请提高你们的门楣，古老的门户，请加大门扉，因为光荣的君王要进入。}\par

10. 然而，即使在天上的万军中，仍有一些天使惊奇不已，对从未见过的如此威仪与光荣惊愕万分，于是他们问：\BibleQuote{谁是这位光荣的君王？} 不过，鉴于天使（和我们一样）也是逐步获得知识，并能够不断进步，他们必然显示出能力和悟性的差异，因为唯一天主超越一切渐进增长的限制，祂从永恒便拥有一切圆满。\par

11. 另一些天使——就是那些曾目睹祂复活、已看见或已认出祂的——回答说：\BibleQuote{就是那位强而有力的上主，那位在战场上大能的上主。}\par

12. 于是，众多天使再次以凯旋的合唱宣告：\BibleQuote{城门，请提高你们的门楣，古老的门户，请加大门扉，因为光荣的君王要进入。}\par

13. 而那些惊愕者再次发问：\BibleQuote{谁是那位光荣的君王？因为我们看见的祂没有俊美，也没有华丽；\BibleRef{依 53:2} 如果那不是祂，那到底\Emph{谁}是这位光荣的君王呢？}\par

14. 那些认识祂的天使回答说：\BibleQuote{\OriginalTerm{万军的上主}{Lord of Sabaoth}，祂就是光荣的君王。} 因此，万军的上主即是圣子。那么，亚略派何以称祂为可错者，而我们却相信祂是万军的上主，正如我们相信圣父是万军的上主一样？当我们发现圣子也被尊称为\BibleQuote{万军的上主}，正如圣父一样时，他们又怎能区分两者的至高权能？因为就在这一段经文中，许多抄本读作：\BibleQuote{万军的上主，祂就是光荣的君王。} 翻译家们将\Quote{万军的上主}有的地方译为\Quote{万军的上主}，有的地方译为\Quote{上主君王}，有的地方译为\Quote{全能的上主}。因此，既然那升天的是圣子，而升天者又是万军的上主，那必然推论出：天主子是全能的！\par

\Emph{没有信德，无人能升天堂；无论如何，那已如此升上去的人将被逐出。因此，信德必须热诚地保守。我们每人内里都有一个天堂，其门必须透过宣认基督的\OriginalTerm{天主性}{Godhead}而打开并高举，这些门不是\OriginalTerm{亚略派}{Arians}所能高举的，也不是那些在世俗事物中寻找圣子的人所能高举的，他们因此必须像\OriginalTerm{玛达肋纳}{Magdalene}一样被遣回宗徒那里，地狱的门不能胜过他们。引用圣经为要显示，主的仆人不可减损其主人的任何尊荣。}\par

15. 那么，我们该做什么呢？我们如何升天呢？那里有掌权者排列，有\OriginalTerm{掌权天使}{principalities}秩序井然地列阵，把守天堂的门，盘问上升的人。除非我宣认基督是全能者，谁将给我通行？天门紧闭——不为任何人开启；不是凡愿意的都能进入，除非他也按照真信德相信。君王的殿廷有卫兵把守。\par

16. 然而，假设有一个不配的人混入，偷越过把守天堂门的掌权天使，坐于主的筵席上；当宴席之主进来，看见一个没有穿着信德婚宴礼服的人，就会把他丢到外面的黑暗里，那里有哀哭和切齿，\BibleRef{玛 22:11}如果他不保守信德与和平。\par

17. 所以，我们应当保守我们所领受的婚宴礼服，不可拒绝基督那原本属于祂的；祂的全能，有天使传报、先知预言、宗徒见证，正如我们以上所已经证明的。\par

18. 或许，先知的言论不仅关乎宇宙天堂的门；因为还有其他的天堂，天主圣言也进入其中，论到这些天堂，经上说：\BibleQuote{我们有一位伟大的司祭，一位大司祭，祂已穿过诸天，就是天主子耶稣。}\BibleRef{希 4:14}那些天堂是什么呢？岂不是先知所说的那\Quote{诸天宣扬天主的光荣}的天堂吗？\par

19. 因为基督站在你灵魂的门前。听祂说话：\BibleQuote{看，我站在门口敲门：谁若听见我的声音而开门，我要进去，同他坐席，他同我一起坐席。}\BibleRef{默 3:20}教会论到祂说：\BibleQuote{我的爱人的声音在敲门。}\BibleRef{歌 5:2}\par

20. 祂站立着，但不是独自一人，因为有天使在祂前面先行，说：\Quote{城门，你们要举起高门；古老的门户，你们要加大开放。}什么门？就是圣咏作者在另一处所吟唱的：\Quote{请给我敞开正义的门。}那么，向基督敞开你们的门吧，好让祂进入你们内——敞开正义之门，贞洁之门，勇德和智慧之门。\par

21. 相信天使的讯息：\Quote{永恒的门户，你们要被举起，光荣的君王要进来，\OriginalTerm{万军}{Sabaoth}的上主。}你的门户是以忠信声音作出的高声宣认；这是主的门，是宗徒切望为他敞开的门，正如他说：\BibleQuote{求天主给我们打开传道的门，好宣讲基督的奥秘。}\BibleRef{哥 4:3}\par

22. 那么，让你们的门向基督敞开，而且不仅敞开，更被高举，如果它确实是永恒的，且不被定罪而毁灭；因为经上记载：\Quote{你们 \Emph{永恒} 的门户，要被举起。}当\OriginalTerm{色辣芬}{seraph}触摸依撒意亚的嘴唇时，门楣为他而举起，他看见了万军的上主。\par

23. 所以，你们的门将被高举，若你们相信天主子是永恒的、全能的、超越一切赞美和理解的，知晓一切过去和未来的事；而如果你们判断祂只有有限的能力和知识，且是下属的，你们就没有举起永恒的门户。\par

24. 所以，愿你们的门被高举，好使基督进入你们，不是亚略派所认为的那种基督——卑微、软弱、奴仆般的——而是具有天主形体的基督，与父同在的基督；愿祂以祂本来的样子进入，超越诸天和万有；愿祂向你们遣发祂的圣神。你们应该相信祂已升天，坐在圣父的右边，这对你们是有益的；因为如果你们以不敬的思想把祂拘羁在受造和尘世的事物中，如果祂不为你们离去，不为你们升天，那么\OriginalTerm{护慰者}{Comforter}便不会来到你们这里，正如基督亲自告诉我们的：\BibleQuote{我若不去，护慰者便不会到你们这里来；我若去了，就要派遣他到你们这里来。}\BibleRef{若 16:7}\par

25. 但若你们在尘世的人中寻找祂，如同玛达肋纳的玛利亚寻找祂那样，就要小心，免得祂对你们说，如同对她说的：\BibleQuote{你别拉住我不放，因为我还没有升到父那里。}\BibleRef{若 20:17}因为你们的门户狭窄——不给我通道——不能被高举，所以我不能进入。\par

26. 所以，你们到我弟兄们那里去——就是到那些永恒的门户那里去，它们一看见耶稣就被高举。伯多禄是一扇\Quote{永恒的门户}，地狱的门不能胜过他。\BibleRef{玛 16:18}若望和雅各伯，即雷霆之子，\BibleRef{谷 3:17}是\Quote{永恒的门户}。教会的门户是永恒的，先知切望宣扬基督的赞美，在那里说：\Quote{好使我在\OriginalTerm{熙雍}{Sion}女子的门口，宣扬你的一切美善。}\par

27. 因此，基督的奥迹是伟大的，连天使也惊愕惶惑。为此，你们有责任朝拜祂，而且身为仆人，不应减损你们主人的尊荣。你们不能以无知为借口，因为祂降来正是为此，好叫你们相信；如果你们不信，祂就不是为你们而降来，不是为你们而受苦。\BibleQuote{假使我没有来，没有教训他们，他们就没有罪；但现在他们对自己的罪，是无可推诿的。恨我的，也恨我的父。}\BibleRef{若 15:22}那么，谁恨基督呢？岂不是那说侮辱祂的话的人吗？——因为爱的本分是给予尊敬，恨的本分则是收回尊敬。恨人的，发出质疑；爱人的，献上尊敬。\par

\Emph{被\OriginalTerm{亚略派}{Arians}误用的\BibleQuote{每一个人的头是基督……而基督的头是天主}这句话，如今转过来攻击他们，以驳斥他们。其次，另一段圣经经文，常被同一异端者用作反对的依据，被引用来表明天主是基督的头，是在基督为人方面，即就祂的人性而言，而他们对那经文\BibleQuote{栽种的和浇灌的都是一样}的反对，显示出他们的不明智。在这些解释之后，圣父在圣子内，圣子在圣父内，以及信友在两者内的教义的意义，就得到阐明。}\par

28. 现在让我们审查\OriginalTerm{亚略派}{Arians}提出的其他异议。他们说，经上记载着：\BibleQuote{每一个人的头是基督，女人的头是男人，基督的头是天主。} \BibleRef{格前 11:3} 他们若愿意，请告诉我他们提出这异议的用意——是要将这四者联合起来，还是分开？假定他们是要联合，并说天主是基督的头，其意义和方式与男人是女人的头相同。看看他们会落入怎样的结论。因为如果这个比较是基于假定的同等地位，而将这四者——女人、男人、基督、天主——视为因具有同一本性而相似，那么女人和天主就会开始归于同一个定义之下。\par

29. 但如果这个结论因其不虔敬而不能令人满意，就让他们按照他们主张的任何原则来区分。如此，如果他们坚持基督与圣父的关系如同女人与男人，那么他们确实宣布基督和天主是同一\OriginalTerm{性体}{substance}，因为女人和男人在肉体方面是同一本性，他们的区别在于性别。但是，鉴于基督和祂的圣父之间没有性别差异，他们就会承认，就本性而言，圣子与圣父是同一且共同的，而否认性别上的差异。\par

30. 这个结论让他们满意吗？或者他们将女人、男人和基督视为同一性体，而将圣父与它们区分开来？这样能合他们的意吗？假设可以，那么就观察他们的归宿。他们必须要么承认自己不仅是\OriginalTerm{亚略派}{Arians}，而且是十足的\OriginalTerm{福提努斯派}{Photinians}，因为他们只承认基督的人性，认为祂只配与人类同列。要么他们必须，无论多么违背自己的倾向，赞同我们的信仰，我们以虔敬和神工的方式持守他们因不虔敬的思路而达到的观点，即基督确实，在祂的\OriginalTerm{神性受生}{divine generation}后，是天主的能力，而在祂取了肉身后，在肉体方面与所有人同一性体，除了祂\OriginalTerm{降生成人}{Incarnation}的特有荣耀，因为祂取了肉身的真实，而非幻像的样式。\par

31. 那么，就让天主成为基督的头，是就人性的条件而言。注意圣经没有说\Emph{圣父}是基督的头，而是说\Emph{天主}是基督的头，因为天主性，作为创造能力，是受造者的头。宗徒说得很好：\BibleQuote{基督的头是天主}；将基督的天主性和祂的肉身都带到我们思想中，也就是说，在提及基督之名时暗示\OriginalTerm{降生成人}{Incarnation}，在提及天主之名时暗示天主性的同一和主权的威严。\par

32. 但是，说就\OriginalTerm{降生成人}{Incarnation}而言天主是基督的头，引向这样的原则：基督作为降生成人者，是人的头，正如宗徒在另一处清楚表明的，他说：\BibleQuote{因为男人是女人的头，正如基督是\OriginalTerm{教会}{Church}的头} \BibleRef{弗 5:23}；而在随后的话中他补充说：\BibleQuote{祂为她舍弃了自己。} \BibleRef{弗 5:25} 所以，在祂降生成人之后，基督是人的头，因为祂的自我牺牲源自祂的降生成人。\par

33. 那么，基督的头是天主，是在考虑祂的奴仆形象，即人性，而非天主性时。但这毫无损于圣子，如果按照祂肉身的真实，祂与人相似，而在祂的天主性方面，祂与父是一体，因为就这点而言，我们并未从祂的主权夺走什么，反而将怜悯归于祂。\par

34. 但是谁能够凭良心否认圣父与圣子的同一\OriginalTerm{天主性}{Godhead}，当我们的主为祂的门徒完成教导时说：\BibleQuote{愿他们合而为一，正如我们一样。} \BibleRef{若 17:11} 这记载为信仰作证，尽管\OriginalTerm{亚略派}{Arians}将其歪曲以适应他们的异端；因为，他们既然不能否认屡次提及的合一，就试图减弱它，以便使圣父与圣子之间存在的天主性的合一，看起来如同人们中间的敬虔和信仰的合一，尽管在人们中间，\OriginalTerm{性体}{nature}的共同性造成了他们的合一。\par

35. 这样，我们以充分的清晰度驳斥了\OriginalTerm{亚略派}{Arians}通常为削弱神圣合一而提出的异议，基于这句经文：\BibleQuote{但栽种的和浇灌的都是一样。} 这段经文，亚略派若是明智，就不会引用来反对我们；因为如果\Person{保禄}和\Person{阿颇罗}在本性和信仰上都是一，他们怎能否认圣父与圣子是一呢？同时，我们确实承认这些不可能在所有方面都是一，因为人间事物不能与神圣事物相比。\par

36. 那么，由于神圣\OriginalTerm{性体}{Substance}的合一，\OriginalTerm{圣言}{Word}与圣父之间不能有分离，在能力上不能有分离，在智慧上不能有分离。再者，圣父在圣子内，如我们常常读到的，然而祂并非如同圣化那缺乏圣化者一般住在圣子内，也不是如同充满虚空一般，因为天主的能力不知虚空。再者，一位的能力不被另一位的能力增加，因为没有两种能力，而是一种能力；\OriginalTerm{天主性}{Godhead}也不容纳天主性，因为没有两个天主性，而是一个天主性。相反，我们将通过从另一位所领受并住在我们内的能力，在\OriginalTerm{基督}{Christ}内合而为一。\par

37. 合一标记是共通的，但\OriginalTerm{天主的本体}{the Substance of God}与\OriginalTerm{人的本体}{substance of man}不同。我们将成为一体，而\Person{圣父}与\Person{圣子}已是一体；我们将藉\OriginalTerm{恩宠}{grace}成为一体，而圣子则是因本体而一体。再者，\OriginalTerm{结合性合一}{unity by conjunction}是一回事，\OriginalTerm{本性合一}{unity by nature}则是另一回事。最后，请注意\WorkTitle{圣经}上记载的话：\BibleQuote{愿他们合而为一，父啊，如同祢在我内，我在祢内。}\BibleRef{若 17:21}\par

38. 现在请注意，祂没有说\BibleQuote{祢在我们内，我们在祢内}，而是说\BibleQuote{祢在我内，我在祢内}，为将自己与受造物区分开来。然后祂又补充说：\BibleQuote{使他们也在我们内}，为在此将祂与父的尊位与我们分开，使我们与父和子的结合显明是基于\OriginalTerm{恩宠}{grace}，而非本性；至于父与子的合一，则应信子并非藉恩宠领受此合一，而是因着祂\OriginalTerm{子的身份}{Sonship}自然拥有的权利。\par

\Emph{异端者所恶意引用的经文，即\BibleQuote{子凭自己什么也不能作}，首先从后续的话加以解释；然后，逐字检视经文，指明若按\OriginalTerm{亚略派}{Arian}的意涵去接受，必然招致不虔或荒谬的指责，证据主要在于世界的创造和\Person{基督}的若干奇迹。}\par

39. 再者，\OriginalTerm{亚略派}{Arians}提出的另一项反对，否认\Person{圣父}与\Person{圣子}的\OriginalTerm{大能}{Power}能够是一而相同的，其根据是祂的话：\BibleQuote{我实实在在告诉你们：子凭自己不能作什么，除非祂看见父所作的。}\BibleRef{若 5:19} 因此他们断言，子凭自己什么也没有作，而且除了祂看见父所作的以外，什么也不能作。\par

40. 噢，对不信者论点的智慧预知，更预备了答覆的方法，加上了随后的话：\BibleQuote{因为凡父所作的，子也照样作}\BibleRef{若 5:19}，这确是接续的话。那么，为何经上写的是\BibleQuote{子作相同的事}，而不是\BibleQuote{类似的事}呢？岂不是要你判断子与父的工程是合一，而非模仿吗？\par

41. 但为反过来考验他们的证据：我要让他们回答一个问题，子是否看见父的工程。我问，祂看见吗？若祂看见，那么祂也作（这些工程）；若祂作（这些工程），就让异端者停止否认祂的全能吧，他们既已承认祂能作一切祂看见父所作的事。\par

42. 但我们又当如何理解\BibleQuote{看见}呢？子有什么必要用肉体的眼睛吗？绝不。若他们这样主张子，那么他们也会推论出父也需要肉体的工作，好让子能看见祂自己将要作的事。\par

43. 再者，\BibleQuote{子凭自己不能作什么}这话是什么意思？让我们提出这问题并辩论。难道有什么是天主的大能与智慧不能成就的吗？请注意，这些乃是\Person{天主圣子}的名字，祂的威能决不是从另一位领受的恩赐；正如祂是\OriginalTerm{生命}{Life}（\BibleRef{若 14:6}），不依赖另一位的赋予生命行动，而是自己赋予他人生命，因为祂就是生命；同样，祂也是\OriginalTerm{智慧}{Wisdom}（\BibleRef{格前 1:24}），不是无知者获得智慧，而是以自己固有的使他人成为智慧；同样，祂也是\OriginalTerm{大能}{Power}（\BibleRef{格前 1:24}），不是因软弱而获得能力的增长，而是祂本身就是大能，并将能力赐予强壮的人。\par

44. 如此，大能怎能好像起誓般宣称：\BibleQuote{我实实在在告诉你们}——意思就是\Quote{我确确实实告诉你们}？\BibleRef{若 5:19} 那么，\Person{主耶稣}，祢确实说了，并确证，甚至重复祢庄严的声明：祢除了看见圣父所作的以外，什么也不能作。祢创造了宇宙。难道祢的圣父造了另一个宇宙，好让祢作为模型吗？因此祢的亵渎者必须承认有两个，或多个宇宙，如同哲学家所主张的那样，从而陷入异教的谬误；或者，如果他们愿意遵循真理，就让他们说，祢所造的，是祢没有借助任何典范造的。\par

45. 主啊，请告诉我，祢何时看见祢的父\OriginalTerm{降生成人}{incarnate}，在海上行走？因为我不知道，而且我认为信父做这种事是亵渎的，因为我知道只有祢取了我们的肉躯。祢何时看见父在婚宴上，把水变成酒？不，但我读过，祢是独一的唯一圣子，由父所生。我曾受教导，只有祢，在\OriginalTerm{降生成人}{Incarnation}的奥迹中，由\OriginalTerm{圣神}{Holy Ghost}及\Person{童贞玛利亚}诞生。所以，我们列举的各项祢的作为，父并没有做，而是祢独自做的，没有依凭祢的父所行任何事工的引导，为了以祢的血赎买世界的\OriginalTerm{救恩}{salvation}，祢从童贞玛利亚的母胎中无玷而出。\par

46. 当他们说\BibleQuote{子凭自己不能作什么}时，他们实在没有排除任何事，以至于有一位亵渎者甚至说：\Quote{祂连一只蚊子也造不出来}，以如此刚愎的亵渎和傲慢侮辱了\OriginalTerm{至高的大能}{Supreme Power}；但也许他们以为祢\OriginalTerm{降生成人的生命}{Incarnate Life}的奥迹是必要的例外。但是，\Person{主耶稣}，请说，圣父没有祢造了什么大地呢？因为没有祢，祂没有造任何高天，因为经上写着：\BibleQuote{因上主的圣言，诸天建立。}\par

47. 然而圣父也没有在祢以外创造大地，因为经上记载着：\BibleQuote{万物是藉着他而造成的；凡受造的，没有一样不是由他而造成的。} 因为如果圣父在祢——\OriginalTerm{天主圣言}{God the Word}——以外造了什么，那么就不是万物都是藉着圣言造成的，\OriginalTerm{圣史}{Evangelist}就是说谎的了。反之，如果万物是藉着圣言造成的，如果凡是先前不存在的，都是藉着祢而开始存在，那么祢自己，确是由祢自己，造了祢未曾见圣父造过的东西；尽管也许我们的对手会求助于\Person{柏拉图}的理论，在祢面前提出哲学家们所臆想的理型，但我们知道，连哲学家自身也已将其推翻了。另一方面，如果祢自己是由祢自己创造了万物，那么不信者们的主张便是徒然的，他们将学业的进步归于那自行提供其技艺教导的万有创造者。\par

48. 但如果异端者否认天和地是祢所造的，就让他们小心，看他们因自己的疯狂将自身投入何等深渊，因为经上记载着：\BibleQuote{愿那没有创造天地的神明灭亡。} \BibleRef{耶 10:11} 那么，\OriginalTerm{亚略}{Arian}啊！那寻获并拯救了已丧亡者的，难道祂还要灭亡吗？不过我们要言归正传。\par

\Emph{继续阐释他曾开始的争议经文，\Person{安博罗削}提出我们断言某事物不能存在的四种理由，并指出前三种理由不适用于基督，从而推断出圣子凭自己不能做什么的唯一原因是祂与圣父在权能上的合一。}\par

49. 在何种意义上圣子凭自己不能做什么呢？让我们探讨祂不能做的是什么。不可能性有多种不同类别。有些事自然不可能，有些事自然可能，却因某种软弱而成为不可能。还有，有些事因力量而成为可能，因不熟练或身体与心灵的软弱而成为不可能。此外，有些事由于不变旨意的法则、坚定意志的持守以及友情的忠信，成为不可能改变的事。\par

50. 为使此点更清楚，让我们藉实例来思考这事。一只鸟不可能学习任何学问或受训于任何技艺：一块石头不可能向任何方向移动，因为它只能因另一物体的运动而被移动。因此，石头凭自身无法移动，无法离开其位置。同样，鹰不能学习人类的学问之道。\par

51. 再举一例，患病的人不可能做强壮者的工作；但在这一情况中，不可能性的理由属于另一种类，因为此人因疾病而不能做他本性有能力做的事。所以在这一情况中，不可能性的原因是疾病，而这类不可能性与第一种不同，因为人是由于身体虚弱而被阻碍，不能有所作为。\par

52. 再者，有第三种不可能的原因。一个人可能本性有能力，他的身体健康也容许他做某项工作，然而他却不能做，因为缺乏技能，或因为他的人生地位使他不够资格；亦即，他缺乏所需学问或身为奴隶。\par

53. 请你们想想，我们列举的这三种不同的不可能原因（且将第四种搁置一旁），哪一种可以恰当地归诸天主圣子的情况呢？祂像石头一样本性无知觉、不移动吗？祂诚然是恶人的绊脚石，信徒的角石；但祂并非无知觉，因有感知之万民的忠诚爱戴寄托于祂。祂不是不动之磐石，\BibleQuote{因为他们饮了伴随他们的神磐石，那磐石就是基督。} \BibleRef{格前 10:4} 那么，圣父的工作并不因本性差异而对基督成为不可能。\par

54. 也许我们可能设想，有些事由于软弱而对祂成为不可能。但祂并不软弱，祂能以权柄之言治好他人的软弱。当祂命瘫子拿起床行走时（\BibleRef{谷 2:11}），祂显出软弱了吗？祂吩咐那人行一项需要健康作为必要条件的事，那时病人甚至还在祈求治疗他的疾病。主是万军的统帅，祂使瞎子看见，使跛者直立，使死者复活（\BibleRef{玛 11:5}），当我们祈祷时祂预先赐下药物的效果，治愈求告祂的人，只要触摸祂的衣边就得以洁净（\BibleRef{谷 6:56}），祂并不软弱。\par

55. 除非，你们这些恶人啊，你们看见祂的创伤时，或许认为那是软弱。诚然，它们是刺穿祂身体的创伤，但那创伤并未显出任何软弱，众生之生命正从那里流出，因此先知说：\BibleQuote{因祂的鞭伤，我们得了痊愈。} \BibleRef{依 53:5} 那么，祂在受创之时并不软弱，难道在祂的至尊权柄上却是软弱的吗？那么我要问，怎么说呢？当祂命令魔鬼，宽恕罪人的过犯时（\BibleRef{路 5:20}），又如何呢？或是当祂向圣父恳求时？\par

56. 在这里，我们的对手或许确实会追问：\Quote{如果圣子一时命令，一时恳求，圣父与圣子怎能是一体呢？} 诚然，他们是一体；也诚然，祂既命令又祈求：然而在祂命令的时刻，祂并不孤单；同样，在祂祈祷的时刻，祂也不软弱。祂不孤单，因为凡圣父所做的，圣子也同样照样做。祂不软弱，因为虽然在肉身内祂为我们的罪而承担了软弱，但那却是为了我们得平安的惩罚（\BibleRef{依 53:5}），临到祂身上，并非祂自身缺乏至尊权能。\par

57. 而且，为使你知道祂是凭人性而祈求，凭天主性而命令，福音中为你记载了祂对\Person{伯多禄}说的话：\BibleQuote{我已为你祈求，叫你的信德不至丧失。}\BibleRef{路 22:32} 对同一位宗徒，另一次他先前曾说：\BibleQuote{祢是基督，永生天主之子。}祂回答说：\BibleQuote{你是伯多禄（磐石），在这磐石上，我要建立我的教会，我要将天国的钥匙交给你。}\BibleRef{玛 16:18} 那么，对于这位祂凭自己的权柄赐予天国，称其为磐石，并以此宣示他为教会的基石的人，难道祂不能坚固他的信德吗？因此，请思考祂祈求的方式和命令的时机。当祂被显示为即将受难时，祂祈求；当祂被相信为天主之子时，祂命令。\par

58. 因此，我们明白，两种不可能性不能提供解释，因为天主的能力既非无感知，也非软弱。那么你们会提出第三种可能，即祂什么也不能做，就像一个没有技艺的学徒若无师傅的指导就什么也不会做，或像一个奴隶若无主人就什么也不能做一样。那么，主耶稣，祢称自己为师傅和主，就是说了谎话，并用祢的话欺骗了祢的门徒：\BibleQuote{你们称我为师傅、主，你们说得对，我本来就是。}\BibleRef{若 13:13} 不，祢这真理，决不会欺骗人，更不会欺骗祢称之为朋友的人。\BibleRef{若 15:14}\par

59. 然而，如果我们的敌人将祢与造物主分离，视祢为笨拙无能，那么让他们看看，他们如何断言祢，即神圣智慧，缺乏技巧；尽管如此，他们却无法割裂祢与圣父所共有的同一性体。其实，工匠与笨拙者之间的差别并非基于本性，而是由于无知；但手艺既不属于圣父，无知也不属于祢，因为不存在无知的智慧。\par

60. 因此，如果圣子没有无感知的属性，也没有软弱、无知和奴性，那么不信者应当深思：圣子就本性和权能而言与圣父是一体的，在运作上祂的能力与圣父并不相悖，因为\BibleQuote{凡父所做的，子也照样做，}因为若不在同一本性的合一中有份，且在工作的方式上也不低劣，就没有人能照样的做同一工作。\par

61. 然而我仍要追问，除非看到圣父在做，圣子到底\Emph{什么}不能做。我要采用愚妄人的方式，从下界的事物中举出一些例子。\BibleQuote{我成了狂妄的人，是你们逼我的。}\BibleRef{格后 12:11} 的确，有什么比争论天主的尊威更愚妄的呢？这与其说是基于信德的虔敬教导，不如说是引发疑问。但是，让争论对付争论，让言语回应他们，而对我们，则是爱，那在天主内的爱，发自纯洁的心、善良的良心和无伪的信德。因此，我不忌讳甚至引进可笑的事来反驳如此虚妄的论点。\par

62. 那么，圣子如何看见圣父呢？马看到一幅画，自然无法模仿。圣子观看圣父并非如此。孩子看到成年人的工作，却无法复制；当然，圣子观看圣父也非如此。\par

63. 那么，如果圣子能凭祂与圣父所共有的同一本性的奥秘能力，以不可见的方式既观看又行动，并凭祂天主性的圆满执行祂旨意的每一决定，我们除了相信圣子由于不可分割的能力合一，除了祂看见圣父所做的以外，什么也不做，因为由于祂无与伦比的爱，圣子凭自己什么也不做，因为祂不愿任何违背圣父旨意的事，这确实不是软弱的证据，而是合一的证据。\par

\Emph{现在考虑第四种不可能性（§49），并指出圣子不做任何圣父不赞同的事，因为祂们之间具有意志与能力的完美合一。}\par

64. 此外，圣子——现在考虑我们的第四前提——并非自行其是，因为祂，神圣的同座者，从未做过任何不符合圣父旨意的事。而且，圣父看到圣子所造的一切，都认为甚好；因为在\BibleBook{创世}{纪}上这样记载：\BibleQuote{天主说：\InnerQuote{要有光！}就有了光。天主见光好。}\BibleRef{创 1:3}\par

65. 那么，圣父在那时曾说：\Quote{要有我所造的那样的光，}还是说\Quote{要有光}——那时光尚不存在；或者圣子询问圣父造了什么样的光？不是的，圣子照自己的旨意造了光，并且完全符合圣父的喜悦，以致圣父认可。这段经文所说的是圣子全新的、独创的工程。\par

66. 再者，如果如亚略派对圣经的解释所主张的那样，圣子造祂所见的是羞辱，而圣经却表明祂造了祂先前未见的事物，并使尚未存在的事物得以存在，那么他们又该怎样论及圣父呢？祂称赞祂所见到的，仿佛祂不能预见将要受造的事物似的？\par

因此，\OriginalTerm{圣子}{Son}看见\OriginalTerm{圣父}{Father}的作为，正如圣父看见圣子的作为；圣父称赞这作为，并非如称赞他人之作，而是认为这是祂自己的作为，因为\BibleQuote{凡圣父所做的，圣子也照样做。}[经上如此记载，为的是]使你们明白，这一作为既是圣父的也是圣子的作为。因此，圣子所作的无非是圣父所认可、所称许、所愿欲的，因为祂的整个存在出于圣父；祂并非如受造之物，常常犯罪，多次违背造物主的旨意，贪恋罪恶并陷入其中。所以，圣子所行的一切，无非是取悦圣父的，因为祂们有一个旨意、一个目的、一个真正的爱、一个行动的果效。\par

再者，为要证明这是出于爱——圣子凭自己不能做什么，除非祂看见圣父所做的——\OriginalTerm{宗徒}{Apostle}在此话之后加上理由：\BibleQuote{凡圣父所做的，圣子也照样做，}\BibleQuote{因为圣父爱圣子。}因此，圣经将圣子不能做的证据，归于不容分隔或分歧的爱之合一。\par

然而，倘若位格在爱内的不可分离性（这确是真的）基于本性的同一，那么，基于同样的理由，他们在行动中也不可分离；当圣子所做的，圣父也做，圣父所做的，圣子也做，圣子所说的，圣父也说时，圣子的作为不可能不与圣父的旨意一致，正如经上所记：\BibleQuote{我的父居住在我内，是祂在说话，我所做的工作，祂自己去做。}\BibleRef{若 14:10}。因为圣父无不凭借祂的能力和上智而行，祂智慧地创造了万物，如经上所记：\BibleQuote{祢以智慧造成万物；}同样，\OriginalTerm{天主圣言}{God the Word}若没有圣父的参与，也一无造成。\par

祂并非离开圣父而工作；祂并非没有圣父的旨意而献上自己，作为为拯救全世界而被宰杀的祭品，承受至圣的苦难；祂并非没有圣父旨意的合作而复活死者。例如，当祂将要复活拉匝禄时，祂举目向天说：\BibleQuote{父啊，我感谢祢，因为祢俯听了我。我知道祢常俯听我，但我说话是为了四周站立的群众，使他们相信祢派遣了我，}\BibleRef{若 11:40}，为的是，虽然祂是按所取的人性、在肉身内说话，祂仍然表达出在旨意和行动上与圣父的合一，因为圣父垂听并看见圣子所愿的一切；因此，圣父也看见圣子的行事，听见祂旨意的表达，因为圣子未曾祈求，却说祂已蒙俯听。\par

再者，我们不能设想，圣父不会俯听圣子旨意所定的一切；为要显示祂常被圣父俯听，不是像仆人，不是像先知，而是作为圣子，祂说：\BibleQuote{我知道祢常俯听我，但我说话是为了四周站立的群众，使他们相信祢派遣了我。}\par

因此，祂是为我们而感谢，免得我们听见圣父和圣子做成同一工作时，以为圣父和圣子是同一位格。此外，为要显示祂的感谢并非出于软弱者所献的酬谢，相反，祂作为\OriginalTerm{天主圣子}{Son of God}，常宣称自己拥有神圣的权威，祂大声说：\BibleQuote{拉匝禄，出来！}这无疑是命令的声音，而非祈祷。\par

\Emph{所着意强调的教理，由这一真理所证实：圣子是圣父的圣言——这圣言，不是我们通常所理解的那种言语，而是生活的、行动的圣言。既然这样，我们就不能否认祂与圣父具有同一的旨意、能力和\OriginalTerm{本体}{Substance}。}\par

然而，回到我们先前所讨论的，完成摆在我们面前的任务。圣子作为\OriginalTerm{圣言}{Word}，执行圣父的旨意。如今，我们通常所理解、所使用的言语，是一种\Emph{言语}。有音节与声音，却与我们内心的思想不相违背；我们以话语的见证，将内心领悟和感受的表露出来，这言语仿佛为我们工作。但我们所说的话语，本身并无直接效力；唯有天主圣言，既非单纯的话语，也非他们所谓的\Quote{内在概念}，却有效地运行，是生活的，并具有治愈的能力。\par

你想知道圣言的本质是什么吗？——请听圣经。\BibleQuote{天主圣言是生活的，是有效力的，比各种双刃的剑还锐利，直穿入灵魂和神魂，关节与骨髓的分离点。}\BibleRef{希 4:12}\par

那么，你听见了天主圣言，还要将祂与圣父的旨意和权能分离吗？你听见祂被称为生活的圣言、治愈的圣言——那么就不要设法将祂与我们口舌的言语相比；因为我们说出的言语，虽无眼可观、无耳可听，却仍能说话，且它所言说的知识，是藉着人性隐藏的奥秘而生效的，那么，一个人若要求天主圣言中的\OriginalTerm{天主性}{Godhead}具有某种形体的视觉和听觉，并认为圣子凭自己什么也不能做，除非祂看见圣父所做的（尽管正如我们所说，在圣父、圣子和圣神内，因着同一本体之故，具有同一旨意，无论行与不行，也具有同一能力），这人怎能逃脱亵渎的指责呢？\par

然而，虽然人在思想和感情上通常各不相同，却仍能就某一命题的意义达成一致；那么，对于天主圣父与圣子，我们该作何想法呢？——尤其是，在天主性的本体中，存在着那为人类之爱所效法者。\par

77. 然而，让我们假设——正如我们的对手所愿——圣子仿佛在复制他所见圣父所行之事的模式。但即便如此，我们也必须承认，这意味着他是\OriginalTerm{同一性体}{same substance}的，因为无人能完全模仿他人的工作，除非他与那人在同一本性内合一。\par

\Emph{异端者的反驳，即圣子不能与圣父同等，因为圣子不能生子，被反弹回提出此说者自身。从人性的事例显明，一个人是否生育后代，与他的能力无关。这尤其应为真，因为否则，圣父自己也得被宣布为缺乏能力。由此可知，我们无权以人性之事判断神性之事，而必须立足圣经的权威，否则我们必须否认圣父或圣子的一切能力。}\par

78. 陛下，某些人提出愚人的抗辩，为要表明圣父与圣子并不同等，说圣父是全能的，因为他生了圣子，但圣子不是全能的，因为他不能生。\par

79. 但看他们的亵渎何等狂妄，他们哲学家的逻辑如何自相矛盾。因为提出此问题必然导致要么他们亲口承认圣子与圣父\OriginalTerm{同永远}{co-eternal}，要么，若他们给圣子的存在设定一个开始，就必定给圣父的能力设定一个开始。因此，当他们否认圣子是全能的时，他们就走向了不虔敬的断言——即圣父因圣子的帮助而开始成为全能的。\par

80. 因为如果圣父因生了圣子而成为全能的，那么，当然，要么圣子与圣父同永远，因为如果圣父是永远全能的，那么圣子也是永远的，要么，如果曾有一段时间没有永恒的圣子，随之曾有一段时间没有全能的圣父。因为他们若声称有一段时期圣子开始存在，他们就滑回到〔错误〕说圣父的能力也不是从永恒就有，而是因圣子的受生而开始存在。这样，他们意图羞辱天主子，却如此增加了他的尊荣，以致似乎使他——违反一切正确信仰——成为圣父之能力的源头，尽管圣子说：\BibleQuote{凡父所有的，都是我的}\BibleRef{若 16:15}——即是说，不是他赋予圣父的，而是他从圣父领受的，因他是圣父所生的圣子，具有子的权利。\par

81. 因此我们宣告圣子为永恒的大能；\BibleRef{罗 1:20} 那么，如果他的能力和天主性是永恒的，他的主权自然也是永恒的。所以，羞辱圣子的就是羞辱圣父，是违背责任与爱德的仇敌与冒犯者。我们该尊敬圣子，圣父在他内喜悦，因为圣父乐意赞美归于圣子，而圣父自己也在他内喜悦。\par

82. 然而，我们要回答他们竭力建立的结论；但我们似乎为个人诉求所驱使，逃避了处理眼前问题的困难。他们说，圣父生了圣子；圣子却不能生。这怎能证明他们不平等呢？生是圣父作为父亲的自然功能，并非他主权能力的必然结果。再者，孝顺的尊重使人与人平等，而非分离他们。而且，我们这些脆弱凡人的经验教导我们，常有弱者有子而强者无子，奴仆有儿女而主人无后，贫者生育而富人无子息。\par

83. 但如果我们的对手说，这也可能是软弱的结果，因为人可能想要生育子女却无能；那么，虽然神性之事不可通过人性之事来判断决定，但他们当明白，在人间亦如在天主那里，有无子女并不依赖或源于其权威能力，而取决于作为父亲的个人属性，且生育不在于我们意志的能力，而取决于我们身体的特质；若是主权之事，那么更强有力的君王就当有更多儿子。因此，有子或无子，与主权并无必然关联。那么，在自然本性上也是如此吗？\par

84. 如果你们〔我的亚略派对手〕将你们所反对的视为本性软弱，并依赖取自人性本性的例子，记住圣父的本性与圣子的本性相同，因此你们要么承认圣子为真子，并因在同一个本性的合一而在圣子的位格上羞辱圣父（因为圣父按本性是天主，圣子亦然；然而宗徒说，\Quote{许多的神}按本性并非如此，只是如此称呼而已）；要么，如果你们否认他是真子，即拥有同一本性，那么他就不是所生的，而如果圣子不是所生的，圣父就没有生他。\par

85. 因此，按照你们的说辞，我们得出的结论是，天主圣父不是全能的，因为他若没有生圣子，而是创造了圣子，他就不能生。但因为圣父是全能的，按你们所持，他是全能者，因他是唯一的存在之创造者，那么他当然生了圣子，而非创造了他。然而，我们应当相信他的话，而非你们的话。他说：\Quote{我生了}，且不止一次，为自己作证他是生育者。\par

86. 那么，基督未曾生育，并不表示本性或权威的软弱，因为生育，如我们已多次说的，与权威的至高无关，而关乎本性中的一种个人特性。因为如果圣父的全能由此构成，即他有一个子，那么他若生了更多的子，他本应更是全能的。\par

87. 难道祂的能力因生一位就耗尽了么？绝非如此。我还要指出，基督也有子女，祂每日生育他们，但这种生育——更准确地说，是再生——关乎个人权柄而非本性，因为义子名分乃是权柄的运用和赋予，而生育是本性的彰显，正如圣经自身所教导的：因为若望说：\BibleQuote{祂在世界中，世界是藉着祂造成的，世界却不认识祂。祂来到自己的领域，自己的人却没有接受祂。但是，凡接受祂的，祂赐给他们权能，好成为天主的子女，就是那些相信祂名字的人。}\par

88. 因此我们说，祂使我们成为天主的子女，乃是祂权柄的功用和行使；而天主的神谕揭示，祂的生育关乎个人的属性，因为天主的智慧说：\BibleQuote{我从至高者的口中出来，}\BibleRef{德 24:5}意思是，不是出于强迫，而是自由的，不是受权柄的束缚，而是按照至高主权和权柄正义的个人能力，在隐秘的出生中诞生的。再者，论到同一智慧，就是主耶稣，圣父在另一处说：\BibleQuote{我在黎明之前，从胎中生了你。}\par

89. 祂说这话，不是要我们想到肉身的母胎，而是要表明真正的生育是祂固有的活动，因为如果我们把这些话理解为由肉身生育，那么（我们便是在暗示）全能圣父怀孕并产难。但绝不可这样！我们不应以这软弱的肉身框架去度量天主的伟大。\Quote{胎}一词代表隐藏的奥秘，是圣父存有的内层圣所，无论天使、总领天使、掌权者、宰制者，还是任何受造之物，都不得进入其间。因为圣子常与圣父同在，并在圣父内——与圣父同在，是凭借永恒天主圣三所固有的、区分而不分离的特性；在圣父内，是因天主性体本质上的合一。\par

90. 那么，这里哪里还有余地让人坐在审判席上审判天主性，质问圣父与圣子——一个生育，另一个不生育？没有人因自己的仆人或婢女生育（或生养）子女而定他们的罪；但那些亚略派的人却因基督没有生育而定祂的罪——他们确实定祂的罪，因为他们暗中将定罪的判决加于祂，剥夺祂的荣耀与尊严。为什么他们没有生育子女的问题，并不会导致已婚夫妇丧失爱意，或否定彼此的优点；但亚略派却因为基督没有生育子嗣，便轻视祂的至高主权。\par

91. 他们问：为什么子不是父？因为，另一方面，父不是子。为什么基督没有生育？正是因为父不是受生的。然而，子并不因为没有成为父而地位稍低；父也不因为没有成为子而地位稍低，因为子说过：\BibleQuote{父所有的一切，都是我的}\BibleRef{若 16:15}——生育真实地包含在父的个人属性中，并不单凭主权而来。\par

92. 可以说，天主圣三的实体，在那区分之中是一个共同的性体，一个不可思议、不可言喻的实体。我们持守父、子、圣神的区分，而非混淆；是区分而无分离；是区分而无复数；因此，我们相信父、子、圣神，每一位都从永恒到永恒地存在于这神圣而奇妙的奥迹中：没有两位父，没有两位子，也没有两位圣神。因为 \BibleQuote{只有一个天主，就是父，万物都出于祂，我们也归于祂；也只有一个主，即耶稣基督，万物藉着祂而有，我们也藉着祂而有。}\BibleRef{格前 8:6} 有一位由父所生的，就是主耶稣，因此祂是独生子。同一位宗徒说过：\BibleQuote{也只有一个圣神。}\BibleRef{格前 12:11} 我们如此相信，如此诵读，如此持守。我们知道区分的事实，却对隐藏的奥秘一无所知；我们不探求原因，只持守赐予我们的外在标记。\par

93. 何等怪诞的邪恶啊！那些对自己生育毫无权力的人，竟敢僭取权力，探求天主的生育！让他们否认吧，否认圣子与圣父同等，因为祂没有生育；让他们否认吧，否认圣子与圣父同等，因为祂有一位父！但是，如果他们以这种论调谈论人，那些有时渴望生育儿子却未能如愿的人，我们会称之为侮辱；同样，如果两人中，一人有儿子，另一人无子，便因此说无子的那人比有子的低一等，我们也会称之为侮辱。我说，即便在单纯的人事上，一个人因为拥有父亲而被轻视，也是怪诞至极。或许，亚略派认为基督处于家庭中一员的地位，因未获自由、不独立于父的权柄、无权管理家业而烦恼。但基督并非处于监护之下；相反，祂废除了一切监护。\par

94. 那么，请他们告诉我们，他们希望这些事如何？——真正的生育，真正的圣子由天主圣父所生，即由父的实体所生，还是由另一实体所生？如果他们说是 \Quote{由父所生，即由天主的实体所生，} 那好极了，因为这样他们便承认圣子是由父的实体所生的。那么，既然是同一实体，当然也具有同一至高能力。反之，如果圣子是由另一实体所生，那么圣父怎么能是全能的，而圣子却不是全能的呢？因为天主若用另一实体造了祂的圣子，那又有什么益处呢？毕竟，众所周知，圣子同样也用另一实体使我们成为天主的子女。因此，圣子要么与圣父同一实体，要么同一至高能力。\par

95. 那么，我们对手的诘难便落空了，因为他们不能审判基督——更确切地说，因为当祂受审时，祂是清楚的。然而，他们既对我们提出这样的问题，就理当被他们自己的判决定罪，因为如果圣子与圣父不平等是由于祂没有生一位子，那么让那些散播这类议论的人，如果他们没有子女，总该承认，他们自己的仆人倒应比他们自己更受看重，因为他们不能与有子女的人平等——反之，如果他们有子女，他们应将这优点归于他们的儿子，而非归于自己。\par

96. 这样，这反对理由便立不住脚：即圣子不能与圣父平等，因为圣父生了圣子，而圣子自己并没有生一位子。正如泉源生出溪流，而溪流却不从自身生出泉源；光产生光辉，而光辉不产生光，但光辉与光的本质却是同一个。\par

\Emph{\OriginalTerm{亚略派}{Arians}提出的各种诡辩，旨在证明圣子有存在之始，现予以审议和驳斥。其根据是，亚略派既没有证明任何东西，或者如果他们证明了什么，那反而证明了反对他们自己的东西，（因为那作为万有之始者，其自身不可能有开端），他们的推论甚至也不适用于人类存在的事实。时间不可能先于那位时间之作者而存在——如果祂确曾一度不存在，那么圣父就曾没有自己的德能与智慧了。再者，我们自己人类的经验也表明，一个人在出生之前即可被说成是存在的。}\par

97. 既然我们的对手未能坚持他们反对基督与圣父平等的理由——这理由是基于基督的受生，那么他们就该看到，他们那尽人皆知的诡辩手法，他们那惯用的歪曲伎俩，是行不通了。他们惯常提出的谜题是：\Quote{圣子怎能与圣父平等？如果祂是子，那么在受生之前祂并不存在。如果祂是存在的，那祂为什么受生？}而那些提出亚略困难的人，却又坚称自己不是\OriginalTerm{亚略派}{Arians}。\par

98. 因此，他们要求我们回答，意在如果我们说：\Quote{圣子在受生之前就已存在，}便用如下巧妙的驳难来应对：\Quote{那么，这样祂在受生之前就是受造的，与其它受造物毫无区别，因为祂在被立为子之前就已开始是受造物了。}他们又追问：\Quote{祂既已存在，为什么要受生？是因为祂不完备，要以后被造得更完备吗？}相反，如果我们回答说圣子在受生之前并不存在，他们就会立即回应：\Quote{那么，借着受生，祂才进入存在，在受生之前并不存在，}从而由此得出：\Quote{圣子曾一度存在，而那时祂并不存在。}\par

99. 但是，让那些提出这难题并企图将真理笼罩在云雾中的人，先告诉我们：圣父施展祂的生养能力，是在时间界限之内，还是在时间界限之外？如果他们说\Quote{在时间界限之内，}那么他们便将他们反对圣子的说法归到了圣父身上，致使圣父似乎也开始成为祂先前所不是的。如果他们回答\Quote{不在这样的界限之内，}那么除了亲自解开他们所提的难题，并承认圣子的受生不受时间的界限和条件所限，因他们否认圣父是那样生的，他们还能有什么呢？\par

100. 那么，如果圣子不是在时间界限内受生，我们便可以断定：在圣子之前不可能存在任何事物，因为祂的存在不受时间限制。事实上，如果在圣子之前真有什么存在，那么即刻便可推知，天上地下的一切都不是在祂内受造的，而宗徒在书信中所写的便是错了，\BibleRef{哥 1:16}；然而如果在祂受生之前什么也没有，我便看不出，一位没有任何在祂之先者的，怎么能被说成在某者之后。\par

101. 与此考虑相联的，还有他们另一个极其亵渎的刁难，其阴险的目的是要混淆纯朴之人的感觉和理解。他们问：凡有终结的事，是否也必曾有一个开始？如果告诉他们有终结的事也必有开始，他们便再追问：圣父是否停止了生祂的圣子？我们若同意这点，他们便推出结论说圣子的受生有开始。你若容许这点，看来便似乎会推出：如果受生有开始，那开始就似乎在受生者内开始；这样，一位先前不曾存在的，便可被称为\Quote{受生者}了——他们的目的便是通过立下这所谓的定论，即圣子有过不存在的时刻，来结束争论。\par

102. 除此以外，还有其它无谓的刁难，如同些口齿伶俐的人信口提出那样。他们说，如果圣子是圣父的\OriginalTerm{圣言}{Word}，那么祂之所以被称为\Quote{受生者}，就因为祂是圣言。但既然祂是圣言，祂就不是一种工程。然而，圣父曾\BibleQuote{以多种方式}发言，\BibleRef{希 1:1}，由此推出祂生了许多子，如果祂所说的是祂的圣言，而不是亲手所造的工程的话。愚蠢的人哪，他们说着话，仿佛不知道发出的话语，与那永远的、自圣父所生的、天主性的圣言——我说生，而不只是说出——之间的区别。在祂内没有音节的组合，而是永恒的、圆满的天主性，以及无尽的生命！\par

接着有另一个亵渎之言，他们借此询问圣父生圣子是由于祂自己的自由意志，还是出于强迫，意图是，如果我们说\Quote{出于祂自己的自由意志}，就显得我们承认圣父的意志先于神圣的受生，并回答说在圣子的存在之前先有了什么，所以圣子不与圣父同为永恒，或者祂像世界万物一样是一个受造物，因为经上写着：\Quote{祂凡祂所愿意的，都做了}，虽然这话不是指圣父和圣子，而是指圣子所造的受造物。另一方面，如果我们回答说圣父是由于强迫而生了圣子，我们就似乎将软弱归给了圣父。\par

但在永恒的受生中，没有任何预先的条件，既没有意志，也没有不情愿，因此我既不能说圣父是出于自己的自由意志生了圣子，也不能说祂是由于强迫生了，因为生并不取决于由意志所决定的可能性，而是似乎在于圣父隐秘存在的某种权利和特质。因为正如圣父不是因祂愿意是善的而成为善的，也不是被迫成为善的，而是超越这些条件——善，即是，由于本性——同样，祂生育能力的发出，既不是出于意志，也不是出于必然性。\par

然而，让我们姑且同意他们的提议。假定受生取决于生者之意志；他们说这一意志行动是何时发生的？如果是在元始，那么显然，子在元始就存在。如果意志是永恒的，那么子也是永恒的。如果意志开始存在，那么天主圣父，在其本来状况中，对自己是如此不悦，以至于改变了自己的状态，也就是说，没有圣子时，祂对自己感到不悦；在圣子内，祂开始感到悦乐。\par

再推究其后果。如果圣父按照人性的方式，产生了生育的欲望，那么祂也经历了人在生产之前所遭遇的一切经验——但我们发现，生育不只是由意志决定，而是一个愿望的对象。\par

他们就这样暴露了自己的不虔，他们主张基督的受生有一个开始，为的是显出那不是圣言永恒的真受生，而是消逝并被遗忘的言语的发声，并且借着引入（众多子的前提），他们可以（据以）否认基督自身拥有神圣属性，最终使祂既不被视为独生子，也不被视为首生子；最后，一旦相信祂的存在有一个开始，就可能认为祂的存在也被命定有一个终结。\par

但天主圣子既没有任何开始，因为祂在元始就已存在，也不会有一个终结，祂是宇宙的起始与终末；因为既然是起始，祂如何能取得并接受那祂早已拥有的，或者是终末，祂如何能有一个终结，祂自己就是万物的终末，好使在那终末中我们有一个无尽的居所？神圣的受生不是发生在时间进程中、并受其限制的事件，因此在它之前没有时间，在它内时间也无立足之地。\par

再者，他们那无目的而又虚妄的问题，即使在针对受造界本身时，也找不到任何可攻击的漏洞；不，实际上，时间性的存在物，在某些情况下，似乎不允许有时间的分割。例如，光产生光芒，但我们既不能想象光芒在光之后才存在，也不能想象光在光芒之前就已经存在，因为哪里有光，哪里就有光芒，哪里有光芒，哪里也有光；因此，我们既不能有光而无光芒，也不能有光芒而无光，因为光在光芒内，光芒在光内。因此，宗徒被教导称子为\BibleQuote{圣父光荣的反映}，\BibleRef{希 1:3}，因为子是祂父之光的反映，由于能力的永恒而共永恒；由于光明的合一而不可分离。\par

如果我们既不能理解也不能分解这些在我们头顶天空中的受造物，就是我们所看见的，我们怎能理解那位我们看不见、超越一切受造存在的天主，正如祂在祂自己受生的至圣所内？我们能将时间作为祂与子之间的阻隔吗，当一切时间都是子的受造物时？\par

因此，让他们停止吧，不要再说什么在祂受生之前子不存在。因为\Quote{之前}这个词是时间的标记，而受生先于一切时间，因此后于某物而来的，不会先于它，作品也不能先于制造者，因为受造物必然从制造它们的工匠那里获得其开始。任何受造物的惯常行为怎能被视为先于其制造者存在，而一切时间都是受造物，每一受造物都是从其创造者获得存在呢？\par

因此，我要进一步审问那些自以为精明的对手，要他们把他们那套理论应用到人的存在上，使他们自圆其说，因为他们用这套理论来贬低天主存在的荣耀，并且远远避开对神圣受生中玄妙奥迹的任何宣认。我要他们在人的生育事实上找到他们反对的依据。关于天主子，他们断言在祂受生之前祂不存在——也就是说，他们这话是论及天主的智慧、能力、圣言，而祂的受生是没有任何先于自身之物的。但是，如果如他们所要我们相信的，曾有一段时间子不存在（这是亵渎的断言），那么，如果后来天主经历了一个生子的过程，就曾有一段时间天主缺乏神圣完美的满全。\par

113. 然而，为向他们显明其反对论点的虚弱和空洞——尽管它与任何神性或人性的真理都无真实关联——我要向他们证明，人在出生之前就已存在。否则，让他们说明，\Person{雅各伯}还在他母亲胎中暗室时，就已欺骗了他的兄长，难道在他出生之前，他并未被选定和命定吗？\BibleRef{创 25:23} 让他们说明，\Person{耶肋米亚}在他出生之前，岂不也正是如此——耶肋米亚，有信息传给他：\BibleQuote{我还没有在母胎内形成你以前，我已认识了你；在你还没有出离母胎以前，我已祝圣了你，选定了你作万民的先知。}\BibleRef{耶 1:5} 还有比这位伟大先知的事例更坚强的证据吗？他在出生之前就被祝圣了，在他成形之前就被认识了。\par

114. 再者，关于\Person{若望}，我要说什么呢？他的圣洁母亲作证：当他还在她胎中时，他便在心灵中感知到主临在，欢喜跳跃，正如我们记得经上如此记载，他母亲说：\BibleQuote{看，你请安的声音一入我耳，胎儿就在我腹中欢喜踊跃。}\BibleRef{路 1:44} 那么，那发预言的，他是存在的，还是不存在？不，他当然是存在的——他当然是存在的，他敬拜了他的创造者；他是存在的，他在母胎中就发言了。因此，\Person{依撒伯尔}充满了他儿子的神灵，\Person{玛利亚}被她的圣神所圣化，因为你可以找到这样的记载：\BibleQuote{胎儿就在她的腹中欢跃，依撒伯尔遂充满了圣神。}\BibleRef{路 1:41}\par

115. 思考每句话的恰当力量。\Person{依撒伯尔}确是首先听到\Person{玛利亚}声音的，但\Person{若望}是先感受到他的主恩宠临在的。预言与预言的和谐，女人与女人的和谐，胎儿与胎儿的和谐，是何等甜美。女人们说出恩宠的话语，胎儿们隐秘地动作，当他们的母亲彼此靠近时，他们便也在进行神秘的爱之交谈；在一个双重的奇迹中，尽管尊荣程度不同，母亲们都在他们小儿的神灵中说预言。我问，成就这奇迹的是谁？岂不是天主子，他使未出生的得以存在吗？\par

116. 因此，你们的反对论点无法与人类存在的真理调和——难道它就能与神圣奥秘相调和吗？你们所谓\Quote{在他受生之前，他不存在}的原则是什么意思？难道父曾有一段时间在进行构思，以致于在子受生之前过去了某些时期？难道他像妇人一样经历分娩阵痛，所以只是这阵痛？你们想怎样？我们何苦探求神圣的奥秘？圣经告诉我天主圣子生成的必要效果，而非如何生成。\par

\Emph{有人根据圣若望所记述，基督因为父而生活，因此不应被视为与父同等。对此异议，答复如下：圣子的生命，就其天主性而言，从未有开始之时；它不依赖任何其他，而该段经文应理解为指他人性的生命，这从他在那里谈及他的身体和血可得到指明。为这段经文提供了两种解释，其中一种指明它指的是基督的人性，第二种则教导他与父同等，以及他与人的相似。那对亚略派施加的斥责，是因他们试图对圣子所施加的侮辱，并且解释了圣子可以说\Quote{因为}父而生活的意义，以及生命与神圣生命的合一。另一个基于圣子祈求父光荣他的异议，被简要驳斥。}\par

118. 有不少人提出这进一步的异议，即经上记载：\BibleQuote{就如那生活的父派遣了我，我因父而生活；照样那吃我的人，也要因我而生活。}\BibleRef{若 6:58} 他们问：\Quote{圣子既说\InnerQuote{我因父而生活}，又怎能与父同等呢？}\par

119. 让那些以此理由反对我们的人首先告诉我们，圣子的生命是什么？难道是父赐予一个缺乏生命者的生命吗？但圣子自己就是生命，又怎会一度缺少生命？正如他所说：\BibleQuote{我是道路、真理、生命。}\BibleRef{若 14:6} 诚然，他的生命是永恒的，正如他的能力是永恒的一样。那么，曾有一段时间（可以这样说）生命自己未曾拥有自己吗？\par

120. 请思考今天所读的关于主耶稣的话：\BibleQuote{他为我们死了，为叫我们不论醒寤或睡眠，都同他一起生活。}\BibleRef{得前 5:10} 那死亡即是生命的他，他的天主性不就是生命吗？因为天主性就是永恒的生命。\par

121. 但是，他的生命果真在父的权下吗？他表明，即便是他肉身的生命，也不在任何其他人的权下，正如我们有记录的：\Quote{我舍掉我的性命，为再取回它来。谁也不能夺去我的性命，而是我甘心情愿舍掉它；我有权舍掉它，我也有权再取回它：这是我由我父所接受的命令。}\par

122. 那么，他的神圣生命是否应被视为依赖另一位的权能，既然他肉身的生命只服从他自己的权能，而不服从任何其他权能？因为若是权能不是统一的，那就会成为另一位的权能。但正如他让我们明白，他舍掉自己的生命是出于他自己的权能，出于他自己的自由意志，同样他也教导我们，因服从父的命令而舍掉它，这显示了他自己的旨意与父的旨意是统一的。\par

123. 那么，既然既不存在某时刻圣子的生命有了开始，也不存在它所屈从的权能，让我们思考他说：\Quote{就如那生活的父派遣了我，我因父而生活}，其意义是什么？让我们尽力阐明他的意思；不，还是让他自己来阐明吧。\par

124. 那么，请注意祂在前面的讲论中说过什么。\BibleQuote{我实实在在告诉你们。}祂先教导你们应如何聆听。\BibleQuote{我实实在在告诉你们：你们若不吃\OriginalTerm{人子}{Son of Man}的肉，不喝祂的血，在你们内便没有生命。}\BibleRef{若 6:54}祂首先申明祂是作为\OriginalTerm{人子}{Son of Man}说话；那么，你们是否认为，祂作为\OriginalTerm{人子}{Son of Man}所说的、关于祂的\OriginalTerm{肉}{Flesh}和\OriginalTerm{血}{Blood}的话，应该应用到祂的天主性上？\par

125. 接着祂又说：\BibleQuote{因为我的肉，真是食粮；我的血，真是饮品。}\BibleRef{若 6:56}你们听见祂谈论自己的\OriginalTerm{肉}{Flesh}和\OriginalTerm{血}{Blood}，你们看到了神圣的抵押——即向我们传达主的死亡之功绩与能力——\BibleRef{若 6:52}你们却侮辱祂的天主性。听听祂自己的话：\BibleQuote{鬼神没有肉和骨头。}\BibleRef{路 24:39}我们每次领受\OriginalTerm{圣事元素}{Sacramental Elements}，就是那些因圣善祈祷的神秘效能而转变为\OriginalTerm{肉}{Flesh}和\OriginalTerm{血}{Blood}的元素，便\Quote{宣示主的死亡}。\par

126. 然后，在呼唤我们注意祂是作为\OriginalTerm{人子}{Son of Man}说话，并一再重复提及祂的\OriginalTerm{肉}{Flesh}和\OriginalTerm{血}{Blood}之后，祂又加上：\BibleQuote{就如那生活的父派遣了我，我因父而生活；照样，那吃我的人，也要因我而生活。}那么他们以为我们应该如何理解这些话呢？——因为这比喻可以呈现为双重比喻。第一个比喻是这样：\Quote{就如那生活的父派遣了我，我因父而生活；}第二个比喻是这样：\Quote{就如那生活的父派遣了我，我因父而生活，照样，那吃我的人，也要因我而生活。}\par

127. 如果我们的对手选择前一种，意思就是：\Quote{正如我是父所派遣，且从父而来，照样我因父而生活。}但祂是以什么身分被派遣并降来呢？岂不是作为\OriginalTerm{人子}{Son of Man}，正如祂自己先前所说的：\BibleQuote{没有人上过天，除了那自天降下的人子。}\BibleRef{若 3:13}那么，正如祂是作为人子被派遣而降来，照样祂作为人子因父而生活。而且，那吃祂的人，既然是吃人子，自己也因人子而生活。这样，祂将自己的\OriginalTerm{降生成人}{Incarnation}的效应与祂的降来作了比较。\par

128. 但如果他们选择第二种方式，我们岂不既推论出圣子与圣父的同等，又推论出祂与人的相似，尽管二者有明确区分？因为\Quote{正如祂自己因父而生活，照样我们也因祂而生活}这句话，除了表明圣子怎样赐人生命，就如圣父在圣子内怎样赐予人性生命之外，还有什么意思呢？\BibleQuote{就如父唤起死者，使他们复生，照样，子也使祂所愿意的人复生，}\BibleRef{若 5:21}正如主自己先前所说的。\par

129. 因此，圣子与圣父的同等，是建立于赐予生命行动的合一上，因为圣子怎样赐生命，圣父也怎样赐生命。所以，要承认祂生命和王权的永恒性。再者，我们与圣子的相似，以及我们在肉身内与祂的某种合一，也被揭示出来，因为就如同圣子在肉身内由圣父赐予生命，照样人也由圣子赐予生命；因为经上这样记载：正如天主使耶稣基督从死者中复活，照样我们这些人，也因圣子而获得生命。\BibleRef{罗 4:24}\par

130. 按照这一解释，那么，永生不仅借着恩宠的慷慨而应用于我们的状况，而且也被宣认为\OriginalTerm{天主性}{Godhead}的专有属性——后者，是因为是天主性赐予生命；前者，是因为人性在基督内被赐予生命。\par

131. 但如果有人把这两种比喻中的任何一种力量应用到基督的\Emph{\OriginalTerm{天主性}{Godhead}}上，那么圣子就被置于与人类相同的立足点，以至于圣子因圣父而生活，正如我们因圣子而生活一样。但是，圣子是以白白恩赐的方式赐予永生，我们却不能这样做。因此，如果祂被置于与我们同等的地位，祂也就不会赐予这恩赐。那么，让亚略的弟子们获得他们信仰的应得报酬——那就是，不能由圣子获得永生。\par

132. 现在我愿进一步探讨。如果我们的对手乐于将此段教导应用于\OriginalTerm{天主性体}{Divine Substance}的永恒性原则，让他们听听第三种解释：我们的主不是明明在说，正如圣父是生活的，照样圣子也生活吗？——谁又能看不出，这里我们必须理解为指生命的合一，因为这同一生命既是圣父的生命，也是圣子的生命？\BibleQuote{就如父在自己内有生命，照样，祂也赐给子在自己内有生命。}\BibleRef{若 5:26}祂赐给——是由于与祂的合一。祂赐给，不是为取走，而是为在圣子内受光荣。祂赐给，不是为让圣父看守它，而是为让圣子拥有它。\par

133. 但亚略派认为，他们必须用祂曾经说过的话来反对这一点：\BibleQuote{我因父而生活。}当然（假设他们将这话视为指祂的天主性），圣子因圣父而生活，因为祂是由圣父所生的圣子——因圣父，因为祂与圣父同一性体——因圣父，因为祂是从圣父心中发出的圣言，因为祂是从圣父而来的，因为祂是由\Quote{父的怀抱}所生的，因为圣父是圣子的生命的泉源和根由。\par

134. 但也许他们会争辩说：\Quote{如果你们认为，子说\InnerQuote{我因父而生活}，是指存在于父与子之间的生命合一，那么，祂说\InnerQuote{那吃我的人，也要因我而生活}，岂不是揭示了子与人类之间的生命合一吗？}\par

135. 正是如此。正如我承认，因天主性体的同一性，圣父圣子内存在着属天生命的合一；同样，除了那些属于天主本性的特恩、或属于我们的主降生成人的效果之外，我断言，圣子因祂人性与我们人性的同一，而与我们分享灵性生命，因为经上说：\BibleQuote{那属天的怎样，凡属天的也怎样。} 再者，正如我们在祂内坐在圣父的右边，这并不意味着我们分享祂的宝座，而是说我们安息在基督的身体内——同样，我要说，我们因身体上的同一而有分于基督的坐位，照样，我们因我们身体与祂身体的同一，而在基督内生活。\par

136. 那么，我不单不怕\BibleQuote{我因圣父生活}这节经文，即使基督说过\BibleQuote{我因圣父的帮助而生活}，我也不应有所畏惧。\par

137. 他们常提出的另一个反对意见，来自这节经文：\BibleQuote{这病不至于死，乃是为天主的光荣，叫天主子因此受到光荣。} \BibleRef{若 11:4} 但圣子不仅借着圣父并由圣父受到光荣，如经上所载：\BibleQuote{圣父啊，光荣我罢！} \BibleRef{若 17:5} 又说：\BibleQuote{如今人子受到了光荣，天主也因祂受了光荣，天主又光荣祂，} \BibleRef{若 13:31} 而且圣父也借着子并由子受到光荣，因为真理曾说过：\BibleQuote{我在地上，已光荣了祢。} \BibleRef{若 17:4}\par

138. 因此，正如圣子是借着圣父受光荣，祂同样也因圣父生活。有些人由这些话引出一个想法，认为[希腊文] \OriginalTerm{\GreekText{δόξα}}{\GreekText{δόξα}} 的意思是\Quote{意见、信念}，而不是\Quote{光荣}，因而他们如此解释说：\Quote{我在地上已给了你们一个\OriginalTerm{\GreekText{δόξα}}{\GreekText{δόξα}}，我完成了祢所托付我作的工；圣父啊，如今求祢也赐给我一个\OriginalTerm{\GreekText{δόξα}}{\GreekText{δόξα}}；}这也就是说：\Quote{我教导人关于祢的信仰，使他们知道祢是真实的天主；求祢也在他们心中，巩固对我的信仰，即我是祢的儿子，也是真天主。}\par

\Emph{亚略派试图从宗徒的教训——即万物\Quote{本于}圣父、\Quote{借着}子——来证明圣父与子之间有某种特别的区别，这一企图被推翻了；因为所引证的经文显示，同样的全能归于圣父与子，这可由多处圣经经文证明，尤其是由\Person{圣保禄}自己的话，异端者竟愚昧地以为这些话仅指圣父而言，虽然它并不能由此证明子的主权有任何减损或低劣。最后，\Quote{本于谁}、\Quote{借着谁}、\Quote{在谁内}这三个短语，显示出不涉及任何（权能上的）差别，且每一个短语对圣三中的每一位都同样适用。}\par

139. 现在我们来看有些人试图用宗徒的证言，来证明圣父与子之间存在权能差异的那种可笑方法，经文写说：\BibleQuote{为我们，却只有一个天主，就是圣父，万物都本于祂，我们也在祂内；并只有一个主，就是耶稣基督，万物都是借着祂，而我们也借着祂。} \BibleRef{格前 8:6} 他们坚称，断言万物\Quote{本于}圣父而\Quote{借着}子，这表示天主尊威的程度有不可谓不大的差别。然而，最清楚不过的是，这里明显给出了圣子全能的理由，因为万物尽管都\Quote{本于}圣父，却也无不\Quote{借着}子。\par

140. 圣父并不在万物\Quote{当中}，因为向祂承认说，\BibleQuote{万物都事奉祢。}子也不算在万物\Quote{当中}，因为\BibleQuote{万物是借着祂造成的，} \BibleRef{若 1:3} 且\BibleQuote{万物都在祂内而存在，祂在万有之上。} \BibleRef{哥 1:17} 因此，子并非在万物\Emph{当中}，而是在万物\Emph{之上}，诚然，按肉身说，祂出自那民族，出自犹太人，但同时祂是万有之上的天主，永远受赞美，祂拥有一个超越一切名字的名号，\BibleRef{斐 2:9} 论到祂说：\BibleQuote{祢将万物都屈伏在祂的脚下。}但在使\Emph{万物}都服从祂时，没有留下一样不顺服的，如宗徒所说。\BibleRef{希 2:8} 然而，假设宗徒的话是指降生成人的主说的；那么，我们怎能怀疑祂\OriginalTerm{神性受生}{Divine Generation}的无与伦比的尊威呢？\par

141. 那么，确实无疑的是，圣父与子之间不能有权能的差异。不但绝无此差异，同一位宗徒更说，万物\Quote{本于}那位使万物得以存在的祂，如下：\BibleQuote{因为万物都本于祂、借着祂、并在祂内。} \BibleRef{罗 11:36}\par

142. 现在，若如他们所设想的，这里唯独指圣父而言，那么，祂不能同时因万物是本于祂就说是全能的，又因万物是借着祂就不说是全能的。按他们自己的论证，他们要么宣称圣父缺乏权能，且非全能；要么至少他们亲口承认，尽管并非情愿，圣子与圣父同样是全能的。\par

143. 不过，让他们自行决定是否愿意将这话理解为指圣父而言。若他们如此决定，那么万物也是\Quote{借着}祂。若他们决定这里指圣子，那么万物既是\Quote{本于}圣子，也是\Quote{本于}圣父。但如果万物也是\Quote{借着}圣父，那么自然就没有理由减损圣子应受的尊荣；而如果万物是\Quote{本于}圣子，圣子就应该像圣父一样受到尊荣。\par

144. 以免对手怀疑我们利用了某节窜入经文的伪造字句，让我们回顾整个段落。宗徒感叹道：\BibleQuote{啊，天主的智慧和知识的富饶，何其高深！祂的决断何其不可测度，祂的道路何其不可探寻！谁曾知道主的心意？谁曾作过祂的谋士？谁曾先给了祂什么，而需要偿还呢？因为万物都出于祂，藉着祂，并在于祂。愿光荣归于祂，直到永远！}\par

145. 那么，他们认为这里所说的是谁——圣父还是圣子？如果是圣父——那么我们回答，圣父并非天主的智慧，因为圣子才是。然而，对于智慧，有什么是不可能的呢？论到她，经上写着：\BibleQuote{她既是全能的、常在的，在她内便使万象更新。}\BibleRef{智 7:27} 因此，我们读到论智慧的话，不是说她接近，而是说她常在。这样，你们便有了\Person{撒罗满}的权威，教导你们智慧的全能与永恒，以及她的美善，因为经上还写着：\BibleQuote{邪恶不能战胜智慧。}\BibleRef{智 7:30}\par

146. 但是言归正传。宗徒说：\BibleQuote{祂的决断何其不可测度！} 现在，如果\BibleQuote{圣父把一切审判交给了圣子，}\BibleRef{若 5:22} 那么似乎圣父指向圣子作为审判者。\par

147. 但现在，为了向我们表明他说的是圣子，而非圣父，\Person{圣保禄}继续说：\BibleQuote{谁曾先给了祂什么？} 因为\BibleQuote{圣父已交给圣子}，但这是承认祂所生者的权利，而非出于慷慨赐予。因此，不可否认的是，圣子蒙受了圣父的给予，正如经上所记：\BibleQuote{我父将一切交给了我，}\BibleRef{玛 11:27} 但说\BibleQuote{谁曾先给了祂什么？}时，宗徒并未否认圣子因其本性已蒙受圣父的恩赐，但他确实表明，圣父与圣子之间，谁也不可被认为先于对方，因为尽管圣父将恩赐给了圣子，却并非像给一个后于自己而开始存在的那样；因为那非受造且不可测度的天主圣三，具有同一的永恒与光荣，既无时间上的先后，也无地位上的等级之分。\par

148. 然而，如果我们更倾向于遵循那些显示为\Quote{\OriginalTerm{谁先给了祂？}{\GreekText{τίς} \GreekText{προσέδωχεν} \GreekText{αὐτῳ}}}的希腊抄本，那么很明显，那位不能加给任何东西的，并不亚于那位圆满完全的。因此，如果宗徒的这一整段经文，在论及圣子时更易理解，我们便看到，关于圣子我们也必须相信万物都出于祂，正如经上所记：\BibleQuote{因为万物都出于祂，藉着祂，并在于祂。}\par

149. 尽管如此，如果他们仍认为这段经文是指圣父，那么我们要记住，正如我们读到万物都是\Emph{出于}祂，我们也读到万物都是\Emph{藉着}祂，也就是说，圣父与圣子的权柄遍及整个受造的宇宙。尽管我们已经通过圣父的全能证明了圣子的全能，但是——因为他们总是醉心于贬低——让他们想想，他们既贬低了圣父也贬低了圣子。因为如果由于万物都是\Emph{藉着}圣子而成就，圣子的能力就受限制，那么我们是否要进一步说，圣父也同样受限制，因为万物也是藉着圣父呢？\par

150. 但是，为了让他们明白这些说法没有任何区别，我要再次证明，无论是万物的\Quote{出于}还是\Quote{藉着}，都是同一位，而且我们读到这两种关系都可用于圣父。因为我们看到：\BibleQuote{天主是忠信的，你们正是藉着祂被召分享祂圣子的共融。} 让我们的对手掂量宗徒这话的意思。我们\Emph{藉着}圣父蒙召——他们对此没有争议；我们\Emph{藉着}圣子受造——他们却将此定为低一等的标记。圣父召叫我们分享祂圣子的共融，这个真理我们理所当然地虔诚接受。圣子创造了万物，而\Person{亚略}的追随者们却想象这不是出于自由意志的决定，而是被迫的、奴隶般的服务！\par

151. 再者，为了更充分理解：既然我们藉着圣父蒙召进入祂圣子的共融，圣父与圣子的权能并无差别，[请注意] 这共融本身就始于圣子，如经上所记：\BibleQuote{从祂的满盈中，我们都领受了，}\BibleRef{若 1:16} 不过，若依循希腊文福音文本，我们应当译作\Quote{\Emph{出于}祂的满盈}。\par

152. 那么，请看，既有\Emph{藉着}圣父的共融，也有\Emph{出于}圣子的共融，然而并非不同的共融，而是同一的。\BibleQuote{而是为使我们与圣父和祂的子耶稣基督共融。}\BibleRef{若一 1:3}\par

153. 还要进一步留意，圣经不仅说我们在圣父和圣子内，也说我们在圣神内拥有同一的共融。宗徒说：\BibleQuote{愿主耶稣基督的恩宠，和天主的爱，以及圣神的共融，与你们众人同在！}\par

154. 现在，我问，那位万物藉着祂而有的，在哪方面显得比那位万物出于祂的低呢？是因为祂被称作\Quote{行动者}吗？但圣父也工作，因为那位说\BibleQuote{我父到现在一直工作，我也工作}\BibleRef{若 5:17} 的所言是真实的。因此，正如圣父工作，圣子也同样工作；所以这位工作的既非能力有限，也非卑微，因为圣父也工作；既然如此，那些圣子与圣父共有的，或圣子由圣父得来的，都不应受轻视，以免异端者在圣子的位格上进一步羞辱圣父。\par

155. 同样不可忽略，为制止亚略异端的争论，同一位圣若望在另一处圣经所写下的话：\BibleQuote{如果你们知道他是正义的，便该知道，凡行正义的，都是由他而生的。}\BibleRef{若一 2:29}。但谁是正义的呢？不是那喜爱正义的主吗？如之前的经文所警告的，我们若没有子，又能靠谁确保永生呢？因此，如果天主子向我们应许了永生，而他是正义的，我们当然是由他而生的。否则，如果我们的对手否认我们是由圣子藉着恩宠而生，他们也同样否认他的正义。\par

156. 因此，你们必须相信，万物都是出于天主子，[正如出于天主圣父，因为正如天主是万有的父，同样地，圣子也是万有的作者和创造者。那么，我们看到，他们这种质疑是徒劳的，因为就圣子而言[如同就圣父而言]，这话也是真实无误的：\BibleQuote{万物都出于他、藉着他、归于他。}\par

157. 我们已指出，万物如何是\Quote{出于}他，同样也如何是\Quote{藉着}他。那么，既然另一处圣经说：\BibleQuote{因为在诸天内的万有，都是在他内建立，并且在他内受造的，他在万有之先已有，万有都赖他而存在}，谁还怀疑万物是\Quote{在他内}呢？\BibleRef{哥 1:16}。因此，你们由他获得恩宠；你们以他为创造者；在他内，你们找到万有的基础。\par

\Emph{在\BibleBook{若望}{福音}中，将圣子比作葡萄树，将圣父比作园丁的比较，必须参照降生成人的奥迹来理解。若参照天主圣子永恒生成的奥迹来理解，便是对圣子加倍的侮辱，使他低于圣保禄，将他降低到与其余人同等的水平，同样地，也使圣父如此，不仅使圣父与同一位宗徒平起平坐，甚至无足轻重。事实上，就其为天主而言，圣子也是园丁，就其人性而言，则是一串葡萄。这是对圣父之卓越性的正确陈述。}\par

158. 还有另一处圣经，我们的对手通常用来反对我们，以证明他们将圣父的天主性与圣子的天主性分裂开来，那就是我们在福音中主的话：\BibleQuote{我是真葡萄树，我父是园丁。}他们说，葡萄树和园丁本性不同，葡萄树在园丁的权下。\par

159. 那么，你们就是要我们相信，就圣子的天主性而言，他就像葡萄树，没有园丁他便算不了什么，可以被弃置甚或被连根拔除。这样，你们便从经文的字面编造谎言，那经文是说我们的主称自己为葡萄树，其用意在于他降生成人的奥迹。不过，如果你们执意要我们按字面争辩，我也承认，甚至我声明，圣子称自己为葡萄树。因为我若不承认他对其子民救恩的保证，我就有祸了！\par

160. 那么，你们打算如何理解\Quote{天主子称自己为葡萄树}这一真理？如果你们就他天主性的本体来解释这句话，并且你们设想圣父与圣子之间的天主性差异，就像园丁与葡萄树之间本性的差异那样大，那你们便是加给圣父与圣子双重的侮辱——对圣子的侮辱在于：如果如你们所声称的，就他的天主性而言，他在园丁之下，那么按同样说法，他必须被视为低于宗徒保禄，因为保禄确实称自己为园丁，正如经上记载的：\BibleQuote{我栽种了，阿颇罗浇灌了，然而使之生长的却是天主。}\BibleRef{格前 3:6}。难道你们愿意保禄比天主子更优越吗？\par

161. 这是第一重侮辱。至于第二重，则在于：如果圣子就其永恒所生的位格而言是葡萄树，那么，当他说了：\BibleQuote{我是葡萄树，你们是枝条，}\BibleRef{若 15:5}那位神性所生者便显得与我们同一本性。但经上记载：\BibleQuote{上主，众神之中谁能像祢？}\BibleRef{出 15:11}；在\BibleBook{圣咏}{集}中又说：\BibleQuote{在云中谁能与上主相比？天主的众子中谁能相似天主。}\par

162. 此外，你们不仅贬抑圣子，也贬抑圣父。因为如果\Quote{园丁}这一称谓在其名称中包含了圣父主权的一切特恩，那么，鉴于保禄也是园丁，你们便将你们否认圣子与之同等的这位宗徒，置于与圣父同等的地位。\par

163. 再者，经上记载：\BibleQuote{然而，栽种的不算什么，浇灌的也不算什么，只在那使之生长的天主。}你们竟将圣父尊威的全部完满安放于一个你们看到的、代表软弱的名称。因为如果栽种的算不了什么，浇灌的也算不了什么，惟有那使之生长的天主[才是一切]，那么，看你们的亵渎意图何为——竟将圣父置于园丁的名衔之下，使他受轻视，并要求另一位天主来促成圣父的劳作之增长。因此，他们邪恶地以为使用\Quote{园丁}这一称谓是在高举天主圣父的尊威，但在这称谓中，天主圣父却被降低到人的层次，仿佛用一个普通的称呼来指称。\par

164. 然而，如果如你们这些异端者所愿，圣父竟被高举于一位其天主性与人类普通境况毫无差别的圣子之上，这又有什么奇怪呢？如果你们认为圣子之所以被称为葡萄树，是就其天主性而言，那么你们便将他视为不仅可能朽坏、易受风雨变化的影响，甚至只分享了人性，因为葡萄树与枝条属于同一本性，以致天主子看起来并非藉着降生成人的奥迹取了我们的血肉，而是完全由血肉生成。\par

165. 但我确实公开承认，祂的肉身虽以一种新颖而奥秘的方式诞生，却与我们拥有相同的本性，且这是我们救恩的保证，而非\OriginalTerm{圣子之生}{Divine Generation}的根源。祂确实是葡萄树，因为祂承受我的苦难，每当迄今为止脆弱的人性依靠祂，便以更新的虔诚结出丰硕的果实，得以成熟。\par

166. 然而，如果园丁的能力吸引你，请告诉我，在先知书中说话的是谁，说：\BibleQuote{上主啊，求你让我知道，我便知道；那时你使我看见他们的作为。我像一只无害的羔羊被牵往宰杀，我却不知道；他们合谋害我，说：\InnerQuote{来，我们将木头丢进他的饼中。}}\BibleRef{耶 11:18} 因为如果圣子在此谈论祂即将来临的\OriginalTerm{降生成人}{Incarnation}的奥秘——因为若以为这些话是指着圣父说，则是亵渎——那么，在更早的经文中说话的就必然是圣子：\BibleQuote{我曾栽种你如一棵结果累累的葡萄树——你怎么竟变成苦树，成了野葡萄？}\BibleRef{耶 2:21}\par

167. 因此你看到，圣子也是园丁——圣子，与圣父同一名号，同一工程，同一尊威及\OriginalTerm{本体}{Substance}。那么，如果圣子既是葡萄树又是园丁，我们便显然推断出葡萄树的意义关系到\OriginalTerm{降生成人}{Incarnation}的奥秘。\par

168. 然而，我们的主不仅称自己为葡萄树——祂也藉先知的口，给自己取了葡萄串的名号——就是当梅瑟奉上主的命令，派侦探到葡萄串山谷的时候。\BibleRef{户 13:24} 那个山谷岂不就是\OriginalTerm{降生成人}{Incarnation}的谦卑与\OriginalTerm{受难}{Passion}的丰饶？我确实认为，祂被称为葡萄串，是因为从埃及带出的葡萄树，即犹太民族，为世界的益处结出了果实。诚然，没有人能将葡萄串理解成\OriginalTerm{圣子之生}{Divine Generation}的表征——若有谁如此理解，他们便别无结论，只能相信葡萄串是从葡萄树生出的。于是，他们愚蠢地将自己拒绝相信属于圣子的事情归于圣父。\par

169. 然而，如今若毫无疑问，天主子之所以被称为葡萄树，正是着眼于并指向祂的\OriginalTerm{降生成人}{Incarnation}，那么你便明白，我们的主在说：\BibleQuote{父比我大。}\BibleRef{若 14:28} 时，所指的是何种隐藏的真理。因为在说了这话之后，祂紧接着说：\BibleQuote{我是真葡萄树，我父是园丁，}好让你知道，父之所以更大，在于祂栽培照料我们主的肉身，如同园丁栽培照料他的葡萄树。此外，我们主的肉身能随年龄增长身材，并能因受苦而受伤，为使全人类得以在这十字架的荫庇下安息，免受尘世享乐的毒热侵袭——十字架上伸展着祂的肢体。\par

\SourceRecord{Translated by H. de Romestin, E. de Romestin and H.T.F. Duckworth.；Nicene and Post-Nicene Fathers, Second Series；Vol. 10.；Edited by Philip Schaff and Henry Wace.；Buffalo, NY: Christian Literature Publishing Co.,；1896.；<http://www.newadvent.org/fathers/34044.htm>.}

% structure: p0001 source=h1 target=section reason=page-title

\section{基督信仰释义，第五卷}

% structure: p0002 source=h3 target=omitted reason=duplicate-page-title

\Emph{谁是忠信而明智的仆人？他的赏报在\Person{伯多禄}的事例中，也在\Person{保禄}的事例中指明了。\Person{安博罗削}渴望遵循\Person{保禄}的指引，希望这本书能加到其他书里，因为它无法包含在前一卷中。然后陈述了讨论的主题，并给出了进行这一讨论的理由。他必须得到宽恕，因为每个仆人都应当为他所受托的钱财被要求交出利息。他们的忠信正是他本人所期望的利息。他若能指望报酬，便是有福的；但他并不那么期盼圣人们的赏报，而是更期望免于惩罚。他敦促所有人都寻求配得此恩。}\par

1. \BibleQuote{那么，谁是忠信而明智的仆人，主人派他管理自己的家人，按时配给他们食粮呢？主人来时，看见他如此行事，那仆人才是有福的。} \BibleRef{玛 24:45} 这个仆人并非无足轻重：他该是一位大人物。让我们想想他可能是谁。\par

2. 这正是\Person{伯多禄}，他由主亲自拣选来喂养祂的羊群，并堪当三次听到这些话：\BibleQuote{你喂养我的羔羊；你牧放我的羊群；你喂养我的羊群。} 于是，通过用信德的食物好好喂养基督的羊群，他抹去了先前跌倒的罪。为此缘故，他三次受命喂养羊群；三次被问是否爱主，好使他能三次宣认那位在祂受难前他曾三次否认的主。\par

3. 那能说这些话的仆人也是有福的：\BibleQuote{我用奶喂养了你们，而没有用饭，因为那时你们还不能吃。} \BibleRef{格前 3:2} 因为他知道如何喂养他们。我们中谁能做到这些呢？我们中谁能真正地说：\BibleQuote{对软弱的人，我就成为软弱的，为赢得那软弱的人}？ \BibleRef{格前 9:22}\par

4. 然而，他虽是这样伟大的人，且被基督拣选来照顾祂的羊群，以坚固软弱的，医治患病的——我说，他在一次规劝之后，便立刻将异端者从所受托的羊栈中赶出去，\BibleRef{铎 3:10} 以免一只迷途之羊的污染像蔓延的疮一样感染整个羊群。他还嘱咐人要避免愚昧的辩论和纷争。\BibleRef{铎 3:9}\par

5. 那么，我们这些无知的人，在久已存在的庄稼地中处于这些新生的莠子之间，又当如何行事呢？\BibleRef{玛 13:25} 我们若缄默不言，便像是让步了；若与他们争辩，又恐怕我们会被视为属血肉的。因为关于这类引起纷争的事，经上记载说：\BibleQuote{主的仆人不可争吵，而应对众人温和，善于教导，凡事忍耐，以温和规劝那些反抗的人。} \BibleRef{弟后 2:24} 又在另一处说：\BibleQuote{若有人想强辩，我们没有这样的惯例，天主的各教会也没有。} \BibleRef{格前 11:16} 为此，我们立意要写些东西，好让我们的著作能不出声地替我们答覆异端者的不虔。\par

6. 因此，我们准备开始这第五卷，\Person{奥古斯都}皇帝陛下。因为第四卷理应以我们对葡萄树的讨论作结，免得我们似乎是把一堆杂乱的主题塞进那卷书，而不是以属灵葡萄园的果实来充满它。另一方面，既然还有如此众多伟大的主题可供讨论，就不宜让信仰的葡萄采收工作半途而废。\par

7. 因此，在第五卷中，我们论述圣父、圣子、圣神不可分割的天主性（但省略了对圣神的全面讨论），因为福音的训导催迫我们，要将那托付给这五卷书的信德五\OriginalTerm{塔冷通}{talent}，如同本金一般，以利息的方式借给人灵；免得主来到时，发现祂的钱埋在地里，便可能对我说：\BibleQuote{可恶懒惰的仆人！你既知道我在没有下种的地方收割，在没有散布的地方聚敛；那么你就该把我的银子交给钱庄里的人，待我回来时，把我的连本带利取回。} \BibleRef{玛 25:26} 或如另一卷书所载：\BibleQuote{等我回来时，可以连本带利取回。} \BibleRef{路 19:23}\par

8. 如果这份长篇大论的冒昧令任何人感到不悦，我请求他们宽恕。我职分的念虑迫使我将自己所领受的托付给他人。\BibleQuote{我们是天上奥秘的管家。} \BibleRef{格前 4:1} 我们是仆人，却并非人人相同。\BibleQuote{但正如主赐给各人的：我栽种，\Person{阿颇罗}浇灌，然而使之生长的却是天主。} \BibleRef{格前 3:5} 因此，每人都该努力，使他能按自己的劳苦领受赏报。\BibleQuote{因为我们是天主的同工，}正如宗徒说的；\BibleQuote{你们是天主的庄田，是天主的建筑物。} \BibleRef{格前 3:9} 所以，那看见自己本金得了利息的人是有福的；那看到自己工作果实的人也是有福的；那\BibleQuote{在信仰的基础上，用金、银、宝石建筑的人，}同样是有福的。\BibleRef{格前 3:12}\par

9. 你们这些听见或读到这些话的人，就是对我们而言的一切。你们是放贷者的利息——是言语的利息，而非金钱的利息；你们是给予农夫的回报；你们是建筑者的金、银和宝石。在你们的功绩里，蕴含着司铎劳苦的主要成果；在你们的灵魂中，闪耀着主教工作的果实；在你们的进步里，闪烁着主的黄金；如果你们持守天主的话，白银就会增加。\BibleQuote{主的言语是纯净的言语，如同在火炉中炼过七次的银子；在地上经过考验，七次精炼。} 因此，你们会使放贷者富有，使农夫出产丰盛；你们将证明总建筑师的手艺。我并不是自夸地说话；因为我渴望的并非自己的利益，而是你们的利益。\par

10. 我多么希望那时能放心地对你们说：\BibleQuote{主，祢给了我五个\OriginalTerm{塔冷通}{talent}，看，我又赚得了另外五个塔冷通；}\BibleRef{玛 25:20}并展示你们诸德的宝贵塔冷通！\BibleQuote{因为我们是在瓦器中存有这宝贝。}\BibleRef{格后 4:7}这些就是主吩咐我们在灵性上用以交易的塔冷通，或者说新约与旧约的两枚钱币，福音中那位\OriginalTerm{撒玛黎雅人}{Samaritan}将它们留给被强盗打伤的人，为治好他的创伤。\BibleRef{路 10:35}\par

11. 弟兄们，我并非出于贪欲而渴求此事，好能管理许多事；你们的长进本身带给我的酬报已足够。但愿我不至被认为不配领受我所得到的！那些为我过于高超的事，就分配给比我更好的人吧。我不要求这些！然而，主啊，愿祢可以说：\BibleQuote{我要给这最后来的，如同给你的一样。}\BibleRef{玛 20:14}让那配得的人领受管辖十座城的权柄吧。\BibleRef{路 19:17}\par

12. 愿他像那写下\OriginalTerm{法律十句话}{Ten Words of the Law}的\Person{梅瑟}一样。愿他像\Person{农}的儿子\Person{若苏厄}一样，他征服了五个国王，并使\OriginalTerm{基贝红人}{Gibeonites}归顺，为预表那位将要来临、与他同名的人，藉着那人的权能，一切肉欲都将被制伏，\OriginalTerm{外邦人}{Gentiles}将要归正，好使他们追随耶稣基督的信仰，而非以往的追求与欲望。愿他像\Person{达味}一样，年轻女子们歌唱着出来迎接他，说：\BibleQuote{\Person{撒乌耳}杀死千千，达味杀死万万。}\BibleRef{撒上 18:7}\par

13. 对我来说，只要不被扔到外面的黑暗中，像那将托付给他的塔冷通埋在地里——可以说就是埋在自己的肉身里——的人一样，也就够了。这正是那\OriginalTerm{会堂长}{ruler of the synagogue}和犹太人其他首领所做的；因为他们将托付给他们的主的话，可说用在了自己肉身的土地上；他们沉溺于肉身的享乐，将天上的托付如同沉入妄自尊大之心的深坑。\par

14. 所以我们不要把主的钱财埋藏在肉身里；也不要将我们的一塔冷通用\OriginalTerm{汗巾}{napkin}包藏起来 \BibleRef{路 19:20}；而要像忠信的兑换银钱的人，以心神和肉身的劳作，以平衡和敏捷的意志，时常称量它，使这话离你不远，就在你的口中，就在你的心里。\BibleRef{申 30:14}\par

15. 这是主的话，这就是那使人得救赎的宝贵塔冷通。这钱币必须时常出现在灵魂的桌案上，好能藉着不断交易，使良币的声音传遍各地，借着它，永生得以购得。\BibleQuote{这就是永生，}全能圣父，祢白白赐予，为使我们认识\BibleQuote{祢，唯一的真天主，和祢所派遣的耶稣基督。}\BibleRef{若 17:3}\par

\Emph{\OriginalTerm{亚略派}{Arians}是何等不虔敬，竟攻击人类幸福之所系。\Person{若望}总是将圣子与圣父合为一体，尤其在他说道：\BibleQuote{使他们认识祢，唯一的真天主，等等}时。因此，在那一处，我们必须也把\BibleQuote{真天主}理解为圣子；因为不能否认祂是天主，也不能说祂是假神，更不能说祂只是名义上的天主。从宗徒的话可以证明这最后一点，我们正确地宣认基督是真天主。}\par

16. \Emph{因此，} 让亚略派看到，他们质疑我们的希望和欲望所归向的，是何等不虔敬。既然他们惯于在这一点上大声叫嚣，说基督与唯一真天主有别，就让我们竭尽所能驳斥他们这亵渎的思想。\par

17. 因为在这一点上，他们更应理解，这正是完美德性的利益和赏报，即这神圣无比、无与伦比的恩赐：我们得以认识基督，与圣父一起，不将圣子与圣父分开；正如圣经也不将他们分开。因为以下的话论及天主的至尊，更倾向于合一而非差异，即认识圣父与圣子，给予我们同样的酬报，同一的尊荣；这赏报唯独那些认识了圣父与圣子的人才能享有。因为正如认识圣父带来永生，认识圣子也同样带来永生。\par

18. 因此，正如那位\OriginalTerm{圣史}{Evangelist}在他敬虔的信德宣认中，一开始就将\OriginalTerm{圣言}{Word}与天主圣父结合在一起，说：\BibleQuote{圣言与天主同在；}\BibleRef{若 1:1}而在这里，当他写下主的话：\BibleQuote{使他们认识祢，唯一的真天主，和祢所派遣的耶稣基督，}\BibleRef{若 17:3}他无疑通过这样的连接，将圣父与圣子结合在了一起，以致无人能将基督作为真天主与圣父的尊威分开，因为结合并不分离。\par

19. 因此，当他说道：\BibleQuote{使他们认识祢，唯一的真天主，和祢所派遣的耶稣基督，}他便终止了\OriginalTerm{萨培里派}{Sabellians}的谬误，也将犹太人——至少那些听他讲话的人——排除在外；如此，前者不至于以为圣父和圣子是同一位，若他没有加上\BibleQuote{基督}，他们可能如此认为；而后者不至于将圣子与圣父分离。\par

20. 但是，我要问，他们为何不愿从已经说过的话中领悟并理解这一点：即正如他宣称父是唯一的真天主，我们同样可以将耶稣基督理解为唯一的真天主？因为若非如此，就无法表达，以免他似乎是在谈论两位神。因为我们并不谈论两位神；然而我们宣认圣子与圣父具有同一的\OriginalTerm{天主性}{Godhead}。\par

21. 因此，我们可否询问，他们凭什么认为在天主性中有区分，以及他们是否否认基督是天主？但他们无法否认。他们否认祂是真天主吗？然而，若他们否认祂是真天主，就让他们说明，他们宣称祂是假天主，还是\OriginalTerm{仅以称号而言的天主}{by appellation}。因为按照圣经，\Quote{天主}一词或指真天主，或仅指\OriginalTerm{称号上的天主}{by appellation}，或指假天主。真天主一如父；\OriginalTerm{称号上的天主}{by appellation}一如诸圣；假天主一如魔鬼和偶像。那么，就让他们说明他们将如何承认及描述天主子。他们是否设想天主之名被假托了；或者，事实上仅仅是天主内居于祂内，恰似仅以\OriginalTerm{称号}{by appellation}而言那样？\par

22. 我不认为他们能说那名号是假托的，因而使自己陷入公然亵渎的恶行；免得他们一方面自投于魔鬼和偶像，另一方面又暗示基督的天主之名是假托而自投于基督。但若他们以为祂被称为天主，是因为有天主性内居于祂内——一如许多圣人那样（因为圣经称那些领受天主圣言的人为神，\BibleRef{若 10:35}）——他们就没有把祂置于其他人之前，反而认为祂可相比；以致他们视祂如同祂准许其他人成为的那样子，正如祂对\Person{梅瑟}说的：\BibleQuote{我已立你作\Person{法郎}的神。}\BibleRef{出 7:1} 因此，在\BibleBook{}{圣咏}中也有话说：\BibleQuote{我曾说过：你们是神。}\par

23. 这些亵渎者的这种想法，\Person{保禄}弃置一旁；因为他说：\BibleQuote{因为虽然有称为神的，或在天上，或在地上}\BibleRef{格前 8:5} 他并没有说\Quote{有神}，而是说\Quote{有称为神的}。但\Quote{基督}，正如所记载的，\Quote{昨日今日常一样。}\BibleRef{希 13:8}\Quote{祂是}，经上说；这就是说，不仅在于名称，也在于真实。\par

24. 经上写得也好：\BibleQuote{祂是昨日今日常一样的}，好使\Person{亚略}的不虔找不到堆积其亵渎的余地。因为他在读到\BibleBook{}{圣咏}第二篇中父对子说：\BibleQuote{你是我的儿子，我今日生了你}，便注意了\Quote{今日}一词，而非\Quote{昨日}，将那论及取我们人性的话归到天主性生育的永恒上去；关于这，\Person{保禄}也在\BibleBook{宗徒}{大事录}中说：\BibleQuote{我们也向你们报告那个恩许，即天主已给我们做子孙的满了这许，叫耶稣复活了，正如圣咏第二篇所记载的：你是我的儿子，我今日生了你。}\BibleRef{宗 13:32}这样，宗徒充满圣神，为摧毁他那狂暴的疯狂，说：\BibleQuote{昨日今日直到永远，常是一样。}\Quote{昨日}是因祂的永恒；\Quote{今日}是因祂摄取了一个人体。\par

25. 所以，基督是，且永远是；因为那自有者常是。基督常是，\Person{梅瑟}论到祂说：\BibleQuote{自有者派遣了我。}\BibleRef{出 3:14} \Person{加俾额尔}曾经是，\Person{辣法耳}曾经是，天使们曾经是；但那些曾经不是的，绝无同等的理由可以说常是。但基督，如我们所读到的，\BibleQuote{\Quote{是}并不是\Quote{是}而又\Quote{非}，在祂内只有\Quote{是}。}\BibleRef{格后 1:19} 因此，唯独天主才有\Quote{是}的特性，祂是永在者。\par

26. 所以，如果他们不敢说祂是\OriginalTerm{称号上的天主}{by appellation}，而说祂是假天主又是一种极度不虔的标志，那么剩下的就是：祂是真天主，与真父并不相异，而是与父同等。祂随心所愿圣化人、使人成义，\BibleRef{罗 9:18}并不是从自身外获取这权能，而是自身内拥有圣化的权能，怎能说祂不是真天主呢？因为宗徒确实称那位按本性就是天主的为真天主，正如经上所载：\BibleQuote{当你们还不认识天主的时候，服事了一些本来不是神的神。}\BibleRef{迦 4:8} 也就是说，那些不能成为真神的，因为这称号绝不因其本性而属于它们。\par

\Emph{既然已经证明子是真天主，且在这一点上不比父低，那么就显示出，当圣经论及父时所用的\OriginalTerm{唯独}{solus}一词，并不排除子；不，这一表述最为适合祂，且只适合祂。天主圣三是唯一的，不是居于万有之中，而是超越万有。子独自行父所行之事，且独享不朽。但我们不能因此就在争论中将祂与父分离。然而，我们可以把这经文理解为降生成人的事。最后，那些想要子与父分离的人，将父排拒于人类救赎的份位之外。}\par

27. 我们在前几卷中已借圣经的章节充分证明，基督是真天主，是的，是极真的天主。因此，如果基督——如所训导的——是真天主，就让我们探究，为什么他们读到唯独父是真天主时，却希望将子与父分开。\par

28. 若他们说唯独父是真天主，他们就不能否认天主圣子独是真理；因为基督就是真理。那么，真理难道就逊于那真者吗？依用词法，人由\Quote{真理}一词而被称为真者，正如智者由智慧而来，义者由正义而来。我们并不认为父与子之间也是如此。因为父一无所缺，父充满真理；而子，正因祂是真理，便与那真者同等。\par

29. 但让他们知道，当他们看到\Quote{独一，}这个词时，圣子绝不可与圣父分离，他们应记得天主曾在先知书中说过：\BibleQuote{我独自展开了诸天。}\BibleRef{依 44:24} 圣父确实未曾离了圣子而独自展开诸天。因为圣子自己，即是天主的智慧，曾说：\BibleQuote{当他布置诸天时，我已在场。}\BibleRef{箴 8:27} 而\Person{保禄}也声明，这句话是指着圣子说的：\Quote{主，你在起初奠定了大地，诸天是你手中的工程。} 因此，无论是圣子创造了诸天——正如宗徒所愿意让人理解的，同时他自己确非离了圣父而独自铺张诸天；或是如\BibleBook{箴}{言}所载：\BibleQuote{主以智慧奠定了大地，以睿智坚定了诸天；}\BibleRef{箴 3:19} 这都证明了：圣父离了圣子并未独自创造诸天，圣子离了圣父也未独自创造。然而，那位铺张诸天的却被称为独一的。\par

30. 为显示我们应当多么清楚地理解\Quote{独一}一词用于圣子（虽然我们绝不该相信他曾做了什么而圣父不知晓），我们在此另有经文，写道：\BibleQuote{他独自展开诸天，步行在海波之上。}\BibleRef{约 9:8} 因为主的福音教导我们，在海上行走的不是圣父，而是圣子，当时\Person{伯多禄}请求他说：\BibleQuote{主，叫我到你那里去！}\BibleRef{玛 14:28} 就连预言本身也为此作证。因为圣\Person{约伯}曾预言主的来临；论到他，他曾确实地说，他将战胜那巨大的\OriginalTerm{里外雅堂}{Leviathan}，\BibleRef{约 41:8} 这也实现了。因为那可怕的\OriginalTerm{里外雅堂}{Leviathan}，即魔鬼，祂在末期藉着自己可钦敬身体所受的苦难，击打了他，打倒了他，将他置于低处。\BibleRef{依 27:1}\par

31. 因此，圣子是唯一的真天主，因为这也归于圣子作为祂独有的权利。对于任何受造物，都不能确切地说他是独一的。一个与其余受造物有联系的受造者，怎能被视作独一而与其他分离呢？因此，人被视为地上一切受造物中的理性存在者，但他并非唯一的理性存在者；因为我们知道天主天上的工程也是理性的，我们承认天使和总领天使是理性的存在者。如果天使是理性的，那么人就不能被称作唯一的理性存在者。\par

32. 但他们说太阳可被称为独一的，因为没有第二个太阳。但太阳本身与星辰有许多共同之处，因为它运行于诸天，它属于那以太的、属天的本质，它是一个受造物，被列入天主的一切工程之中。它与万物一同侍奉天主，与万物一同赞颂祂，与万物一同赞美祂。因此，不能确切地说它是独一的，因为它并未与其余物分开。\par

33. 因此，既然没有任何受造物能与圣父、圣子和圣神的天主性体相比，这天主性体是独一的，并非在万物之中，而是在万物之上（我们关于圣神的声明暂被搁置）；正如圣父被称为唯一的真天主，因为他与别物毫无共通之处；同样，圣子也是唯一真天主的肖像，祂独为圣父的手，祂独为天主的德能与智慧。\par

34. 因此，圣子独自做圣父所行的事；因为经上记载：\BibleQuote{凡我所做的，祂也照样做。}\BibleRef{若 5:19} 既然圣父与圣子的工程是一体的，那么论及圣父和圣子时，说天主单独工作，是恰当的；因此，当我们论及造物主时，我们既承认圣父，也承认圣子。因为确实，当\Person{保禄}说：\BibleQuote{敬拜受造物，胜过敬拜造物主，}\BibleRef{罗 1:25} 时，他既没有否认圣父是造物主——万物都出于祂，也没有否认圣子是造物主——万物都由祂造成。\BibleRef{罗 11:36}\par

35. 这与经上所载的并不矛盾：\BibleQuote{那独享不死的，}\BibleRef{弟前 6:16} 因为那位在自己内有生命的，怎能没有不死呢？祂因本性而有此生命，因本质而有之；祂拥有它，并非作为一种暂时的恩宠，而是借着祂永恒的天主性体。祂不是像仆人一样以赠予的方式拥有，而是由于祂生自父的特有权利，作为同永恒的圣子。祂也如同圣父一样拥有。\BibleQuote{正如父在自己内有生命，同样祂也赐给子在自己内有生命。}\BibleRef{若 5:26} 正如父怎样有生命，祂也同样赐给了子。你们已经知道祂怎样赐予，这并非出于恩宠的慷慨赠与，而是祂出生的奥秘。既然圣父与圣子之间的生命没有差异，怎能认为只有圣父独享不死，而圣子却没有呢？\par

36. 因此，让他们明白，在这段经文中，圣子不应与圣父分离，圣父是唯一的真天主。因为他们不能证明圣子不是唯一而真实的天主，尤其是正如我所说的，在这里也可以推断出基督也是真实而唯一的天主；或者这段经文至少可以部分地理解为指圣父和圣子的天主性体，部分地指基督的降生成人：因为知识若不承认耶稣基督从永恒就是\OriginalTerm{独生天主}{only-begotten God}、真天主之子，并且按肉躯由童贞女所生，便不算圆满。这正是这位圣史在别处教导我们的，他说：\BibleQuote{凡明认耶稣基督是在肉身内降来的神，便是出于天主。}\BibleRef{若一 4:2}\par

37. 最后，这段经文整体教导我们，在此节经文中提及降生成人的圣事并非不当。因为经上记载说：\BibleQuote{父啊！时候来到了，光荣你的圣子罢！}\BibleRef{若 17:1}所以，当祂说时候已到，并祈求受光荣时，怎能设想祂不是按取了我们肉身而说的呢？因为天主性没有固定的时限，永光也不需要光荣。因此，在唯一的真天主，就是圣父内，我们也按天主性的合一，理解唯一的真天主子。而在耶稣基督的名字内，这名是祂由童贞玛利亚诞生时所领受的，我们承认降生成人的圣事。\par

38. 但是，如果他们读到父是唯一的真天主时想要分离圣子，我想当他们读到关于圣子降生成人的话：\BibleQuote{这石头就是你们匠人所弃而不用的，却成了屋角的基石；}并且又读到：\BibleQuote{在天下人间，没有赐下别的名字，使我们赖以得救的。}\BibleRef{宗 4:11}时，他们便会设想父被排除在赐予我们救恩的恩惠之外。然而，没有父就没有救恩，没有圣子也没有永生。\par

\Emph{对于亚略派的反驳，即认为实体的合一引入了两位神，回答是，更多位神更可能从实体的差异中推断出来。此外，他们的指责反归于自己。多重的差异正是为何两个男人不能被称为一个人的原因，尽管所有人个别地都称之为男人，此时指的是本性上的合一。在他们内有单独一个本性，但在天主圣三位格内，却完全是合一。因此，圣子不可与圣父分离，尤其因为他们不敢否认当受敬拜的是祂。}\par

39. 但亚略派坚持以下说法：如果你说，正如父是唯一的真天主，子也是如此，并且承认父与子同属一个实体，那么你引入的不是一位天主，而是两位。因为那些同属一个实体的，似乎不是一位天主，而是两位天主。正如人们说两个人或两只羊或更多，而一个人和一只羊不被称为两个人或两只羊，而是一个人一只羊。\par

40. 这就是亚略派所说的；他们想用这种狡猾的论证捕捉头脑较简单的人。然而，如果我们阅读\WorkTitle{圣经}，就会发现多个的概念更常出现在那些具有不同且相异实体，即 \OriginalTerm{异质}{\GreekText{ἑτερούσια}} 的事物上。这在\Person{撒罗满}的著作中有所阐明，在那一节中他说：\BibleQuote{有三件事令我不可思议，不，是四件我所不知道的事：鹰在天空飞翔之道，蛇在岩石爬行之道，船在海中航行之道，少年人的道路。}\BibleRef{箴 30:18}鹰、船和蛇不属于同一种类和本性，而是属于可区分且不同的实体，但它们却是三样。因此，根据\WorkTitle{圣经}的见证，他们明白自己的论证反对了自己。\par

41. 因此，当他们说父与子的实体不同，且他们的天主性可区分时，他们自己断言存在两位天主。但我们，当承认父与子，并宣布他们仍同属一个天主性时，我们说没有两位天主，而是只有一位天主。我们凭主的话确立了这一点。因为凡是多个，就存在本性或意志和工作上的差异。最后，为使他们因自己的见证而被驳倒，提到两个人：尽管他们因出生权利而具有同一本性，但在时间、思想、工作和处所上，他们是分开的；所以，一个人不能在两个的意义和数目上被提及；因为有差异就没有合一。但天主被称为一位，而圣父、圣子、圣神的荣耀与圆满便这样表达出来。\par

42. 实际上，合一的真理如此真实，以至于当只指出人性或人肉身的本性时，就用\Quote{人}这个单数词来指称许多人，如经上所载：\Quote{上主是我的助佑，我无所畏惧，人能把我怎样？}也就是说，不是指一个人的个体，而是指同一肉身、人本性中的同一脆弱。经上接着又说：\Quote{信赖上主，胜过信赖世人。}这里也不是指某一个具体的人，而是指普遍的境况。然后，紧接着在论及许多人时又补充说：\Quote{信赖上主，胜过信赖王侯。}正如我们已经说过的，当提及\Quote{人}时，是指存在于众人之间的共同本性上的合一；而当提及\Quote{王侯}时，则指出他们不同权力之间的某种区别。\par

43. 在人中间，或在人内，存在某一方面上的合一，或是在爱中，或是在欲望中，或是在肉身中，或是在虔诚中，或是在信仰中。然而，那按照天主性的荣耀包容万有的普遍合一，却唯独属于圣父、圣子和圣神。\par

44. 因此，主在指出人之间的差异时，他们没有任何能趋向不可分之实体合一性的共同点，说：\BibleQuote{在你们的法律上也记载着：两个人的见证是真实的。}\BibleRef{若 8:17}但尽管祂说过：\BibleQuote{两个人的见证是真实的，}当祂论及自己和父的见证时，祂并没有说：\Quote{我们的见证是真实的，因为这是两位神的见证；}而是说：\BibleQuote{我是为自己作证，派遣我来的父也为我作证。}\BibleRef{若 8:18}更早祂还说：\BibleQuote{即使我判断，我的判断也是真实的，因为我不是独自一个，而是有我，还有派遣我来的父。}\BibleRef{若 8:16}这样，无论在这处还是那处，祂都同时指出父与子，但既没有暗示多位，也没有分离他们天主性实体的合一。\par

45. 因此，显而易见，所有同属一个本质的，不可能被割裂，即使它并非单一，而是一体。我所说的单一，即希腊人所谓的 \OriginalTerm{单一性}{\GreekText{μονοτής}}。单一与位格有关；一体则与本性有关。那些本质不同的，通常不称为单独一个，而称为多个；这一点虽已由先知的证言所证明，但宗徒自己也明确地这样说：\BibleQuote{因为虽然有被称为神的，或在天上，或在地上，} \BibleRef{格前 8:5} 那么，你岂不见那些本质不同、不属同一本性真实的，被称为\Quote{神}吗？但圣父与圣子，因同属一个本质，并非两个天主，而是 \BibleQuote{一个天主，就是圣父，万物都出于他；一个主耶稣基督，万物藉他而有。} \BibleRef{格前 8:6} \BibleQuote{一个天主，}他说，\BibleQuote{和一个主耶稣基督；}而前面又说：\BibleQuote{一个天主，不是两个天主；}然后说：\BibleQuote{一个主，不是两个主。} \BibleRef{格前 8:4}\par

46. 因此，复数性被排除，但一体性并未破坏。不过，一方面，当我们读到主耶稣时，正如我已说过的，我们不会将圣父排除在统御的权能之外，因为这一权能是他与圣子共有的；另一方面，当我们读到惟一的真天主，即圣父时，我们也不能将圣子排除在惟一真天主的权能之外，因为这一权能也是他与圣父共有的。\par

47. 让他们说说自己的感受或想法，当我们读到：\BibleQuote{你要朝拜上主你的天主，惟独事奉他。} \BibleRef{玛 4:10} 难道他们认为基督不应受朝拜，也不应受事奉吗？但如果那位朝拜了他的\Place{客纳罕}妇人，\BibleRef{玛 15:25} 得以获得她所求的，而宗徒\Person{保禄}，在自己书信的开端就自认是基督的仆人，得以成为宗徒，\BibleQuote{不是由于人，也不是藉着人，而是藉着耶稣基督，} \BibleRef{迦 1:1} 就让他们说说他们以为当有何结论。他们是宁愿与\Person{亚略}一同结盟行奸，通过否认基督是惟一的真天主，来表明他们认为基督既不应受朝拜，也不应受事奉呢？还是宁愿与\Person{保禄}同行，因为保禄在事奉和朝拜基督时，并未在言语和心里否认那惟一的真天主，即以虔敬的事奉所承认的那一位呢？\par

\Emph{异端者反对说，基督向他的圣父献上朝拜。但相反，我们指出，这必须归因于他的人性，这从对该段经文的检验中便可清楚看出。然而，这也为他的天主性提供了新的见证，正如我们在基督所做的其他行动中常见的那样。}\par

48. 但若有人因经上记载说：\BibleQuote{你们朝拜你们所不认识的，我们朝拜我们所认识的，} \BibleRef{若 4:22} 便说圣子朝拜天主圣父，那么他当考虑这话是在何时说的，是对谁说的，又是回应谁的请求。\par

49. 在本章前数节中，并非没有理由地记载，耶稣因行路疲倦，便坐下休息，并向一位\Place{撒玛黎雅}妇人要水喝；\BibleRef{若 4:6} 因为他那时是以人的身份说话；因为作为天主，他既不可能疲倦，也不可能口渴。\par

50. 因此，当这妇人称他为犹太人，并以为他是先知时，他以犹太人的身份回答她，并以属神的方式教导法律的奥迹：\BibleQuote{你们朝拜你们所不认识的，我们朝拜我们所认识的。} \BibleQuote{我们，}他说；因为他将自己与人联结。但他如何与人联结呢？岂不是按血肉吗？而为表明他是以成为肉身者的身份回答，他又加上：\BibleQuote{因为救恩是出自犹太人。} \BibleRef{若 4:22}\par

51. 但紧接着，他便放下了人的情感，说：\BibleQuote{然而时候要到，且现在就是，那些真正朝拜的人，将以心神以真理朝拜父。} \BibleRef{若 4:23} 他并没有说：\BibleQuote{我们将要朝拜。} 如果他分担我们的服从，他定会这样说。\par

52. 当我们读到玛利亚朝拜他，\BibleRef{玛 28:9} 我们应当明白，在同一本性之下，他既以仆人的身份朝拜，又以主的身份受朝拜，这是不可能的；更好说，作为人，他在人中间朝拜，而作为主，他受仆人们的朝拜。\par

53. 因此，我们在\OriginalTerm{降生成人}{Incarnation}的\OriginalTerm{奥迹}{sacrament}的光照下，诵读并相信许多事。但在我们人性本有的情感中，也可以看见天主的尊威。耶稣因行路而疲倦，为的是使疲倦者得以安歇；当他将要给口渴者属神的饮料时，自己却想喝水；当他将要给饥饿者救恩的食粮时，自己却饥饿；他死去，为能复活；他被埋葬，为能再起；他悬挂在可怕的木架上，为坚固恐惧中的人们；他以密厚的黑暗遮蔽苍天，为能赐予光明；他使大地震动，为使其稳固；他掀动海洋，为使其平息；他开启死人的坟墓，为显示它们是活人的居所；他由童贞女所生，为使人相信他由天主所生；他佯装不知，为教导无知者知晓；作为犹太人，经上说，他朝拜，为使圣子得以作为真天主受朝拜。\par

\Emph{\Person{安博罗削}回答那些紧抓住主对载伯德儿子们的母亲所说的话的人，他指出，这些话是出于慈爱说的，因为基督不愿使她悲伤。作者还提出了这种温柔态度的充分理由。主宁愿将赐予那请求的权力交给圣父，而不愿宣告那是不可能的。然而，基督的这个答复，并不有损于他，这一点从他亲口说的话，以及将这些话与其他经文比较，都得到了证明。}\par

54. \Quote{何以，}他们说，\Quote{天主子如何能是惟一的真天主，与父相似，当他自己对\Person{载伯德}的儿子们说：\InnerQuote{你们固然要饮我的爵；但坐在我的右边或左边，不是我可以给的，而是我父给谁预备了，就给谁？}}？ \BibleRef{玛 20:23} 那么，这便是你们所期望的，关于天主性不平等之证明；尽管在其中，你们更应敬畏主的慈爱，并朝拜他的恩宠；如果，也就是说，你们能领悟天主德能与智慧的深邃奥秘的话。\par

55. 要想想那位为她儿子们、偕同她儿子们提出这请求的妇人。她是一位母亲，因着对儿子们荣誉的焦虑，尽管她的渴望有些过于无节制，却仍可获得宽恕。她是一位年事已高、热忱虔敬、失去安慰的母亲；那时，她本可得到健壮子女的帮助与扶持，却任由孩子们离开她，宁愿自己的儿子因追随基督而得的赏报，胜过自己的欢愉。因为我们读到，他们一蒙主召唤，立刻舍下网和父亲，跟随了祂。\BibleRef{玛 4:22}\par

56. 于是，她多少顺着母亲热忱的虔敬，恳求救主说：\BibleQuote{求祢叫我这两个儿子在祢的国里，一个坐在祢右边，一个坐在祢左边。}\BibleRef{玛 20:21} 这虽是个过失，却是出于母亲慈爱的过失；因为母亲的心不懂得忍耐。她虽热切渴望得到所求，但这渴望是可宽恕的，因为她贪图的不是金钱，而是恩宠。她的请求并非不知羞耻，因为她不是为自己着想，而是为她的孩子着想。请细思这位母亲，思考她吧。\par

57. 但是，如果你认为全能的天主圣父对其独生圣子的爱无关紧要，那么父母对其子女的感情在你看来微不足道，也就不足为奇了。天地的主（按照我们取了血肉之躯和灵魂德能的说法）感到羞惭——我说，祂感到羞惭，并且，用祂自己的话说，心乱不安，以致不愿拒绝一位为自己儿子们请求的母亲，甚至不愿拒绝让她分享祂自己的座位。你有时主张永恒天主的亲生圣子站着服侍，有时又说祂同坐如同侍从，就是说，不是因为尊威同等，而是因为这是圣父的安排；你竟不给与那位真实天主的圣子，明明连祂自己都不愿拒绝给人的恩惠。\par

58. 因为祂想到了那母亲的爱，她用儿子们的赏报来安慰自己的晚年，尽管被母亲的渴望所折磨，却忍受着那些她最心爱的表记不在身边的苦楚。\par

59. 也请想想这个女人，也就是这较软弱的性别，主尚未以自己的苦难坚强她。我说，请想想这第一位女人\Person{厄娃}的后裔，她承受着那传给所有人的无节制情欲的遗毒，沉沦其下；而且，主尚未以自己的宝血救赎她，也尚未以自己的宝血洗净那植入所有人心中的、对无限尊荣的贪求，甚至超越正当界限。因此，这女人因着遗传的犯罪倾向而犯了罪。\par

60. 一个母亲为她的子女争取优先地位（这远比她为自己争取要好得多），这有什么可奇怪的呢？因为我们读到，连宗徒们自己也曾彼此争论，谁应该居首位。\BibleRef{路 22:24}\par

61. 因此，这位医生不该用羞辱的斥责去伤害一位失去一切的母亲，去伤害一颗受苦的心灵，以免当请求已被提出，却被傲慢地拒绝后，她因自己的祈求被指为不合理而悲伤。\par

62. 最后，主知道母亲的爱当受尊重，便没有回答那女人，而是回答她的儿子们，说：\BibleQuote{你们能饮我将要饮的爵吗？} 他们回答说：\BibleQuote{我们能。} 耶稣对他们说：\BibleQuote{我的爵你们固然要饮；但坐在我的右边和左边，不是我可以给的，而是我父给谁预备了，就给谁。} \BibleRef{玛 20:22}\par

63. 主是多么忍耐和良善啊！祂的智慧何其深奥，祂的慈爱何其美好！因为祂想表明，门徒们所求的并非小事，而是他们不能获得的，祂便将自己特有的权利保留给圣父尊荣，毫不怕减损自己的权利：\BibleQuote{祂虽具有天主的形体，并没有以自己与天主同等，为应当把持不舍，}\BibleRef{斐 2:6} 同时，祂也爱自己的门徒，（因为经上记载，\BibleQuote{祂爱他们到底}），\BibleRef{若 13:1} 因此祂不愿显得拒绝祂所爱的人所求的事；我说，这位善而圣的主，宁愿隐藏自己的一些特权，也不愿放弃丝毫祂的爱。\BibleQuote{因为爱德是含忍的，慈祥的，爱德不嫉妒，不求己益，} \BibleRef{格前 13:4}\par

64. 最后，为使你明白，祂说：\BibleQuote{不是我可以给的，} 这话并非软弱的表现，而是体贴的表示；请注意，当载伯德的儿子们没有偕同母亲去请求时，祂丝毫没有提到父；因为经上这样记载：\BibleQuote{不是我可以给的，而是给谁预备了，就给谁。} \BibleRef{谷 10:40} 圣史\Person{玛尔谷}是这样陈述的。但当母亲为儿子们请求时，正如我们在\BibleBook{玛窦}{福音}中所见，祂却说：\BibleQuote{不是我可以给的，而是我父给谁预备了，就给谁。} \BibleRef{玛 20:23} 这里祂加上了：\BibleQuote{我父，} 因为母亲的情感需要更大的体贴。\par

65. 但是，如果他们认为，祂说：\BibleQuote{我父给谁预备了，就给谁，} 是将更大的权能归于圣父，或是从自己身上减损了什么；就请他们说说，他们是否认为，由于圣子在福音中论到圣父说：\BibleQuote{父不审判任何人，} \BibleRef{若 5:22} 便也减损了圣父的权能。\par

66. 但如果我们以为，圣父把一切审判都交给了圣子，以致于祂自己没有了审判——因为祂拥有审判，并且不能失去神圣尊威因其本性所拥有的——这样的想法是不虔敬的，那么，我们同样应该认为，设想圣子不能赐予人所能赚得的，或任何受造物所能领受的，也是不虔敬的；尤其是祂自己曾说：\BibleQuote{我要往我父那里去，你们因我的名无论向祂求什么，我必要践行。}\BibleRef{若 14:12} 因为如果圣子不能赐予圣父所能赐予的，那么真理就是说谎了，并且不能行因祂的名向圣父所求的事。因此，祂并没有说：\BibleQuote{这是我父所预备的，} 为叫人只向圣父祈求。因为祂已声明，凡向圣父所求的一切，祂都要赐予。最后，祂没有说：\BibleQuote{你们无论向我求什么，我必要践行；} 而是说：\BibleQuote{你们因我的名无论向祂求什么，我必要践行。}\par

\Emph{为了更充分地回答上述反对意见，他主张，这一请求，若非本身不可能，基督原本是可以应允的；尤其是因为圣父已将一切审判交给了祂；这恩赐我们应当理解为毫无缺陷地赐予的。然而，他证明这一请求必须算为不可能的事。为使它真的可能，他教导说，基督的回答必须按照祂的人性来理解，并通过对经文的阐释来表明这一点。最后，他再次确认了他对于基督坐席之不可能性所作的答复。}\par

67. 我现在问他们，是否认为载伯德的妻子和儿子们所提出的请求，对于人的处境，或对任何受造物而言，是可能的还是不可能的呢？如果这请求是可能的，那么那位从无中创造万有的主，怎么会没有能力赐给祂的宗徒在祂右边和左边的座位呢？或者，那位由圣父赐予一切审判权的主，怎么会不能判断人的功劳呢？\par

68. 我们很清楚祂是以何种方式赐予的；因为那从虚无中创造万有的圣子，怎么会如同有所匮乏似地接受呢？祂岂不拥有对祂所造之物本性的审判权吗？圣父将一切审判都交给了圣子，\BibleQuote{为使众人都尊敬子如同尊敬父一样。}\BibleRef{若 5:23} 因此，增长的并非是圣子的权能，而是我们对这权能的认识；我们所学到的，并不给祂的存在增添什么，而只为我们自己的益处；好使我们藉着认识天主子，而能获得永生。\par

69. 因此，我们对天主子的认识，关系到祂的尊荣，但这是我们的益处，而非祂的益处；如果有人以为天主的权能因这尊荣而增加，那么他也必须相信天主圣父能得增加；因为正如圣子一样，父也因我们认识祂而受到光荣：正如圣经记载圣子的话：\BibleQuote{我在地上，已光荣了你。}\BibleRef{若 17:4} 所以，如果所祈求的无论如何是可能的，那么圣子当然是有权能应允的。\par

70. 如果他们以为可能，就让他们指出，在人或其他受造物中，有谁坐在天主的右边或左边。因为圣父对圣子说：\BibleQuote{你坐在我的右边。} 因此，若有人坐在圣子的右边，那么圣子（用人能理解的话来说）就介于自己与圣父之间坐着了。\par

71. 那么，人向祂所请求的，是一件为人不可能的事。但祂不愿说人不能与祂同坐；因为祂愿意祂天主性的光荣被遮蔽，并在祂复活之前不显明出来。\BibleRef{玛 17:9} 因为在此之前，当祂在光荣中显于祂的侍者梅瑟和厄里亚之间时，祂曾嘱咐门徒们，不要把所看见的事告诉任何人。\par

72. 因此，如果人或其他受造物不能赚取此事，那么圣子不因祂未将圣父未曾赐给人或其他受造物的，赐给祂的宗徒，就显得权能较小。不然，就让他们说说看，祂赐给了他们中的哪一个。肯定没有赐给天使；圣经论及他们说，所有天使都站在宝座周围。\BibleRef{默 7:11} 同样，加俾额尔也说他侍立着，正如经上说：\BibleQuote{我是站在天主面前的加俾额尔。}\BibleRef{路 1:19}\par

73. 那么，祂既没有将这赐给天使，也没有赐给那些朝拜坐宝座者的长老们；因为他们不是坐在尊威的座位上，而是如圣经所说，在宝座周围；因为还有另外二十四个座位，正如我们在\BibleBook{若望}{默示录}中所看到的：\BibleQuote{我看见那二十四位长老坐在座位上。}\BibleRef{默 4:4} 在\BibleBook{}{福音}中，主自己也说：\BibleQuote{当人子坐在祂光荣的宝座上时，你们也要坐在十二宝座上，审判以色列十二支派。}\BibleRef{玛 19:28} 祂并没有说可以将自己宝座的位分赐给宗徒，而是说还有另外的十二个宝座；然而，我们不应当以为这是指实际的坐席，而是指属神恩宠的美满结局。\par

74. 最后，在\BibleBook{}{列王纪}中，米加雅先知说：\BibleQuote{我看见上主以色列的天主坐在祂的宝座上，天上的万军都侍立在祂左右。}\BibleRef{列上 22:19} 那么，当天使侍立在上主天主的左右，当天上的万军都侍立时，人又怎能坐在天主的右边或左边呢？这些人所领受的德行报酬，乃是相似天使，正如主所说的：\BibleQuote{你们将如天上的天使一样。}\BibleRef{玛 22:30} \BibleQuote{如天使一样，} 祂说，而不是 \BibleQuote{超过天使。}\par

75. 那么，如果圣父所赐的不比圣子多，圣子所赐的也定然不比圣父少。因此，圣子决不可能小于圣父。\par

76. 然而，假设人们有可能获得他们所渴望的；那么，当祂说：\BibleQuote{但坐在我的右边和左边，不是我可以给你们的}\BibleRef{玛 20:23}，这是什么意思？\Quote{我的}指什么？前面祂说：\BibleQuote{你们必要饮我的爵}；随后祂又加上：\BibleQuote{不是我可以给你们的}。前面祂说\Quote{我的}，接下来祂又说\Quote{我的}。祂没有改变。因此，先前的经文告诉我们祂为什么说\Quote{我的}。\par

77. 因为有一个妇人以人的身份请求祂，让她的儿子们坐在祂的右边和左边；既然她以人的身份请求祂，主也就像纯粹的人一样，论到祂的受难回答说：\BibleQuote{你们能饮我将要饮的爵吗？}\BibleRef{玛 20:22}\par

78. 因此，由于祂是依照肉身谈论祂身体的受难，祂愿意表明：依照肉身，祂给我们留下了忍受苦难的榜样和模范；但以祂作为人的身份，祂不能赐给他们分享天上的宝座。这就是为什么祂说：\Quote{不是我的}；正如在另一处祂说：\BibleQuote{我的教训不是我的}\BibleRef{若 7:16}。祂说，这不是按我的肉身说的；因为属神的话不属于肉身。\par

79. 然而，祂多么清楚地显示了祂对所爱门徒的慈爱，祂首先说：\BibleQuote{你们要饮我的爵吗？}因为祂虽不能赐予他们所寻求的，却先给他们提供了别的东西，好能在拒绝他们之前，先说出祂将要分配给他们的；好让他们明白，未能如愿的原因，更多在于他们对祂请求的不当，而不是他们的主不愿施恩。\par

80. 祂说：\BibleQuote{你们必要饮我的爵}；就是说：\Quote{我不会拒绝你们去受我的肉身将要承受的苦难。因为我作为人所接受的一切，你们都可以效法。我已赐给你们苦难的胜利，就是十字架的产业。\InnerQuote{但坐在我的右边和左边，不是我可以给你们的。}}祂没有说：\Quote{不是我可以给的}，而是说：\Quote{不是我可以给你们的}；这意味着，不是祂缺少能力，而是祂的受造物缺乏功德。\par

81. 或者，以另一种方式来理解这句话：\Quote{不是我可以给你们的}，就是说：\Quote{这不是我的事，因为我本是来教人谦逊的；这不是我的事，因为我本是来服事人，而不是受人服事；这不是我的事，因为我施行的是正义，而不是偏爱。}\par

82. 然后，论到父，祂加上一句：\BibleQuote{是为谁预备的}，为表明父也不是单单应允请求，而是按功德赏报；因为天主是\OriginalTerm{不看情面}{respecter of persons}的。\BibleRef{宗 10:34}为此，宗徒也说：\BibleQuote{因为他所预选的人，也预定他们}\BibleRef{罗 8:29}。祂并不是在\OriginalTerm{预知}{foreknow}他们之前就\OriginalTerm{预定}{predestinate}他们，而是预定了那些祂预知其功德的人的赏报。\par

83. 因此，那位妇人的确受到制止，她向主所请求的是不可能的，仿佛是一种特殊的恩宠；而主却凭祂自愿的恩赐，不仅赐给了两位宗徒，也赐给了所有门徒那些祂判定要赐给圣徒的事物；而且无需任何人祈求，正如经上所载：\BibleQuote{你们也要坐在十二宝座上，审判以色列十二支派。}\BibleRef{玛 19:28}\par

84. 因此，尽管我们可以认为那请求是可能的，却没有理由进行错误的攻击。然而，当我读到\OriginalTerm{色辣芬}{seraphim}侍立\BibleRef{依 6:2}时，我怎能想象人能坐在天主子的右边或左边呢？主坐在\OriginalTerm{革鲁宾}{cherubim}之上，正如经上说：\BibleQuote{坐于革鲁宾之上的，求你显示自己。}宗徒又怎能坐在革鲁宾之上呢？\par

85. 我得出这个结论，不是凭我自己的心思，而是凭我主亲口所说的话。因为主后来在将宗徒托付给父时，说道：\BibleQuote{父啊！你所赐给我的人，我愿我在哪里，他们也同我在一起。}\BibleRef{若 17:24}如果祂认为父会将神圣的宝座赐给人，祂就会说：\Quote{我愿我坐在哪里，他们也同我一起坐。}但祂说：\Quote{我愿他们同我在一起}，而不是\Quote{他们同我一起坐}；祂说\Quote{我在哪里}，而不是\Quote{如同我一样}。\par

86. 接下来便是：\BibleQuote{使他们看见我的光荣。}这里祂也没有说：\Quote{使他们拥有我的光荣}，而是\Quote{使他们看见}光荣。因为仆人是看见，主是拥有；达味也曾教导我们说：\BibleQuote{使我得见上主的福乐。}主自己在福音中也启示了，说：\BibleQuote{心里洁净的人是有福的，因为他们要看见天主。}\BibleRef{玛 5:8}\Quote{他们要看见}，而不是\Quote{他们要同天主坐在革鲁宾之上}。\par

87. 所以，他们要停止依照\OriginalTerm{天主性}{Godhead}而轻视天主子，免得他们也轻视圣父。因为谁若对圣子持有错误的信仰，便不能正确地思想圣父；谁若对圣神持有错误的想法，便不能正确地思想圣子。因为在同一尊严、同一光荣、同一爱、同一尊威之处，凡你认为了三位中的任何一位而应减损的，也就同样从其他所有位格减损了。因为那能被你分割并瓜分为不同部分的，永远不可能具有圆满性。\par

\Emph{有人对以下经文提出异议：\BibleQuote{你爱了他们，如同爱了我一样。}为消除这一异议，他首先指出亚略派解释的不虔敬；然后将这些话与其他经文比较；最后，将整段经文纳入考量。由此他得出结论：基督的派遣，虽然要按肉身来接受，却无损于祂。在证明了这一点之后，他便阐述这神圣的派遣是如何发生的。}\par

88. \Person{奥古斯都}皇帝啊，有些人渴望否认\OriginalTerm{天主性体}{divine Substance}的单一性，力图贬低圣父与圣子的爱，因为经上记载说：\BibleQuote{祢爱了他们，如同祢爱了我一样。}\BibleRef{若 17:23}但他们这样说时，除了在天主圣子与人之间采用一种相似性比较之外，还做了什么？\par

89. 人果真能像圣子那样蒙天主所爱吗？父在圣子内所喜悦的。\BibleRef{玛 3:17}祂本身是蒙喜悦的；我们则是藉着祂。凡天主在人身上看到按祂自己肖像所造的自己的圣子，祂便藉着圣子接纳他们进入儿子的恩宠。因此，我们怎样经过相似而趋于相似，也怎样藉着\OriginalTerm{圣子的受生}{Generation of the Son}而蒙召成为嗣子。天主性体的永恒之爱是一回事，恩宠之爱是另一回事。\par

90. 如果他们开始就经上记载的话进行辩论：\BibleQuote{祢爱了他们，如同祢爱了我一样，}并认为这里意在比较，那么他们就必须认为以下的话也是以比较的方式说的：\BibleQuote{你们应当慈悲，就像你们的父是慈悲的一样；}\BibleRef{路 6:36}在另一处又说：\BibleQuote{你们应当是成全的，如同你们的父是成全的一样。}\BibleRef{玛 5:48}但如果说祂是在圆满的光荣中成全，我们则不过是在我们内德增长中的成全。圣子也按那永存之爱的圆满为父所爱，但在我们内，恩宠的增长赢得天主的爱。\par

91. 那么，你既看见天主如何将恩宠赐予人，难道你还想割裂圣父与圣子之间那本性的、不可分割的爱吗？当你注意到提及至尊的统一性时，难道你仍力图使这些话语落空吗？\par

92. 请思考这整段经文，看看祂是从什么立场说话的；因为你听见祂说：\BibleQuote{父啊！请在祢面前光荣我，赐给我在世界未有以前，我与祢共有的光荣罢！}\BibleRef{若 17:5}看看祂如何从第一个人的立场说话。因为祂在那祈求中为我们恳求那些，身为人的祂记得在堕落之前在乐园中曾赐予的事物，正如祂在受难时对右盗论及此事所说的：\BibleQuote{我实在告诉你：今天你就要与我一同在乐园里。}\BibleRef{路 23:43}这就是世界未有以前的光荣。但祂用\Quote{世界}一词代替\Quote{人}，正如你也有这话：\BibleQuote{看！全世界都追随祂去了；}\BibleRef{若 12:19}又说\BibleQuote{为叫世界知道是祢派遣了我。}\BibleRef{若 17:21}\par

93. 然而，为叫你认识伟大的天主，就是赋予生命且全能的圣子，祂加上了一种祂尊威的证明说：\BibleQuote{我的一切都是祢的，祢的一切都是我的。}\BibleRef{若 17:10}祂拥有一切，而你却歪曲祂被派遣的事实，来委屈祂吗？\par

94. 但是，如果你不接受祂按肉身被派遣的真理，一如宗徒所论及的，\BibleRef{罗 8:3}却仅凭一个字眼就下判断反对它，好使你能说下级惯常由上级派遣；那么对于圣子被派遣到人中间这一事实，你将给出什么答复呢？因为如果你认为被派遣的低于派遣他的，那么你也必须知道，下级曾派遣了上级，上级也曾被派遣到下级那里。因为\Person{多俾亚}派遣了总领天使\Person{辣法耳}，\BibleRef{多 9:3}一位天使被派到\Person{巴郎}那里，\BibleRef{户 22:22}天主圣子被派到犹太人中间。\par

95. 难道天主圣子低于祂被派往的犹太人吗？论到祂记载说：\BibleQuote{最后祂打发自己的儿子到他们那里去，说：他们会敬重我的儿子。}\BibleRef{玛 21:37}请注意，祂先提及仆人，然后提及儿子，为叫你明白，天主，按祂天主性的大能而为独生子，与仆人既无共同的名分，也无共同的命运。祂被派遣是为受敬重，不是为与家仆相比。\par

96. 祂加上\Quote{我的}一词，这样我们便能相信，祂来，不是作为众人中的一位，也不是作为具有较低本性或次要权能的一位，而是作为出自真实者的真实者，作为\OriginalTerm{圣父性体的肖像}{Image of the Father's Substance}。\par

97. 然而，假设被派遣的低于派遣他的。那么基督就低于\Person{比拉多}了；因为比拉多把祂送到了\Person{黑落德}那里。但一个字眼并不损害祂的权能。圣经既说祂由父所派遣，也说祂由一位长官所派遣。\par

98. 因此，如果我们明智地持守与天主圣子相称的事，我们就当如此理解祂的被派遣：天主圣言，从祂尊威深处的\OriginalTerm{不可理解且不可言喻的奥迹}{incomprehensible and ineffable mystery}中，将自己交付给我们的理智予以理解，好使我们能领会祂，不但在祂\Quote{\OriginalTerm{虚己}{emptied}}的时候，也在祂住在我们内的时候，正如经上记载：\BibleQuote{我要在他们内居住。}\BibleRef{格后 6:16}在另一处也说，天主说：\BibleQuote{来，我们下去，混乱他们的语言。}\BibleRef{创 11:7}诚然，天主从不从天降下；因为祂说：\BibleQuote{我充满天地。}\BibleRef{耶 23:24}然而，当天主圣言进入我们心中时，祂好似降下，一如先知所说：\BibleQuote{预备主的道，修直祂的途径！}\BibleRef{依 40:3}我们应这样做，为能如祂自己所应许的，祂与圣父同来，并住在我们这里。\BibleRef{若 14:23}那么，很明显祂是怎样来的。\par

\Emph{基督，就其为真天主圣子而言，没有主，只有在祂是人的时候才有；祂的话证明了这一点，在话中祂时而称呼父，时而称呼主。圣经中的一节经文就堵住了多少异端的口！我们必须区分属于基督作为天主圣子与作为达味之子的那些事物。因为只有在后一种称呼下，我们才将祂是仆人的属性归于祂。最后，他指出，许多经文除非是指\OriginalTerm{降生成人}{Incarnation}，否则不能应用。}\par

99. \Emph{因此}，也很清楚祂如何称祂所认识为父的那位为\Quote{主}。因为祂说：\BibleQuote{父啊！天地的主宰！我称谢你，} \BibleRef{玛 11:25} 那原初的智慧先提到自己的父，然后宣称祂是受造物的主。为此，主在祂的福音中指明，凡有真正子嗣之处，便没有主仆关系，祂说：\BibleQuote{论到基督，你们的意见如何？祂是谁的儿子呢？他们回答说：达味的。耶稣对他们说：那么，达味如何因圣神感动称祂为主呢？说：\InnerQuote{上主对我主说：你坐在我右边……}？} 然后祂加上说：\BibleQuote{如果达味因圣神称祂为主，祂怎么又是达味的儿子呢？没有人能回答祂一句话。} \BibleRef{玛 22:42}\par

100. 主为顾及亚略派，在作此见证时，对信德是何等关心啊！因为祂没有说：\Quote{圣神称祂为主，}而是说\Quote{达味因圣神说话；}为使人相信，正如祂按肉身是达味之子，照祂的天主性，祂也是达味的主和天主。那么，你可见用于父子关系和主仆关系的称谓之间是有区别的。\par

101. 主恰当地论及自己的父，而论及天地的主宰；这样，当你读到父和主时，可理解为：父乃是圣子之父，主乃是受造物之主。前者基于本性的要求，后者基于治理的权威。因为祂取了奴仆的形象，所以称祂为主，是因祂服从了服事；在形象上是天主，与天主同等，但在肉身的形象上却是奴仆：因为服事是肉身应做的，而主性是天主性应得的。因此，宗徒也说：\BibleQuote{我们的主耶稣基督的天主和父，光荣的父，} \BibleRef{格后 1:3} 也就是说，称祂为取了人性的天主，又称祂为光荣的父。难道天主有两个儿子，基督和光荣吗？当然不是。因此，如果天主子只有一个，就是基督，那么基督就是光荣。你为何竭力贬低那身为父之光荣的呢？\par

102. 那么，如果圣子是光荣，父也是光荣（因为光荣的父不能不是光荣），光荣便没有分离，光荣乃是一。所以，光荣属于自身特有的本性，而主性则属于所取肉身的服事。因为如果肉身服从义人的灵魂，如经上记载：\BibleQuote{我痛打我的身体，使之服服帖帖，} \BibleRef{格前 9:27} 那么，它岂不更该服从天主性？论到这天主性，经上说：\BibleQuote{万物都要事奉你。}\par

103. 主藉一个问题就同时驳斥了撒伯里乌派、福提努斯派和亚略派。因为当祂说上主对我主讲话时，撒伯里乌便被排除了，他主张父与子是同一位。福提努斯也被排除了，他认为祂仅仅是个人；因为除了天主，无人能作达味王的主，正如经上记载：\BibleQuote{你要敬畏上主你的天主，惟独事奉祂。} \BibleRef{申 6:13} 难道那在法律下执政的先知会违反法律吗？亚略也被排除了，他听到子坐在父的右边；因此，他若用人的方式来论证，便驳倒了自己，使他亵渎论证的毒液回流到自己身上。因为他在以人间习性的类比来解释父与子的不平等时（二者都偏离了真理），反而把被他轻视的那位置于首位，承认自己听到坐在右边的那位是首位的。摩尼教徒也被排除了，因为他并不否认按肉身说，祂是达味之子；当瞎子呼求说：\BibleQuote{耶稣，达味之子，可怜我们吧！} \BibleRef{玛 20:30} 祂因他们的信德而喜悦，就站住治好了他们。但摩尼教徒否认这适用于祂的永恒性，如果祂只被那些虚伪的人称为达味之子。\par

104. 因为\Quote{天主子}驳斥了厄俾敖派，\Quote{达味之子}驳斥了摩尼教徒；\Quote{天主子}驳斥了福提努斯，\Quote{达味之子}驳斥了马西翁；\Quote{天主子}驳斥了萨莫撒塔的保禄，\Quote{达味之子}驳斥了华伦提努斯；\Quote{天主子}驳斥了亚略和撒伯里乌，他们是异教错谬的继承者。\Quote{达味之主}驳斥了犹太人，他们虽亲眼看见降生成人的天主子，却以不虔之疯狂相信祂仅仅是个人。\par

105. 但在教会的信德中，同一位既是天主父之子，又是达味之子。因为天主降生成人的奥迹，是拯救整个受造界的救恩，正如经上记载：\BibleQuote{为了使祂无天主，为人人尝到死味；} \BibleRef{希 2:9} 也就是说，为使一切受造物不用自己受苦，却藉着主天主性的血得以救赎，如经上另一处所说：\BibleQuote{一切受造物都将脱离败坏的控制。} \BibleRef{罗 8:21}\par

106. 按照天主性体被称为子是一回事，按照所取人性的收养被称为子则是另一回事。因为，按照\OriginalTerm{天主性的生发}{the divine Generation}，圣子与天主父同等；而按照所取肉身的收养，祂是天主父的奴仆。因为经上说：\BibleQuote{祂却使自己空虚，取了奴仆的形体，} \BibleRef{斐 2:7} 然而圣子是同一的。另一方面，按照祂的光荣，祂是圣祖达味的主，但按血统却是他的儿子，祂并没有舍弃自己的任何事物，反而为自己获取了由收养进入我们人类而来的权利。\par

107. 祂不仅因降生自\Person{达味}而具有人的身份并为此服役，也因祂的名号而如此，正如经上所载：\BibleQuote{我拣选了我的仆达味；} 又另处载道：\BibleQuote{看，我必要派遣我的仆人旭日，他的名字是旭日。} \BibleRef{匝 3:8} 圣子自己也曾说：\BibleQuote{上主从母胎中形成我作祂的仆人，曾对我说过：你作我的仆人，这是一件大事。看，我立你为人民的盟约，为万民的光明，为作我的救恩直到地极。} \BibleRef{依 49:5} 这话若非对\Person{基督}，又是对谁说的呢？祂虽具有天主的形像，却空虚了自己，取了奴仆的形像。\BibleRef{斐 2:6} 但除了那具有完满天主性的，谁还能有天主的形像呢？\par

108. 那么，你们要明白\Quote{祂取了仆人的形像}是什么意思。那是指祂完全且圆满地取了人性的一切完美，并完全地取了服从。因此，在第三十篇圣咏中记载：\BibleQuote{你使我的脚立于宽广之地。我成了我一切仇敌的耻辱。求你使你的面容光照你的仆人。}\Quote{仆人}是指祂在那人性中而被圣化的人；是指祂在那人性中而被傅油的人；是指祂在那人性中而生于法律之下，生于童贞女的人；简言之，是指祂在那位格中而有母亲的人，正如经上所载：\BibleQuote{上主，我是你的仆人，我是你的仆人，你婢女的儿子；} 又载：\BibleQuote{我极度卑微，遭受磨难。}\par

109. 除了那藉着服从而来解救众人的\Person{基督}，谁又是极度卑微、遭受磨难的呢？\BibleQuote{正如因一人的悖逆，众人都成了罪人；照样，因一人的服从，众人也都要成为义人。} \BibleRef{罗 5:19} 是谁领受了救恩之杯？是\Person{基督}，这位大司祭，而非从未担任过司祭职、也未曾受苦难的\Person{达味}。是谁献上了感恩祭？\par

110. 但这还不够；再看另一处：\BibleQuote{求你保护我的灵魂，因为我是圣的。} \Person{达味}是指自己说这话吗？不，而是那位也说过：\BibleQuote{你绝不会将我遗弃在阴府，也绝不让你的圣者见到腐朽。} 的那一位说的。因此，这两处经文都是指同一位而说的。\par

111. 祂又接着说：\BibleQuote{拯救你的仆人吧；} 再稍后：\BibleQuote{将你的力量赐给你的仆人，赐给你婢女的儿子；} 又在别处，即\BibleBook{厄则克耳}{书}中：\BibleQuote{我要为他们兴起一个牧人，他必牧放他们，就是我的仆人达味；他要牧放他们，作他们的牧人。我上主将作他们的天主，我的仆人达味作他们中间的元首。} \BibleRef{则 34:23} 现在，\Person{叶瑟}之子\Person{达味}早已逝世。因此，这是指\Person{基督}而说的，祂为了我们，以人的形像成为婢女之子；因为按祂的天主性诞生而言，祂无母，只有父：祂不是世间欲望的果实，而是天主的永恒大能。\par

112. 因此，当我们读到主说：\BibleQuote{我的时候还没有到；} \BibleRef{若 7:8} 以及 \BibleQuote{我还有一点点时间与你们同在；} 和 \BibleQuote{我就要回到派遣我来的那里去；} \BibleRef{若 7:33} 以及 \BibleQuote{现今人子受到了光荣；} \BibleRef{若 13:31} 时，我们应将这一切归于\OriginalTerm{降生成人的奥迹}{sacramentum Incarnationis}。但当我们读到：\BibleQuote{天主也在祂身上受到了光荣，天主也光荣了祂；} \BibleRef{若 13:31} 时，这里还有什么疑惑呢？圣子由圣父光荣，圣父由圣子光荣。\par

113. 接着，为了阐明\OriginalTerm{一体}{Unitatis}的信德和\OriginalTerm{天主圣三的合一}{Unionis Trinitatis}，祂还说自己将受到圣神的光荣，正如经上所载：\BibleQuote{祂要由我领受，并传告你们。} \BibleRef{若 16:14} 因此，圣神也光荣圣子。那么，祂又怎会说：\BibleQuote{如果我自己光荣自己，我的光荣便算不了什么。} \BibleRef{若 8:54} 呢？难道圣子的光荣算不了什么吗？除非我们将这话应用于祂的肉身上，否则这样说便是亵渎；因为圣子是以人的身份说话，因为与天主性相比，肉身并无光荣可言。\par

114. 他们应停止那些邪恶的反对意见，这些意见只会反弹到他们自己的虚伪之上。因为他们说，经上记载：\BibleQuote{现今人子受到了光荣。} 我不否认经上记载：\BibleQuote{人子受到了光荣。} 但让他们看看接下来的是什么：\BibleQuote{天主也在祂身上受到了光荣。} 我可以为人子辩解，但祂却不为祂的父辩解；因为圣父并没有取了我们的人性。我可以辩解，但我并不使用这辩解。祂没有辩解，却受到了诬蔑。我既可按字面意义理解，也可将至论及肉身的话应用于肉身上。虔诚的心分辨哪些话是按肉身说的，哪些是按天主性说的。不虔敬的心却将一切涉及肉身卑微的话，都转向对天主性的侮辱。\par

\Emph{圣人驳斥那些以犹太式思维反对\Quote{因父及子及圣神之名}词序的人，反驳说圣子也常被置于圣父之前；不过他首先指出，对此反对意见他自己先前已给出了回答。}\par

115. 为何亚略派人士，仿效犹太人的方式，如此错谬且无耻地解释天主的圣言，甚至竟敢说圣父的能力是一样，圣子的能力是一样，圣神的能力又是一样，既然经上记载：\BibleQuote{你们要去使万民成为门徒，因父及子及圣神之名给他们授洗}？他们为何只因词序的缘故，就对天主的能力加以区分呢？\par

116. 虽然我在前书中已为尊威与名字的统一提供了这同一见证，但如果他们以此作为争论的根基，我仍能凭圣经的证言主张：在许多地方，子是首先被提及的，而父则在祂之后被提及。因此，难道因为子的名字放在前面，就如亚略派所愿的，仅是字词的偶然，父就次于子了吗？天主绝不允许，我说，天主绝不允许。信德不知这种次序；它不知父与子的尊荣被分割。我未曾读过、听过、发现过天主内有任何不同的等级。我从未读过有第二位，或第三位天主。我读到的是第一位天主，\BibleRef{依 44:6} 我听到的是第一位和唯一的天主。\par

117. 如果我们过分注重次序，那么子就不该坐在父的右边，也不该称自己为第一和开端。那位圣史若以\Quote{圣言}而非\Quote{天主}开篇，便是错误的，因为他说：\Quote{在起初已有圣言，圣言与天主同在。} \BibleRef{若 1:1} 因为，按人的惯常次序，他本该先提父。那位宗徒也不晓得他们的次序，他说：\Quote{耶稣基督的仆人保禄，蒙召作宗徒，被选拔为天主的福音服务；} \BibleRef{罗 1:1} 在另一处又说：\Quote{愿主耶稣基督的恩宠，和天主的爱，以及圣神的相通。} \BibleRef{格后 13:14} 如果遵循字词的次序，他把子放在前面，父放在后面。但字词的次序常常改变；因此在论及天主父及其子时，你们不该质疑次序或等级，因为在神性内没有统一的分裂。\par

\Emph{亚略派公然与异教徒站在一起，攻击\Quote{那信我的，不是信我}等语。这段经文的确切含义被揭示出来；为阻止我们认为主禁止我们信仰祂，它说明祂有时以天主身份说话，有时以人身份说话。在列举了那信德的各种结果之例子后，他指出其他一些经文也必须同样理解。}\par

118. 最后，为表明他们不是基督徒，他们否认我们应当信仰基督，说经上记载：\Quote{那信我的，不是信我，而是信那派遣我来的。} \BibleRef{若 12:44} 我一直等着这认信；你们为何用诡辩哄骗我？我早知必须与异教徒争辩。不，他们确实皈依了，你们却没有。如果他们相信，那（圣洗）圣事是稳妥的；你们领受了它，却毁坏了它，抑或许从未领受，而是自始便是虚假的。\par

119. 他们说，经上写着：\Quote{那信我的，不是信我，而是信那派遣我来的。} 但是请看下文，并看天主子愿意如何被人看见；因为接着说：\Quote{谁看见我，就是看见派遣我来的，} \BibleRef{若 12:45} 因为父在子内被看见。这样，祂解释了先前所说的，即那承认父的，就信了子。因为谁不认识子，也不认识父。凡否认子的，就没有父；但谁明认子，就有父也有子。\BibleRef{若一 2:23}\par

120. 那么，\Quote{不是信我}是什么意思呢？这就是说，不是信你们能以形体感知的，也不是仅仅信你们所见的人。因为祂已说明，我们不应只信一个人，而是要使你们相信耶稣基督自身既是天主又是人。因此，基于这两重理由，祂说：\Quote{我不是由我自己来的；} \BibleRef{若 7:28} 又说：\Quote{我是元始，正是我给你们所讲的。} \BibleRef{若 8:25} 作为人，祂不是由自己来的；作为天主子，祂不以人为元始；而是，\Quote{我是，}祂说，\Quote{正是那\InnerQuote{我给你们所讲的元始}。我所说的话，也不是属人的，而是属神的。}\par

121. 也不可相信祂否认了我们应该信仰祂，因为祂自己说：\Quote{好使信我的，不住在黑暗中；} \BibleRef{若 12:46} 在另一处又说：\Quote{因为这是我父的旨意：凡看见子，并信从子的，必获永生；} \BibleRef{若 6:40} 又说：\Quote{你们信天主，也当信我。} \BibleRef{若 14:1}\par

122. 所以，谁也不可只接受子而不要父，因为我们读到的是子。子有父，但不是时间意义上的，也不是因祂的苦难，或因祂的受孕，或靠恩宠。我读到的是祂的\Quote{生发}，而不是祂的\Quote{受孕}。父说：\Quote{我生了；}祂没有说：\Quote{我造了。}子在祂神圣永恒的生发中，不称天主为创造者，而称为父。\par

123. 祂时而又以人的身份展现自己，时而在天主的尊威中；时而宣称自己与父有同一的神性，时而承担人性肉身的一切脆弱；时而说祂没有自己的教训，时而说祂不寻求自己的旨意；时而指出祂的证据不真，时而说它是真的。因为祂自己说过：\Quote{我若为我作证，我的证据不足凭信。} \BibleRef{若 5:31} 后来祂又说：\Quote{我即使为我自己作证，我的证据也是真实的。} \BibleRef{若 7:14}\par

124. 主耶稣，祢的证据怎么会不真呢？难道那个相信祢的人，虽然他悬在十字架上，为他承认的罪行受罚，却丢弃了强盗应有的报应，而获得了无辜者的赏报吗？\BibleRef{路 23:41}\par

125. 保禄受骗了吗？他因信而恢复了视力；\BibleRef{宗 9:12} 而他先前在不信时是失明的。\par

126. \Person{农的儿子若苏厄}在认出天军统帅时犯了错误吗？\BibleRef{苏 5:13} 但他一相信，就立即得胜，被认定堪当在信德的战斗中凯旋。再者，他没有率领武装队伍出战，也没有用攻城锤或其他战争机械摧毁敌人城墙的壁垒，而是靠着司祭的七支号角声。如此，号角的响声和司祭的标记终结了一场残酷的战争。\par

127. 一个妓女看见这事；她在城被摧毁时已失去任何安全得救的希望，但因她的信德已得胜，就在窗户上系了一条朱红线，如此举起了她信德的标志和\OriginalTerm{主的受难}{Lord's Passion}的旗帜；\BibleRef{苏 2:18} 好使那将要救赎世界的\OriginalTerm{神秘的血}{mystic blood}的象征，存留于记忆中。因此，在外，\Person{若苏厄}的名字是战斗者的胜利标志；在内，\OriginalTerm{主的受难}{Lord's Passion}的象征是危急者的救恩标志。因此，因为\Person{辣哈布}明白了天上的奥秘，主在\BibleBook{}{圣咏}中说：\BibleQuote{我要记念认识我的辣哈布和巴比伦。}\par

128. 主啊，那么，你的见证如何不真实呢，除非那是按照人的脆弱而给予的？因为\BibleQuote{人人都是说谎者。}\par

129. 最后，为了证明他说话是凭人性，他说：\BibleQuote{派遣我来的父，他为我作证。}\BibleRef{若 8:18} 但他作为天主的见证是真实的，正如他自己说的：\BibleQuote{我的证据是真实的：因为我知道我从哪里来，往哪里去；你们却不知道我从哪里来，往哪里去。你们只凭肉身判断。}\BibleRef{若 8:14} 因此，那些认为基督没有作证能力的人，不是凭\OriginalTerm{天主性}{Godhead}，而是凭\OriginalTerm{人性}{manhood}在判断。\par

130. 因此，当你听到\BibleQuote{信的人，不是信我；}或：\BibleQuote{派遣我来的父，他给了我命令；}\BibleRef{若 12:49} 你现在已经知道应当把这些话归于何处。最后，他揭示了那命令是什么，说：\BibleQuote{我舍掉我的性命，为再取回它来。谁也不能夺去我的性命，而是我甘心情愿舍掉它。}\BibleRef{若 10:17} 那么，你看到，这话是要表明他有全权舍掉或取回他的性命；正如他还说过：\BibleQuote{我有权舍掉它，我也有权再取回它来：这是我由我父所接受的命令。}\BibleRef{若 10:18}\par

131. 因此，无论是一个命令，还是如一些拉丁抄本所载的一个指示，都绝不是赐给身为天主的他，而是赐给降生成人的他，关乎他通过受难受死将要赢得的胜利。\par

\Emph{我们必须把基督被说成不凭自己说话的事实，归于他的人性。在解释了为何说作为天主的他听见并看见父是恰当的之后，他通过大量的证据，确凿地证明天主子不是受造物。}\par

132. 难道我们竟把天主子贬低到如此卑微的地步，以致若非听见，他就不知道如何行动或说话吗？难道因为经上记载：\BibleQuote{我不凭自己说话}和下文：\BibleQuote{父怎样对我说，我就怎样说}\BibleRef{若 12:50}，我们就要设想一个固定的行动或说话的度量被分配给他吗？然而，那些话是指肉身的服从，或者是指向\OriginalTerm{合一}{Unity}的信德。因为很多博学的人承认，子听见，父藉着\OriginalTerm{性体}{Nature}的合一与子说话；因为子通过意志的结合，知道父所愿意的，他似乎就是听见了。\par

133. 这里指的不是个人的职责，而是一个不可分割的合作定旨。因为这并不表示任何实际听见话语，而是意志与能力的合一，这既存在于父内，也存在于子内。他在另一处指出，这也存在于圣神内，说：\BibleQuote{因为他不凭自己说话，只要他听见什么，就说什么}\BibleRef{若 16:13}，好叫我们明白：圣神说什么，子也说什么；子说什么，父也说什么。因为在\OriginalTerm{天主圣三}{Trinity}内，只有一个意念和一种运作模式。因为，正如父在子内被看见，不是以形体显现，而是在\OriginalTerm{天主性}{Godhead}的合一中，同样，父也在子内说话，不是用尘世的声音，不是用人的声响，而是在他们工作的合一中。所以当他曾说：\BibleQuote{住在我内的父，他说话；我所做的工作，他做}\BibleRef{若 14:10} 之后，他又加上：\BibleQuote{你们要相信我：我在父内，父在我内；不然，你们至少该因那些工作而相信我}\BibleRef{若 14:17}。\par

134. 这是按照全部圣经所教导的我们理解的方式；但那些不愿对天主怀正确想法的\OriginalTerm{亚略异端者}{Arians}，可以用一个正合适他们应得惩罚的例子驳倒他们；以便他们不要以属肉体的方式相信一切，因为连他们自己也不能以肉眼看到他们的祖宗魔鬼的工作。所以主就指着他们的同胞犹太人说：\BibleQuote{你们所做的，是你们由你们的祖宗那里所看见的}\BibleRef{若 8:38}；他们受责备，不是因为他们看见了魔鬼的工作，而是因为他们行了魔鬼的意愿，因为不可见的魔鬼按着他们自己的邪恶在他们内运行罪恶。我们写下这些，正如宗徒所做的，是因为这些背信者的愚妄。\BibleRef{弟后 3:9}\par

135. 然而，我们已通过圣经的例证充分证明了：父住在子内，子似乎从父听见他所说的话，乃是\OriginalTerm{神圣尊威}{divine majesty}合一的特性。除了藉着知道对父与子应存同样的敬意之外，我们怎能理解尊威的合一呢？因为有什么能比宗徒所说的荣耀的主被钉在十字架上这句话更恰当呢？\BibleRef{格前 2:8}\par

136. 因此，圣子是荣耀的天主和荣耀的主，但荣耀不受制于受造物；所以圣子不是受造物。\par

137. 圣子是\OriginalTerm{圣父实体的肖像}{Image of the Father's Substance}；\BibleRef{希 1:3} 但所有受造物都与那神圣的\OriginalTerm{实体}{Substance}不相似，然而父之子却与天主相似；因此圣子不是受造物。\par

圣子认为自己与天主同等，并非僭越；\BibleRef{斐 2:6} 但任何受造物都不与天主同等，而圣子却是同等的；因此圣子不是受造物。\par

一切受造物都是可变的；而天主子却不是可变的；因此天主子不是受造物。\par

一切受造物按其本性的能力都会遭遇善恶的偶然事件，也感受其消逝；但就天主子而言，在祂的天主性内，没有任何事物能从祂消逝或给祂增添什么；因此天主子不是受造物。\par

天主的一切工程都必受审判；\BibleRef{训 12:14} 但天主子却不受审判；因为祂自己施行审判；因此天主子不是受造物。\par

最后，为使你们明白这合一，救主论到祂的羊时说：\BibleQuote{谁也不能从我手中把他们夺去。我那赐给我羊群的父，超越万有，谁也不能从我父手里将他们夺去。我与父原是一体。}\BibleRef{若 10:28}\par

所以圣子如同圣父那样赋予生命。\BibleQuote{就如父唤起死者，使他们复生，照样，子也随自己的旨意使人复生。}\BibleRef{若 5:21} 所以圣子如同圣父那样唤起，也如同圣父那样保全。祂在恩宠上既不相异，在权能上又如何相异呢？同样，圣子也不毁灭，正如圣父不毁灭。因此，为避免有人以为有两位天主，或想象权能上有分别，祂便说祂与圣父是一体。一个受造物怎能这样说呢？因此天主子不是受造物。\par

治理与服事不是一回事；而基督既是君王，又是君王之子。因此天主子不是仆人。然而，一切受造物都是服事的。但天主子使仆人成为天主的子女，祂却不服事人。因此天主子不是仆人。\par

\Emph{祂借着那个往远方去为自己取得王位的富人的比喻，确证了先前所说的话；并指明，当圣子将王国交与圣父时，我们断不可以贬抑的方式看待圣父使万物屈伏于祂的事实。在此世，我们是基督的王国，并在基督的王国里；将来，我们将在天主的王国里，那时天主圣三将一同为王。}\par

祂以属神的方式讲述了那个富人的比喻：这富人往远方去，为取得王位，还要回来，\BibleRef{路 19:12} 以此描述祂自己的天主性实体和人性。因为祂原在天主性的圆满中是富有的，且是永恒的君王、永恒君王之子，却为了我们成为贫穷的，虽然祂本是富有的；我是指，祂借着取了一个身体，往异乡去了，因为祂踏上世人的道路，如同踏上一段陌生的旅程，来到这世界，为从我们当中为自己预备一个王国。\par

因此，耶稣来到这地上，为从我们这里为自己取得一个王国，祂对我们说：\BibleQuote{天主的国就在你们中间。}\BibleRef{路 17:21} 这就是基督所领受的王国，这就是祂交与圣父的王国。因为祂本是永恒的君王，又如何为自己取得一个王国呢？\Quote{因此，人子来，是为取得一个王国，而后返回。}犹太人不愿意承认祂，论到他们，祂说：\BibleQuote{至于那些不愿意我作王统治他们的人，把他们押到这里来，在我面前杀掉。}\BibleRef{路 19:27}\par

让我们遵循圣经的脉络。那要来者，将把王国交与天主圣父；当祂交出王国之后，祂自己也要屈伏于那使万物屈伏于祂的，好叫天主成为万物中的万有。\BibleRef{格前 15:24} 如果天主子以人子的身分领受了王国，那么，祂也必以人子的身分交出所领受的。如果祂以人子的身分交出，那么作为人子，祂确然是在血肉之躯的条件下承认自己的屈伏，而非在天主性的尊威中。\par

难道你们还要反对并轻慢祂，因为天主使万物屈伏于祂，而你们却听到人子将王国交与天主，并且读过，一如我们在先前书卷中所说的：\BibleQuote{凡不是派遣我的父所吸引的人，谁也不能到我这里来，而我在末日要叫他复活。}\BibleRef{若 6:44} 如果我们按字面遵循，更该看到并留意彼此给予的尊荣的合一：圣父使万物屈伏于圣子，而圣子将王国交与圣父。现在请说说，哪一样更大：是交出，还是使人复活？我们岂不按人的方式，说交出是一种服事，而使人复活是一种能力吗？然而，圣子既向圣父交出，圣父也向圣子交出。圣子使人复活，圣父也使人复活。在权能合一之处，让他们去编造亵渎的分裂的谎言吧。\par

那么，就让圣子将祂的王国交与圣父吧。祂所交出的王国，并非基督失去了，而是增长了。我们就是这王国，因为经上对我们说：\BibleQuote{天主的国就在你们中间。}\BibleRef{路 17:21} 并且我们就是这王国，先属于基督，后属于圣父；如经上所载：\BibleQuote{除非经过我，谁也不能到父那里去。}\BibleRef{若 14:6} 当我在路上时，我属于基督；当我通过以后，我便属于圣父；但处处都是经由基督，又处处都属祂权下。\par

在基督的王国里是件美事，好使基督与我们同在；正如祂自己说的：\BibleQuote{看，我同你们天天在一起，直到今世的终结。}\BibleRef{玛 28:20} 但与基督同在更为美好：\BibleQuote{因为离开此世与基督同在，实在好得多。}\BibleRef{斐 1:23} 虽然我们在此世仍处于罪恶之下，基督却与我们同在，是要 \BibleQuote{因一人的服从，众人成为义人。}\BibleRef{罗 5:19} 如果我脱离了此世的罪恶，我便开始与基督同在。因此祂说：\BibleQuote{我必再来接你们到我那里去；}\BibleRef{若 14:3} 并在稍后说：\BibleQuote{好使我在那里，你们也在那里。}\BibleRef{若 14:3}\par

因此，我们如今尚在肉身之时，是在基督的统治之下，还未脱去祂所取的形象，就是祂\Quote{空虚自己}时所取的形象。但当我们得见祂在创世之前的荣耀时，我们将在天主的国里，那里有圣祖和先知，正如经上记载：\BibleQuote{你们将看见亚巴郎、依撒格和雅各伯，并众先知在天主的国里；}\BibleRef{路 13:28} 这样，我们将对天主有更深的认识。\par

但在圣子的国里，圣父也临在为王；在圣父的国里，圣子也临在为王：因为圣父在圣子内，圣子在圣父内；圣子住在谁内，圣父也住在谁内；圣父住在谁内，圣子也住在谁内，正如经上记载：\BibleQuote{我和我父要到他那里去，并在他那里作我们的住所。}\BibleRef{若 14:23} 因此，既有一个住所，也就有一个国。的确，圣父的国与圣子的国是同一个，以致圣父领受圣子所交付的，圣子也并不失去圣父所领受的。因此，在这一个国里，有着权能的合一。所以，谁也不可在圣父与圣子间分裂天主性。\par

\Emph{怀抱着学习的心，要了解服从基督的意义，在提出并否定了种种关于服从的观念后，他再度引述了宗徒的话语；从而结束了异端者对此事的亵渎性见解。此处所论的服从，既然是未来的事，便不可能与天主性有关，因为圣父与圣子的意志，向来是全然和谐的。再者，就圣子的天主性而言，万有确已服属于祂；但之所以说万有尚未服属于祂，乃是因为并非所有的人都服从祂的诫命。然而，当万有都服属于祂之后，基督便也要在他们内服属于父，而圣父的工程也就得以圆满。}\par

然而，如果天主这一个名号和权柄，既是圣父和圣子所共有的，因为天主子也是真天主，也是永恒的君王，那么天主子在其天主性内，是并不服属的。那么，\Person{奥古斯都}皇帝，让我们思想一下，我们又当如何看待祂的服属呢？\par

天主子是如何服属的呢？像受造物服属于虚幻一样吗？可是，若这样去想象天主性\OriginalTerm{实体}{Substance}，那便是亵渎了。\par

或者，像一切受造物服属于天主子一样吗？因为经上确实记载说：\Quote{你将一切放在祂脚下}？可是，基督并不服属于自己。\par

或者，像女人对男人一样吗？正如我们读到：\BibleQuote{作妻子的，应当服从自己的丈夫；}\BibleRef{弗 5:22} 又说：\BibleQuote{女人要在沉静中受教，事事服从}\BibleRef{弟前 2:11}？可是，将男人比作圣父，或将女人比作天主子，那是亵渎的。\par

或者，如伯多禄所说：\BibleQuote{你们要服从一切人}\BibleRef{伯前 2:13}？可是基督当然不是这样服属的。\par

或者，如保禄写道：\BibleQuote{你们要怀着敬畏基督的心，互相顺从。}\BibleRef{弗 5:21}？可是，基督的服属，既不是靠祂自己的敬畏，也不是靠另一位基督的敬畏。因为基督只有一位。然而，请注意这些话的力度：我们既服从圣父，也敬畏基督。\par

那么，我们如何理解祂的服从呢？不如我们重温宗徒所写的整章经文，免得给人留下我们故意遗漏了什么，或有意削弱经文力量的印象？他说：\BibleQuote{如果我们只在今生寄望于基督，我们就是众人中最可怜的了。但是，基督从死者中实在复活了，做了死者的初果。}\BibleRef{格前 15:19} 你看，他是如何讨论基督复活的问题的。\par

他说：\BibleQuote{因为死亡既因一人而来，死者的复活也因一人而来；就如在亚当内，众人都死了，照样在基督内，众人也要复活；不过各人要依照自己的次序：首先是初果的基督，然后是在基督再来时属于基督的人，再后才是结局；那时，基督将消灭一切率领者、一切掌权者和大能者，把自己的王权交于天主父。因为基督必须为王，直到把一切仇敌屈伏在他的脚下。最后被毁灭的仇敌便是死亡；因为天主使万物都屈伏在他的脚下。既然说万物都已屈伏了，显然那使万物屈伏于他的不在其内。万物都屈伏于他以后，子自己也要屈伏于那使万物屈伏于自己的父，好叫天主成为万物之中的万有。}\BibleRef{格前 15:21-28} 同样，这位宗徒在\BibleBook{希伯来}{书}中也说：\BibleQuote{但是，我们至今还没有看见一切全服从于他}\BibleRef{希 2:8}。我们已聆听了宗徒的整篇论述。\par

那么，我们如何论述祂的服从呢？\OriginalTerm{萨培里派}{Sabellians}和\OriginalTerm{马西昂派}{Marcionites}声称，基督对天主父的这种服从，将会是圣子被\OriginalTerm{再吸收}{re-absorbed}到圣父之内。如果圣言的服从意味着天主圣言将被吸收到圣父之内；那么，所有服从圣父与圣子者，也都将被吸收到圣父与圣子之内，好使天主成为万物并内在于祂的一切受造物中。但这说法是愚昧的。因此，再吸收并非服从。因为，另有其他事物被造而服从，即那些受造物，而另有一位是它们所服从的对象。让那些宣讲严酷再吸收的人缄默吧。\par

那些因不能证明天主的圣言和天主的智慧可被再吸收，便归咎服从的软弱于祂的天主性，并援引经上的话说：\BibleQuote{万物都屈服于祂以后，子自己也要屈服于那使万物屈服于自己的父}\BibleRef{格前 15:28}。但愿这些人也缄默吧。\par

163. 我们由此看出，\WorkTitle{圣经}陈述祂尚未屈服，但这事将要发生：因此现在，圣子尚未屈服于天主圣父。那么，你们说圣子将以什么方式屈服呢？若是在祂的\OriginalTerm{天主性}{Godhead}内，祂并非不服从，因为祂并非与圣父相悖；祂也不是被屈服，因为祂不是奴仆，而是属于自己父的唯一圣子。最后，当祂创造天、形成地时，祂施展了能力和仁爱。因此，在\Person{基督}的天主性内，没有像奴仆那样的\OriginalTerm{屈服}{subjection}。但若没有屈服，那么意志便是自由的。\par

164. 但如果他们以为圣子的屈服就是：圣父使万物与祂的旨意和谐一致，那么让他们明白，这其实是不分离能力的证明。因为祂们旨意的合一不是在时间中开始的，而是永远存在的。但哪里有恒常的旨意合一，哪里就不会有暂时屈服的软弱。因为如果祂因本性而被屈服，祂就会永远处于屈服之中；但既然祂被称为在时间中被屈服，那么这种屈服必定是担任了一种职务，而不是永远软弱的证明：特别是天主永恒的能力不能在时间中改变自己的状态，治理的权利也不能在时间中归于圣父。因为如果圣子将来会如此改变，以至于在祂的天主性中被屈服，那么按照你们的解释，如果天主圣父将来获得了更多能力，并使圣子在祂的天主性中屈服于自己，那么在此期间，圣父在此时就应被认为是次等的。\par

165. 但圣子有什么过错，要让我们相信祂以后会在天主性里被屈服呢？难道祂作为人，擅自取得了坐在圣父右边的权利，或者违背圣父的旨意，僭取了圣父的宝座的特权吗？但祂自己说：\BibleQuote{因为我常作祂所喜悦的事。}\BibleRef{若 8:29} 因此如果圣子在凡事上都使圣父喜悦，为什么祂要被屈服呢，既然祂以前未曾被屈服？\par

166. 那么让我们看看，那种屈服不是天主性的屈服，而是我们在\Person{基督}的敬畏中的屈服，这是如此充满恩宠和奥秘的真理。为此，我们再次思量宗徒的话：\BibleQuote{万物都屈服于他以后，圣子自己也要屈服于那使万物屈服于自己的天主圣父，好叫天主成为万物之中的万有。}那么你们怎么说呢？难道万物现在不已经屈服于祂吗？难道圣人们的歌队尚未屈服吗？难道那些曾在地上侍奉祂的天使尚未屈服吗？\BibleRef{玛 4:11} 难道那些被派遣给\Person{玛利亚}预言主降临的总领天使尚未屈服吗？难道所有天军尚未屈服吗？难道\OriginalTerm{革鲁宾}{cherubim}和\OriginalTerm{色辣芬}{seraphim}尚未屈服？难道朝拜和赞美祂的\OriginalTerm{上座者}{thrones}、\OriginalTerm{宰制者}{dominions}和\OriginalTerm{异能者}{powers}尚未屈服吗？\par

167. 那么，它们怎样才能屈服呢？就是主自己所说的方式。\BibleQuote{背起我的轭，跟我学罢！}\BibleRef{玛 11:29} 负轭的不是凶猛的人，而是谦卑温柔的人。这显然不是对人卑下的屈服，而是光荣的：\BibleQuote{一听到耶稣的名字，无不屈膝叩拜，凡上天、地上和地下的一切；一切唇舌无不明认耶稣基督是主，以光荣天主圣父。}\BibleRef{斐 2:10} 但是之所以先前万物没有屈服，是因为他们还没有接受天主的智慧，还没有把\OriginalTerm{圣言}{Word}的柔和之轭放在心灵上仿佛负于颈项。\BibleQuote{凡接受他的，}正如经上所载，\BibleQuote{他给他们权能，好成为天主的子女。}\BibleRef{若 1:12}\par

168. 难道会有人说，因为许多人信了，\Person{基督}现在已经屈服了吗？当然不会。因为基督的屈服不在于少数人，而在于人类。正如若我内的肉欲仍相反圣神，圣神相反肉欲，\BibleRef{迦 5:17} 我就似乎不被屈服，尽管我已部分被制服；同样，因为整个教会是基督的一个身体，只要人类仍有分歧，我们就分裂了基督。因此，基督尚未屈服，因为祂的肢体尚未被屈服。但当我们成为，不再有许多肢体，而是一个心神，那时祂也将屈服，好使藉着祂的屈服，\BibleQuote{天主成为万物之中的万有。}\par

169. 但正如\Person{基督}尚未屈服，天主的工作也尚未完成；因为天主子说过：\BibleQuote{我的食物就是承行派遣我者的旨意，完成他的工程。}\BibleRef{若 4:34} 既然我自身尚未完善，圣父的工作在我内尚未完成，那么圣子的屈服在我内当然仍属未来吗？我，使天主的工作未完成，是否也使天主子屈服呢？但那不是过失之事，而是恩宠之事。因为就我们被屈服而言，我们屈服于法律，屈服于恩宠，是为我们的益处，而非为天主性的益处。因为从前，如宗徒自己所言，肉性的智慧与天主为敌，因为它\BibleQuote{不服从天主的法律，}\BibleRef{罗 8:7} 但现在却藉着\Person{基督}的\OriginalTerm{苦难}{Passion}屈服了。\par

\Emph{他继续讨论所进入的难题，并教导说基督并非屈服，而只是按肉身而言。然而，基督在肉身屈服时，仍给出了祂天主性的证据。他驳斥基督在这方面屈服的观点。确实，祂所摄取的人性，已经在我们内屈服，同样，我们的人性也在祂那完全相同的人性中被提升。最后，我们获知，基督的那同一屈服将在何时发生。}\par

170. 然而，免得有人挑剔，请看圣经在天主默感下是何等谨慎。因为它既向我们显示\Person{基督}在何事上服从了天主，又教导我们在何事上他使宇宙服从了自己。所以它说：\BibleQuote{现今我们还没有看见万物全属他权下。}\BibleRef{希 2:8} 因为我们看见\Person{耶稣}因着死亡的痛苦，变得比天使卑微一点。\BibleRef{希 2:9} 因此它显示，他在取我们的肉躯时变得卑微。那么，他在取我们的肉躯时，那肉躯使他能叫万有服从自己，而他自己在那肉躯内则被置于天主圣父的权下，这一切既已如此，又有什么能阻止他公开显示其服从呢？\par

171. 那么，让我们想想他的服从。他说：\BibleQuote{父啊！你如果愿意，请给我免去这杯罢！但不要随我的意愿，惟照你的旨意成就吧。}\BibleRef{路 22:42} 因此那服从将是依照所取人性的；正如我们所读到的：\BibleQuote{他贬抑自己，听命至死。}\BibleRef{斐 2:8} 所以服从就是听命；听命就是死亡；死亡就是所取的人性；因此那服从就是所取之人性的服从。这样，天主性内绝无软弱，反而有一种虔诚职责的如此履行。\par

172. 看，我如何不怕他们的企图。他们声称他必须服从天主圣父，我却说他曾服从他的母亲\Person{玛利亚}。因为论及\Person{若瑟}和\Person{玛利亚}时记载：\BibleQuote{他就同他们下去，来到纳匝肋，属他们管辖。}\BibleRef{路 2:51} 然而，如果他们这样想，请他们说，天主性何以服属于人。\par

173. 切莫让人说祂被置于服从地位这一事实对祂有所损害，因为祂既被称为仆人，又被宣称为被钉死，又被说成已经死亡，这些都没有给祂带来任何损害。因为祂死时，祂仍是活着的；祂服从时，祂仍在为王；祂被埋葬时，祂又复活了。祂让自己服从人的权下，可是在另一时候祂又宣称自己是永恒荣耀的主。祂站在审判官面前，却又为自己要求坐在天主右边宝座上，作为永远的审判者。因为经上这样记载：\BibleQuote{从此以后，你们将看见人子坐在大能者的右边，乘着天上的云彩降来。}\BibleRef{玛 26:64} 祂被犹太人鞭打，却又命令天使；祂由\Person{玛利亚}所生，生于法律之下\BibleRef{迦 4:4}；祂在\Person{亚巴郎}之前，却在法律之上。在十字架上，祂受万有的崇敬；太阳逃离；大地震动；天使静默。元素可曾看见祂的\OriginalTerm{生成}{Generation}？但它们却害怕看见祂的苦难。它们既不能忍受祂肉身的服从，又怎能接受在祂内那应受钦崇的天主性服从呢？\par

174. 但既然圣父、圣子和圣神是同一性体，圣父自然不会服从自己。因此，在那圣子与圣父合而为一的方面，圣子也不会处于服从地位；免得由于天主性的一体，显得圣父也服从于圣子。所以，正如在十字架上，被置于服从之下的不是圆满的天主性，而是我们的软弱，同样，将来圣子也将在分享我们的人性时服从于圣父，好使肉身的私欲被制服之后，内心不再眷恋财富、野心或享乐；而是让天主成为我们的一切，只要我们按照祂的肖像和模样生活，尽我们所能，透过万有达到这一境界。\par

175. 那么，这益处已从个体传给了团体；因为祂在自己肉躯内驯服了全人类肉躯的性情。因此，按宗徒的话：\BibleQuote{我们怎样带了那属于土的肖像，也要怎样带那属于天上的肖像。}\BibleRef{格前 15:49} 这事显然只有在内在的人身上才能实现。因此，\BibleQuote{戒绝这一切，}那就是我们读到的：\BibleQuote{忿怒、暴戾、恶意、诟骂和污秽的言谈，}\BibleRef{哥 3:8} 正如下文他接着说的：\BibleQuote{你们原已脱去了旧人和他的作为，却穿上了新人，这新人即是照创造他者的肖像而更新，为获得知识的。}\BibleRef{哥 3:9}\par

176. 为使你知道，当他说\Quote{天主成为万物之中的万有}时，并没有将\Person{基督}与天主圣父分开，他也对哥罗森人说：\Quote{在这一点上，已没有希腊人或犹太人，受割损的或未受割损的，野蛮人、叔提雅人、奴隶、自由人的分别，而只有是一切并在一切内的\Person{基督}。}\BibleRef{哥 3:11} 同样，他对格林多人说：\Quote{天主成为万物之中的万有}，他在此句话中包含了\Person{基督}与天主圣父的结合与平等，因为圣子并不与圣父分开。正如圣父在万有内作一切，同样\Person{基督}也在万有内作一切。那么，如果\Person{基督}也在万有内作一切，祂就不是在天主性的荣耀上被置于服从，而是在我们内被置于服从。但是，祂如何在我们内被置于服从呢？岂不是和祂变得比天使卑微的方式一样，也就是在祂身体的圣事中吗？因为一切万物，自起初便事奉它们的造物主，但似乎尚未在这一点上服属于祂。\par

177. 但如果你问，祂如何在我们内被置于服从，祂自己就指示给我们，说：\BibleQuote{我在监里，你们来探望了我；我患病，你们看顾了我；凡你们对我这些最小兄弟中的一个所做的，就是对我做的。}\BibleRef{玛 25:36} 你听到祂患病和软弱，并不感动；你听到祂在服从中，却感动了，虽然祂正是在那作为服从而患病和软弱的形式内，也是在为我们的缘故成为罪和诅咒者的那一位内。\par

178.~因此，正如祂并非为自己而是为我们成了罪和诅咒，同样，祂并非为自己而是为我们，在我们内受到服从；在祂永恒的本性里，祂既不受服从，也不受诅咒。\BibleQuote{凡被悬在木架上的，都是可咒骂的。}\BibleRef{迦 3:13} 祂受了咒骂，因为祂承担了我们的咒骂；祂也在服从之下，因为祂把我们的服从担在自己身上，但这服从是在取了奴仆的形体时，而非在天主的光荣内；这样，当祂在肉身内分担我们的软弱时，祂就使我们因祂的大能而在天主性上有份。然而无论在这一点上，或在另一点上，我们都与基督\OriginalTerm{天上的受生}{Generation}没有任何本性的共融，在基督内的天主性也没有任何服从。但是正如宗徒说过，因着那身为我们救恩保证的肉身，我们与祂一同坐在天上，\BibleRef{弗 2:6} 尽管肯定不是我们自己坐，同样，祂也因着取了我们的本性，而被称为在我们内服从。\par

179.~那么，正如我们已说过的，谁竟如此狂妄，竟以为在天主父的右边有一个尊位归给祂，而这一位却是生祂的父按肉身赐予基督的，即一个具有天上同等权能的座位？天使们朝拜祂，难道你竟企图以不敬的妄断推翻天主的宝座吗？\par

180.~你说经上写道：\BibleQuote{当我们死在罪恶中的时候，便叫我们与基督一同活过来；你们得救就是靠着恩宠。他又使我们同他一起复活，在基督耶稣内使我们和他一同坐在天上。}\BibleRef{弗 2:5} 我承认经上是这样写的；但经上并不是说，天主容许人坐在祂的右边，而只是说，在基督的位格内坐在那里。因为祂是万有的根基，是教会的头，\BibleRef{弗 5:23} 我们按肉身的共同人性，在祂内才配得了天上宝座的权利。因为肉身因分享基督——即是天主的基督——而受到尊荣，全人类的性体因分享这肉身而受到尊荣。\par

181.~因此，正如我们因在肉身性体上的共融，而坐在祂内，同样，那因取了我们的肉身而为我们成为诅咒的（因为诅咒不能落在可赞美的天主子身上），我要说，祂也因着万民的服从，将在我们内服从于父；当外邦人信仰了，犹太人承认了那位他们曾钉死的；当\OriginalTerm{摩尼教徒}{Manichæan}朝拜了那位他们曾不相信已降生成人的；当\OriginalTerm{亚略派}{Arian}宣认了那位他们曾否认的全能者；最后，当天主的智慧、他的正义、和平、仁爱、复活在万有内时。基督将藉着祂自己的工程，并藉着诸德行的种种形式，在我们内服从于父。那时，当邪恶舍弃，罪债终结，一种心神开始在万民心中，在一切事上契合于天主时，天主就会成为万有之中的万有。\BibleRef{格前 15:28}\par

\Emph{他简要地再次提出同一争论点，并从父与子天主性权能的合一中精明地得出结论：凡论及子的服从，都仅指祂的人性。他更以存在于双方同等的爱为证据，进一步确证这一点。}\par

182.~那么，让我们对整件事的结论简短总结如下。权能的合一排除了任何贬抑性服从的想法。祂交出权能，以及祂作为征服者赢得的对死亡的胜利，并没有减损祂的权能。服从带来服从。基督自己取了服从，甚至服从到取了我们的肉身，以及十字架，以赢取我们的救恩。因此，工作在哪里，工作的创始者也就在哪里。所以，当万物因着基督的服从，都服从于基督，使所有都因祂的名字屈膝时，那时祂自己将成为万有之中的万有。因为如今，既然并非所有都相信，也就不是所有都看似服从。但当所有都相信并承行了天主的旨意，那时基督就成为万有之中的万有。而当基督是万有之中的万有时，天主也就是万有之中的万有；因为父永远住在子内。那么，那位救赎了软弱者的，怎能被显示为软弱的呢？\par

183.~而且，以免你偶然因为经上记载天主将万物置于祂的脚下，而归因于子的软弱；你要学习，是祂自己将万物归于自己的权下，因为经上写道：\BibleQuote{我们的家乡原是在天上，我们也从那里等候救主、主耶稣基督，他将按那能使万物屈服于自己的大能，改变我们卑贱的身体，相似他光荣的身体。}\BibleRef{斐 3:20} 因此，你已经学到，凭祂天主性的大能，祂能使万物屈服于自己。\par

184.~现在学习，按肉身论，祂如何承受万物的服从，如经上所记：\BibleQuote{他在基督身上所施展的强有力的大能，使他由死者中复活，叫他在天上坐在自己右边，超乎一切率领者、掌权者、异能者、宰制者，以及一切现世及来世可称呼的名号以上；又将万有置于他的脚下。}\BibleRef{弗 1:20} 所以，按肉身论，万物在服从中被赐予了祂；按此，祂也从死者中复活，既在祂人性的灵魂上，也在祂理性的服从上。\par

185.~经上写道：\BibleQuote{我的灵魂确实服从于天主吗？} 许多人对这节经文作出了高贵的解释；他说的是灵魂，不是天主性，灵魂不是光荣。为了使我们知道，主藉先知论及我们人性的被收养，他加上说：\BibleQuote{你们攻击一个人要到几时？} 正如他也在福音中说：\BibleQuote{你们为什么图谋杀害我这个说实话的人呢？}\BibleRef{若 8:40} 他又加上说：\BibleQuote{然而他们想要拒绝我的赎价，他们狂渴奔驰，他们口虽祝福，心却诅咒。} 因为当犹达斯把那价钱拿回来时，\BibleRef{玛 27:4} 犹太人却不接受它，在狂妄的渴中奔驰，因为他们拒绝了属灵饮料的恩宠。\par

这就是对\Quote{服从}的虔敬解释，因为既然这是主受难的职分，便要在所遭受的事上，在我们内成为服从的。我们可问这是为什么？为的是\BibleQuote{无论是天使，掌权者，崇高，深幽，现世的事物，是将来来的事物，或是任何其他受造物，都不能使我们与天主的爱相隔绝，就是在基督耶稣内的爱。}\BibleRef{罗 8:38} 由此可见，从上述所言，没有任何受造物被排除在外；而是无论何种，都被列在他上面所提及的那些之中。\par

同时，我们也必须思量这些话：他先说：\BibleQuote{谁能使我们与基督的爱相隔绝？}\BibleRef{罗 8:35} 然后写道：\BibleQuote{无论是死亡，生活，或是任何其他受造物，都不能使我们与天主的爱相隔绝，就是在基督耶稣内的爱。}由此可见，天主的爱与基督的爱是相同的。因此，他论及天主的爱时写\Quote{就是在基督耶稣内的}，并非没有理由，免得你误以为天主的爱与基督的爱是分开的。然而，爱不分隔任何事物，永恒的天主性没有不能做的，没有什么是真理所不知的，或能欺骗正义，或能逃过智慧的鉴察。\par

\Emph{亚略派人士受圣神借达味之口所谴责：因为他们竟敢限制基督的知识。他们为证明此事所引用的那节经文，绝非没有遭人篡改的嫌疑。但为纠正此事，我们必须注意\Quote{子}一字。因为知识不可能不临于身为圣子的基督，因为就是智慧；认知也不会缺少任何部分，因为创造了万物。那创造了诸世代的主，不可能不知道未来，更不用说审判之日了。这样的知识，无论事关重大还是微小，都不可否认子拥有，同样也不可否认圣神拥有。最后，将提供各种证据，好使我们从中推论出，基督确实拥有这种知识。}\par

\Emph{因此}我们应当知道，说这种话的人，是被圣神所咒骂和谴责的。因为除了亚略派的首要人物之外，先知还能谴责谁呢？因为他们竟说圣子不认识时辰和岁月。因为天主没有不知道的事；而基督，诚然是至高者基督，就是天主，因为是的\BibleQuote{在万有之上的天主}。\BibleRef{罗 9:5}\par

请看圣达味对这类人多么惊骇，因为他们限制了圣子的知识。因为经上这样记载：\BibleQuote{他们不像别人一样受苦，也不像别人一样受鞭笞；因此骄傲攫住了他们；他们被自己的邪恶和亵渎所覆盖；他们的邪恶仿佛因肥胖而挺出；他们走入了心中的思念。}真的谴责了那些认为应该用人心思念的角度来看待天主之事的人。因为天主不受安排或顺序的约束；我们甚至看到，那些人间常见、在人类历史中经常发生的事，也并非总是按某种既定规律安排而发生，反而常常以某些隐秘莫测的方式突然发生。\par

\BibleQuote{他们思想}，说，\BibleQuote{并且说邪恶的话，他们说出反对至高者的话，他们开口攻击上天。}由此可见，谴责那些声称有权按照我们人性的样式来安排天上奥秘的人，视之为犯了恶毒的亵渎之罪。\par

他们还说：\BibleQuote{天主怎能知道？至高者岂有知识？}亚略派人士不正是天天附和这话，说一切知识不能存在于基督内吗？因为他们说，声称自己不知道那日子和时辰。他们岂不正是在说：怎能知道呢？他们一面主张除了听见和看见的之外一无所知，一面又用一种亵渎的解释，把关乎天主性统一的道理用来削弱的能力。\par

他们说，经上记载着：\BibleQuote{至于那日子和那时刻，除了父以外，谁也不知道，连天上的天使和子都不知道。}\BibleRef{谷 13:32} 首先，古老的希腊文抄本并未包含\Quote{子也不知道}一句话。但若那些篡改神圣圣经的人，也伪造了这段经文，倒也不足为奇。这句话被插入的原因十分明显，因为只要引用它，就能展开这样的亵渎。\par

然而，假若圣史的确是如此记述的。\Quote{子}之名包含两种本性。因为也被称为人子，因此，在与我们所取人性相关的无知中，似乎不知道那要来的审判之日。因为圣子怎能不知道那日子呢，既然一切智慧和知识的宝藏，都藏在祂内？\BibleRef{哥 2:3}\par

那么我要问，是凭自己的存有，还是偶然地拥有这知识呢？因为一切知识来到我们身上，或者是透过本性，或者是透过学习。本性供应的，比如马能跑，鱼能游，因为它们不用学习就会。另一方面，人是通过学习才会游泳的，因为除非学习，否则就不会。因此，既然本性使无灵的动物能做、能知它们所未学之事，你又为何要对圣子妄加论断，说祂的知识是来自教导还是来自本性呢？如果是来自教导，那么就不是以智慧的身分受生了，是逐渐才开始成为完美的，而不是恒常如此。但如果是来自本性，那么在起初就是完美的，从父而来时就是完美的；因此，不需要任何对未来的预知。\par

因此，他并非不知道那些日子；因为\OriginalTerm{天主的智慧}{Wisdom of God}不可能部分知道、部分不知道。那位创造万物的怎能对局部无知呢？因为知道比起创造来是更小的事。我们知道许多我们不能创造的事物，我们也不是都以同样的方式知道，而是部分知道。因为一个乡下人知道风的力量和星辰的运行是一种方式——城市居民知道它们又是另一种方式——而领航员则用第三种方式。尽管不是所有人都知道一切，他们仍被称为知道；但惟有他，即创造万物的那位，完全知道一切。领航员知道大角星持续多少更次，他会发现猎户座怎样的升起，但他对昴星团和其他星辰的联系，或它们的数目与名称却一无所知，正如那位\BibleQuote{数点众星、一一称呼它们名号者}所知道的；确实，他工作的威能无法瞒过他自己。\par

那么，你希望天主子如何造了这些呢？如同一个没有感觉的印记戒指吗？但圣父在智慧中创造一切，就是说，他藉着圣子创造一切，圣子正是\OriginalTerm{天主的德能与智慧}{Virtue and Wisdom of God}。\BibleRef{格前 1:24} 但对于这样的智慧，知晓他自己作为的能力和原因是合宜的。因此，万物的创造者不能对他所造的无知——或者对他自己所赐予的毫无知识。因此，他知道他所造的那日子。\par

但你说他知道现在，却不知道未来。虽然这是个愚昧的想法，但为了让你从圣经的根据上满意，你要知道，他不仅创造了过去的，也创造了将来的，正如经上所记：\BibleQuote{他造成了将来的事。}\BibleRef{依 45:11} 圣经另处又说：\BibleQuote{藉着他造成了诸世代，他是他光荣的辉耀，是他本体的真像。}\BibleRef{希 1:2} 而世代有过去、现在和未来。那么将来的是如何造成的呢？难道不是因为他主动的大能和知识在自己内包含了一切世代的数目吗？正如他叫那不存在的成为存在的，\BibleRef{罗 4:17} 他也使将来的事如同已存在一样。它们不能不成为现实。他预定要存在的事物，必然将会存在。因此，那造成将要来临之事者，按它们将存在的方式知道它们。\par

如果关于世代我们相信这一点，关于审判的日子我们就更应相信了，其根据是天主子对它拥有知识，因为它已是他所造的了。因为经上记载：\BibleQuote{按照你的法令，日子仍要延续。}他没有只说\Quote{日子延续，}而是说\Quote{仍要延续，}好使将来的事受他的法令管理。难道他不知道自己所命令的吗？\BibleQuote{那栽种耳朵的，难道他自己听不见？那形成眼睛的，难道他自己看不见？}\par

然而，让我们看一看，是否可能有什么大事超出了造物主的知识；或者至少让他们选择，他们愿意认为那是什么伟大而超群的事，还是什么微小而卑微的事。若是微小而卑微的，按我们的说法，不知道那无价值且琐碎的事，这并不算损失。因为知道最伟大的事是能力的标记，而去关注那价值较小的，反倒似乎显得是较低劣的工作。这样，他就免于挑剔，却不至于失去他的能力。\par

但如果他们认为知道审判的日子是件重大而重要的事：那么让他们说出什么比天主圣父更伟大或更美善。他认识天主圣父，正如他自己所说的：\BibleQuote{除了子和子所愿意启示的人以外，没有人认识父。}\BibleRef{玛 11:27} 我说，他认识父，却不认识那日子吗？那么你们相信他启示了父，却无法启示那日子吗？\par

其次，因为你们划分等级，将父置于子之前，又将子置于圣神之前，请告诉我，圣神是否知道审判的日子？因为关于他，此处没有任何记载。你们完全予以否认。但如果我向你们证明他知道呢？因为经上记载：\BibleQuote{可是天主藉着圣神将这一切启示给我们了，因为圣神通透一切，就连天主的深奥事理他也洞悉。}\BibleRef{格前 2:10} 因此，因为他通透天主的深奥事理，既然天主知道审判的日子，圣神也知道。因为他知道天主所知道的一切，正如宗徒所说的：\BibleQuote{除了人内里的心神外，有谁能知道那人的事呢？同样，除了天主圣神外，谁也不能明了天主的事。}\BibleRef{格前 2:11} 因此，你要小心，免得要么因否认圣神知道，而否认父知道；（因为天主的事，天主圣神也知道，但圣神不知道的事，就不是天主的事）。要么因承认天主圣神知道，却否认天主子知道，而将圣神置于子之前，这与你们自己的声明相悖。但在这一点上犹豫不决，不仅亵圣，且属愚昧。\par

现在想一想，知识是如何获取的，让我们指出，子自己证明了他是知道那日子的。因为我们所知道的，我们或是通过提到时间、地点、征兆、人物，或是通过给出它们的次序，使之显明。那么，那位描述了审判的时刻和地点、以及征兆与情形的主，怎会不知道审判的日子呢？\par

所以你们有这样的记载：\BibleQuote{在这一天，那在屋顶上的，不要下来取他屋里的东西；那在田地里的，同样不要回来。}\BibleRef{路 17:31} 他对未来的危险后果知道得如此深远，竟至于向处于危险中的人指出了安全的方法。\par

那位曾亲口说人子是安息日的主的，怎会不知道那日子呢？\BibleRef{玛 12:8}\par

205. 主也在别处划定了一个地点，当祂的门徒把圣殿的建筑指给祂看时，祂对他们说：\BibleQuote{你们看见这一切吗？我实在告诉你们：将来这里决没有一块石头留在另一块石头上，而不被拆毁的。}\BibleRef{玛 24:2}\par

206. 当门徒问祂关于征兆的问题时，祂回答说：\BibleQuote{你们要谨慎，不要受欺骗！因为将有许多人假冒我的名字来说：我就是基督；}\BibleRef{路 21:8} 随后祂又说：\BibleQuote{将到处有大地震、饥荒和瘟疫；将有恐怖的现象，天上要有巨大的异兆。}\BibleRef{路 21:11} 这样，祂既描述了人，也描述了征兆。\par

207. 祂以怎样的方式预言军队将围绕\Place{耶路撒冷}，外邦人的时期何时满全，以及这些事的次序——这一切都藉着福音话语的见证向我们揭示出来。因此，祂知道一切。\par

\Emph{基督为了我们的益处，不愿揭示审判的日子。这一点由主的其他话语以及保禄著作中一段并非不相似的经文清楚表明。其他一些似乎将同样的无知归于圣父的章节，被引用来回应那些急切想知道基督为什么回答门徒，仿佛祂不知道的人。从这些章节中，盎博罗削反驳他们说，如果他们在圣父身上承认无知和无能，他们就必须承认圣子与圣父具有同一的\OriginalTerm{实体}{Substance}；除非他们宁愿控告圣子说谎；因为欺骗既不属于祂，也不属于圣父，但上述章节指出了两者的合一。}\par

208. 但我们要问，祂不愿说出那时间的原因是什么。如果我们追问，就会发现这并非出于无知，而是出于智慧。因为知道这事对我们没有益处；正是为了我们对即将来临的审判的时刻茫然无知，我们才能时刻保持警惕，立于德行的瞭望台上，从而远离罪恶的习惯；免得主的日子在我们作恶时突然临于我们。因为知道未来对我们没有益处，反该心存畏惧；因为经上记载：\BibleQuote{不可心高气傲，反应畏惧。}\BibleRef{罗 11:20}\par

209. 因为如果祂清楚地说出了那日子，祂似乎就只为最接近审判的那一代人规定了生活的法则，而较早时代的义人便会更加疏忽，罪人也会更无忧无虑。因为奸淫者除非日日惧怕惩罚，就不能停止犯奸淫的欲望；盗贼除非知道刑罚日日悬在头上，也不能离开他所栖身的林中藏匿之处。因为不洁通常驱使他们放纵，而恐惧却终觉惹人厌烦。\par

210. 因此我说，知道这事对我们没有益处；反而，不知道对我们有益，藉着无知我们可以畏惧，藉着警惕我们可以改正，正如祂自己所说：\BibleQuote{你们应准备妥当，因为你们不料想的时辰，人子就来了。}\BibleRef{玛 24:44} 因为士兵若不知道战争临近，就不知道怎样在营中警戒。\par

211. 因此，另一次主自己被宗徒们问及（是的，因为他们不像\Person{亚略}那样理解，而是相信天主子知道未来。因为除非他们相信这一点，他们决不会问这问题）——我说，主在被问及何时将国度复兴给\Place{以色列}时，并没有说祂不知道，而是说：\BibleQuote{不是你们应当知道的时候或年期，那是父以自己的权柄所定的。}\BibleRef{宗 1:7} 注意祂说的话：不是你们应当知道的！再读一次，\Quote{不是你们。} 祂说 \Quote{你们}，而不是 \Quote{为我}，因为此时祂不是按自己的完美说话，而是按对人性和我们灵魂有益的方式说话。所以祂说 \Quote{你们}，而不是 \Quote{为我}。\par

212. 宗徒也效法了这个例子：\BibleQuote{弟兄们，关于时候与季节，}他说，\BibleQuote{不需要给你们写什么。}\BibleRef{得前 5:1} 因此，甚至宗徒本人，基督的仆人，也没有说他不知道那些季节，而是说不需要教导人们；因为他们应当时常用属灵的武器武装起来，使基督的德能在每个人身上彰显出来。但当主说：\BibleQuote{那些时候，是父以自己的权柄所定的，}\BibleRef{宗 1:7} 祂当然不可能不分享父的知识，因为祂在权柄上也绝非没有份。因为权柄源于智慧和德能；而基督兼具这两者。\par

213. 但你问，为什么祂不以知者的身份拒绝门徒，却不肯说；为什么祂反而声明，无论是天使还是子都不知道？\BibleRef{谷 13:32} 我也要问你，为什么天主在\BibleBook{创}{世纪}中说：\BibleQuote{我现在要下去，看看他们所行的是否全如达到我耳中的呼声一样；如果不是，我也要知道。}\BibleRef{创 18:21} 为什么圣经也论及天主说：\BibleQuote{上主遂降下，要看看世人所造的城和塔。}\BibleRef{创 11:5} 为什么先知在\BibleBook{圣}{咏集}中又说：\BibleQuote{上主由高天俯视世人，察看有无寻觅天主的智者}？就好像在前一处，如果天主没有降下，在后一处，如果祂没有俯视，祂就会对人的作为或功过一无所知似的。\par

214. 在\BibleBook{路加}{福音}中你也有同样的情形，因为圣父说：\BibleQuote{我要怎样办呢？我要派遣我的爱子；或许他们会尊敬他。}\BibleRef{路 20:13} 在\BibleBook{玛窦}{福音}和\BibleBook{马尔谷}{福音}中记载：\BibleQuote{后来他派了自己的儿子去，说：他们会尊敬我的儿子。}\BibleRef{玛 21:37} 在一部福音中说：\BibleQuote{或许他们会尊敬我的儿子；}\BibleRef{谷 12:6} 并且显得犹豫，仿佛不知道；因为这是犹豫者的口吻。但在另外两部福音中祂说：\BibleQuote{他们会尊敬我的儿子；}\BibleRef{玛 21:37}；\BibleRef{谷 12:6} 就是说，祂声明他们将表示尊敬。\par

215. 但天主既不会疑惑，也不会受骗。因为只有对未来无知者才会疑惑；而受骗者，则是预言了一件事，却发生了另一件事。然而，有什么比圣经记述圣父论及圣子时所说的一件事，而同部圣经却证实发生了另一件事更明显的呢？圣子受鞭打，遭戏弄，被钉十字架，且死去了。祂在肉身中所受的苦，比先前被派定的那些仆人更为惨重。圣父是受骗了，还是对此无知，抑或不能施以援手？但那真实者不会犯错；因为经上记载：\BibleQuote{天主是忠信的，绝不撒谎。} \BibleRef{铎 1:2} 那知晓一切者，怎会无知？那能行一切者，又怎会有所不能？\par

216. 然而，如果圣父或是不知，或是无能（因为你们宁愿承认圣父不知，也不愿承认圣子知道），你们由此正可看出，圣子与圣父是同一性体；因为按你们愚妄的想法，圣子与圣父一样，并非知晓万事，也并非无所不能。我并不急于或冒失地赞扬圣子，以至胆敢说圣子所能做的比圣父更多；因为我不在圣父与圣子之间区分能力。\par

217. 但或许你们会说，圣父并未如此说，而是圣子论及圣父时犯了错。这样，你们便断定圣子不仅软弱，还犯了亵渎与撒谎之罪。然而，如果你们在论及圣父的事上不信圣子，那么在论及那事上，你们也不可信祂。因为如果祂说圣父疑惑，仿佛不知将要发生何事，以图欺骗我们，那么祂说自己不知未来，也是为了欺骗我们。祂在自己所行的事上遮掩无知，远比在祂论及圣父时所预言的事上，因结果与预言相反而显得受骗，更能让人接受。\par

218. 但圣父既未受骗，圣子也未欺骗。圣经惯常如此说话，正如我已举出的例子，以及许多其他例证所表明的那样，即天主假装不知道祂所知道的事。由此证明了在圣父与圣子内存在着\OriginalTerm{天主性的一体}{unity of Godhead}，以及品格的一致；因为，如同天主圣父隐藏祂所知之事，同样，作为天主之像的圣子在这方面也隐藏祂所知之事。\par

\Emph{他欲解释主对宗徒们的答复，便将所领受的回答归于基督的温良。随后，当别人提供另一理由时，他承认那是真实的；因为主是出于祂人性情感而说出那话。由此他推论，圣父与圣子的知识是相等的，圣子并不低于圣父。在将一段说圣子较低下的经文，与另一段宣称圣子同等的经文并置之后，他谴责\OriginalTerm{亚略派}{Arians}在判断圣子时的鲁莽，并指出，当他们邪恶地将圣子贬为低下时，祂却恰当地自称为磐石。}\par

219. 因此，我们受教导，天主子并非不知未来。如果他们承认这一点，那么我——现在为了答复祂为何宣称，不是天使，也不是子，而是唯有父知道——也要记起祂在此处对门徒们惯常的爱，以及祂的恩宠，这恩宠以其频繁出现本当为众人所知。因为主对祂的门徒满怀深情，当他们向祂询问祂认为对他们无益得知的事时，祂宁愿显得对所知的有所不知，也不愿拒绝回答。祂更乐意提供对我们有益的事，而非炫耀祂自己的能力。\par

220. 然而，有些人不像我这般胆怯。因为我宁愿敬畏天主的深奥，也不愿自作聪明。然而，有些人依据这句话：\BibleQuote{耶稣在年龄、智慧和天主及人前的恩爱上，渐渐增长，} \BibleRef{路 2:52} 便大胆地说，按照祂的天主性，祂确实不可能不知未来，但祂在摄取我们人性的状态下，以人子的身分说，祂在受难之前有所不知。因为当祂提到子时，祂并非像是在论及另一位；因为祂自己就是我们的主，天主子，也是童贞女之子。但祂用一个涵盖二者的词语，引导我们的心灵，以至于祂作为人子，由于采纳了我们的无知与知识的增长，可以被理解为尚未完全知晓一切。因为知晓未来并非我们的事。这样，在祂进展的状态中，祂似乎是无知的。因为按照天主性，祂怎会进展呢？在祂内居住着圆满的\OriginalTerm{天主性}{Godhead}。\BibleRef{哥 2:9} 或者，天主子怎会有什么不知道的呢？祂曾说：\BibleQuote{你们为什么心里思念恶事呢？} \BibleRef{玛 9:4} 经文论及祂时说：\BibleQuote{耶稣却知道他们的心意，} \BibleRef{路 6:8} 祂又怎会不知呢？\par

221. 这是别人所说的，但我——回到我先前的论点，我曾指出经上论及圣父说：\BibleQuote{他们或许会尊敬我的儿子，}——我认为这段经文之所以如此记载，是为了使圣父在谈论人时，也像是带着人的情感说话。但我更倾向于认为，那位与世人同行、度着人的生活、并取了我们的血肉的圣子，也摄取了我们的情感；好使祂能借着我们的无知，说祂不知道，尽管没有什么是祂不知道的。因为虽然祂在肉身的真实中看似是一个人，但祂却是生命、光明，并有能力从祂身上发出，\BibleRef{路 6:19} 以祂威能的大能治愈受伤者。\par

222. 那么，你们看，这事已为你们解决了，因为圣子的话被归因于祂完全摄取了我们的人性状态，而关于圣父的记载也是如此，为使你们停止对圣子的苛责。\par

223. 那时，\Person{天主子}一无所不知，因为\Person{圣父}一无所不知。但如果正如我们现在所断定的，\Person{天主子}并非一无所知，就让他们说说他们希望祂在哪方面显得低下。如果\Person{天主}生了一个比自己低的子，祂就给了祂较少的。如果祂给了较少的，祂要么是想给得少，要么是只能给得少。但\Person{圣父}既不软弱也不忌妒，因为在子之前，既无意志也无能力。试问，祂在哪方面低下呢？祂拥有的一切，正如同\Person{圣父}拥有的一样。祂因着祂的\OriginalTerm{生成}{Generation}的权利，从\Person{圣父}那里领受了一切，\BibleRef{若 16:15}，并以祂尊威的光荣完全彰显了\Person{圣父}。\par

224. 他们说，经上记载：\BibleQuote{因为父比我大。} \BibleRef{若 14:28} 经上也记载：\BibleQuote{他虽具有天主的形体，却没有以自己与天主同等为僭越。} \BibleRef{斐 2:6} 经上又记载犹太人想要杀害他，因为他说自己是天主子，使自己与天主平等。\BibleRef{若 5:18} 经上记载：\BibleQuote{我与父原是一体。} \BibleRef{若 10:30} 他们读出\Quote{一}，没有读出\Quote{多}。那么，在同一个本性内，他能既是低下又是同等的吗？不是的，前者针对他的天主性，后者针对他的肉躯。\par

225. 他们说他是低下的：我请问，谁测量过，谁如此心高气傲，竟将\Person{圣父}和\Person{圣子}置于他的审判座前，以决断谁更大？\Person{达味}说：\BibleQuote{我的心不做慢，我的眼不高视。} \Person{达味}王在人世间的事上，不敢心高气傲，但我们却在对抗\Person{天主}奥秘的事上，高抬我们的心。谁能决断\Person{天主子}的事呢？\OriginalTerm{上座者}{Thrones}、\OriginalTerm{宰制者}{Dominions}、\OriginalTerm{天使}{Angels}、\OriginalTerm{异能者}{Powers}吗？但\OriginalTerm{总领天使}{Archangels}侍立服侍他，\OriginalTerm{革鲁宾}{Cherubim}和\OriginalTerm{色辣芬}{Seraphim}在他前供职赞美他。在读到\Person{圣父}自己认识\Person{圣子}，却不审判他时，谁还敢决断\Person{天主子}呢？\BibleQuote{除了父外，没有人认识子。} \BibleRef{玛 11:27} 经上说\Quote{认识}，而不是\Quote{审判}。认识是一回事，审判是另一回事。\Person{圣父}自身内拥有知识。\Person{圣子}没有比自己更高的权力。还有：\BibleQuote{除了子外，没有人认识父；} 而且他自己认识父，如同父认识他一样。\par

226. 但你们说，他说过他是低下的，他说过他是\Quote{石头}。你们说的更多，却仍不虔敬地攻击他。我说的较少，却怀着敬意增加他的尊荣。你们说他低下，却又承认他高于\OriginalTerm{天使}{angels}。我说他低于\OriginalTerm{天使}{angels}，却并不减损他的尊荣；因为我不否认他的\OriginalTerm{天主性}{Godhead}，但我宣扬他的\OriginalTerm{慈悯}{pity}。\par

\Emph{这位\OriginalTerm{圣人}{Saint}转向\Person{天主圣父}，解释他为何不嘲笑\Person{圣子}比\Person{圣父}低下，然后他宣称轮不到他去测量\Person{天主子}，因为测量\Place{耶路撒冷}——甚至可能只是作为人的\Person{基督}——的事交给了\OriginalTerm{天使}{angel}。他说，\Person{亚略}已显明自己是\Person{撒殚}的模仿者。对\OriginalTerm{天主性的生成}{divine Generation}进行讨论是鲁莽的事。既然\Person{依撒意亚}已给出了人类生成的如此大的标记，我们就不应该在神圣事物上进行比附。最后，他通过摆出\WorkTitle{圣经}的各种例子，表明我们应当何等小心地避免\Person{亚略}式的傲慢。}\par

227. 全能的\Person{圣父}啊，我现在含着泪向你倾诉。我实在轻易地称你为不可亲近的、不可理解的、不可估量的；但我不敢说你的\Person{圣子}低于你自己。因为当我读到他是你光荣的光辉，是你本体的肖像，\BibleRef{希 1:3} 我害怕，当我说你本体的肖像是低下的时，我会显得在说你那位\Person{圣子}是肖像的本体是低下的；因为你\OriginalTerm{天主性}{Godhead}的完整全然在\Person{圣子}内。我时常读到，也诚心相信，你和你的\Person{圣子}及\Person{圣神}是无际、无量、不可估量、不可言喻的。因此我不能评估你，好能以衡量你。\par

228. 就算这样，是我怀着鲁莽孟浪的精神想要测量你吗？我请问，我从何处测量你呢？先知看见一根\OriginalTerm{麻绳}{line of flax}，\OriginalTerm{天使}{angel}用它来测量\Place{耶路撒冷}。是\OriginalTerm{天使}{angel}在测量，不是\Person{亚略}。而且他测量的是\Place{耶路撒冷}，不是\Person{天主}。甚至可能\OriginalTerm{天使}{angel}都不能测量\Place{耶路撒冷}，因为那是一位人。因此经上记载：\BibleQuote{我举目观望，看见一个人，手中有麻绳。} \BibleRef{则 40:3} 他是一位人，因为如此显示了将来要取的\OriginalTerm{身体}{body}的\OriginalTerm{预像}{type}。他是一位人，关于他经上说：\BibleQuote{在我以后来的一个人，我连给他解鞋带都不配。} \BibleRef{若 1:27} 所以，\Person{基督}以\OriginalTerm{预像}{type}的方式测量了\Place{耶路撒冷}。\Person{亚略}却测量\Person{天主}。\par

229. 甚至\Person{撒殚}也伪装成\OriginalTerm{光明的天使}{angel of light}；\BibleRef{格后 11:14} 那么，如果\Person{亚略}模仿他的\OriginalTerm{始作俑者}{Author}，竟攫取不该做的事，又有什么奇怪呢？虽然他的父\OriginalTerm{魔鬼}{devil}并未在自己身上这样做，那人却以难以容忍的亵渎，将那关于\Person{天主}\OriginalTerm{奥秘}{divine secrets}和\OriginalTerm{天上生成}{heavenly Generation}之奥迹的知识归于自己。因为\OriginalTerm{魔鬼}{devil}承认了真正的\Person{天主子}，\Person{亚略}却否认\Person{祂}。\par

230. 那么，既然我不能测量你，全能的\Person{圣父}啊，我岂能不亵渎地讨论你\OriginalTerm{生成}{Generation}的奥秘呢？当那位由你所生的他自己说：\BibleQuote{凡父所有的一切，都是我的；} \BibleRef{若 16:15} 我还能说在你和你的\Person{圣子}之间有什么多或什么少吗？\BibleQuote{谁立了我作人间事物的判官和分施者呢？} 这是\Person{子}说的，\BibleRef{路 12:14} 而我们却要在\Person{圣父}和\Person{圣子}之间进行分配并作出审判吗？一份应有的敬虔之心，即使在分配遗产时也避免\OriginalTerm{仲裁人}{arbiters}。而我们竟要成为\OriginalTerm{仲裁人}{arbiters}，在你和你\Person{圣子}之间分配\OriginalTerm{非受造本体}{uncreated Substance}之光荣吗？\par

231. \BibleQuote{这一世代，}经上说，\BibleQuote{是个邪恶的世代；它寻求征兆，除了约纳先知的征兆外，必不给它任何征兆。}\BibleRef{路 11:29} 因此，\OriginalTerm{天主性}{Godhead}的征兆并未赐下，所赐的只是\OriginalTerm{降生成人}{Incarnation}的征兆。因此，当先知要论及降生成人时，他说：\BibleQuote{你们向上主求一个征兆吧！}而当国王说：\BibleQuote{我不求，也不愿试探上主}时，得到的答复是：\BibleQuote{看，贞女将怀孕生子。} 因此，我们无法看见天主性的征兆，难道我们还要寻求它的量度吗？哀哉！我真是可怜！我们竟敢放肆地议论那位我们不能以相称的祈祷向祂祈求的天主！\par

232. 让\OriginalTerm{亚略派}{Arians}看看他们所做的是什么。我竟非法地将祢、圣父啊，与祢的受造物相比较，说祢比一切更大。若如亚略所主张的，祢大于祢的圣子，那我便作了邪恶的判断。那判断将首先针对祢。因为除非经过比较，便无法作出选择；除非先对某一位作出判断，也不能把某人置于另一人之前。\par

233. 我们指着天起誓是不合法的，但判断天主却是合法的。然而祢却将一切审判都交给了祢的圣子。\par

234. 若翰害怕为主施洗，若翰阻止祂说：\BibleQuote{我该当受祢的洗，祢却反来就我吗？}\BibleRef{玛 3:4} 难道我竟该把基督置于我的判断之下么？\par

235. 梅瑟推辞司祭职，伯多禄试图逃避所要求的职务服从；而亚略竟考察天主的深奥事理吗？但亚略并非圣神。不，这话甚至是对亚略和所有人说的：\BibleQuote{不要寻求那些超越你的事。}\BibleRef{德 3:22}\par

236. 梅瑟被阻止看见天主的面；\BibleRef{出 33:23} 亚略却有幸在暗中窥见。梅瑟和亚郎在祂的司祭中。那曾在荣耀中与主一同显现的梅瑟，那时也仅看见了天主背影的显现；亚略却面对面地完全看见了天主！但经上说：\BibleQuote{没有人能看见我的面而能生活的。}\par

237. 保禄也论及人认知的有限：\BibleQuote{我们现在所知道的，只是局部的；我们作先知所讲的，也只是局部的。}\BibleRef{格前 13:9} 亚略却说：\BibleQuote{我认识天主完全是全面的，不是局部的。} 这样，保禄反不如亚略，被选的器皿只知道局部，而丧亡的器皿却知道全部。他说：\BibleQuote{我知道一个人——或在身内，或在身外，我不知道，天主知道——他被提到乐园里，听到了不可言传的话。}\BibleRef{格后 13:3} 保禄被提到三层天上，却不自知；亚略生活在污秽中，却认识天主。保禄论自己说：\BibleQuote{天主知道；} 亚略论天主却说：\BibleQuote{我知道。}\par

238. 但亚略并未被提到天上，尽管他追随了那以可憎的狂傲擅敢僭越属神之事者，那曾说：\BibleQuote{我要将我的宝座高举在云彩之上；我要与至高者相平。}\BibleRef{依 14:14} 因为正如他所说：\BibleQuote{我要与至高者相平，} 同样，亚略也希望至高无上的圣子显得像他自己一样；他不以永恒的\OriginalTerm{天主性}{Godhead}光荣敬拜圣子，却以肉躯的软弱来度量祂。\par

\SourceRecord{Translated by H. de Romestin, E. de Romestin and H.T.F. Duckworth.；Nicene and Post-Nicene Fathers, Second Series；Vol. 10.；Edited by Philip Schaff and Henry Wace.；Buffalo, NY: Christian Literature Publishing Co.,；1896.；<http://www.newadvent.org/fathers/34045.htm>.}
